\documentclass{article}

% Numeric, compressed citations (matches the original splncs04 numeric style).
\PassOptionsToPackage{numbers,compress}{natbib}

% ---------------------------------------------------------------
% NeurIPS 2026 style file
% Use [preprint] / [final] options as appropriate; default = anonymized submission.
\usepackage[preprint]{neurips_2026}

% ---------------------------------------------------------------
% Standard packages
\usepackage[utf8]{inputenc}
\usepackage[T1]{fontenc}
\usepackage{hyperref}
\usepackage{url}
\usepackage{booktabs}
\usepackage{amsfonts}
\usepackage{nicefrac}
\usepackage{microtype}
\usepackage{xcolor}

% ---------------------------------------------------------------
% Paper-specific packages
\usepackage{graphicx}
\usepackage{mathtools,amssymb}
\usepackage{derivative}
\usepackage{algorithm}
\usepackage{algpseudocode}
\usepackage[capitalize]{cleveref}
\usepackage{pgfplots}\pgfplotsset{compat=1.18}
\usepgfplotslibrary{groupplots}% multi-panel rows via panel_profile_group.tex
\usepgfplotslibrary{fillbetween}% bootstrap CI bands in panel_profile_paper*.tex
\usetikzlibrary{matrix}        % required by per-cell legend.tex
\usepackage{subcaption}
% --- deps for figures/panel_pose_budget.tex (pose error vs minimal samples) ---
\definecolor{tabblue}{HTML}{1F77B4}
\definecolor{tabgreen}{HTML}{2CA02C}
\definecolor{tabred}{HTML}{D62728}
\pgfdeclareplotmark{hexagon}{%
  \pgfpathmoveto{\pgfpointpolar{0}{\pgfplotmarksize}}%
  \pgfpathlineto{\pgfpointpolar{60}{\pgfplotmarksize}}%
  \pgfpathlineto{\pgfpointpolar{120}{\pgfplotmarksize}}%
  \pgfpathlineto{\pgfpointpolar{180}{\pgfplotmarksize}}%
  \pgfpathlineto{\pgfpointpolar{240}{\pgfplotmarksize}}%
  \pgfpathlineto{\pgfpointpolar{300}{\pgfplotmarksize}}%
  \pgfpathclose\pgfusepathqfillstroke}

% ---------------------------------------------------------------
% Figure fragments live in figures/ alongside this file.  Generator output
% is copied here so the build is self-contained (see figures/ for the
% panel_*_group.tex fragments and the shared style files).
% Shared per-method pgfplots styles (ours_ransac, ours_msac, ...) used by
% every Profile panel.  Sensitivity styles are shared separately below.
% paper_styles_shared.tex — colour palette + per-method named
% styles for the paper-mode panels.  Identical across all
% cells/experiments.  \input once into your paper preamble
% (or copy verbatim).  Standalone preview.tex \input's this
% before any panel; main.tex should do the same once globally.

% --- Colour palette (\providecolor: paper-side override wins) ---
\providecolor{ColorGaU}{HTML}{42D4F4}
\providecolor{ColorIGMLE}{HTML}{4363D8}
\providecolor{ColorIGRobust}{HTML}{3CB44B}
\providecolor{ColorJeffreys}{HTML}{E6194B}
\providecolor{ColorMAGSAC}{HTML}{F58231}
\providecolor{ColorMSAC}{HTML}{FFE119}
\providecolor{ColorOracle}{HTML}{808080}
\providecolor{ColorRANSAC}{HTML}{FFB6C1}
\providecolor{ColorgauPerPairBest}{HTML}{42D4F4}
\providecolor{ColormagsacPerPairBest}{HTML}{F58231}
\providecolor{ColormsacPerPairBest}{HTML}{FFE119}
\providecolor{ColorransacPerPairBest}{HTML}{FFB6C1}

% --- Per-method addplot styles ---
\pgfplotsset{ours_ransac/.style={color=ColorRANSAC, line width=1.0pt, solid, mark options={solid}}}
\pgfplotsset{ours_ransac_marker/.style={ours_ransac, mark=x, mark size=2pt, mark options={solid}}}
\pgfplotsset{ours_msac/.style={color=ColorMSAC, line width=1.0pt, solid, mark options={solid}}}
\pgfplotsset{ours_msac_marker/.style={ours_msac, mark=diamond*, mark size=2pt, mark options={solid}}}
\pgfplotsset{ours_gau/.style={color=ColorGaU, line width=1.0pt, solid, mark options={solid}}}
\pgfplotsset{ours_gau_marker/.style={ours_gau, mark=triangle*, mark size=2pt, mark options={solid}}}
\pgfplotsset{ours_magsac/.style={color=ColorMAGSAC, line width=1.0pt, solid, mark options={solid}}}
\pgfplotsset{ours_magsac_marker/.style={ours_magsac, mark=star, mark size=2pt, mark options={solid}}}
\pgfplotsset{ours_jeffreys/.style={color=ColorJeffreys, line width=1.0pt, solid, mark options={solid}}}
\pgfplotsset{ours_jeffreys_marker/.style={ours_jeffreys, mark=pentagon, mark size=2pt, mark options={solid}}}
\pgfplotsset{ours_ig_robust/.style={color=ColorIGRobust, line width=1.0pt, solid, mark options={solid}}}
\pgfplotsset{ours_ig_robust_marker/.style={ours_ig_robust, mark=o, mark size=2pt, mark options={solid}}}
\pgfplotsset{ours_ig_mle/.style={color=ColorIGMLE, line width=1.0pt, solid, mark options={solid}}}
\pgfplotsset{ours_ig_mle_marker/.style={ours_ig_mle, mark=square, mark size=2pt, mark options={solid}}}
\pgfplotsset{ours_oracle/.style={color=black, line width=1.0pt, solid}}

% Numbers the prose reads from its data (generated by ../scripts/paper_numbers.jl).
% generated by scripts/paper_numbers.jl from configs/paper_numbers.toml; do not edit
\newcommand{\BaselinesStdVtwoLo}{0.04}
\newcommand{\BaselinesStdVtwoHi}{0.38}
\newcommand{\StdGapVtwoLo}{2}
\newcommand{\StdGapVtwoHi}{29}
\newcommand{\ConvRatioPT}{128}
\newcommand{\ConvRatioETH}{8}


% ---------------------------------------------------------------
% Math commands
\DeclareMathOperator{\RSS}{RSS}
\DeclareMathOperator{\Cov}{Cov}
\DeclareMathOperator{\smax}{smax}
\DeclareMathOperator{\blockdiag}{blockdiag}
\DeclareMathOperator{\diag}{diag}

\makeatletter
% Auto-generate \bA..\bZ and \ba..\bz as \mathbf shortcuts.
% Skip lowercase 'f' — \bf is the deprecated LaTeX2e font command.
% \edef captures the current value of \@tmp at definition time;
% \noexpand keeps \mathbf as a token in the definition body.
\@for\@tmp:={A,B,C,D,E,F,G,H,I,J,K,L,M,N,O,P,Q,R,S,T,U,V,W,X,Y,Z,%
              a,b,c,d,e,g,h,i,j,k,l,m,n,o,p,q,r,s,t,u,v,w,x,y,z}\do{%
  \expandafter\edef\csname b\@tmp\endcsname{\noexpand\mathbf{\@tmp}}%
}
% Auto-generate \cA..\cZ as \mathcal shortcuts.
\@for\@tmp:={A,B,C,D,E,F,G,H,I,J,K,L,M,N,O,P,Q,R,S,T,U,V,W,X,Y,Z}\do{%
  \expandafter\edef\csname c\@tmp\endcsname{\noexpand\mathcal{\@tmp}}%
}
% Auto-generate \balpha..\bomega and \bGamma..\bOmega as \boldsymbol Greek
% shortcuts.  Pattern: \b<greekname> -> \boldsymbol{\<greekname>}.
% Skip 'eta': \b+eta = \beta would clash with the standard Greek letter.
\@for\@tmp:={alpha,beta,gamma,delta,epsilon,varepsilon,zeta,theta,vartheta,%
              iota,kappa,lambda,mu,nu,xi,pi,varpi,rho,varrho,sigma,varsigma,tau,%
              upsilon,phi,varphi,chi,psi,omega,%
              Gamma,Delta,Theta,Lambda,Xi,Pi,Sigma,Upsilon,Phi,Psi,Omega}\do{%
  \expandafter\edef\csname b\@tmp\endcsname{%
    \noexpand\boldsymbol{\expandafter\noexpand\csname\@tmp\endcsname}%
  }%
}
\makeatother

% Cases the loops cannot reach
\newcommand{\bzero}{\mathbf{0}}             % digit, not a letter
\newcommand{\btSigma}{\widetilde{\bSigma}}  % composite (tilde + bold Sigma)
\newcommand{\NN}{\mathcal{N}}               % readable alias for \cN (normal dist.)
\newcommand{\Reals}{\mathbb{R}}
\newcommand{\IG}{\mathrm{IG}}

% Inlier/outlier densities
\newcommand{\pin}{p_{\mathrm{in}}}
\newcommand{\pout}{p_{\mathrm{out}}}
% Jacobian notation: use \pdv{f}!{x} from the derivative package (renders ∂_x f)

% ---------------------------------------------------------------
\title{RANSAC Scoring Done Right}

\author{%
  James Pritts \quad Felix Seegräber \quad Kevin Köser \\
  Marine Data Science, Institut für Informatik \\
  Kiel University, Germany \\
  \texttt{jpr@informatik.uni-kiel.de} \quad
  \texttt{fse@informatik.uni-kiel.de} \quad
  \texttt{kk@informatik.uni-kiel.de} \\
}

\begin{document}

\maketitle

\begin{abstract}
	The most widely used RANSAC variants score candidate models by
	counting inliers or summing per-point scores that saturate beyond
	a residual threshold. Every such score requires a user-supplied
	parameter that is a function of the inlier scale, which must
	itself be estimated from contaminated data.  We remove this
	dependence by reversing the usual order of inference: rather than
	estimating the scale and then scoring against it, we marginalize
	$\sigma$ analytically in closed form under a conjugate
	Inverse-Gamma prior for a fixed inlier partition, then optimize
	over partitions.  A single closed-form expression spans the
	non-informative Jeffreys limit and informative empirical-Bayes
	priors, so the same score adapts across data-rich and data-scarce
	regimes without any change to the algorithm.  The proposed RANSAC
	score is the first in which the inlier scale is genuinely absent
	from the formula.  The score admits $O(N \log N)$ computation via
	sort-and-sweep.  On a benchmark of nearly $70\,000$ image pairs
	spanning different two-view estimation problems and both
	engineered and learned feature pipelines, the proposed score
	exceeds the state of the art (RANSAC, MSAC, GaU,
	MAGSAC++): it stays nearly flat under threshold miscalibration
	where baselines degrade, reaches near-optimal accuracy from as few
	as two validation pairs where baselines need $\ConvRatioPT\times$
	more on PhotoTourism, and tightens its prior regularization as validation data
	grows scarce.\footnote{Code will be released upon acceptance.}
\end{abstract}

% =====================================================================
\section{Introduction}
\label{sec:intro}
% =====================================================================

RANSAC~\cite{Fischler1981} scores candidate geometric models by
evaluating how well they explain the observed data in the presence
of outliers.  Every widely used scoring function, whether inlier
counting (RANSAC), truncated squared error (MSAC~\cite{Torr1998}),
or the Gaussian-Uniform marginal likelihood
(MLESAC~\cite{Torr2000}), depends on a hyperparameter that
has an implicit dependence on the inlier scale.  The user
must either supply this quantity or estimate it from
contaminated data, neither of which is straightforward in
practice.

Shekhovtsov~\cite{Shekhovtsov2025} provides a unifying analysis
of these scoring functions through the lens of a single
probabilistic model: the Gaussian-Uniform (GaU) mixture, in which
each inlier residual follows a Gaussian density and each
outlier is uniform.  Within this framework, MSAC is the
\emph{profile} likelihood (maximizing over latent inlier labels)
and MLESAC is the \emph{marginal} likelihood (summing over
labels); the two differ only in replacing a hard~$\max$ with a
soft maximum.  Both are parameterized by a decision
threshold~$\tau$ that is itself a quantile of the inlier scale,
$\tau = q\,\sigma$ for a fixed quantile factor $q$: tuning $\tau$
is equivalent to tuning $\sigma$ in different units.  We call
scores of this form \emph{threshold-based}.

MAGSAC~\cite{Barath2019} and MAGSAC++~\cite{Barath2020} were
introduced specifically to remove this dependence by
``marginalizing'' over~$\sigma$ under a uniform prior
$\sigma \sim \mathcal{U}[0, \bar\sigma]$.
Shekhovtsov~\cite{Shekhovtsov2025} deconstructs the derivation
and identifies three errors (an inlier density that assigns
zero likelihood to the ground-truth model, an event
probability equated with a density value, and IRLS weights
chosen first and integrated into a score rather than derived
from one) that cancel by numerical coincidence, leaving a
score numerically equivalent to GaU at fixed $\bar\sigma$.
The hyperparameter the user actually sets is
$\tau = \kappa\,\bar\sigma$ for a truncation quantile $\kappa$,
so MAGSAC++ does not in fact remove the inlier-scale parameter:
it replaces $\sigma$ with a proxy hyperparameter in disguise
(\Cref{sec:supp-unified}).  Moreover, Shekhovtsov's experiments
show that when the threshold is properly tuned, \emph{all} of
these scoring functions (MSAC, GaU, MAGSAC++, and even a
learned per-scene score) perform identically for model
selection.  Only vanilla inlier counting is measurably worse.
The bottleneck is not the choice of scoring function but the
dependence on a well-chosen inlier-scale quantile.  We
extend this finding: the threshold-based baselines (with the
exception of vanilla RANSAC) not only tie at their respective
optima but exhibit the same characteristic degradation as the
hyperparameter is perturbed and the same dependence on
validation-set size (Figs.~\ref{fig:profile-grid},
\ref{fig:sensitivity-grid}).  The proposed score breaks this
pattern, evidencing a structural difference from the
threshold-based family.

We propose a scoring framework that \emph{genuinely} eliminates the
inlier scale.  The key idea is to reverse the usual elimination order.
The scores above first remove the latent inlier/outlier labels (by
summation in the marginal likelihood or by maximization in the profile
likelihood), leaving a score that still depends on a quantile of
$\sigma$.  We instead marginalize $\sigma$ analytically under a
conjugate Inverse-Gamma prior given a fixed inlier partition; the
non-informative Jeffreys prior arises as a limiting case of this
Inverse-Gamma family, so the same closed form covers both informative
and non-informative priors.  A single score therefore adapts to data
availability without any change to the algorithm: an empirical-Bayes
prior fit from a small validation set when labelled data exists, and
the Jeffreys limit, which needs no validation data at all, when it
does not.  We then optimize over partitions, which requires only a
sort over residuals.

The resulting \emph{scale-free marginal score} is a closed-form
expression involving the residual sum of squares, the inlier
count, and the outlier domain half-width~$a$
(\Cref{sec:scoring}), and admits
$O(N\log N)$ computation via a sort-and-sweep algorithm.
Crucially, $a$ is \emph{not} a relabelled inlier scale: it
parameterizes the outlier component of the mixture and is
not tied by any quantile relation to the inlier scale.  After marginalization the inlier scale
is genuinely absent from the formula; this is the structural
sense in which our score is the first in the RANSAC family to
properly marginalize $\sigma$, and, where the per-pair inlier
count is large enough for the data to dominate the prior, the
score is markedly less sensitive to $a$ than threshold-based
scores are to their hyperparameter.

We validate these claims at scale on nearly $70\,000$ image
pairs spanning homography, essential-, and fundamental-matrix
estimation with both engineered (RootSIFT) and learned
(SuperPoint+LightGlue) feature pipelines (\Cref{sec:experiments}).
The score matches the best-tuned threshold-based baselines at
their optima while being far more robust to threshold
miscalibration where the data informs the scale and reaching
near-optimal accuracy from as few as two validation pairs, where
the baselines need $\ConvRatioPT\times$ more on PhotoTourism.

% =====================================================================
\section{Probabilistic Model for Robust Fitting}
\label{sec:partition}
% =====================================================================

Consider $N$ observations $\{\bx_i\}_{i=1}^N$ and a candidate
geometric model~$\theta \in \Reals^{n_\theta}$.  Each observation
imposes $d_g$ algebraic constraints
$\bg(\bx_i,\theta) = \bzero$; the minimal sample size is
$m = \lceil n_\theta/d_g \rceil$
(\Cref{tab:model-dims} in the supplement lists
$(n_\theta, d_g, m)$ for common models).
We adopt the standard two-component
mixture~\cite{Torr2000,Shekhovtsov2025}: a latent label
$z_i \in \{0,1\}$ marks observation~$i$ as inlier ($z_i{=}1$)
or outlier ($z_i{=}0$) with prior $\Pr(z_i{=}1) = \gamma$, and
the residual $\br_i \in \Reals^{d_g}$ (an algebraic residual
vector, e.g.\ a Sampson-type residual) has conditional density
\begin{equation}
	p(\br_i \mid z_i, \theta, \sigma)
	= \begin{cases}
		\pin(\br_i;\sigma) \coloneqq
		(2\pi\sigma^2)^{-d_g/2}\,
		\exp\!\bigl({-}\|\br_i\|^2/2\sigma^2\bigr)
		 & z_i = 1, \\[3pt]
		\pout \coloneqq (2a)^{-d_g}
		 & z_i = 0,
	\end{cases}
	\label{eq:component-densities}
\end{equation}
a Gaussian inlier density and a uniform outlier density on
$[-a,a]^{d_g}$; the half-width~$a$ has the same units
as~$\sigma$ (e.g.\ pixels for a Sampson-type distance).
By independence, the joint density of residuals and labels
$\bz = (z_1,\ldots,z_N)$ factorizes as
\begin{equation}
	p(\br, \bz \mid \theta, \sigma)
	= \prod_{i=1}^{N}
	\bigl[\gamma\, \pin(\br_i;\sigma)\bigr]^{z_i}\,
	\bigl[(1{-}\gamma)\, \pout\bigr]^{1-z_i}.
	\label{eq:joint}
\end{equation}
Maximizing~\eqref{eq:joint} over labels classifies $z_i^*=1$
iff $\gamma\,\pin(\br_i;\sigma) > (1{-}\gamma)\,\pout$;
collecting positive labels defines the inlier set
$\cI \coloneqq \{i : z_i^* = 1\}$ with complement
$\cO = \{1,\ldots,N\}\setminus\cI$, and $n_I = |\cI|$,
$n_O = |\cO|$.  This per-point rule is the profile-likelihood
labeling at fixed~$\sigma$; in \Cref{sec:scoring} we instead
treat the partition~$\cI$ as a free variable to optimize.
Given a partition, the likelihood factorizes as
\begin{equation}
	p(\br \mid \theta, \cI, \sigma)
	= \prod_{i \in \cI}
	\frac{1}{(2\pi\sigma^2)^{d_g/2}}\,
	\exp\!\Bigl({-}\frac{\|\br_i\|^2}{2\sigma^2}\Bigr)
	\;\cdot\;
	\frac{1}{(2a)^{n_O d_g}}\,.
	\label{eq:joint-partition}
\end{equation}

% =====================================================================
\section{Scale-Free Marginal Likelihood Scoring}
\label{sec:scoring}
% =====================================================================

The likelihood~\eqref{eq:joint-partition} depends on the noise
scale~$\sigma$ through the Gaussian density.  Together with the
discrete labels~$\bz$, the joint~\eqref{eq:joint} has two nuisance
quantities besides the model~$\theta$; scoring a candidate model
requires eliminating both.
Standard approaches eliminate the labels first---by summation
(marginal likelihood, e.g.\ MLESAC~\cite{Torr2000}) or
maximization (profile likelihood, e.g.\
MSAC~\cite{Torr1998})---but the resulting scores still depend
on~$\sigma$
~\cite{Shekhovtsov2025}.
In practice this means every threshold-based score
requires~$\sigma$ to be estimated from data or supplied by the
user as a hyperparameter.

% ---------------------------------------------------------------------
\subsection{Reversing the Elimination Order}
\label{sec:reverse-order}

We take a different path.  Rewrite the
joint~\eqref{eq:joint} by grouping labels into a partition
$\cI = \{i : z_i = 1\}$:
\begin{equation}
	p(\br, \bz \mid \theta, \sigma)
	= \gamma^{n_I}(1{-}\gamma)^{n_O}
	\;\cdot\; p(\br \mid \theta, \cI, \sigma),
	\label{eq:joint-partition-decomp}
\end{equation}
where $p(\br \mid \theta, \cI, \sigma)$ is the partition-conditioned
likelihood~\eqref{eq:joint-partition} and the prefactor depends only
on~$|\cI|$.  Rather than eliminating labels first (which
entangles~$\sigma$ with the per-point sums or maxima), we
\emph{reverse the order}: for a \emph{fixed} partition~$\cI$,
marginalize~$\sigma$ in closed form, then optimize over~$\cI$.

Concretely, for a fixed partition~$\cI$ we marginalize~$\sigma$
under a conjugate Inverse-Gamma prior
$\pi(\sigma^2) = \IG(\alpha_0, \beta_0)$,
defining the \emph{scale-marginalized log-likelihood}
\begin{equation}
	S(\theta, \cI)
	\;\coloneqq\; \log \int_0^\infty
	p(\br \mid \theta, \cI, \sigma)\;
	\pi(\sigma^2)\, d\sigma^2.
	\label{eq:S-def}
\end{equation}
The Inverse-Gamma is the conjugate prior for the variance of
a zero-mean Gaussian, so the integral admits a closed form for
any $(\alpha_0, \beta_0) \ge 0$.  The non-informative
Jeffreys prior~$\pi(\sigma^2) \propto 1/\sigma^2$~\cite{Torr2002}
arises as the limit $(\alpha_0, \beta_0) \to (0,0)$ and
introduces no scale assumption; it is improper but its
posterior is proper whenever the effective degrees of freedom
$\nu = n_I\,d_g - n_\theta > 0$---a condition enforced by
Algorithm~\ref{alg:score}.  We carry $(\alpha_0, \beta_0)$
through the derivation symbolically; the Jeffreys special case
$(\alpha_0,\beta_0)=(0,0)$ is recovered in Section~\ref{sec:score},
and two empirical-Bayes estimators for the informative case are
deferred to Appendix~\ref{sec:supp-prior-estimation}; a
posterior and pseudo-observation reading of $(\alpha_0,\beta_0)$
is given in Appendix~\ref{sec:supp-pseudo-obs}.
The full score from~\eqref{eq:joint-partition-decomp} also
includes the label prior
$\gamma^{n_I}(1{-}\gamma)^{n_O}$.  For general~$\gamma$ this
adds $n_I\log\!\tfrac{\gamma}{1-\gamma}$ (plus a constant) to
the score, biasing the inlier/outlier boundary.  Since our goal is to avoid
such hyperparameters, we adopt the principled non-informative
choice $\gamma = \tfrac{1}{2}$, for which the prefactor becomes
$2^{-N}$---constant across all partitions---and drops out of the
optimization.\footnote{One could instead marginalize~$\gamma$
under a conjugate Beta prior, exactly as we marginalize~$\sigma$:
$\int_0^1 \gamma^{n_I}(1{-}\gamma)^{n_O}\,\pi(\gamma)\,d\gamma$
is closed-form and, depending only on the count~$n_I$, preserves
the sort-and-sweep optimality and $O(N\log N)$ cost.  We fix
$\gamma=\tfrac12$ instead so that the label prior is
partition-independent and the consensus decision is driven solely
by fit quality and outlier volume rather than a separately
estimated inlier rate, which itself depends on the partition.}

Our overall model score is then
\begin{equation}
	S^*(\theta) \coloneqq \max_{\cI} \; S(\theta, \cI).
	\label{eq:profile-I}
\end{equation}
This is a \emph{profile-marginal} hybrid: profile over
labels~$\bz$ (equivalently partitions~$\cI$), marginal
over~$\sigma$.  By contrast, the standard marginal likelihood
(e.g.\ MLESAC~\cite{Torr2000}) and profile likelihood
(e.g.\ MSAC~\cite{Torr1998}) both eliminate labels for
\emph{fixed}~$\sigma$, so the resulting score still
requires~$\sigma$ as input.  A fully Bayesian score would
marginalize \emph{both} $\bz$ and~$\sigma$, but the resulting
sum $\log \sum_{\cI} \exp S(\theta,\cI)$ has $2^N$ terms and is
intractable.  Our approach avoids this combinatorial sum while
still eliminating~$\sigma$ analytically.  The key enabler is
that the Gaussian likelihood over the inlier
partition~\eqref{eq:joint-partition} admits a closed-form
$\sigma$-integral, which we evaluate in \Cref{sec:supp-score-deriv}.

% ---------------------------------------------------------------------
\subsection{The Scale-Free Marginal Score}
\label{sec:score}

Marginalizing $\sigma^2$ under the conjugate Inverse-Gamma prior
admits a closed form; accounting for the $n_\theta$ parameters
consumed by fitting~$\theta$ (which reduce the effective degrees
of freedom from $n_I\,d_g$ to $\nu = n_I\,d_g - n_\theta$) and
taking logarithms yields the \emph{scale-free marginal score}
(\Cref{sec:supp-score-deriv} gives the full derivation):
\begin{equation}
	\boxed{
		S(\theta, \cI)
		= \log\Gamma\!\Bigl(\frac{\nu}{2} + \alpha_0\Bigr)
		- \Bigl(\frac{\nu}{2} + \alpha_0\Bigr)
		\log\!\Bigl(\frac{\RSS_\cI}{2} + \beta_0\Bigr)
		- \frac{\nu}{2}\log(2\pi)
		- n_O\, d_g\,\log(2a),
	}
	\label{eq:score}
\end{equation}
where $\nu = n_I\,d_g - n_\theta$
and $\RSS_\cI = \sum_{i\in\cI} q_i$ with $q_i = \|\br_i\|^2$
the per-point squared residual.  The best model maximizes
$S$ jointly over $\theta$ and~$\cI$.  All four terms
of~\eqref{eq:score} depend on the partition~$\cI$:
\begin{enumerate}
	\item $\log\Gamma(\nu/2 + \alpha_0)$: \emph{Inlier reward}---grows
	      super-linearly with the effective degrees of freedom plus
	      the prior pseudo-observation count $2\alpha_0$, encouraging
	      larger consensus sets.
	\item $-(\nu/2 + \alpha_0)\log(\RSS_\cI/2 + \beta_0)$: \emph{Fit
		      quality}---penalizes large residuals among the inliers,
	      smoothed by the prior pseudo-RSS $2\beta_0$ to prevent
	      collapse when the data alone would drive $\RSS_\cI \to 0$.
	      Writing $\hat s^2 \coloneqq (\RSS_\cI + 2\beta_0)/(\nu + 2\alpha_0)$
	      for the pseudo-observation variance estimator (data RSS plus
	      prior pseudo-RSS, divided by data plus prior DOF; equivalently
	      $1/\mathbb{E}[\sigma^{-2}\mid\cI]$ under the posterior), this term
	      favors partitions with small~$\hat s$.
	\item $-(\nu/2)\log(2\pi)$: Gaussian normalizer---a mild
	      per-inlier penalty that partially offsets the
	      $\log\Gamma$ reward.
	\item $-n_O\,d_g\,\log(2a)$: \emph{Outlier penalty}---each
	      outlier incurs a cost of $d_g\log(2a)$.  Adding a marginal
	      point to~$\cI$ is beneficial only if its squared residual is
	      small enough that the fit quality improvement outweighs the
	      lost outlier penalty.
\end{enumerate}

\paragraph{Noninformative limit (Jeffreys prior).}
At $(\alpha_0, \beta_0) = (0,0)$ the IG prior collapses to the
improper Jeffreys prior $\pi(\sigma^2) \propto 1/\sigma^2$, and
the score~\eqref{eq:score} simplifies to
\begin{equation}
	S_J(\theta, \cI)
	= \log\Gamma\!\Bigl(\frac{\nu}{2}\Bigr)
	- \frac{\nu}{2}\,\log(\pi\,\RSS_\cI)
	- n_O\, d_g\,\log(2a),
	\label{eq:score-jeffreys}
\end{equation}
where the combination
$(\nu/2)\log(2\pi) + (\nu/2)\log(\RSS_\cI/2) =
	(\nu/2)\log(\pi\,\RSS_\cI)$ has been merged.  This is the
``no scale assumption'' specialization: every term is data-only,
without any pseudo-observation contribution from the prior.
$S_J$ is thus the natural baseline against which an informative
prior adds value only when the data are scale-poor.

\paragraph{Pseudo-observation reading.}\label{par:pseudo-obs}
The conjugate posterior of~$\sigma^2$ given the inlier residuals
is
\begin{equation}
  \sigma^2 \mid \cI \;\sim\;
  \IG\!\Bigl(
  \alpha_0 + \tfrac{\nu}{2},\;
  \beta_0 + \tfrac{\RSS_\cI}{2}
  \Bigr),
  \label{eq:ig-posterior}
\end{equation}
the integrand of the scale integral~\eqref{eq:marg-sigma} read
as an Inverse-Gamma kernel.  It reads the prior parameters as
fictitious data: the prior contributes $2\alpha_0$
\emph{pseudo-residuals} carrying $2\beta_0$ of \emph{pseudo-RSS},
i.e.\ a prior pseudo-scale $\beta_0/\alpha_0$.  Thus $\alpha_0$
controls \emph{how much} the prior pulls and $\beta_0/\alpha_0$
\emph{toward what scale}, yielding three regimes:
\begin{itemize}
  \item \emph{Data-dominated} ($\nu \gg 2\alpha_0$): the prior is
        negligible and the score reduces to Jeffreys~\eqref{eq:score-jeffreys}.
  \item \emph{Prior-dominated} ($\nu \ll 2\alpha_0$): the partition's
        RSS contributes little and the score behaves as if $\sigma^2$
        were locked at $\beta_0/\alpha_0$.
  \item \emph{Balanced} ($\nu \sim 2\alpha_0$): the prior smooths the
        per-partition scale by $(\nu+2\alpha_0)/\nu$ toward $\beta_0/\alpha_0$.
\end{itemize}
This adaptivity---data-driven where inliers are plentiful,
prior-regularized where they are scarce---underlies the
data-rich/data-scarce behavior in \Cref{sec:experiments}
(further detail in \Cref{sec:supp-pseudo-obs}).

% ---------------------------------------------------------------------
\subsection{The Sort-and-Sweep Algorithm}

For a fixed~$\theta$, the per-point scores $q_i = \|\br_i\|^2$
depend only on~$\theta$, not on the partition.  Sorting them
$q_{(1)} \leq \cdots \leq q_{(N)}$, the optimal inlier set is
always a prefix $\cI^* = \{(1),\ldots,(k^*)\}$: swapping any
included point for an excluded one with a smaller residual
lowers~$\RSS_\cI$ at fixed~$n_I$ and strictly raises the score.
Since $\RSS_\cI = \sum_{j=1}^{k} q_{(j)}$ is then a prefix sum, a
running accumulator over every prefix with positive degrees of
freedom ($\nu = k\,d_g - n_\theta > 0$) evaluates the full score
in $O(N\log N)$; the Jeffreys limit ($\alpha_0{=}\beta_0{=}0$) is the
same algorithm without the prior pseudo-counts.
Algorithm~\ref{alg:score} summarizes the
procedure.  Within RANSAC, the
sweep is applied to each candidate model generated from a
minimal sample~\cite{Fischler1981}; the model with the highest
score~$S^*$ is selected.

\begin{algorithm}[h]
	\caption{Sort-and-sweep computation of the
		scale-free marginal score~\eqref{eq:score}.
		The algorithm returns the best-scoring prefix
		$\cI^* = \{(1),\ldots,(k^*)\}$ in $O(N\log N)$ time.}
	\label{alg:score}
	\begin{algorithmic}[1]
		\Require per-point scores
		$q_1,\ldots,q_N$;
		codimension~$d_g$;
		model parameters~$n_\theta$;
		outlier half-width $a$;
		prior parameters $(\alpha_0, \beta_0)$
		\Ensure best score $S^*$; optimal inlier count $k^*$
		\State Sort so that
		$q_{(1)} \leq q_{(2)} \leq \cdots \leq q_{(N)}$
		\State $k_{\min} \gets \min\!\bigl(\lfloor n_\theta/d_g \rfloor,\, N\bigr)$
		\Comment{largest prefix with $\nu \leq 0$}
		\State $\RSS \gets \sum_{j=1}^{k_{\min}} q_{(j)}$
		\Comment{seed the zero-DOF prefix}
		\State $S^* \gets -\infty$, \quad $k^* \gets 0$
		\For{$k = k_{\min}+1, \ldots, N$}
		\State $\RSS \gets \RSS + q_{(k)}$
		\State $\nu \gets k\,d_g - n_\theta$
		\Comment{positive DOF guaranteed}
		\State $S \gets \log\Gamma\!\bigl(\nu/2 + \alpha_0\bigr)
			- \bigl(\tfrac{\nu}{2}+\alpha_0\bigr)\log\!\bigl(\tfrac{\RSS}{2}+\beta_0\bigr)
			- \tfrac{\nu}{2}\log(2\pi)
			- (N{-}k)\,d_g\,\log(2a)$
		\If{$S > S^*$}
		\State $S^* \gets S$, \quad $k^* \gets k$
		\EndIf
		\EndFor
	\end{algorithmic}
\end{algorithm}

\paragraph{Implicit threshold.}
Although the score eliminates~$\sigma$ explicitly, the sweep's
optimal boundary defines an \emph{implicit} threshold~$\tau^*$
that, in the large-sample limit, satisfies the MAP condition
$\pin(\tau^*;\hat s) = \pout$ with $\hat s = \sqrt{\RSS_{k^*}/\nu}$,
giving the closed form
\begin{equation}
  \tau^* = \hat s\,\sqrt{2\log\!\bigl(a\sqrt{2/\pi}/\hat s\bigr)}
  \qquad\bigl(\approx 2\text{--}3\,\hat s \ \text{for}\ a/\hat s \in [10,100]\bigr),
\end{equation}
the usual $2$--$3\sigma$ range.  This matches the GaU threshold at an
effective prior $\gamma_{\mathrm{eff}} = \tfrac12$, so the
non-informative $\gamma = \tfrac12$ (\Cref{sec:reverse-order}) is
self-consistent.  Crucially $\tau^*$ \emph{adapts per hypothesis}
through~$\hat s$, whereas the baselines apply one global~$\tau$;
in the high-$\nu$ limit our score is the Bayes-optimal
Gaussian--Uniform ranking objective and threshold-based scoring
its $\sigma$-fixed specialization (\Cref{sec:implicit-threshold}).

% =====================================================================
\section{Experiments}
\label{sec:experiments}
% =====================================================================

State-of-the-art RANSAC pipelines achieve performance through
the combined effect of multiple components: guided sampling
(PROSAC~\cite{ChumMatas2005}, NAPSAC~\cite{Myatt2002napsac}),
degeneracy testing
(DEGENSAC~\cite{Chum2005degensac}, QDEGSAC~\cite{FrahmPollefeys2006qdegsac}),
local optimization
(LO-RANSAC~\cite{Chum2003,Lebeda2012}), and
varied termination criteria~\cite{Schoenberger2025stopping}.
Many of these components are domain-specific (e.g., guided
sampling for two-view matching~\cite{Barath2019pnapsac}) and do
not transfer to other model classes.  Proper
ablation of the score from the rest is rarely performed, so a
method's apparent advantage may stem from other components
rather than the scoring strategy~\cite{Shekhovtsov2025}.  We follow~\cite{Shekhovtsov2025} in isolating the score, with
some adjustments called out in the following paragraphs.  The
takeaway is that every method evaluates the same per-pair pool
of minimal-sample hypotheses, differing only in its scoring
function and its single hyperparameter.

We evaluate each scoring method along two
complementary axes.
The \emph{Profile} experiment (\Cref{sec:profile}) sweeps each
method's hyperparameter and measures \emph{robustness to
	hyperparameter value}, the flatness of the error under
perturbation of the optimal parameter on the validation set.
The \emph{Sensitivity} experiment (\Cref{sec:sensitivity-expt})
varies only the validation-set size and measures the smallest
validation set sufficient to reach near-optimal test
performance.

\subsection{Evaluation Protocol}
\label{sec:eval-protocol}

For each image pair we sample $1\,000$ two-view geometric
relations (homography, essential, or fundamental matrix,
depending on the dataset) from pre-extracted feature
correspondences.  Each method scores this shared model pool;
the highest-scored model is evaluated against the ground-truth
transformation by \emph{orientation error}, which is the
$\mathrm{SO}(3)$ angular distance between the recovered and
oracle rotations.  The rotation is recovered per model class:
directly from the essential matrix, and via the known
calibration for the fundamental matrix ($\bF\!\to\!\bE$) and the
homography (by decomposing $\bH$).  Both Profile and Sensitivity use this same
metric.  Datasets group image pairs by scene.  The Sensitivity
experiment (\Cref{sec:sensitivity-expt}) draws its
validation/test split hierarchically by scene on HEB---resampling
sites, then pairs within each, to test cross-scene
generalization---and by a flat pair-level split on PhotoTourism
and ETH3D; the Profile bands (\Cref{sec:profile}) are pair-level
on all three (\Cref{tab:datasets}).

Shekhovtsov~\cite{Shekhovtsov2025} showed that the four
threshold-based scoring functions we compare---RANSAC~\cite{Fischler1981},
MSAC~\cite{Torr1998}, GaU~\cite{Torr2000}, and
MAGSAC++~\cite{Barath2020}---can each be reparameterized with a
single hyperparameter.  Our three proposed marginal-likelihood variants (Jeffreys,
IG-Robust, IG-MLE) likewise share a single hyperparameter, the
outlier half-width from~\eqref{eq:score}.  We write~$\kappa$ for
each method's single hyperparameter generically (the half-width
$a$ for our marginal scores, the threshold quantile for the
baselines), normalized to its per-dataset optimum~$\kappa^\star$
in all plots.  However, the physical interpretation of the role of the
hyperparameter in each family is different, and we should not
expect them to share the same numeric range nor do we constrain
them to.

Shekhovtsov~\cite{Shekhovtsov2025} hand-sets per-method search
intervals and shows that even a generous 2+~decade interval
based on a reasonable guess can disadvantage some methods,
since the interval acts as a weak hyperprior on the
hyperparameter.  We instead derive every method's interval
from data: we draw a single validation set per dataset and
evaluate each method's test-set orientation error on a coarse
log-spaced grid; the minimum-error grid point becomes the
anchor~$\kappa^\star$ (if $\kappa^\star$ lies on a boundary, the search extends
by two more decades).  The two-decade window
$[\kappa^\star/10,\,\kappa^\star\cdot 10]$ is the per-method, per-dataset
interval used for hyperparameter selection in both the
\emph{Profile} (\Cref{sec:profile}) and \emph{Sensitivity}
(\Cref{sec:sensitivity-expt}) experiments.  To parallelize
scoring we sample 100 points per decade uniformly in log-space
within this window (200 points total); hyperparameter
selection in both experiments is done over this fixed grid.
All seven methods use the same resolution and the same
two-decade width.

Our empirical-Bayes variants (IG-MLE and IG-Robust) need
ground-truth annotations---the relative pose and a set of
ground-truth point matches per pair---during the validation
phase.  Both fit the Inverse-Gamma prior on the noise
variance~$\sigma^2$ from per-pair noise-variance
estimates~$\hat\sigma^2_p$ on the validation
set~(\Cref{sec:supp-per-pair-sigma}).  \emph{IG-MLE} fits the
prior by type-II maximum likelihood: since
$\sigma^2\!\sim\!\IG(\alpha,\beta) \iff
1/\sigma^2\!\sim\!\mathrm{Gamma}(\alpha,1/\beta)$, it fits a
Gamma to $\{1/\hat\sigma^2_p\}$ and reads off
$(\alpha_0,\beta_0)$, capturing both the location and the
dispersion of the $\hat\sigma^2_p$ but remaining sensitive to
outlier per-pair estimates.  \emph{IG-Robust} instead fixes
$\alpha_0 = 1$ (two pseudo-observations) and sets
$\beta_0 = \alpha_0\,\mathrm{median}(\hat\sigma^2_p)$, trading
dispersion information for robustness to the heavy
$\hat\sigma^2_p$ tails that arise when match-set cardinality
approaches the model's degrees of freedom (demonstrated on
HEB~\cite{Barath2023} in \Cref{sec:sensitivity-expt}).  Each
$\hat\sigma^2_p$ is computed from residuals on that pair's GT
matches under the GT pose; how those annotations are obtained
is dataset-specific and detailed in \Cref{sec:supp-datasets}.

% ---------------------------------------------------------------------
\subsection{Datasets}
\label{sec:datasets}

We evaluate on three datasets, one per two-view geometric
model class, with model proposals generated from both a
hand-crafted and a state-of-the-art learned feature detector
(nearly $70\,000$ pairs in total), spanning three axes of
variation: dataset, geometric problem, and feature pipeline.
\emph{PhotoTourism} (PT)~\cite{Jin2021,RansacTutorial2020}
supplies 21 landmark scenes (44\,211 pairs) with
RootSIFT~\cite{Arandjelovic2012} features for essential-matrix
($\bE$) estimation; \emph{ETH3D}~\cite{Schoeps2017} supplies 13
sub-scenes (3\,039 pairs) with learned
SuperPoint\,+\,LightGlue~\cite{DeTone2018,Lindenberger2023}
features for fundamental-matrix ($\bF$) estimation; and the
\emph{Homography Estimation Benchmark} (HEB)~\cite{Barath2023}
supplies 10 architectural sites (21\,999 pairs) with
SIFT~\cite{Lowe2004} features for homography ($\bH$)
estimation.  The empirical-Bayes variants (IG-MLE, IG-Robust)
additionally consume per-pair ground-truth matches during
validation (\Cref{sec:eval-protocol}).  Full dataset
provenance, annotation and feature-pipeline details, and the
HEB plane-assignment caveat that surfaces in some HEB panels
below are given in \Cref{sec:supp-datasets}
(\Cref{tab:datasets}).

% ---------------------------------------------------------------------
\subsection{Profile}
\label{sec:profile}

The Profile experiment plots each method's median rotation
error across all pairs in the dataset as a function of its
hyperparameter, swept exhaustively over the two-decade
interval described in \Cref{sec:eval-protocol}.  This is not a deployable tuning protocol: it is
computationally expensive (every grid point evaluated on the
full pool), and a user is unlikely to have ground-truth
annotations for so many examples of their problem
setting (the Sensitivity experiments of
\Cref{sec:sensitivity-expt} outline a pragmatic protocol for
tuning the estimators).  But
the curve shape is informative on two fronts: each method's
minimum is its achievable lower bound on the dataset, and the
rate at which error climbs as the hyperparameter is perturbed
off-optimum over two decades is a direct readout of
sensitivity to miscalibration.  Validation~$=$~test is fine
here because no model is selected: every grid point is plotted
on its own, so the ``tune on val, report on val'' leakage mode
does not apply.

% All three Profile cells now use production data.

\begin{figure}[t]
	\centering
	% Groupplot row: y-axis annotations rendered once on the leftmost cell.
	% Per-subplot bbox = 0.32\linewidth, with 0.3cm horizontal sep between
	% cells.  All three group files share the same scaffold; the x-axis,
	% tick density, and panel_axis defaults are set on the groupplot.
	% paper_styles.tex — cell-/criterion-specific bits.
% Standalone preview: \input{paper_styles_shared}\input{paper_styles}
% BEFORE \input{panel_profile_paper}.  In main.tex,
% \input{paper_styles_shared} once globally, omit this file,
% and \renewcommand the \Panel* macros per panel-input as needed.

\providecommand{\PanelXLabel}{$\kappa / \kappa^\star$ \ \ (factor relative to method-specific optimum)}
\providecommand{\PanelYLabel}{median $e_{\mathbf{R}}$ (deg)}
\providecommand{\PanelXTicks}{-1,-0.5,0,0.5,1}
\providecommand{\PanelXTickLabels}{$\tfrac{1}{10}$,$\tfrac{1}{\sqrt{10}}$,1,$\sqrt{10}$,10}
\providecommand{\PanelArgminSize}{1.5pt}

% Named axis style — panels opt-in via `panel_axis` in their
% axis options.  Legends/insets/color-bars stay clean.
% NOTE: visual properties only — NO width/height.  Layout (size,
% subfigure dimensions, groupplot geometry) is the paper's call;
% baking width/height here would silently override paper-side
% sizing via `every axis/.append style{}` last-wins.  Standalone
% preview supplies the default size in `preview.tex`.
\pgfplotsset{panel_axis/.style={
    grid style={gray!20},
    label style={font=\footnotesize},
    tick label style={font=\scriptsize},
}}
%
	\pgfplotsset{panel_axis/.append style={
				label style={font=\scriptsize},
				title style={font=\scriptsize},
				tick label style={font=\tiny}}}%
	\renewcommand{\PanelXLabel}{$\kappa / \kappa^\star$}%
	\begin{tikzpicture}
		\begin{groupplot}[panel_axis,
				width=0.26\linewidth, height=3.5cm, scale only axis,
				group style={group size=3 by 1, horizontal sep=0.8cm,
						ylabels at=edge left},
				xmin=-1, xmax=1, xtick={\PanelXTicks},
				xticklabels={$\tfrac{1}{10}$,$\tfrac{1}{\sqrt{10}}$,$1$,$\sqrt{10}$,$10$},
				xlabel={\PanelXLabel}, ylabel={\PanelYLabel}]
			% Per-cell ymin = each cell's GT-Oracle so all three black Oracle
			% lines sit at the same visual height (bottom of plot).  Cells (a)
			% and (b) cap ymax=2 to clip the high-error failure tail; cell (c)
			% HEB auto-scales the upper bound since its data spans 2-24 deg.
			% Per-cell explicit ymin/ymax giving the GT-Oracle (black line) a
			% visible cushion below.  Each ymin set so Oracle sits ~5\% above
			% the panel bottom; ymax clipped to 2 on (a)+(b), HEB auto-ymax
			% (broader range, can't share scale).
			\let\origNGP\nextgroupplot%
			% Per-cell ymin tied to each cell's own Oracle so all three black
			% Oracle lines sit at ~5\% above panel bottom (consistent visual
			% anchor across the row).  ymax clipped to 1.5 on (a)+(b), 15 on (c).
			\renewcommand\nextgroupplot[1][]{\origNGP[#1, ymin=0.23, ymax=1.5,
					ytick={0.25,0.5,0.75,1.0,1.25,1.5}]}%
			% panel_profile_group.tex — \nextgroupplot fragment.
% Host preamble MUST load: \usepgfplotslibrary{groupplots},
%                          \usepgfplotslibrary{fillbetween}.
% Must be \input inside a \begin{groupplot}[...] with these
% options set at the GROUPPLOT level (not per-cell):
%   panel_axis, xmin=-1, xmax=1, xtick={\PanelXTicks},
%   xlabel={\PanelXLabel}, ylabel={\PanelYLabel},
%   ylabels at=edge left, yticklabels at=edge left,
%   group style={group size=N by 1, horizontal sep=...}
% That way pgfplots' native `ylabels at=edge left` shows the y-
% annotations only on the leftmost cell.  Per-cell options here
% are LIMITED to `title=` so the row's wrapper retains all axis
% control.  See `profile_row_template` for the canonical wrapper.

\nextgroupplot[title={PhotoTourism / RootSIFT / \(\bE\)}]
  \addplot[name path=U_ransac_grp, draw=none, mark=none, forget plot] coordinates {(-0.9999999999999999,1.6909721043675183) (-0.9899497487437182,1.6909721043675183) (-0.9798994974874373,1.6909721043675183) (-0.9698492462311556,1.6909721043675183) (-0.959798994974874,1.6909721043675183) (-0.9497487437185931,1.6909721043675183) (-0.9396984924623114,1.6909721043675183) (-0.9296482412060298,1.5891558175564955) (-0.9195979899497488,1.5891558175564955) (-0.9095477386934673,1.5891558175564955) (-0.8994974874371856,1.5891558175564955) (-0.8894472361809047,1.5891558175564955) (-0.879396984924623,1.5891558175564955) (-0.8693467336683414,1.5891558175564955) (-0.8592964824120605,1.5175084602983941) (-0.8492462311557789,1.5175084602983941) (-0.8391959798994972,1.5175084602983941) (-0.8291457286432162,1.5175084602983941) (-0.8190954773869346,1.5175084602983941) (-0.8090452261306529,1.5175084602983941) (-0.798994974874372,1.4756798049629363) (-0.7889447236180904,1.4756798049629363) (-0.7788944723618088,1.4756798049629363) (-0.7688442211055279,1.4756798049629363) (-0.7587939698492462,1.4756798049629363) (-0.7487437185929646,1.4756798049629363) (-0.7386934673366836,1.4258778680945918) (-0.7286432160804021,1.4258778680945918) (-0.7185929648241204,1.4258778680945918) (-0.7085427135678395,1.4258778680945918) (-0.6984924623115578,1.379869003348285) (-0.6884422110552761,1.379869003348285) (-0.6783919597989952,1.379869003348285) (-0.6683417085427136,1.379869003348285) (-0.658291457286432,1.379869003348285) (-0.6482412060301507,1.3435028922126229) (-0.6381909547738694,1.3435028922126229) (-0.6281407035175877,1.3435028922126229) (-0.6180904522613064,1.3435028922126229) (-0.6080402010050252,1.311011292683501) (-0.5979899497487435,1.311011292683501) (-0.5879396984924622,1.311011292683501) (-0.5778894472361809,1.2876179726824595) (-0.5678391959798992,1.2876179726824595) (-0.5577889447236181,1.2876179726824595) (-0.5477386934673368,1.2876179726824595) (-0.537688442211055,1.2596496080255468) (-0.5276381909547738,1.2596496080255468) (-0.5175879396984925,1.2596496080255468) (-0.5075376884422108,1.2435739503168637) (-0.4974874371859296,1.2435739503168637) (-0.4874371859296484,1.21625543280089) (-0.4773869346733666,1.21625543280089) (-0.46733668341708545,1.21625543280089) (-0.4572864321608041,1.196732022181084) (-0.44723618090452244,1.196732022181084) (-0.4371859296482412,1.196732022181084) (-0.42713567839195987,1.1808041710762522) (-0.41708542713567825,1.1808041710762522) (-0.4070351758793969,1.1659570040627787) (-0.3969849246231153,1.1659570040627787) (-0.386934673366834,1.1466633022831638) (-0.3768844221105528,1.1466633022831638) (-0.3668341708542711,1.1466633022831638) (-0.35678391959798983,1.134352884840779) (-0.34673366834170855,1.134352884840779) (-0.33668341708542693,1.1224095649739727) (-0.32663316582914564,1.1224095649739727) (-0.31658291457286436,1.109674148053414) (-0.3065326633165827,1.095113985788568) (-0.29648241206030146,1.095113985788568) (-0.2864321608040202,1.0897750058853681) (-0.27638190954773856,1.0897750058853681) (-0.2663316582914572,1.0831273832439363) (-0.25628140703517593,1.0831273832439363) (-0.24623115577889426,1.0710518436978667) (-0.236180904522613,1.072558174985197) (-0.22613065326633172,1.072558174985197) (-0.2160804020100501,1.0651276662742173) (-0.20603015075376882,1.05589390790282) (-0.19597989949748754,1.05589390790282) (-0.1859296482412058,1.0535425398759208) (-0.1758793969849246,1.0488131683232533) (-0.16582914572864338,1.0429631176116076) (-0.15577889447236168,1.0429631176116076) (-0.14572864321608034,1.035920108289624) (-0.13567839195979914,1.0380711166669316) (-0.12562814070351738,1.0377803964592858) (-0.11557788944723618,1.0377803964592858) (-0.10552763819095484,1.0246927671157169) (-0.09547738693467323,1.02511455751653) (-0.08542713567839195,1.0236488097147503) (-0.07537688442211071,1.0211461288538433) (-0.06532663316582901,1.0212233832692466) (-0.055276381909547756,1.016123800414358) (-0.04522613065326643,1.015591814130799) (-0.0351758793969848,1.0120503778250052) (-0.025125628140703526,1.0183564661298654) (-0.01507537688442226,1.0173200815105734) (-0.0050251256281406455,1.021379997373677) (0.00502512562814067,1.0212538893972751) (0.015075376884422353,1.015754191856432) (0.02512562814070361,1.0111788613332822) (0.035175879396984876,1.0160774691514711) (0.04522613065326653,1.0201197805642628) (0.05527638190954779,1.0193313843358829) (0.06532663316582907,1.01508773078444) (0.07537688442211074,1.0193023155659475) (0.08542713567839198,1.024969592700672) (0.09547738693467332,1.022669676194305) (0.10552763819095497,1.0277283383105815) (0.11557788944723618,1.0238656910342865) (0.12562814070351755,1.031103954939435) (0.13567839195979917,1.0364524679241545) (0.14572864321608045,1.0428958546030156) (0.1557788944723617,1.0505176046126898) (0.16582914572864335,1.0543404418642155) (0.17587939698492464,1.0548344191110912) (0.1859296482412059,1.0651166069670206) (0.19597989949748756,1.0703200264186796) (0.20603015075376882,1.0793262646963708) (0.21608040201005016,1.0806725092872849) (0.2261306532663318,1.0945962709605268) (0.23618090452261306,1.1063815533613999) (0.24623115577889432,1.1139153418474224) (0.256281407035176,1.122274781442547) (0.2663316582914573,1.1279644373739242) (0.27638190954773856,1.1400788732579668) (0.28643216080402023,1.1484504350355103) (0.29648241206030146,1.1591178130504185) (0.3065326633165827,1.1767810291049947) (0.31658291457286436,1.1875706003201825) (0.3266331658291457,1.2074451736480785) (0.336683417085427,1.2170053744788625) (0.34673366834170866,1.2315933557872802) (0.35678391959798994,1.2474630747213296) (0.36683417085427117,1.265318932072715) (0.3768844221105528,1.2765899166406818) (0.38693467336683407,1.3007085122515551) (0.39698492462311535,1.3174310106637308) (0.40703517587939697,1.3415421316981064) (0.4170854271356783,1.3529887504197153) (0.4271356783919599,1.3723625988641275) (0.4371859296482412,1.3929796340659557) (0.44723618090452255,1.4119660559606926) (0.45728643216080417,1.4322687338764468) (0.46733668341708545,1.4642923024684136) (0.47738693467336674,1.4858872635035965) (0.4874371859296484,1.5106901946783635) (0.4974874371859296,1.543610095855644) (0.5075376884422109,1.5672861412084116) (0.5175879396984926,1.5941332391510032) (0.5276381909547738,1.6146776987410636) (0.5376884422110552,1.6376843923150688) (0.5477386934673368,1.6769289532825455) (0.5577889447236181,1.7002742069070913) (0.5678391959798994,1.7281372126898973) (0.577889447236181,1.7578988053989641) (0.5879396984924623,1.7836531679325793) (0.5979899497487435,1.8260392357016708) (0.6080402010050252,1.8618603874748745) (0.6180904522613064,1.8908208388901018) (0.6281407035175878,1.9205892197654462) (0.6381909547738694,1.9632034555495035) (0.6482412060301507,1.9980724323351664) (0.6582914572864319,2.0405907489173445) (0.6683417085427136,2.0764808943442357) (0.6783919597989949,2.120219834098518) (0.6884422110552761,2.1537179316022117) (0.6984924623115578,2.196432077184056) (0.708542713567839,2.2253877596279543) (0.7185929648241204,2.2633927788814097) (0.7286432160804021,2.314849335717572) (0.7386934673366833,2.3584101373090762) (0.7487437185929646,2.411194455861169) (0.7587939698492463,2.4482047263595983) (0.7688442211055275,2.5176705729169964) (0.7788944723618089,2.5673403427820443) (0.7889447236180904,2.621639552914029) (0.7989949748743718,2.675232330643897) (0.8090452261306534,2.7240323796182766) (0.8190954773869347,2.775898550575101) (0.8291457286432159,2.8382021547780396) (0.8391959798994976,2.8906789741111916) (0.8492462311557789,2.9631447034320644) (0.8592964824120601,3.026346074256863) (0.8693467336683418,3.0853058717783206) (0.8793969849246231,3.161612742825711) (0.8894472361809044,3.2253875651601285) (0.899497487437186,3.2931655483583047) (0.9095477386934673,3.3502495738011002) (0.9195979899497486,3.417612229348817) (0.9296482412060302,3.4928205188511097) (0.9396984924623115,3.5618884416183545) (0.9497487437185927,3.6244365195918458) (0.9597989949748744,3.695396303757313) (0.9698492462311558,3.765149358504189) (0.979899497487437,3.831984868826202) (0.9899497487437187,3.9057952595014664) (0.9999999999999999,3.9813151391614294) };
  \addplot[name path=L_ransac_grp, draw=none, mark=none, forget plot] coordinates {(-0.9999999999999999,1.6228702158656743) (-0.9899497487437182,1.6228702158656743) (-0.9798994974874373,1.6228702158656743) (-0.9698492462311556,1.6228702158656743) (-0.959798994974874,1.6228702158656743) (-0.9497487437185931,1.6228702158656743) (-0.9396984924623114,1.6228702158656743) (-0.9296482412060298,1.53105462295071) (-0.9195979899497488,1.53105462295071) (-0.9095477386934673,1.53105462295071) (-0.8994974874371856,1.53105462295071) (-0.8894472361809047,1.53105462295071) (-0.879396984924623,1.53105462295071) (-0.8693467336683414,1.53105462295071) (-0.8592964824120605,1.4670694705403486) (-0.8492462311557789,1.4670694705403486) (-0.8391959798994972,1.4670694705403486) (-0.8291457286432162,1.4670694705403486) (-0.8190954773869346,1.4670694705403486) (-0.8090452261306529,1.4670694705403486) (-0.798994974874372,1.4243354443046448) (-0.7889447236180904,1.4243354443046448) (-0.7788944723618088,1.4243354443046448) (-0.7688442211055279,1.4243354443046448) (-0.7587939698492462,1.4243354443046448) (-0.7487437185929646,1.4243354443046448) (-0.7386934673366836,1.3771982722816307) (-0.7286432160804021,1.3771982722816307) (-0.7185929648241204,1.3771982722816307) (-0.7085427135678395,1.3771982722816307) (-0.6984924623115578,1.3347343774015354) (-0.6884422110552761,1.3347343774015354) (-0.6783919597989952,1.3347343774015354) (-0.6683417085427136,1.3347343774015354) (-0.658291457286432,1.3347343774015354) (-0.6482412060301507,1.2961674646612487) (-0.6381909547738694,1.2961674646612487) (-0.6281407035175877,1.2961674646612487) (-0.6180904522613064,1.2961674646612487) (-0.6080402010050252,1.2681523950889673) (-0.5979899497487435,1.2681523950889673) (-0.5879396984924622,1.2681523950889673) (-0.5778894472361809,1.24550675314646) (-0.5678391959798992,1.24550675314646) (-0.5577889447236181,1.24550675314646) (-0.5477386934673368,1.24550675314646) (-0.537688442211055,1.2195134787005086) (-0.5276381909547738,1.2195134787005086) (-0.5175879396984925,1.2195134787005086) (-0.5075376884422108,1.204561994940896) (-0.4974874371859296,1.204561994940896) (-0.4874371859296484,1.1783218046684039) (-0.4773869346733666,1.1783218046684039) (-0.46733668341708545,1.1783218046684039) (-0.4572864321608041,1.1555529109359957) (-0.44723618090452244,1.1555529109359957) (-0.4371859296482412,1.1555529109359957) (-0.42713567839195987,1.1424544327175663) (-0.41708542713567825,1.1424544327175663) (-0.4070351758793969,1.1306587524432876) (-0.3969849246231153,1.1306587524432876) (-0.386934673366834,1.1108266774042785) (-0.3768844221105528,1.1108266774042785) (-0.3668341708542711,1.1108266774042785) (-0.35678391959798983,1.098160182705975) (-0.34673366834170855,1.098160182705975) (-0.33668341708542693,1.0832690795445306) (-0.32663316582914564,1.0832690795445306) (-0.31658291457286436,1.0726993469120862) (-0.3065326633165827,1.0603662481925913) (-0.29648241206030146,1.0603662481925913) (-0.2864321608040202,1.0556238649780836) (-0.27638190954773856,1.0556238649780836) (-0.2663316582914572,1.0492120596362553) (-0.25628140703517593,1.0492120596362553) (-0.24623115577889426,1.039052198102483) (-0.236180904522613,1.0411396594872184) (-0.22613065326633172,1.0411396594872184) (-0.2160804020100501,1.033042184839221) (-0.20603015075376882,1.0245604332379565) (-0.19597989949748754,1.0245604332379565) (-0.1859296482412058,1.020845525923711) (-0.1758793969849246,1.016748831923281) (-0.16582914572864338,1.0127331321114836) (-0.15577889447236168,1.0127331321114836) (-0.14572864321608034,1.0025202261722843) (-0.13567839195979914,1.0031875950622147) (-0.12562814070351738,1.0039328428899916) (-0.11557788944723618,1.0039328428899916) (-0.10552763819095484,0.9916300361578693) (-0.09547738693467323,0.9932355153975468) (-0.08542713567839195,0.9915591357629603) (-0.07537688442211071,0.9880257294675258) (-0.06532663316582901,0.9891780716820913) (-0.055276381909547756,0.9852381149372873) (-0.04522613065326643,0.984143339369196) (-0.0351758793969848,0.980894154921554) (-0.025125628140703526,0.9877365799557513) (-0.01507537688442226,0.9868081372362226) (-0.0050251256281406455,0.990116312109567) (0.00502512562814067,0.9893888190146579) (0.015075376884422353,0.984409232702232) (0.02512562814070361,0.9809916913046003) (0.035175879396984876,0.985888062400034) (0.04522613065326653,0.9889102426483025) (0.05527638190954779,0.9888793465967036) (0.06532663316582907,0.9840232872472533) (0.07537688442211074,0.9882896020401291) (0.08542713567839198,0.9929090774308797) (0.09547738693467332,0.9923250801405904) (0.10552763819095497,0.995854967082158) (0.11557788944723618,0.9935650314610665) (0.12562814070351755,1.0002450303288817) (0.13567839195979917,1.005343486623324) (0.14572864321608045,1.0102038144958052) (0.1557788944723617,1.0177802685244055) (0.16582914572864335,1.0222098529291428) (0.17587939698492464,1.0252761108940245) (0.1859296482412059,1.034824436466052) (0.19597989949748756,1.0413392179212242) (0.20603015075376882,1.0472048398606486) (0.21608040201005016,1.0486197908494341) (0.2261306532663318,1.06149810850254) (0.23618090452261306,1.074060042818284) (0.24623115577889432,1.0815722488892905) (0.256281407035176,1.08795923110296) (0.2663316582914573,1.0932408368642053) (0.27638190954773856,1.1069513373930469) (0.28643216080402023,1.1164093080202342) (0.29648241206030146,1.1254234982741986) (0.3065326633165827,1.1407259730016646) (0.31658291457286436,1.1519736430388057) (0.3266331658291457,1.172763267111082) (0.336683417085427,1.182917685211891) (0.34673366834170866,1.1929316286778229) (0.35678391959798994,1.2093541187733845) (0.36683417085427117,1.2308143012536303) (0.3768844221105528,1.2418366388223865) (0.38693467336683407,1.2665208484791561) (0.39698492462311535,1.2806195757992134) (0.40703517587939697,1.2998722256324229) (0.4170854271356783,1.3135040500241228) (0.4271356783919599,1.3350682400636962) (0.4371859296482412,1.353000741928061) (0.44723618090452255,1.373331056348969) (0.45728643216080417,1.389233963778128) (0.46733668341708545,1.4224576254325731) (0.47738693467336674,1.4408574243198002) (0.4874371859296484,1.4648901606844795) (0.4974874371859296,1.4982485017927645) (0.5075376884422109,1.5217753840525752) (0.5175879396984926,1.5530894897642238) (0.5276381909547738,1.572510992446153) (0.5376884422110552,1.5932716325689167) (0.5477386934673368,1.6268309937375283) (0.5577889447236181,1.6524902707101) (0.5678391959798994,1.6775957667655421) (0.577889447236181,1.7092804590976327) (0.5879396984924623,1.7344928061107756) (0.5979899497487435,1.7783432952501266) (0.6080402010050252,1.8136957868126795) (0.6180904522613064,1.8369387370065018) (0.6281407035175878,1.8632775137329245) (0.6381909547738694,1.9098691269564076) (0.6482412060301507,1.9411090342051684) (0.6582914572864319,1.9801701528280016) (0.6683417085427136,2.020844585132592) (0.6783919597989949,2.06248411523878) (0.6884422110552761,2.0961164461914645) (0.6984924623115578,2.135863616100738) (0.708542713567839,2.1674273773826003) (0.7185929648241204,2.2049581885349143) (0.7286432160804021,2.249381923680015) (0.7386934673366833,2.2915771938725022) (0.7487437185929646,2.3405637769007916) (0.7587939698492463,2.3842518872330865) (0.7688442211055275,2.4417877257100598) (0.7788944723618089,2.4916916428477824) (0.7889447236180904,2.5491625915796923) (0.7989949748743718,2.6034697432106593) (0.8090452261306534,2.6552223325459057) (0.8190954773869347,2.696257903485896) (0.8291457286432159,2.7568392649447664) (0.8391959798994976,2.8112307259237395) (0.8492462311557789,2.8802802889580272) (0.8592964824120601,2.937073578409465) (0.8693467336683418,2.9912038726880157) (0.8793969849246231,3.074848297008952) (0.8894472361809044,3.1395269236421757) (0.899497487437186,3.195379269714282) (0.9095477386934673,3.2622026866984863) (0.9195979899497486,3.3287043257094617) (0.9296482412060302,3.3968982463912054) (0.9396984924623115,3.4613166363237142) (0.9497487437185927,3.5312158755739653) (0.9597989949748744,3.5963882677195977) (0.9698492462311558,3.6617104049474856) (0.979899497487437,3.7326113276188457) (0.9899497487437187,3.809473527470033) (0.9999999999999999,3.880147413393358) };
  \addplot[ours_ransac, fill, fill opacity=0.18, draw=none, mark=none, forget plot] fill between [of=U_ransac_grp and L_ransac_grp];
  \addplot+[ours_ransac, mark=none] coordinates {(-0.9999999999999999,1.6569211601165963) (-0.9899497487437182,1.6569211601165963) (-0.9798994974874373,1.6569211601165963) (-0.9698492462311556,1.6569211601165963) (-0.959798994974874,1.6569211601165963) (-0.9497487437185931,1.6569211601165963) (-0.9396984924623114,1.6569211601165963) (-0.9296482412060298,1.5601052202536028) (-0.9195979899497488,1.5601052202536028) (-0.9095477386934673,1.5601052202536028) (-0.8994974874371856,1.5601052202536028) (-0.8894472361809047,1.5601052202536028) (-0.879396984924623,1.5601052202536028) (-0.8693467336683414,1.5601052202536028) (-0.8592964824120605,1.4922889654193714) (-0.8492462311557789,1.4922889654193714) (-0.8391959798994972,1.4922889654193714) (-0.8291457286432162,1.4922889654193714) (-0.8190954773869346,1.4922889654193714) (-0.8090452261306529,1.4922889654193714) (-0.798994974874372,1.4500076246337905) (-0.7889447236180904,1.4500076246337905) (-0.7788944723618088,1.4500076246337905) (-0.7688442211055279,1.4500076246337905) (-0.7587939698492462,1.4500076246337905) (-0.7487437185929646,1.4500076246337905) (-0.7386934673366836,1.4015380701881113) (-0.7286432160804021,1.4015380701881113) (-0.7185929648241204,1.4015380701881113) (-0.7085427135678395,1.4015380701881113) (-0.6984924623115578,1.3573016903749102) (-0.6884422110552761,1.3573016903749102) (-0.6783919597989952,1.3573016903749102) (-0.6683417085427136,1.3573016903749102) (-0.658291457286432,1.3573016903749102) (-0.6482412060301507,1.3198351784369358) (-0.6381909547738694,1.3198351784369358) (-0.6281407035175877,1.3198351784369358) (-0.6180904522613064,1.3198351784369358) (-0.6080402010050252,1.289581843886234) (-0.5979899497487435,1.289581843886234) (-0.5879396984924622,1.289581843886234) (-0.5778894472361809,1.2665623629144598) (-0.5678391959798992,1.2665623629144598) (-0.5577889447236181,1.2665623629144598) (-0.5477386934673368,1.2665623629144598) (-0.537688442211055,1.2395815433630277) (-0.5276381909547738,1.2395815433630277) (-0.5175879396984925,1.2395815433630277) (-0.5075376884422108,1.22406797262888) (-0.4974874371859296,1.22406797262888) (-0.4874371859296484,1.197288618734647) (-0.4773869346733666,1.197288618734647) (-0.46733668341708545,1.197288618734647) (-0.4572864321608041,1.1761424665585398) (-0.44723618090452244,1.1761424665585398) (-0.4371859296482412,1.1761424665585398) (-0.42713567839195987,1.1616293018969093) (-0.41708542713567825,1.1616293018969093) (-0.4070351758793969,1.1483078782530332) (-0.3969849246231153,1.1483078782530332) (-0.386934673366834,1.1287449898437212) (-0.3768844221105528,1.1287449898437212) (-0.3668341708542711,1.1287449898437212) (-0.35678391959798983,1.116256533773377) (-0.34673366834170855,1.116256533773377) (-0.33668341708542693,1.1028393222592516) (-0.32663316582914564,1.1028393222592516) (-0.31658291457286436,1.0911867474827501) (-0.3065326633165827,1.0777401169905796) (-0.29648241206030146,1.0777401169905796) (-0.2864321608040202,1.0726994354317259) (-0.27638190954773856,1.0726994354317259) (-0.2663316582914572,1.0661697214400958) (-0.25628140703517593,1.0661697214400958) (-0.24623115577889426,1.0550520209001748) (-0.236180904522613,1.0568489172362077) (-0.22613065326633172,1.0568489172362077) (-0.2160804020100501,1.0490849255567192) (-0.20603015075376882,1.0402271705703883) (-0.19597989949748754,1.0402271705703883) (-0.1859296482412058,1.037194032899816) (-0.1758793969849246,1.0327810001232671) (-0.16582914572864338,1.0278481248615456) (-0.15577889447236168,1.0278481248615456) (-0.14572864321608034,1.0192201672309542) (-0.13567839195979914,1.0206293558645732) (-0.12562814070351738,1.0208566196746387) (-0.11557788944723618,1.0208566196746387) (-0.10552763819095484,1.008161401636793) (-0.09547738693467323,1.0091750364570384) (-0.08542713567839195,1.0076039727388553) (-0.07537688442211071,1.0045859291606845) (-0.06532663316582901,1.005200727475669) (-0.055276381909547756,1.0006809576758227) (-0.04522613065326643,0.9998675767499975) (-0.0351758793969848,0.9964722663732796) (-0.025125628140703526,1.0030465230428083) (-0.01507537688442226,1.002064109373398) (-0.0050251256281406455,1.005748154741622) (0.00502512562814067,1.0053213542059665) (0.015075376884422353,1.000081712279332) (0.02512562814070361,0.9960852763189412) (0.035175879396984876,1.0009827657757526) (0.04522613065326653,1.0045150116062826) (0.05527638190954779,1.0041053654662933) (0.06532663316582907,0.9995555090158466) (0.07537688442211074,1.0037959588030383) (0.08542713567839198,1.0089393350657758) (0.09547738693467332,1.0074973781674477) (0.10552763819095497,1.0117916526963697) (0.11557788944723618,1.0087153612476765) (0.12562814070351755,1.0156744926341583) (0.13567839195979917,1.0208979772737392) (0.14572864321608045,1.0265498345494104) (0.1557788944723617,1.0341489365685477) (0.16582914572864335,1.0382751473966791) (0.17587939698492464,1.0400552650025578) (0.1859296482412059,1.0499705217165363) (0.19597989949748756,1.055829622169952) (0.20603015075376882,1.0632655522785097) (0.21608040201005016,1.0646461500683595) (0.2261306532663318,1.0780471897315334) (0.23618090452261306,1.090220798089842) (0.24623115577889432,1.0977437953683564) (0.256281407035176,1.1051170062727536) (0.2663316582914573,1.1106026371190647) (0.27638190954773856,1.1235151053255068) (0.28643216080402023,1.1324298715278722) (0.29648241206030146,1.1422706556623086) (0.3065326633165827,1.1587535010533296) (0.31658291457286436,1.169772121679494) (0.3266331658291457,1.1901042203795802) (0.336683417085427,1.1999615298453767) (0.34673366834170866,1.2122624922325516) (0.35678391959798994,1.228408596747357) (0.36683417085427117,1.2480666166631726) (0.3768844221105528,1.2592132777315341) (0.38693467336683407,1.2836146803653556) (0.39698492462311535,1.299025293231472) (0.40703517587939697,1.3207071786652647) (0.4170854271356783,1.333246400221919) (0.4271356783919599,1.3537154194639118) (0.4371859296482412,1.3729901879970083) (0.44723618090452255,1.3926485561548307) (0.45728643216080417,1.4107513488272874) (0.46733668341708545,1.4433749639504934) (0.47738693467336674,1.4633723439116983) (0.4874371859296484,1.4877901776814215) (0.4974874371859296,1.5209292988242042) (0.5075376884422109,1.5445307626304934) (0.5175879396984926,1.5736113644576135) (0.5276381909547738,1.5935943455936084) (0.5376884422110552,1.6154780124419927) (0.5477386934673368,1.651879973510037) (0.5577889447236181,1.6763822388085956) (0.5678391959798994,1.7028664897277197) (0.577889447236181,1.7335896322482984) (0.5879396984924623,1.7590729870216775) (0.5979899497487435,1.8021912654758987) (0.6080402010050252,1.837778087143777) (0.6180904522613064,1.8638797879483018) (0.6281407035175878,1.8919333667491853) (0.6381909547738694,1.9365362912529556) (0.6482412060301507,1.9695907332701674) (0.6582914572864319,2.010380450872673) (0.6683417085427136,2.048662739738414) (0.6783919597989949,2.091351974668649) (0.6884422110552761,2.124917188896838) (0.6984924623115578,2.166147846642397) (0.708542713567839,2.1964075685052773) (0.7185929648241204,2.234175483708162) (0.7286432160804021,2.2821156296987937) (0.7386934673366833,2.3249936655907892) (0.7487437185929646,2.3758791163809803) (0.7587939698492463,2.4162283067963424) (0.7688442211055275,2.479729149313528) (0.7788944723618089,2.5295159928149133) (0.7889447236180904,2.5854010722468606) (0.7989949748743718,2.639351036927278) (0.8090452261306534,2.689627356082091) (0.8190954773869347,2.7360782270304984) (0.8291457286432159,2.797520709861403) (0.8391959798994976,2.8509548500174655) (0.8492462311557789,2.921712496195046) (0.8592964824120601,2.981709826333164) (0.8693467336683418,3.038254872233168) (0.8793969849246231,3.1182305199173315) (0.8894472361809044,3.182457244401152) (0.899497487437186,3.2442724090362933) (0.9095477386934673,3.3062261302497933) (0.9195979899497486,3.3731582775291393) (0.9296482412060302,3.4448593826211575) (0.9396984924623115,3.5116025389710344) (0.9497487437185927,3.5778261975829055) (0.9597989949748744,3.6458922857384555) (0.9698492462311558,3.7134298817258373) (0.979899497487437,3.782298098222524) (0.9899497487437187,3.8576343934857498) (0.9999999999999999,3.9307312762773936) };
  \addplot+[ours_ransac, mark=none, line width=0.4pt, dashed, forget plot] coordinates {(-1.0,0.39590751000620505) (1.0,0.39590751000620505)};
  \addplot+[ours_ransac, mark=*, only marks, mark size=\PanelArgminSize, forget plot] coordinates {(0.02512562814070361,0.9960852763189412)};
  \addplot[name path=U_msac_grp, draw=none, mark=none, forget plot] coordinates {(-0.9999999999999999,1.2891997594754352) (-0.9899497487437185,1.2784362707335657) (-0.979899497487437,1.271031776355509) (-0.9698492462311556,1.2657506810676706) (-0.9597989949748744,1.2583285454041047) (-0.9497487437185927,1.2536754969114583) (-0.9396984924623114,1.2496813377208762) (-0.9296482412060302,1.2424568000149796) (-0.9195979899497485,1.231549466567015) (-0.9095477386934673,1.223573473306512) (-0.8994974874371859,1.2167071814749442) (-0.8894472361809044,1.2078698286358933) (-0.879396984924623,1.2021920352925994) (-0.8693467336683417,1.1963601755320163) (-0.8592964824120601,1.1892415184140481) (-0.8492462311557789,1.1806205697987542) (-0.8391959798994976,1.1715252513164882) (-0.8291457286432159,1.1651051912137196) (-0.8190954773869346,1.1605182927855242) (-0.8090452261306534,1.1537145109785187) (-0.7989949748743717,1.1467842877593857) (-0.7889447236180904,1.141446847040968) (-0.7788944723618088,1.1346792127567797) (-0.7688442211055275,1.1276049562427495) (-0.7587939698492462,1.1227500301594733) (-0.7487437185929645,1.117279617901227) (-0.7386934673366833,1.1109013550523719) (-0.7286432160804021,1.1025250251400684) (-0.7185929648241204,1.096301738101253) (-0.708542713567839,1.092658919156792) (-0.6984924623115578,1.0897421556634972) (-0.6884422110552761,1.0812602049674365) (-0.6783919597989949,1.0743765085709174) (-0.6683417085427136,1.0689950051360582) (-0.6582914572864319,1.0642402940020286) (-0.6482412060301507,1.056514311328577) (-0.6381909547738693,1.0519303875175279) (-0.6281407035175877,1.0461602176271452) (-0.6180904522613064,1.0414226877955863) (-0.6080402010050252,1.0388088516578273) (-0.5979899497487435,1.0346547136489366) (-0.5879396984924623,1.029148377449902) (-0.577889447236181,1.025013728237725) (-0.5678391959798993,1.0215230659177874) (-0.5577889447236181,1.0161122280516053) (-0.5477386934673368,1.011174104941593) (-0.5376884422110552,1.0057240663205735) (-0.5276381909547738,1.0032781490838065) (-0.5175879396984925,1.0011117248663768) (-0.5075376884422109,0.9971653232008512) (-0.4974874371859296,0.9924662312373591) (-0.4874371859296483,0.9892901733819116) (-0.4773869346733667,0.984325772080093) (-0.4673366834170854,0.9809200003569061) (-0.4572864321608041,0.9778469365848947) (-0.4472361809045225,0.9749202176588366) (-0.43718592964824116,0.9729011934507265) (-0.4271356783919599,0.9689171886596151) (-0.41708542713567825,0.9658947089403709) (-0.40703517587939697,0.9610981482995784) (-0.3969849246231153,0.9577711348034943) (-0.386934673366834,0.955589543860085) (-0.3768844221105528,0.9509328127018132) (-0.3668341708542711,0.9483794377532552) (-0.3567839195979898,0.9447628257782191) (-0.34673366834170855,0.9410920345626429) (-0.33668341708542693,0.9384629381513581) (-0.32663316582914564,0.9389287424959073) (-0.3165829145728644,0.9367042523624562) (-0.3065326633165827,0.9332630713607821) (-0.29648241206030146,0.931752519144885) (-0.2864321608040201,0.928271566633152) (-0.27638190954773856,0.9273443942315249) (-0.2663316582914572,0.9241393880615835) (-0.256281407035176,0.9207952217810862) (-0.24623115577889426,0.9208777801069582) (-0.236180904522613,0.9193170549493531) (-0.22613065326633172,0.9189701039422736) (-0.2160804020100501,0.9173870347222062) (-0.20603015075376882,0.9158191484148324) (-0.19597989949748754,0.9148098108448155) (-0.1859296482412059,0.914091398688043) (-0.1758793969849246,0.9127588765915414) (-0.1658291457286433,0.9120210487195981) (-0.15577889447236168,0.9123386330832094) (-0.14572864321608034,0.909595068013954) (-0.13567839195979914,0.9100760320525085) (-0.12562814070351747,0.9074693832405041) (-0.11557788944723618,0.9084097727464194) (-0.10552763819095491,0.9090627194968443) (-0.09547738693467323,0.9082546766817567) (-0.08542713567839195,0.9072589456770779) (-0.07537688442211066,0.9065765972151212) (-0.06532663316582907,0.9067138200441784) (-0.055276381909547756,0.9070164798105566) (-0.04522613065326648,0.9064231202495787) (-0.03517587939698485,0.9077822300137551) (-0.025125628140703526,0.9084738990914858) (-0.01507537688442226,0.9092069910147625) (-0.005025125628140596,0.9100680944576721) (0.00502512562814067,0.9107242127440285) (0.015075376884422353,0.9116636528310808) (0.02512562814070361,0.9132977001929897) (0.035175879396984876,0.9150906946073637) (0.04522613065326653,0.9164765876428268) (0.05527638190954779,0.917672595628414) (0.06532663316582907,0.9205163755921756) (0.07537688442211074,0.921541926351161) (0.08542713567839198,0.9248365626962106) (0.09547738693467332,0.9275932417538505) (0.10552763819095497,0.9305461002712786) (0.11557788944723618,0.9340631493655062) (0.12562814070351755,0.9371801001188212) (0.13567839195979917,0.9392696180261697) (0.14572864321608045,0.943494745627188) (0.1557788944723617,0.9477320957138453) (0.16582914572864335,0.9516081321184144) (0.17587939698492464,0.9549488468660696) (0.1859296482412059,0.9579177712653116) (0.19597989949748756,0.9620935240019263) (0.20603015075376882,0.9681210545853766) (0.2160804020100501,0.971749055467549) (0.2261306532663318,0.9774608299121532) (0.236180904522613,0.9829847002189644) (0.24623115577889437,0.9866768758713541) (0.256281407035176,0.9917764435886315) (0.2663316582914573,0.9975083154391029) (0.2763819095477385,1.004037159425296) (0.2864321608040202,1.0108916824054575) (0.2964824120603015,1.018481889845843) (0.3065326633165828,1.0265770672726202) (0.31658291457286436,1.032871604069452) (0.3266331658291457,1.0404391814191343) (0.336683417085427,1.047840758058315) (0.34673366834170855,1.0543045543036758) (0.35678391959798994,1.062747133792231) (0.36683417085427117,1.0721660038728713) (0.3768844221105528,1.0818362146818017) (0.38693467336683407,1.0920792380809115) (0.39698492462311535,1.0994667484265506) (0.407035175879397,1.107624044789018) (0.4170854271356783,1.1166485662134793) (0.4271356783919599,1.1234550879803777) (0.4371859296482412,1.1338738455142345) (0.44723618090452255,1.1443851654665311) (0.45728643216080417,1.1548612561906058) (0.46733668341708545,1.167429519832745) (0.4773869346733668,1.1811267994809758) (0.4874371859296484,1.1972652443945941) (0.49748743718592964,1.2119976610263312) (0.5075376884422109,1.223866738629563) (0.5175879396984926,1.2389540778557306) (0.5276381909547738,1.2542352271593027) (0.5376884422110552,1.2687263492291823) (0.5477386934673367,1.28341103711699) (0.5577889447236181,1.3016806056251513) (0.5678391959798994,1.32063052600556) (0.5778894472361807,1.3400407518058683) (0.5879396984924622,1.3547865031639108) (0.5979899497487438,1.3701660919074852) (0.6080402010050251,1.3897971814686456) (0.6180904522613064,1.4077729937431425) (0.6281407035175879,1.4264188432032137) (0.6381909547738692,1.4427325718104613) (0.6482412060301507,1.463175695574507) (0.6582914572864321,1.48571630767926) (0.6683417085427135,1.5082681572812708) (0.678391959798995,1.5313365914842194) (0.6884422110552765,1.5573311291534198) (0.6984924623115577,1.577193172741981) (0.7085427135678392,1.6003019833775585) (0.7185929648241206,1.626200430952471) (0.7286432160804021,1.6514585849762602) (0.7386934673366833,1.6768190416409592) (0.7487437185929647,1.7014873872969942) (0.7587939698492462,1.7334538057187103) (0.7688442211055275,1.7661100512459298) (0.778894472361809,1.7959471515362724) (0.7889447236180904,1.8273985429741157) (0.7989949748743718,1.858167753895191) (0.8090452261306531,1.8896415477025221) (0.8190954773869347,1.927565724564456) (0.8291457286432159,1.9589492412876015) (0.8391959798994975,1.9936657236539708) (0.8492462311557789,2.0285115626456376) (0.8592964824120601,2.059310354582195) (0.8693467336683417,2.096868207535749) (0.8793969849246231,2.1333604660839565) (0.8894472361809044,2.1756387609970798) (0.8994974874371858,2.2174696817094333) (0.9095477386934673,2.2589482463778268) (0.9195979899497488,2.308941227157617) (0.9296482412060301,2.35330726881097) (0.9396984924623115,2.394306817113885) (0.949748743718593,2.4302295003889474) (0.9597989949748743,2.4715763400524655) (0.9698492462311558,2.51610231556773) (0.9798994974874372,2.5558778078479647) (0.9899497487437184,2.602395922418266) (0.9999999999999999,2.652697095304312) };
  \addplot[name path=L_msac_grp, draw=none, mark=none, forget plot] coordinates {(-0.9999999999999999,1.245305915727793) (-0.9899497487437185,1.2370637422933972) (-0.979899497487437,1.231256766336625) (-0.9698492462311556,1.226379665788947) (-0.9597989949748744,1.2185737350922579) (-0.9497487437185927,1.2134188201947549) (-0.9396984924623114,1.2100666313640824) (-0.9296482412060302,1.2038225690474194) (-0.9195979899497485,1.1926610708595229) (-0.9095477386934673,1.184232379670979) (-0.8994974874371859,1.1776116757933524) (-0.8894472361809044,1.1680661571911606) (-0.879396984924623,1.1623653091450354) (-0.8693467336683417,1.1563688082492196) (-0.8592964824120601,1.1503025825965039) (-0.8492462311557789,1.1434069696178888) (-0.8391959798994976,1.1354722539063438) (-0.8291457286432159,1.1288004554596291) (-0.8190954773869346,1.1240352900200394) (-0.8090452261306534,1.1172216470685778) (-0.7989949748743717,1.1115806215858914) (-0.7889447236180904,1.1067565466243356) (-0.7788944723618088,1.0998812722550448) (-0.7688442211055275,1.093359046577474) (-0.7587939698492462,1.0872492895314243) (-0.7487437185929645,1.081667869379067) (-0.7386934673366833,1.0756846287217308) (-0.7286432160804021,1.0682414077078624) (-0.7185929648241204,1.062375589984029) (-0.708542713567839,1.0586109235436043) (-0.6984924623115578,1.0558671141575313) (-0.6884422110552761,1.0472449057703488) (-0.6783919597989949,1.0411993425641213) (-0.6683417085427136,1.0370176201473782) (-0.6582914572864319,1.033426622637329) (-0.6482412060301507,1.0269208212094614) (-0.6381909547738693,1.0220994400844203) (-0.6281407035175877,1.0166258690630947) (-0.6180904522613064,1.0109956447658506) (-0.6080402010050252,1.007555393531204) (-0.5979899497487435,1.001989170880087) (-0.5879396984924623,0.996351778657105) (-0.577889447236181,0.9926089858674653) (-0.5678391959798993,0.9883478485845313) (-0.5577889447236181,0.9831872360572385) (-0.5477386934673368,0.9791798581442592) (-0.5376884422110552,0.9743130498736885) (-0.5276381909547738,0.9721801622397929) (-0.5175879396984925,0.9706160791340428) (-0.5075376884422109,0.9673770359957089) (-0.4974874371859296,0.9636633009312616) (-0.4874371859296483,0.9598240737497151) (-0.4773869346733667,0.9541531034031138) (-0.4673366834170854,0.950340273131955) (-0.4572864321608041,0.9470969765871056) (-0.4472361809045225,0.9442347721656766) (-0.43718592964824116,0.9420665045023896) (-0.4271356783919599,0.9382012675438907) (-0.41708542713567825,0.9347266314798159) (-0.40703517587939697,0.9299633781877774) (-0.3969849246231153,0.9270442991004648) (-0.386934673366834,0.9250901958187459) (-0.3768844221105528,0.9210517680954357) (-0.3668341708542711,0.9185333468618285) (-0.3567839195979898,0.9156142318896295) (-0.34673366834170855,0.9125465950245887) (-0.33668341708542693,0.9104524797669051) (-0.32663316582914564,0.9108981080421751) (-0.3165829145728644,0.9088037617933997) (-0.3065326633165827,0.9061639552068093) (-0.29648241206030146,0.9041817853406644) (-0.2864321608040201,0.9015275455962557) (-0.27638190954773856,0.9009781642500057) (-0.2663316582914572,0.8979340051036996) (-0.256281407035176,0.8940043537482059) (-0.24623115577889426,0.8946409607926461) (-0.236180904522613,0.8928538309611805) (-0.22613065326633172,0.8921212063417644) (-0.2160804020100501,0.890703655668851) (-0.20603015075376882,0.8896618035506088) (-0.19597989949748754,0.8886058275814452) (-0.1859296482412059,0.8881400703605982) (-0.1758793969849246,0.8871891513113087) (-0.1658291457286433,0.8857908736697664) (-0.15577889447236168,0.8859246872883207) (-0.14572864321608034,0.8831097486365498) (-0.13567839195979914,0.8836346056622584) (-0.12562814070351747,0.8801481776510736) (-0.11557788944723618,0.8812377887112123) (-0.10552763819095491,0.8818559077138594) (-0.09547738693467323,0.8811156156237463) (-0.08542713567839195,0.8800739785320784) (-0.07537688442211066,0.8797094061235433) (-0.06532663316582907,0.8798580810631875) (-0.055276381909547756,0.880723782327464) (-0.04522613065326648,0.880461388699213) (-0.03517587939698485,0.8820964759336006) (-0.025125628140703526,0.8827704038570663) (-0.01507537688442226,0.8834853585597707) (-0.005025125628140596,0.884536393912356) (0.00502512562814067,0.8850499232514868) (0.015075376884422353,0.8861971706101778) (0.02512562814070361,0.8871387802490048) (0.035175879396984876,0.8891626062269137) (0.04522613065326653,0.8902864603607669) (0.05527638190954779,0.8917440532355789) (0.06532663316582907,0.8942533669747691) (0.07537688442211074,0.895261135195813) (0.08542713567839198,0.8977930017670074) (0.09547738693467332,0.9003636547339503) (0.10552763819095497,0.9028020909780862) (0.11557788944723618,0.9056038015662885) (0.12562814070351755,0.9079489733117682) (0.13567839195979917,0.9096898138326499) (0.14572864321608045,0.9135081674300611) (0.1557788944723617,0.9169424180281673) (0.16582914572864335,0.9208298416557341) (0.17587939698492464,0.9241933052046241) (0.1859296482412059,0.9273143192489243) (0.19597989949748756,0.9311744149525157) (0.20603015075376882,0.9372273650101116) (0.2160804020100501,0.9407471245973116) (0.2261306532663318,0.9467987126830377) (0.236180904522613,0.9526493913272662) (0.24623115577889437,0.9571001166874825) (0.256281407035176,0.962551094037113) (0.2663316582914573,0.9678087963290389) (0.2763819095477385,0.9734334542069036) (0.2864321608040202,0.9800187881483325) (0.2964824120603015,0.9861460779809936) (0.3065326633165828,0.9929942943555619) (0.31658291457286436,0.9986549075860741) (0.3266331658291457,1.0067704266327462) (0.336683417085427,1.0136642649314804) (0.34673366834170855,1.0208127885029599) (0.35678391959798994,1.0288195245687017) (0.36683417085427117,1.037107884592822) (0.3768844221105528,1.0447095765188945) (0.38693467336683407,1.0534761617060364) (0.39698492462311535,1.0617628084195512) (0.407035175879397,1.0724226002189006) (0.4170854271356783,1.0832368109214365) (0.4271356783919599,1.0895429367022518) (0.4371859296482412,1.0999006615070406) (0.44723618090452255,1.1108607596239521) (0.45728643216080417,1.1205619148532004) (0.46733668341708545,1.1321549473690224) (0.4773869346733668,1.1435707593205142) (0.4874371859296484,1.157396772835849) (0.49748743718592964,1.1724794246935808) (0.5075376884422109,1.1854681452648155) (0.5175879396984926,1.1995851553921117) (0.5276381909547738,1.2136675812433422) (0.5376884422110552,1.229707527487782) (0.5477386934673367,1.2432070903429846) (0.5577889447236181,1.257928553145358) (0.5678391959798994,1.274575072669479) (0.5778894472361807,1.2930466415708874) (0.5879396984924622,1.309515857905582) (0.5979899497487438,1.3274315759902537) (0.6080402010050251,1.3418485738987238) (0.6180904522613064,1.3610584865355113) (0.6281407035175879,1.3813896377068342) (0.6381909547738692,1.3963333433108522) (0.6482412060301507,1.4152620467589454) (0.6582914572864321,1.434112073103729) (0.6683417085427135,1.456394621806676) (0.678391959798995,1.476992905631243) (0.6884422110552765,1.5028992156529342) (0.6984924623115577,1.5255555886512617) (0.7085427135678392,1.544682263799689) (0.7185929648241206,1.5674572633512374) (0.7286432160804021,1.5944510904098137) (0.7386934673366833,1.62024778528771) (0.7487437185929647,1.647588080872938) (0.7587939698492462,1.6723077708112069) (0.7688442211055275,1.7019459713216165) (0.778894472361809,1.7343140486679434) (0.7889447236180904,1.7623198241480946) (0.7989949748743718,1.7922596655603806) (0.8090452261306531,1.8229008230034136) (0.8190954773869347,1.8559468411669013) (0.8291457286432159,1.8897691136056105) (0.8391959798994975,1.9257076816319667) (0.8492462311557789,1.9559753749841287) (0.8592964824120601,1.98862984672845) (0.8693467336683417,2.022997508391593) (0.8793969849246231,2.062152024004592) (0.8894472361809044,2.098189207386975) (0.8994974874371858,2.137590138636244) (0.9095477386934673,2.1782934581354088) (0.9195979899497488,2.218889201434565) (0.9296482412060301,2.261491637493692) (0.9396984924623115,2.309130441272559) (0.949748743718593,2.3545933523709306) (0.9597989949748743,2.38968925744742) (0.9698492462311558,2.432660744902979) (0.9798994974874372,2.4737187047259224) (0.9899497487437184,2.5179241818704474) (0.9999999999999999,2.5565662369258355) };
  \addplot[ours_msac, fill, fill opacity=0.18, draw=none, mark=none, forget plot] fill between [of=U_msac_grp and L_msac_grp];
  \addplot+[ours_msac, mark=none] coordinates {(-0.9999999999999999,1.2672528376016141) (-0.9899497487437185,1.2577500065134815) (-0.979899497487437,1.251144271346067) (-0.9698492462311556,1.2460651734283088) (-0.9597989949748744,1.2384511402481813) (-0.9497487437185927,1.2335471585531066) (-0.9396984924623114,1.2298739845424793) (-0.9296482412060302,1.2231396845311995) (-0.9195979899497485,1.212105268713269) (-0.9095477386934673,1.2039029264887455) (-0.8994974874371859,1.1971594286341483) (-0.8894472361809044,1.187967992913527) (-0.879396984924623,1.1822786722188174) (-0.8693467336683417,1.176364491890618) (-0.8592964824120601,1.169772050505276) (-0.8492462311557789,1.1620137697083215) (-0.8391959798994976,1.153498752611416) (-0.8291457286432159,1.1469528233366744) (-0.8190954773869346,1.1422767914027818) (-0.8090452261306534,1.1354680790235483) (-0.7989949748743717,1.1291824546726386) (-0.7889447236180904,1.1241016968326518) (-0.7788944723618088,1.1172802425059123) (-0.7688442211055275,1.1104820014101118) (-0.7587939698492462,1.1049996598454488) (-0.7487437185929645,1.099473743640147) (-0.7386934673366833,1.0932929918870513) (-0.7286432160804021,1.0853832164239654) (-0.7185929648241204,1.079338664042641) (-0.708542713567839,1.0756349213501981) (-0.6984924623115578,1.0728046349105143) (-0.6884422110552761,1.0642525553688926) (-0.6783919597989949,1.0577879255675193) (-0.6683417085427136,1.0530063126417182) (-0.6582914572864319,1.0488334583196788) (-0.6482412060301507,1.0417175662690192) (-0.6381909547738693,1.037014913800974) (-0.6281407035175877,1.03139304334512) (-0.6180904522613064,1.0262091662807185) (-0.6080402010050252,1.0231821225945157) (-0.5979899497487435,1.0183219422645118) (-0.5879396984924623,1.0127500780535035) (-0.577889447236181,1.0088113570525952) (-0.5678391959798993,1.0049354572511593) (-0.5577889447236181,0.9996497320544219) (-0.5477386934673368,0.9951769815429261) (-0.5376884422110552,0.990018558097131) (-0.5276381909547738,0.9877291556617996) (-0.5175879396984925,0.9858639020002098) (-0.5075376884422109,0.98227117959828) (-0.4974874371859296,0.9780647660843104) (-0.4874371859296483,0.9745571235658134) (-0.4773869346733667,0.9692394377416034) (-0.4673366834170854,0.9656301367444305) (-0.4572864321608041,0.9624719565860002) (-0.4472361809045225,0.9595774949122566) (-0.43718592964824116,0.957483848976558) (-0.4271356783919599,0.9535592281017529) (-0.41708542713567825,0.9503106702100934) (-0.40703517587939697,0.9455307632436779) (-0.3969849246231153,0.9424077169519796) (-0.386934673366834,0.9403398698394154) (-0.3768844221105528,0.9359922903986244) (-0.3668341708542711,0.9334563923075418) (-0.3567839195979898,0.9301885288339243) (-0.34673366834170855,0.9268193147936158) (-0.33668341708542693,0.9244577089591316) (-0.32663316582914564,0.9249134252690412) (-0.3165829145728644,0.922754007077928) (-0.3065326633165827,0.9197135132837957) (-0.29648241206030146,0.9179671522427747) (-0.2864321608040201,0.9148995561147039) (-0.27638190954773856,0.9141612792407653) (-0.2663316582914572,0.9110366965826415) (-0.256281407035176,0.907399787764646) (-0.24623115577889426,0.9077593704498021) (-0.236180904522613,0.9060854429552668) (-0.22613065326633172,0.905545655142019) (-0.2160804020100501,0.9040453451955286) (-0.20603015075376882,0.9027404759827206) (-0.19597989949748754,0.9017078192131304) (-0.1859296482412059,0.9011157345243206) (-0.1758793969849246,0.8999740139514251) (-0.1658291457286433,0.8989059611946822) (-0.15577889447236168,0.8991316601857651) (-0.14572864321608034,0.8963524083252519) (-0.13567839195979914,0.8968553188573835) (-0.12562814070351747,0.8938087804457888) (-0.11557788944723618,0.8948237807288159) (-0.10552763819095491,0.8954593136053518) (-0.09547738693467323,0.8946851461527515) (-0.08542713567839195,0.8936664621045781) (-0.07537688442211066,0.8931430016693322) (-0.06532663316582907,0.893285950553683) (-0.055276381909547756,0.8938701310690103) (-0.04522613065326648,0.8934422544743958) (-0.03517587939698485,0.8949393529736779) (-0.025125628140703526,0.895622151474276) (-0.01507537688442226,0.8963461747872666) (-0.005025125628140596,0.897302244185014) (0.00502512562814067,0.8978870679977576) (0.015075376884422353,0.8989304117206293) (0.02512562814070361,0.9002182402209973) (0.035175879396984876,0.9021266504171387) (0.04522613065326653,0.9033815240017968) (0.05527638190954779,0.9047083244319964) (0.06532663316582907,0.9073848712834723) (0.07537688442211074,0.908401530773487) (0.08542713567839198,0.911314782231609) (0.09547738693467332,0.9139784482439004) (0.10552763819095497,0.9166740956246824) (0.11557788944723618,0.9198334754658973) (0.12562814070351755,0.9225645367152947) (0.13567839195979917,0.9244797159294098) (0.14572864321608045,0.9285014565286246) (0.1557788944723617,0.9323372568710063) (0.16582914572864335,0.9362189868870743) (0.17587939698492464,0.9395710760353468) (0.1859296482412059,0.942616045257118) (0.19597989949748756,0.946633969477221) (0.20603015075376882,0.9526742097977441) (0.2160804020100501,0.9562480900324303) (0.2261306532663318,0.9621297712975955) (0.236180904522613,0.9678170457731153) (0.24623115577889437,0.9718884962794183) (0.256281407035176,0.9771637688128723) (0.2663316582914573,0.9826585558840709) (0.2763819095477385,0.9887353068160998) (0.2864321608040202,0.995455235276895) (0.2964824120603015,1.0023139839134183) (0.3065326633165828,1.009785680814091) (0.31658291457286436,1.0157632558277632) (0.3266331658291457,1.0236048040259402) (0.336683417085427,1.0307525114948977) (0.34673366834170855,1.0375586714033178) (0.35678391959798994,1.0457833291804663) (0.36683417085427117,1.0546369442328467) (0.3768844221105528,1.063272895600348) (0.38693467336683407,1.072777699893474) (0.39698492462311535,1.0806147784230509) (0.407035175879397,1.0900233225039593) (0.4170854271356783,1.099942688567458) (0.4271356783919599,1.1064990123413148) (0.4371859296482412,1.1168872535106376) (0.44723618090452255,1.1276229625452416) (0.45728643216080417,1.137711585521903) (0.46733668341708545,1.1497922336008837) (0.4773869346733668,1.162348779400745) (0.4874371859296484,1.1773310086152216) (0.49748743718592964,1.192238542859956) (0.5075376884422109,1.2046674419471892) (0.5175879396984926,1.2192696166239212) (0.5276381909547738,1.2339514042013224) (0.5376884422110552,1.2492169383584821) (0.5477386934673367,1.2633090637299873) (0.5577889447236181,1.2798045793852546) (0.5678391959798994,1.2976027993375194) (0.5778894472361807,1.316543696688378) (0.5879396984924622,1.3321511805347463) (0.5979899497487438,1.3487988339488695) (0.6080402010050251,1.3658228776836847) (0.6180904522613064,1.384415740139327) (0.6281407035175879,1.403904240455024) (0.6381909547738692,1.4195329575606568) (0.6482412060301507,1.4392188711667262) (0.6582914572864321,1.4599141903914945) (0.6683417085427135,1.4823313895439734) (0.678391959798995,1.5041647485577312) (0.6884422110552765,1.530115172403177) (0.6984924623115577,1.5513743806966214) (0.7085427135678392,1.5724921235886238) (0.7185929648241206,1.5968288471518541) (0.7286432160804021,1.622954837693037) (0.7386934673366833,1.6485334134643346) (0.7487437185929647,1.674537734084966) (0.7587939698492462,1.7028807882649586) (0.7688442211055275,1.7340280112837732) (0.778894472361809,1.7651306001021079) (0.7889447236180904,1.7948591835611052) (0.7989949748743718,1.8252137097277858) (0.8090452261306531,1.8562711853529679) (0.8190954773869347,1.8917562828656787) (0.8291457286432159,1.924359177446606) (0.8391959798994975,1.9596867026429687) (0.8492462311557789,1.9922434688148831) (0.8592964824120601,2.0239701006553226) (0.8693467336683417,2.059932857963671) (0.8793969849246231,2.097756245044274) (0.8894472361809044,2.1369139841920273) (0.8994974874371858,2.1775299101728387) (0.9095477386934673,2.2186208522566178) (0.9195979899497488,2.263915214296091) (0.9296482412060301,2.307399453152331) (0.9396984924623115,2.351718629193222) (0.949748743718593,2.392411426379939) (0.9597989949748743,2.4306327987499428) (0.9698492462311558,2.4743815302353545) (0.9798994974874372,2.5147982562869435) (0.9899497487437184,2.5601600521443566) (0.9999999999999999,2.6046316661150737) };
  \addplot+[ours_msac, mark=none, line width=0.4pt, dashed, forget plot] coordinates {(-1.0,0.5913201531976) (1.0,0.5913201531976)};
  \addplot+[ours_msac, mark=*, only marks, mark size=\PanelArgminSize, forget plot] coordinates {(-0.07537688442211066,0.8931430016693322)};
  \addplot[name path=U_gau_grp, draw=none, mark=none, forget plot] coordinates {(-0.9999999999999999,1.420863656338869) (-0.9899497487437182,1.4081055808873697) (-0.9798994974874373,1.3985286629834075) (-0.9698492462311556,1.3904547092124269) (-0.959798994974874,1.3808861618750794) (-0.9497487437185931,1.370553913915141) (-0.9396984924623114,1.3569273538680042) (-0.9296482412060298,1.3475188949891876) (-0.9195979899497488,1.338123791389929) (-0.9095477386934673,1.3305740795189205) (-0.8994974874371856,1.3230703414678713) (-0.8894472361809047,1.3125766100523637) (-0.879396984924623,1.3024855146081893) (-0.8693467336683414,1.2920877886522846) (-0.8592964824120605,1.2782134871290163) (-0.8492462311557789,1.2658902496922622) (-0.8391959798994972,1.2572855153108804) (-0.8291457286432162,1.2492591330742726) (-0.8190954773869346,1.242423769788627) (-0.8090452261306529,1.2343107448066806) (-0.798994974874372,1.2282030795504773) (-0.7889447236180904,1.2188530915498517) (-0.7788944723618088,1.2134732430461928) (-0.7688442211055279,1.2080298328330736) (-0.7587939698492462,1.2028923381631302) (-0.7487437185929646,1.1926753984615064) (-0.7386934673366836,1.1846440774943057) (-0.7286432160804021,1.174987816389305) (-0.7185929648241204,1.1658765008928096) (-0.7085427135678395,1.1580869271524368) (-0.6984924623115578,1.1520037507526906) (-0.6884422110552761,1.1453178333695204) (-0.6783919597989952,1.140088885773585) (-0.6683417085427136,1.1327977447829658) (-0.658291457286432,1.1270309025333693) (-0.6482412060301507,1.1216794334110305) (-0.6381909547738694,1.1149388805576848) (-0.6281407035175877,1.107085462739125) (-0.6180904522613064,1.10014965069351) (-0.6080402010050252,1.0924728122759932) (-0.5979899497487435,1.0871573953184013) (-0.5879396984924622,1.0807437085709335) (-0.5778894472361809,1.0750247743474008) (-0.5678391959798992,1.0720709754577835) (-0.5577889447236181,1.0642571023301237) (-0.5477386934673368,1.0595193445854556) (-0.537688442211055,1.0527854438484006) (-0.5276381909547738,1.0495469953866199) (-0.5175879396984925,1.0439960138362612) (-0.5075376884422108,1.0389586904591015) (-0.4974874371859296,1.034146800078275) (-0.4874371859296484,1.0303620843543544) (-0.4773869346733666,1.0268133294746968) (-0.46733668341708545,1.021671991572669) (-0.4572864321608041,1.018507433742704) (-0.44723618090452244,1.0136128467985446) (-0.4371859296482412,1.008541923407279) (-0.42713567839195987,1.0028921734698617) (-0.41708542713567825,0.9976946487663747) (-0.4070351758793969,0.9934479736773685) (-0.3969849246231153,0.9905282148020528) (-0.386934673366834,0.9861556845907196) (-0.3768844221105528,0.9820299549084656) (-0.3668341708542711,0.9771526932062151) (-0.35678391959798983,0.9740510904365614) (-0.34673366834170855,0.9699325145392016) (-0.33668341708542693,0.9683261671255109) (-0.32663316582914564,0.9657982116389577) (-0.31658291457286436,0.9620532972651114) (-0.3065326633165827,0.958605592344254) (-0.29648241206030146,0.9563756310717547) (-0.2864321608040202,0.9532987207845943) (-0.27638190954773856,0.9493874579618932) (-0.2663316582914572,0.9468219531624721) (-0.25628140703517593,0.944745451082108) (-0.24623115577889426,0.9411299668973975) (-0.236180904522613,0.9384875780423135) (-0.22613065326633172,0.9352886903368905) (-0.2160804020100501,0.9323061251859452) (-0.20603015075376882,0.9291922665518227) (-0.19597989949748754,0.9275295389663617) (-0.1859296482412058,0.9255371227825512) (-0.1758793969849246,0.9243182613111217) (-0.16582914572864338,0.922076505089218) (-0.15577889447236168,0.9220133963724828) (-0.14572864321608034,0.9194564287897834) (-0.13567839195979914,0.918977911281889) (-0.12562814070351738,0.9168189502287875) (-0.11557788944723618,0.9144908093789799) (-0.10552763819095484,0.9125566620671914) (-0.09547738693467323,0.9113124504807364) (-0.08542713567839195,0.910243083819665) (-0.07537688442211071,0.9093166297301872) (-0.06532663316582901,0.9092214421603732) (-0.055276381909547756,0.9091921420707355) (-0.04522613065326643,0.9089414084873785) (-0.0351758793969848,0.9090908976370113) (-0.025125628140703526,0.9072985234158223) (-0.01507537688442226,0.906288333438176) (-0.0050251256281406455,0.9060729594270296) (0.00502512562814067,0.9055949839913817) (0.015075376884422353,0.9053748064890239) (0.02512562814070361,0.9050559424044059) (0.035175879396984876,0.9043588549096648) (0.04522613065326653,0.9050474947718716) (0.05527638190954779,0.9054762388911595) (0.06532663316582907,0.9057241055866263) (0.07537688442211074,0.9074523769638138) (0.08542713567839198,0.9077499142140667) (0.09547738693467332,0.9081758349455036) (0.10552763819095497,0.9090089170182052) (0.11557788944723618,0.9105127127866112) (0.12562814070351755,0.9114747478363628) (0.13567839195979917,0.9132186727320541) (0.14572864321608045,0.9139524644960759) (0.1557788944723617,0.9160064526343022) (0.16582914572864335,0.9190550557644869) (0.17587939698492464,0.9211350399518516) (0.1859296482412059,0.9234032466994407) (0.19597989949748756,0.9264786876495638) (0.20603015075376882,0.930574030052644) (0.21608040201005016,0.933586802299084) (0.2261306532663318,0.9375012482249004) (0.23618090452261306,0.9400948455604902) (0.24623115577889432,0.9428182865467476) (0.256281407035176,0.945083635766375) (0.2663316582914573,0.949265911626299) (0.27638190954773856,0.955365439888629) (0.28643216080402023,0.9618506131894502) (0.29648241206030146,0.966868482394758) (0.3065326633165827,0.9718356915808036) (0.31658291457286436,0.9788145310299698) (0.3266331658291457,0.9833071853463522) (0.336683417085427,0.9898103093438441) (0.34673366834170866,0.9974594932559777) (0.35678391959798994,1.0054705196537768) (0.36683417085427117,1.0121927250582787) (0.3768844221105528,1.0182423598625232) (0.38693467336683407,1.0274848567354247) (0.39698492462311535,1.0357554211480269) (0.40703517587939697,1.044129419909666) (0.4170854271356783,1.0524492938078756) (0.4271356783919599,1.0617706068211963) (0.4371859296482412,1.0734067181377096) (0.44723618090452255,1.0851611201212787) (0.45728643216080417,1.0984307948418672) (0.46733668341708545,1.1074191235838606) (0.47738693467336674,1.116586199353466) (0.4874371859296484,1.1261480749180839) (0.4974874371859296,1.1377130338213073) (0.5075376884422109,1.1491392049625415) (0.5175879396984926,1.1633030347742275) (0.5276381909547738,1.1780012953308645) (0.5376884422110552,1.1904112022408362) (0.5477386934673368,1.208499436439268) (0.5577889447236181,1.224319415023073) (0.5678391959798994,1.2393387568259364) (0.577889447236181,1.2570729354348784) (0.5879396984924623,1.2695220721731173) (0.5979899497487435,1.2856339580071823) (0.6080402010050252,1.3042969373300033) (0.6180904522613064,1.3223328292899788) (0.6281407035175878,1.344276249870674) (0.6381909547738694,1.3548355830115433) (0.6482412060301507,1.3679075106563452) (0.6582914572864319,1.3844237736754756) (0.6683417085427136,1.4043643207948024) (0.6783919597989949,1.4187168945664732) (0.6884422110552761,1.4331784109530765) (0.6984924623115578,1.4507451024016953) (0.708542713567839,1.4657095367106585) (0.7185929648241204,1.4811618346370297) (0.7286432160804021,1.49893596724715) (0.7386934673366833,1.5112737642544605) (0.7487437185929646,1.5255780779807033) (0.7587939698492463,1.5379405890441524) (0.7688442211055275,1.5504351986017426) (0.7788944723618089,1.5619262927766324) (0.7889447236180904,1.5714247843307756) (0.7989949748743718,1.5794384917242303) (0.8090452261306534,1.5912074777348824) (0.8190954773869347,1.600690945466247) (0.8291457286432159,1.6101243567014842) (0.8391959798994976,1.6179067965979081) (0.8492462311557789,1.627661384535894) (0.8592964824120601,1.6335476287594652) (0.8693467336683418,1.6390806693873181) (0.8793969849246231,1.6464866932800957) (0.8894472361809044,1.6522201015832194) (0.899497487437186,1.6572126923045218) (0.9095477386934673,1.6639811560112334) (0.9195979899497486,1.6689990454309074) (0.9296482412060302,1.6752880705120512) (0.9396984924623115,1.6783445193601714) (0.9497487437185927,1.6823430895965583) (0.9597989949748744,1.6865936169644056) (0.9698492462311558,1.6904615942958163) (0.979899497487437,1.6939398062894442) (0.9899497487437187,1.6976989602775316) (0.9999999999999999,1.7011357059570449) };
  \addplot[name path=L_gau_grp, draw=none, mark=none, forget plot] coordinates {(-0.9999999999999999,1.3721167899435363) (-0.9899497487437182,1.358689788071045) (-0.9798994974874373,1.3515585084501693) (-0.9698492462311556,1.3442766411300313) (-0.959798994974874,1.3359401399682715) (-0.9497487437185931,1.3249940346611537) (-0.9396984924623114,1.3114283407109293) (-0.9296482412060298,1.3029465967217426) (-0.9195979899497488,1.2941655029888777) (-0.9095477386934673,1.2858585224710681) (-0.8994974874371856,1.2776710629083619) (-0.8894472361809047,1.2673721256003398) (-0.879396984924623,1.2568373318966133) (-0.8693467336683414,1.2465249785034433) (-0.8592964824120605,1.232650425396311) (-0.8492462311557789,1.2239564902786757) (-0.8391959798994972,1.2176126474791886) (-0.8291457286432162,1.2108065631377138) (-0.8190954773869346,1.2032141202429798) (-0.8090452261306529,1.1932533213894247) (-0.798994974874372,1.1872207068464784) (-0.7889447236180904,1.1792671479316967) (-0.7788944723618088,1.17408700548538) (-0.7688442211055279,1.1671178379348055) (-0.7587939698492462,1.160193867673247) (-0.7487437185929646,1.1511194032431797) (-0.7386934673366836,1.1441842667475433) (-0.7286432160804021,1.1369911954801701) (-0.7185929648241204,1.1277908338660503) (-0.7085427135678395,1.1203602819781617) (-0.6984924623115578,1.1151075495928893) (-0.6884422110552761,1.109278378104623) (-0.6783919597989952,1.104551182423627) (-0.6683417085427136,1.096470022671909) (-0.658291457286432,1.0908331991143356) (-0.6482412060301507,1.0842589811254837) (-0.6381909547738694,1.0778554426531064) (-0.6281407035175877,1.0697282371652013) (-0.6180904522613064,1.063102007187705) (-0.6080402010050252,1.0564570533633095) (-0.5979899497487435,1.0527581249601867) (-0.5879396984924622,1.0468231436708408) (-0.5778894472361809,1.0409715513177944) (-0.5678391959798992,1.0380853163441) (-0.5577889447236181,1.0311624436363427) (-0.5477386934673368,1.0267696364631793) (-0.537688442211055,1.0210792604831476) (-0.5276381909547738,1.018150482735334) (-0.5175879396984925,1.0129441710458165) (-0.5075376884422108,1.0077610375935482) (-0.4974874371859296,1.0020338708371455) (-0.4874371859296484,0.9971547240796804) (-0.4773869346733666,0.993839800818635) (-0.46733668341708545,0.9889922848745025) (-0.4572864321608041,0.9862696775659519) (-0.44723618090452244,0.9811444857902007) (-0.4371859296482412,0.9762741726570832) (-0.42713567839195987,0.9712350236326391) (-0.41708542713567825,0.9660789813160937) (-0.4070351758793969,0.9624592478028065) (-0.3969849246231153,0.9600821817276403) (-0.386934673366834,0.9558700774415867) (-0.3768844221105528,0.9500821832693315) (-0.3668341708542711,0.9445731471146044) (-0.35678391959798983,0.9415706765929013) (-0.34673366834170855,0.9377613313879446) (-0.33668341708542693,0.9367085453780918) (-0.32663316582914564,0.9342301461718648) (-0.31658291457286436,0.9311156653914417) (-0.3065326633165827,0.927451556776139) (-0.29648241206030146,0.925297388736588) (-0.2864321608040202,0.9227666975146304) (-0.27638190954773856,0.9193090246360819) (-0.2663316582914572,0.9174547262506746) (-0.25628140703517593,0.9155751208987702) (-0.24623115577889426,0.9120884767369583) (-0.236180904522613,0.9095225176235878) (-0.22613065326633172,0.9066706021378176) (-0.2160804020100501,0.9037480689883619) (-0.20603015075376882,0.9007337567432018) (-0.19597989949748754,0.8992336834338143) (-0.1859296482412058,0.8969888095509638) (-0.1758793969849246,0.8956562342633673) (-0.16582914572864338,0.8934137854665262) (-0.15577889447236168,0.8936596321723533) (-0.14572864321608034,0.8918305455780596) (-0.13567839195979914,0.891107553686874) (-0.12562814070351738,0.8888082644729841) (-0.11557788944723618,0.8872211520709137) (-0.10552763819095484,0.8852938145944134) (-0.09547738693467323,0.8838403556203065) (-0.08542713567839195,0.8831847204112647) (-0.07537688442211071,0.882446880977125) (-0.06532663316582901,0.8819144286346708) (-0.055276381909547756,0.8817241683178522) (-0.04522613065326643,0.8814516590218374) (-0.0351758793969848,0.8817310212705223) (-0.025125628140703526,0.8796367973348981) (-0.01507537688442226,0.878654037012338) (-0.0050251256281406455,0.8785662198676447) (0.00502512562814067,0.8782732418946977) (0.015075376884422353,0.8782650498515441) (0.02512562814070361,0.8783858518700961) (0.035175879396984876,0.8784731607429959) (0.04522613065326653,0.8788872711643585) (0.05527638190954779,0.8794239439040143) (0.06532663316582907,0.8794211271301338) (0.07537688442211074,0.8804232506709863) (0.08542713567839198,0.8806156931380401) (0.09547738693467332,0.8816967356390262) (0.10552763819095497,0.88275688085965) (0.11557788944723618,0.8839828688898005) (0.12562814070351755,0.8853588080401621) (0.13567839195979917,0.8868496176290757) (0.14572864321608045,0.887572836990072) (0.1557788944723617,0.8896045557917253) (0.16582914572864335,0.8923190111303135) (0.17587939698492464,0.8939812700185821) (0.1859296482412059,0.8960694080514474) (0.19597989949748756,0.8982903605515408) (0.20603015075376882,0.9014212799828437) (0.21608040201005016,0.903982891492742) (0.2261306532663318,0.907587981620889) (0.23618090452261306,0.910017498208188) (0.24623115577889432,0.9130976803422395) (0.256281407035176,0.915407425585054) (0.2663316582914573,0.919355602034592) (0.27638190954773856,0.9241403510295468) (0.28643216080402023,0.9300616530654227) (0.29648241206030146,0.9344508595507128) (0.3065326633165827,0.9394058406062679) (0.31658291457286436,0.9474754536463272) (0.3266331658291457,0.9531045585306956) (0.336683417085427,0.9601995651265268) (0.34673366834170866,0.9665538983943855) (0.35678391959798994,0.9733838943799503) (0.36683417085427117,0.9796768729391712) (0.3768844221105528,0.9850190354448024) (0.38693467336683407,0.9920055506510667) (0.39698492462311535,1.0003286864185041) (0.40703517587939697,1.0088223277599475) (0.4170854271356783,1.0175708888711026) (0.4271356783919599,1.0274055281070018) (0.4371859296482412,1.0371977660711398) (0.44723618090452255,1.046503216379788) (0.45728643216080417,1.0594355261614707) (0.46733668341708545,1.0706555919129626) (0.47738693467336674,1.0811833709222731) (0.4874371859296484,1.0912130857830515) (0.4974874371859296,1.1030386016949913) (0.5075376884422109,1.1143859667108682) (0.5175879396984926,1.1259496721337712) (0.5276381909547738,1.138712022557468) (0.5376884422110552,1.1507811996570774) (0.5477386934673368,1.1659062872960102) (0.5577889447236181,1.183303154824953) (0.5678391959798994,1.1978690084778567) (0.577889447236181,1.214505467886704) (0.5879396984924623,1.2298917783208887) (0.5979899497487435,1.2446027599785814) (0.6080402010050252,1.2598989305269077) (0.6180904522613064,1.2751913827982713) (0.6281407035175878,1.294150157948988) (0.6381909547738694,1.3074789925348964) (0.6482412060301507,1.3239075863539596) (0.6582914572864319,1.3372822469622703) (0.6683417085427136,1.3560998430746234) (0.6783919597989949,1.374571081224127) (0.6884422110552761,1.3886449936684107) (0.6984924623115578,1.4040526552378358) (0.708542713567839,1.416520564773876) (0.7185929648241204,1.430786248158949) (0.7286432160804021,1.444776090834263) (0.7386934673366833,1.4566739127208514) (0.7487437185929646,1.4718594340366666) (0.7587939698492463,1.4831312359121376) (0.7688442211055275,1.4944996460122135) (0.7788944723618089,1.509414115680499) (0.7889447236180904,1.521232864520736) (0.7989949748743718,1.5290654112809716) (0.8090452261306534,1.5387085031021328) (0.8190954773869347,1.5467889138864022) (0.8291457286432159,1.5543585123536015) (0.8391959798994976,1.5608218896046346) (0.8492462311557789,1.570398556755419) (0.8592964824120601,1.5769386829315974) (0.8693467336683418,1.5838066145705711) (0.8793969849246231,1.5914903437813266) (0.8894472361809044,1.5973840229211405) (0.899497487437186,1.6024731902628775) (0.9095477386934673,1.6080694357226237) (0.9195979899497486,1.6129915329660522) (0.9296482412060302,1.6190069407802208) (0.9396984924623115,1.6222947845700446) (0.9497487437185927,1.6270250947306077) (0.9597989949748744,1.632037471931667) (0.9698492462311558,1.6366895182846937) (0.979899497487437,1.6400895960138115) (0.9899497487437187,1.6434982951843198) (0.9999999999999999,1.646847078256166) };
  \addplot[ours_gau, fill, fill opacity=0.18, draw=none, mark=none, forget plot] fill between [of=U_gau_grp and L_gau_grp];
  \addplot+[ours_gau, mark=none] coordinates {(-0.9999999999999999,1.3964902231412026) (-0.9899497487437182,1.3833976844792073) (-0.9798994974874373,1.3750435857167884) (-0.9698492462311556,1.367365675171229) (-0.959798994974874,1.3584131509216755) (-0.9497487437185931,1.3477739742881474) (-0.9396984924623114,1.3341778472894668) (-0.9296482412060298,1.3252327458554651) (-0.9195979899497488,1.3161446471894034) (-0.9095477386934673,1.3082163009949943) (-0.8994974874371856,1.3003707021881166) (-0.8894472361809047,1.2899743678263518) (-0.879396984924623,1.2796614232524013) (-0.8693467336683414,1.269306383577864) (-0.8592964824120605,1.2554319562626637) (-0.8492462311557789,1.244923369985469) (-0.8391959798994972,1.2374490813950345) (-0.8291457286432162,1.2300328481059932) (-0.8190954773869346,1.2228189450158033) (-0.8090452261306529,1.2137820330980527) (-0.798994974874372,1.2077118931984778) (-0.7889447236180904,1.1990601197407742) (-0.7788944723618088,1.1937801242657864) (-0.7688442211055279,1.1875738353839396) (-0.7587939698492462,1.1815431029181886) (-0.7487437185929646,1.171897400852343) (-0.7386934673366836,1.1644141721209245) (-0.7286432160804021,1.1559895059347376) (-0.7185929648241204,1.14683366737943) (-0.7085427135678395,1.1392236045652993) (-0.6984924623115578,1.13355565017279) (-0.6884422110552761,1.1272981057370717) (-0.6783919597989952,1.122320034098606) (-0.6683417085427136,1.1146338837274374) (-0.658291457286432,1.1089320508238525) (-0.6482412060301507,1.102969207268257) (-0.6381909547738694,1.0963971616053956) (-0.6281407035175877,1.088406849952163) (-0.6180904522613064,1.0816258289406075) (-0.6080402010050252,1.0744649328196514) (-0.5979899497487435,1.069957760139294) (-0.5879396984924622,1.0637834261208872) (-0.5778894472361809,1.0579981628325976) (-0.5678391959798992,1.0550781459009417) (-0.5577889447236181,1.0477097729832332) (-0.5477386934673368,1.0431444905243175) (-0.537688442211055,1.036932352165774) (-0.5276381909547738,1.033848739060977) (-0.5175879396984925,1.0284700924410388) (-0.5075376884422108,1.0233598640263248) (-0.4974874371859296,1.0180903354577102) (-0.4874371859296484,1.0137584042170174) (-0.4773869346733666,1.010326565146666) (-0.46733668341708545,1.0053321382235858) (-0.4572864321608041,1.002388555654328) (-0.44723618090452244,0.9973786662943727) (-0.4371859296482412,0.9924080480321811) (-0.42713567839195987,0.9870635985512505) (-0.41708542713567825,0.9818868150412342) (-0.4070351758793969,0.9779536107400875) (-0.3969849246231153,0.9753051982648465) (-0.386934673366834,0.9710128810161531) (-0.3768844221105528,0.9660560690888985) (-0.3668341708542711,0.9608629201604098) (-0.35678391959798983,0.9578108835147313) (-0.34673366834170855,0.9538469229635731) (-0.33668341708542693,0.9525173562518013) (-0.32663316582914564,0.9500141789054113) (-0.31658291457286436,0.9465844813282766) (-0.3065326633165827,0.9430285745601965) (-0.29648241206030146,0.9408365099041713) (-0.2864321608040202,0.9380327091496123) (-0.27638190954773856,0.9343482412989875) (-0.2663316582914572,0.9321383397065733) (-0.25628140703517593,0.9301602859904391) (-0.24623115577889426,0.9266092218171779) (-0.236180904522613,0.9240050478329507) (-0.22613065326633172,0.9209796462373541) (-0.2160804020100501,0.9180270970871536) (-0.20603015075376882,0.9149630116475123) (-0.19597989949748754,0.913381611200088) (-0.1859296482412058,0.9112629661667575) (-0.1758793969849246,0.9099872477872445) (-0.16582914572864338,0.9077451452778721) (-0.15577889447236168,0.9078365142724181) (-0.14572864321608034,0.9056434871839215) (-0.13567839195979914,0.9050427324843815) (-0.12562814070351738,0.9028136073508858) (-0.11557788944723618,0.9008559807249468) (-0.10552763819095484,0.8989252383308024) (-0.09547738693467323,0.8975764030505214) (-0.08542713567839195,0.8967139021154649) (-0.07537688442211071,0.8958817553536561) (-0.06532663316582901,0.895567935397522) (-0.055276381909547756,0.8954581551942938) (-0.04522613065326643,0.895196533754608) (-0.0351758793969848,0.8954109594537668) (-0.025125628140703526,0.8934676603753602) (-0.01507537688442226,0.892471185225257) (-0.0050251256281406455,0.8923195896473372) (0.00502512562814067,0.8919341129430397) (0.015075376884422353,0.891819928170284) (0.02512562814070361,0.891720897137251) (0.035175879396984876,0.8914160078263303) (0.04522613065326653,0.8919673829681151) (0.05527638190954779,0.8924500913975869) (0.06532663316582907,0.89257261635838) (0.07537688442211074,0.8939378138174) (0.08542713567839198,0.8941828036760534) (0.09547738693467332,0.8949362852922649) (0.10552763819095497,0.8958828989389276) (0.11557788944723618,0.8972477908382058) (0.12562814070351755,0.8984167779382625) (0.13567839195979917,0.9000341451805649) (0.14572864321608045,0.9007626507430739) (0.1557788944723617,0.9028055042130138) (0.16582914572864335,0.9056870334474002) (0.17587939698492464,0.9075581549852169) (0.1859296482412059,0.9097363273754441) (0.19597989949748756,0.9123845241005523) (0.20603015075376882,0.9159976550177439) (0.21608040201005016,0.918784846895913) (0.2261306532663318,0.9225446149228947) (0.23618090452261306,0.9250561718843391) (0.24623115577889432,0.9279579834444935) (0.256281407035176,0.9302455306757145) (0.2663316582914573,0.9343107568304455) (0.27638190954773856,0.9397528954590879) (0.28643216080402023,0.9459561331274364) (0.29648241206030146,0.9506596709727354) (0.3065326633165827,0.9556207660935357) (0.31658291457286436,0.9631449923381485) (0.3266331658291457,0.9682058719385239) (0.336683417085427,0.9750049372351854) (0.34673366834170866,0.9820066958251816) (0.35678391959798994,0.9894272070168636) (0.36683417085427117,0.995934798998725) (0.3768844221105528,1.0016306976536629) (0.38693467336683407,1.0097452036932457) (0.39698492462311535,1.0180420537832655) (0.40703517587939697,1.0264758738348068) (0.4170854271356783,1.035010091339489) (0.4271356783919599,1.044588067464099) (0.4371859296482412,1.0553022421044247) (0.44723618090452255,1.0658321682505334) (0.45728643216080417,1.078933160501669) (0.46733668341708545,1.0890373577484116) (0.47738693467336674,1.0988847851378696) (0.4874371859296484,1.1086805803505677) (0.4974874371859296,1.1203758177581493) (0.5075376884422109,1.1317625858367049) (0.5175879396984926,1.1446263534539993) (0.5276381909547738,1.1583566589441663) (0.5376884422110552,1.1705962009489568) (0.5477386934673368,1.1872028618676391) (0.5577889447236181,1.203811284924013) (0.5678391959798994,1.2186038826518966) (0.577889447236181,1.2357892016607912) (0.5879396984924623,1.249706925247003) (0.5979899497487435,1.2651183589928818) (0.6080402010050252,1.2820979339284555) (0.6180904522613064,1.298762106044125) (0.6281407035175878,1.319213203909831) (0.6381909547738694,1.3311572877732198) (0.6482412060301507,1.3459075485051524) (0.6582914572864319,1.360853010318873) (0.6683417085427136,1.380232081934713) (0.6783919597989949,1.3966439878953) (0.6884422110552761,1.4109117023107436) (0.6984924623115578,1.4273988788197656) (0.708542713567839,1.4411150507422672) (0.7185929648241204,1.4559740413979894) (0.7286432160804021,1.4718560290407066) (0.7386934673366833,1.483973838487656) (0.7487437185929646,1.498718756008685) (0.7587939698492463,1.510535912478145) (0.7688442211055275,1.522467422306978) (0.7788944723618089,1.5356702042285657) (0.7889447236180904,1.5463288244257558) (0.7989949748743718,1.554251951502601) (0.8090452261306534,1.5649579904185076) (0.8190954773869347,1.5737399296763246) (0.8291457286432159,1.5822414345275428) (0.8391959798994976,1.5893643431012714) (0.8492462311557789,1.5990299706456566) (0.8592964824120601,1.6052431558455313) (0.8693467336683418,1.6114436419789446) (0.8793969849246231,1.6189885185307111) (0.8894472361809044,1.62480206225218) (0.899497487437186,1.6298429412836997) (0.9095477386934673,1.6360252958669286) (0.9195979899497486,1.6409952891984798) (0.9296482412060302,1.647147505646136) (0.9396984924623115,1.650319651965108) (0.9497487437185927,1.654684092163583) (0.9597989949748744,1.6593155444480363) (0.9698492462311558,1.663575556290255) (0.979899497487437,1.6670147011516279) (0.9899497487437187,1.6705986277309257) (0.9999999999999999,1.6739913921066054) };
  \addplot+[ours_gau, mark=none, line width=0.4pt, dashed, forget plot] coordinates {(-1.0,0.6049128040341839) (1.0,0.6049128040341839)};
  \addplot+[ours_gau, mark=*, only marks, mark size=\PanelArgminSize, forget plot] coordinates {(0.035175879396984876,0.8914160078263303)};
  \addplot[name path=U_magsac_grp, draw=none, mark=none, forget plot] coordinates {(-1.0,1.2674475426092127) (-0.9899497487437188,1.2582205188474993) (-0.9798994974874371,1.2517712602501199) (-0.9698492462311558,1.246102136343721) (-0.9597989949748745,1.2385913495033716) (-0.9497487437185929,1.2331042299530128) (-0.9396984924623116,1.225680861366827) (-0.9296482412060304,1.2190037155766758) (-0.9195979899497487,1.2124557953373951) (-0.9095477386934675,1.2080982848865152) (-0.8994974874371862,1.2022977317793535) (-0.8894472361809045,1.1934371546773554) (-0.8793969849246231,1.185472556202828) (-0.8693467336683416,1.1775578398192055) (-0.8592964824120602,1.1699027027822877) (-0.849246231155779,1.1624595886058269) (-0.8391959798994973,1.1570485537331376) (-0.8291457286432161,1.1505844885538832) (-0.8190954773869348,1.1462133760438338) (-0.8090452261306531,1.1401019748951815) (-0.7989949748743719,1.135487434693508) (-0.7889447236180906,1.1319005712550683) (-0.7788944723618089,1.1259721337702182) (-0.7688442211055276,1.1210471269539632) (-0.7587939698492464,1.1148037562410638) (-0.7487437185929647,1.1067515690917327) (-0.7386934673366835,1.1004931337056736) (-0.7286432160804022,1.0936692232593712) (-0.7185929648241205,1.0880066527225514) (-0.7085427135678393,1.082268577130273) (-0.6984924623115579,1.0780680351520653) (-0.6884422110552763,1.0734369719050296) (-0.678391959798995,1.0698460073024938) (-0.6683417085427138,1.0658804440703251) (-0.6582914572864321,1.0608583057654588) (-0.6482412060301509,1.0556915849474737) (-0.6381909547738696,1.050802444819158) (-0.6281407035175879,1.0459007490833772) (-0.6180904522613067,1.0404196369884988) (-0.608040201005025,1.036345166808934) (-0.5979899497487438,1.0324246841330091) (-0.5879396984924625,1.0292036683701513) (-0.5778894472361807,1.0247987979678745) (-0.5678391959798995,1.0199935821520245) (-0.5577889447236182,1.0150974675836963) (-0.5477386934673365,1.0107566597563555) (-0.5376884422110553,1.0062950872577259) (-0.527638190954774,1.0023381577264057) (-0.5175879396984923,1.0001609968813108) (-0.507537688442211,0.9972541957704482) (-0.4974874371859298,0.9942343994854256) (-0.4874371859296481,0.9904151131494983) (-0.4773869346733669,0.9865546686051719) (-0.4673366834170855,0.9829293300762622) (-0.45728643216080395,0.9797573297622135) (-0.44723618090452266,0.9755264952904873) (-0.4371859296482414,0.9719051055440204) (-0.42713567839195976,0.9676772681936852) (-0.4170854271356784,0.9641112039592963) (-0.4070351758793972,0.9609820540181113) (-0.39698492462311546,0.95948267414718) (-0.38693467336683424,0.9579156028921185) (-0.37688442211055295,0.9556334328665863) (-0.3668341708542713,0.9527603286453665) (-0.35678391959799,0.9484916163660811) (-0.3467336683417084,0.947323827900883) (-0.33668341708542704,0.9448200696192477) (-0.3266331658291458,0.9417011298943339) (-0.31658291457286414,0.939949759496284) (-0.3065326633165829,0.9360439430879745) (-0.2964824120603017,0.9338440285082474) (-0.2864321608040199,0.9306688450103232) (-0.2763819095477387,0.927966895577563) (-0.2663316582914574,0.9272340599455312) (-0.2562814070351757,0.9241949396642856) (-0.2462311557788945,0.9230643972963657) (-0.23618090452261317,0.9214565293693474) (-0.2261306532663315,0.9202991186099247) (-0.21608040201005027,0.9186352076412788) (-0.206030150753769,0.9164241730300393) (-0.19597989949748731,0.9150561270427556) (-0.18592964824120603,0.9137414734386989) (-0.17587939698492477,0.9123423446381611) (-0.1658291457286431,0.9110212453210778) (-0.15577889447236187,0.9105979269127966) (-0.14572864321608053,0.9094339765631013) (-0.13567839195979886,0.9080262311079891) (-0.12562814070351766,0.9080799993223456) (-0.115577889447236,0.9065891710336405) (-0.10552763819095472,0.905731881012296) (-0.0954773869346734,0.9063887450512477) (-0.08542713567839177,0.9071302239414748) (-0.07537688442211048,0.9059091552805207) (-0.06532663316582918,0.9057682022221986) (-0.05527638190954754,0.9046060612676919) (-0.045226130653266215,0.9040859696718303) (-0.03517587939698501,0.902073894634677) (-0.02512562814070327,0.9019427426298876) (-0.015075376884422061,0.9032389878281295) (-0.005025125628140791,0.9036324828877471) (0.005025125628140861,0.9039079154111624) (0.015075376884422167,0.904562542148228) (0.025125628140703425,0.9046352639728233) (0.035175879396985056,0.9059327653902112) (0.045226130653266354,0.9067913101462985) (0.05527638190954762,0.9071034449739506) (0.06532663316582932,0.9077272502904891) (0.0753768844221105,0.9081012911566609) (0.08542713567839183,0.9094841349073143) (0.09547738693467347,0.9103532038999124) (0.10552763819095481,0.9117605594813128) (0.11557788944723604,0.912702265226347) (0.1256281407035177,0.9142847107147146) (0.13567839195979894,0.9168866436733325) (0.14572864321608067,0.920153702367462) (0.15577889447236196,0.921725502247203) (0.16582914572864316,0.9246657289329656) (0.17587939698492483,0.9279568630631194) (0.18592964824120614,0.9298922598214246) (0.19597989949748737,0.9325493507732502) (0.20603015075376907,0.9356353402750066) (0.21608040201005033,0.9377990908349856) (0.22613065326633164,0.9400341037440341) (0.23618090452261328,0.9422632110397041) (0.24623115577889454,0.945730407436147) (0.25628140703517577,0.9502260442641942) (0.2663316582914575,0.95669050335572) (0.2763819095477388,0.9604153880168763) (0.28643216080401995,0.9662984739797785) (0.2964824120603017,0.9713107542440383) (0.30653266331658297,0.9759491203753178) (0.31658291457286425,0.9808113943824026) (0.326633165829146,0.9852156860747491) (0.33668341708542715,0.9900986400820367) (0.34673366834170866,0.9970163301293413) (0.3567839195979901,1.004034540762216) (0.36683417085427134,1.010856535665086) (0.3768844221105529,1.0149778361633928) (0.3869346733668343,1.0234158019683026) (0.3969849246231155,1.031652839117006) (0.407035175879397,1.040265444191095) (0.41708542713567853,1.0463470008341837) (0.4271356783919598,1.0556524077856004) (0.4371859296482412,1.0630695092888347) (0.44723618090452266,1.0732225908452235) (0.45728643216080417,1.0842644067324299) (0.46733668341708545,1.0946327659284425) (0.47738693467336696,1.1038802807530044) (0.4874371859296484,1.1120349869573072) (0.49748743718592964,1.1218167643588262) (0.5075376884422111,1.1325613313526541) (0.5175879396984926,1.141089954843092) (0.5276381909547738,1.15289890143942) (0.5376884422110554,1.1664070939156879) (0.5477386934673368,1.1784249696881912) (0.5577889447236181,1.1907972458949985) (0.5678391959798995,1.2059538339308602) (0.5778894472361811,1.2207783303469042) (0.5879396984924625,1.2365499733939156) (0.5979899497487438,1.2522959800407458) (0.6080402010050252,1.266893822955316) (0.6180904522613068,1.280195761809094) (0.628140703517588,1.2975740317901563) (0.6381909547738694,1.3195132252331827) (0.6482412060301509,1.33696302863826) (0.6582914572864321,1.3523674176955003) (0.6683417085427136,1.3659302761183396) (0.678391959798995,1.3851257065200915) (0.6884422110552766,1.404694553071569) (0.6984924623115579,1.4208341955391781) (0.7085427135678393,1.4415177609097287) (0.7185929648241208,1.4646092789989793) (0.7286432160804022,1.4915732460833349) (0.7386934673366835,1.5161054435781716) (0.748743718592965,1.5357611862922889) (0.7587939698492463,1.5580722775363032) (0.7688442211055277,1.5774597767374636) (0.7788944723618092,1.6010237923315764) (0.7889447236180904,1.626686772449587) (0.7989949748743719,1.6531929005996582) (0.8090452261306534,1.6793187074226785) (0.8190954773869349,1.7054090520999525) (0.8291457286432162,1.7350934008584133) (0.8391959798994977,1.7644819789433823) (0.8492462311557791,1.7950159734880198) (0.8592964824120604,1.8255900846583066) (0.8693467336683418,1.857523128999578) (0.8793969849246233,1.8849464785622945) (0.8894472361809046,1.916273631203379) (0.899497487437186,1.9464940785000744) (0.9095477386934675,1.9800758402164091) (0.9195979899497488,2.016073137725173) (0.9296482412060302,2.052454475271634) (0.9396984924623117,2.0887200982785488) (0.9497487437185932,2.1242171045702776) (0.9597989949748745,2.161744812310917) (0.9698492462311559,2.207793048010002) (0.9798994974874374,2.2492042146544438) (0.9899497487437187,2.292105783971373) (1.0000000000000002,2.336094395299332) };
  \addplot[name path=L_magsac_grp, draw=none, mark=none, forget plot] coordinates {(-1.0,1.2256768855394322) (-0.9899497487437188,1.2175054674462724) (-0.9798994974874371,1.212160773884746) (-0.9698492462311558,1.205661842085745) (-0.9597989949748745,1.1978412396637728) (-0.9497487437185929,1.19313591296807) (-0.9396984924623116,1.1866356204740802) (-0.9296482412060304,1.1788962721784675) (-0.9195979899497487,1.1708034667964584) (-0.9095477386934675,1.164401810895748) (-0.8994974874371862,1.1579001043783277) (-0.8894472361809045,1.1494903134569732) (-0.8793969849246231,1.1440238091253132) (-0.8693467336683416,1.1385787352476915) (-0.8592964824120602,1.1324212105985543) (-0.849246231155779,1.1255300197761717) (-0.8391959798994973,1.1197585415737534) (-0.8291457286432161,1.1140892681012) (-0.8190954773869348,1.110170415412262) (-0.8090452261306531,1.1050482927601752) (-0.7989949748743719,1.1002118447520493) (-0.7889447236180906,1.0963402776521864) (-0.7788944723618089,1.0895050883381165) (-0.7688442211055276,1.0840006068021193) (-0.7587939698492464,1.0773553914772762) (-0.7487437185929647,1.0702477778198316) (-0.7386934673366835,1.064033554529477) (-0.7286432160804022,1.059108409592266) (-0.7185929648241205,1.054976446279252) (-0.7085427135678393,1.050737380270876) (-0.6984924623115579,1.0471246949828008) (-0.6884422110552763,1.0422351133002292) (-0.678391959798995,1.0381613472793938) (-0.6683417085427138,1.0344486947213263) (-0.6582914572864321,1.0298775130564755) (-0.6482412060301509,1.024831957572957) (-0.6381909547738696,1.0204656951323834) (-0.6281407035175879,1.0167092112706508) (-0.6180904522613067,1.0109066915656681) (-0.608040201005025,1.0060260858735015) (-0.5979899497487438,1.0019368104049597) (-0.5879396984924625,0.9984275412256475) (-0.5778894472361807,0.9944111935307945) (-0.5678391959798995,0.9898701593593098) (-0.5577889447236182,0.9853746232930013) (-0.5477386934673365,0.9814004904909673) (-0.5376884422110553,0.9760906443810478) (-0.527638190954774,0.9719820295178757) (-0.5175879396984923,0.9695696505812165) (-0.507537688442211,0.9671304166450256) (-0.4974874371859298,0.9643122919075328) (-0.4874371859296481,0.9606081678480245) (-0.4773869346733669,0.9559248679851416) (-0.4673366834170855,0.9514294996915217) (-0.45728643216080395,0.9479881385580775) (-0.44723618090452266,0.9440039640574373) (-0.4371859296482414,0.9408354486769084) (-0.42713567839195976,0.9365532276774635) (-0.4170854271356784,0.9328816457642319) (-0.4070351758793972,0.9296684150820229) (-0.39698492462311546,0.9283718312996581) (-0.38693467336683424,0.9264710301197563) (-0.37688442211055295,0.9245783078680538) (-0.3668341708542713,0.9219874792072025) (-0.35678391959799,0.9183752162030298) (-0.3467336683417084,0.9174522813828131) (-0.33668341708542704,0.915327429135621) (-0.3266331658291458,0.9125728414046044) (-0.31658291457286414,0.9108225554068594) (-0.3065326633165829,0.9077205176255064) (-0.2964824120603017,0.9054989064136065) (-0.2864321608040199,0.9030349232671425) (-0.2763819095477387,0.8998455667780328) (-0.2663316582914574,0.8992047518349336) (-0.2562814070351757,0.8961354357666219) (-0.2462311557788945,0.8952556851320381) (-0.23618090452261317,0.8941993434367107) (-0.2261306532663315,0.8933832941794198) (-0.21608040201005027,0.8914259847177203) (-0.206030150753769,0.889309748077626) (-0.19597989949748731,0.887845340273257) (-0.18592964824120603,0.8864257709915803) (-0.17587939698492477,0.8852704614916029) (-0.1658291457286431,0.8835407674502391) (-0.15577889447236187,0.8836733222789612) (-0.14572864321608053,0.8820670349633847) (-0.13567839195979886,0.8806827210317051) (-0.12562814070351766,0.880939152321262) (-0.115577889447236,0.8792323877426333) (-0.10552763819095472,0.8783488318478511) (-0.0954773869346734,0.8789674259939132) (-0.08542713567839177,0.8798109909748438) (-0.07537688442211048,0.8786000586940059) (-0.06532663316582918,0.8787231563808912) (-0.05527638190954754,0.8773081079459838) (-0.045226130653266215,0.8766818029898648) (-0.03517587939698501,0.8752372162137104) (-0.02512562814070327,0.8752651811152607) (-0.015075376884422061,0.8771218724978884) (-0.005025125628140791,0.8771866419897767) (0.005025125628140861,0.8776608358289508) (0.015075376884422167,0.8784864120801068) (0.025125628140703425,0.8787421316730158) (0.035175879396985056,0.8795831574138089) (0.045226130653266354,0.8801280828283364) (0.05527638190954762,0.8808989423454736) (0.06532663316582932,0.8817268455326567) (0.0753768844221105,0.8820204619069058) (0.08542713567839183,0.8834526534351824) (0.09547738693467347,0.8844043783458856) (0.10552763819095481,0.886390411368792) (0.11557788944723604,0.8870851771193864) (0.1256281407035177,0.8883950547699669) (0.13567839195979894,0.8912536313925863) (0.14572864321608067,0.894339691172515) (0.15577889447236196,0.895347644033765) (0.16582914572864316,0.8980775928520057) (0.17587939698492483,0.900894785304508) (0.18592964824120614,0.9026563174229615) (0.19597989949748737,0.9050159454737626) (0.20603015075376907,0.9075664729580348) (0.21608040201005033,0.9093472556649052) (0.22613065326633164,0.9112007471372474) (0.23618090452261328,0.9130752224249578) (0.24623115577889454,0.9162198637686614) (0.25628140703517577,0.9198396476159388) (0.2663316582914575,0.9250821824861) (0.2763819095477388,0.9283030979423267) (0.28643216080401995,0.9329921692766506) (0.2964824120603017,0.9386535520495284) (0.30653266331658297,0.9444496272787862) (0.31658291457286425,0.9507911626453782) (0.326633165829146,0.9562092260785774) (0.33668341708542715,0.9615077977206711) (0.34673366834170866,0.9673585638198914) (0.3567839195979901,0.972934884630271) (0.36683417085427134,0.979492926971439) (0.3768844221105529,0.9838030708748926) (0.3869346733668343,0.990387327953153) (0.3969849246231155,0.9974608355335653) (0.407035175879397,1.0054735629286184) (0.41708542713567853,1.0115015862025305) (0.4271356783919598,1.021751011647494) (0.4371859296482412,1.029591319317777) (0.44723618090452266,1.0378241664562735) (0.45728643216080417,1.0469557535705516) (0.46733668341708545,1.0566389665338163) (0.47738693467336696,1.068118023981474) (0.4874371859296484,1.0781816464387786) (0.49748743718592964,1.0876373356277822) (0.5075376884422111,1.097909932632962) (0.5175879396984926,1.1077445301926498) (0.5276381909547738,1.1186775596885243) (0.5376884422110554,1.129597457384185) (0.5477386934673368,1.1401554419206343) (0.5577889447236181,1.1518892769977924) (0.5678391959798995,1.165634790535728) (0.5778894472361811,1.1822893753159183) (0.5879396984924625,1.196601174741281) (0.5979899497487438,1.2101417436026727) (0.6080402010050252,1.227638541498285) (0.6180904522613068,1.2409501316412443) (0.628140703517588,1.2558509661354273) (0.6381909547738694,1.2747227990504568) (0.6482412060301509,1.29058332620223) (0.6582914572864321,1.3069389538740992) (0.6683417085427136,1.3240364468002783) (0.678391959798995,1.3399152079681846) (0.6884422110552766,1.3592766514927033) (0.6984924623115579,1.3787052034269622) (0.7085427135678393,1.397182398492025) (0.7185929648241208,1.4162643827455135) (0.7286432160804022,1.4395111130712839) (0.7386934673366835,1.465494937821116) (0.748743718592965,1.4836456475262765) (0.7587939698492463,1.5058912123527064) (0.7688442211055277,1.527612408771695) (0.7788944723618092,1.54921671352992) (0.7889447236180904,1.5730518309684285) (0.7989949748743719,1.5995031763086223) (0.8090452261306534,1.624023284737218) (0.8190954773869349,1.648642889942103) (0.8291457286432162,1.6746970909118974) (0.8391959798994977,1.7036757799886137) (0.8492462311557791,1.7373276084404616) (0.8592964824120604,1.765441732448306) (0.8693467336683418,1.7945151723785375) (0.8793969849246233,1.821758229379212) (0.8894472361809046,1.8467744753110462) (0.899497487437186,1.879051328971451) (0.9095477386934675,1.9164286926857104) (0.9195979899497488,1.9452240329936676) (0.9296482412060302,1.9818197027842352) (0.9396984924623117,2.020069780855603) (0.9497487437185932,2.054322003665807) (0.9597989949748745,2.0909206945715413) (0.9698492462311559,2.126253977174189) (0.9798994974874374,2.1683105337854784) (0.9899497487437187,2.207713884897435) (1.0000000000000002,2.245398985253429) };
  \addplot[ours_magsac, fill, fill opacity=0.18, draw=none, mark=none, forget plot] fill between [of=U_magsac_grp and L_magsac_grp];
  \addplot+[ours_magsac, mark=none] coordinates {(-1.0,1.2465622140743224) (-0.9899497487437188,1.2378629931468859) (-0.9798994974874371,1.231966017067433) (-0.9698492462311558,1.225881989214733) (-0.9597989949748745,1.2182162945835722) (-0.9497487437185929,1.2131200714605415) (-0.9396984924623116,1.2061582409204537) (-0.9296482412060304,1.1989499938775716) (-0.9195979899497487,1.1916296310669268) (-0.9095477386934675,1.1862500478911315) (-0.8994974874371862,1.1800989180788406) (-0.8894472361809045,1.1714637340671643) (-0.8793969849246231,1.1647481826640707) (-0.8693467336683416,1.1580682875334485) (-0.8592964824120602,1.151161956690421) (-0.849246231155779,1.1439948041909993) (-0.8391959798994973,1.1384035476534455) (-0.8291457286432161,1.1323368783275416) (-0.8190954773869348,1.128191895728048) (-0.8090452261306531,1.1225751338276784) (-0.7989949748743719,1.1178496397227786) (-0.7889447236180906,1.1141204244536274) (-0.7788944723618089,1.1077386110541674) (-0.7688442211055276,1.1025238668780413) (-0.7587939698492464,1.09607957385917) (-0.7487437185929647,1.0884996734557821) (-0.7386934673366835,1.0822633441175753) (-0.7286432160804022,1.0763888164258186) (-0.7185929648241205,1.0714915495009016) (-0.7085427135678393,1.0665029787005744) (-0.6984924623115579,1.062596365067433) (-0.6884422110552763,1.0578360426026294) (-0.678391959798995,1.0540036772909438) (-0.6683417085427138,1.0501645693958257) (-0.6582914572864321,1.0453679094109671) (-0.6482412060301509,1.0402617712602154) (-0.6381909547738696,1.0356340699757707) (-0.6281407035175879,1.031304980177014) (-0.6180904522613067,1.0256631642770835) (-0.608040201005025,1.0211856263412178) (-0.5979899497487438,1.0171807472689844) (-0.5879396984924625,1.0138156047978995) (-0.5778894472361807,1.0096049957493345) (-0.5678391959798995,1.0049318707556671) (-0.5577889447236182,1.0002360454383488) (-0.5477386934673365,0.9960785751236614) (-0.5376884422110553,0.9911928658193868) (-0.527638190954774,0.9871600936221406) (-0.5175879396984923,0.9848653237312637) (-0.507537688442211,0.9821923062077369) (-0.4974874371859298,0.9792733456964792) (-0.4874371859296481,0.9755116404987614) (-0.4773869346733669,0.9712397682951568) (-0.4673366834170855,0.9671794148838919) (-0.45728643216080395,0.9638727341601455) (-0.44723618090452266,0.9597652296739623) (-0.4371859296482414,0.9563702771104644) (-0.42713567839195976,0.9521152479355743) (-0.4170854271356784,0.9484964248617641) (-0.4070351758793972,0.9453252345500671) (-0.39698492462311546,0.943927252723419) (-0.38693467336683424,0.9421933165059374) (-0.37688442211055295,0.9401058703673201) (-0.3668341708542713,0.9373739039262845) (-0.35678391959799,0.9334334162845555) (-0.3467336683417084,0.932388054641848) (-0.33668341708542704,0.9300737493774344) (-0.3266331658291458,0.9271369856494691) (-0.31658291457286414,0.9253861574515717) (-0.3065326633165829,0.9218822303567404) (-0.2964824120603017,0.919671467460927) (-0.2864321608040199,0.9168518841387329) (-0.2763819095477387,0.9139062311777979) (-0.2663316582914574,0.9132194058902324) (-0.2562814070351757,0.9101651877154537) (-0.2462311557788945,0.9091600412142019) (-0.23618090452261317,0.9078279364030291) (-0.2261306532663315,0.9068412063946723) (-0.21608040201005027,0.9050305961794995) (-0.206030150753769,0.9028669605538326) (-0.19597989949748731,0.9014507336580063) (-0.18592964824120603,0.9000836222151396) (-0.17587939698492477,0.898806403064882) (-0.1658291457286431,0.8972810063856584) (-0.15577889447236187,0.8971356245958789) (-0.14572864321608053,0.895750505763243) (-0.13567839195979886,0.8943544760698471) (-0.12562814070351766,0.8945095758218038) (-0.115577889447236,0.8929107793881369) (-0.10552763819095472,0.8920403564300735) (-0.0954773869346734,0.8926780855225804) (-0.08542713567839177,0.8934706074581593) (-0.07537688442211048,0.8922546069872633) (-0.06532663316582918,0.8922456793015449) (-0.05527638190954754,0.8909570846068379) (-0.045226130653266215,0.8903838863308475) (-0.03517587939698501,0.8886555554241937) (-0.02512562814070327,0.8886039618725742) (-0.015075376884422061,0.890180430163009) (-0.005025125628140791,0.8904095624387619) (0.005025125628140861,0.8907843756200566) (0.015075376884422167,0.8915244771141674) (0.025125628140703425,0.8916886978229196) (0.035175879396985056,0.89275796140201) (0.045226130653266354,0.8934596964873175) (0.05527638190954762,0.8940011936597121) (0.06532663316582932,0.8947270479115729) (0.0753768844221105,0.8950608765317833) (0.08542713567839183,0.8964683941712484) (0.09547738693467347,0.897378791122899) (0.10552763819095481,0.8990754854250524) (0.11557788944723604,0.8998937211728667) (0.1256281407035177,0.9013398827423408) (0.13567839195979894,0.9040701375329594) (0.14572864321608067,0.9072466967699885) (0.15577889447236196,0.908536573140484) (0.16582914572864316,0.9113716608924857) (0.17587939698492483,0.9144258241838137) (0.18592964824120614,0.9162742886221931) (0.19597989949748737,0.9187826481235064) (0.20603015075376907,0.9216009066165207) (0.21608040201005033,0.9235731732499454) (0.22613065326633164,0.9256174254406407) (0.23618090452261328,0.9276692167323309) (0.24623115577889454,0.9309751356024042) (0.25628140703517577,0.9350328459400665) (0.2663316582914575,0.94088634292091) (0.2763819095477388,0.9443592429796015) (0.28643216080401995,0.9496453216282146) (0.2964824120603017,0.9549821531467834) (0.30653266331658297,0.960199373827052) (0.31658291457286425,0.9658012785138904) (0.326633165829146,0.9707124560766632) (0.33668341708542715,0.9758032189013539) (0.34673366834170866,0.9821874469746164) (0.3567839195979901,0.9884847126962435) (0.36683417085427134,0.9951747313182625) (0.3768844221105529,0.9993904535191427) (0.3869346733668343,1.0069015649607278) (0.3969849246231155,1.0145568373252856) (0.407035175879397,1.0228695035598567) (0.41708542713567853,1.0289242935183571) (0.4271356783919598,1.0387017097165472) (0.4371859296482412,1.0463304143033059) (0.44723618090452266,1.0555233786507485) (0.45728643216080417,1.0656100801514907) (0.46733668341708545,1.0756358662311294) (0.47738693467336696,1.0859991523672392) (0.4874371859296484,1.0951083166980429) (0.49748743718592964,1.1047270499933042) (0.5075376884422111,1.115235631992808) (0.5175879396984926,1.124417242517871) (0.5276381909547738,1.1357882305639722) (0.5376884422110554,1.1480022756499364) (0.5477386934673368,1.1592902058044128) (0.5577889447236181,1.1713432614463954) (0.5678391959798995,1.1857943122332941) (0.5778894472361811,1.2015338528314112) (0.5879396984924625,1.2165755740675983) (0.5979899497487438,1.2312188618217093) (0.6080402010050252,1.2472661822268005) (0.6180904522613068,1.2605729467251692) (0.628140703517588,1.2767124989627918) (0.6381909547738694,1.2971180121418198) (0.6482412060301509,1.313773177420245) (0.6582914572864321,1.3296531857847997) (0.6683417085427136,1.344983361459309) (0.678391959798995,1.362520457244138) (0.6884422110552766,1.3819856022821362) (0.6984924623115579,1.3997696994830702) (0.7085427135678393,1.419350079700877) (0.7185929648241208,1.4404368308722464) (0.7286432160804022,1.4655421795773094) (0.7386934673366835,1.4908001906996438) (0.748743718592965,1.5097034169092827) (0.7587939698492463,1.5319817449445048) (0.7688442211055277,1.5525360927545793) (0.7788944723618092,1.5751202529307482) (0.7889447236180904,1.5998693017090078) (0.7989949748743719,1.6263480384541402) (0.8090452261306534,1.6516709960799483) (0.8190954773869349,1.6770259710210278) (0.8291457286432162,1.7048952458851554) (0.8391959798994977,1.734078879465998) (0.8492462311557791,1.7661717909642407) (0.8592964824120604,1.7955159085533063) (0.8693467336683418,1.8260191506890577) (0.8793969849246233,1.8533523539707533) (0.8894472361809046,1.8815240532572126) (0.899497487437186,1.9127727037357627) (0.9095477386934675,1.9482522664510598) (0.9195979899497488,1.9806485853594202) (0.9296482412060302,2.0171370890279348) (0.9396984924623117,2.054394939567076) (0.9497487437185932,2.0892695541180424) (0.9597989949748745,2.126332753441229) (0.9698492462311559,2.1670235125920954) (0.9798994974874374,2.208757374219961) (0.9899497487437187,2.249909834434404) (1.0000000000000002,2.2907466902763804) };
  \addplot+[ours_magsac, mark=none, line width=0.4pt, dashed, forget plot] coordinates {(-1.0,0.6106970255362396) (1.0,0.6106970255362396)};
  \addplot+[ours_magsac, mark=*, only marks, mark size=\PanelArgminSize, forget plot] coordinates {(-0.02512562814070327,0.8886039618725742)};
  \addplot[name path=U_jeffreys_grp, draw=none, mark=none, forget plot] coordinates {(-1.0,0.9173069836951669) (-0.9899497487437185,0.9167844912824815) (-0.9798994974874372,0.9170907958172967) (-0.9698492462311558,0.9173492745795743) (-0.9597989949748743,0.9175842266327544) (-0.9497487437185931,0.9175253938012318) (-0.9396984924623115,0.9173493561220654) (-0.9296482412060301,0.9168711811234399) (-0.9195979899497487,0.9171905155955551) (-0.9095477386934674,0.9167419911072363) (-0.8994974874371859,0.9157560737696947) (-0.8894472361809044,0.9156482500207511) (-0.8793969849246231,0.9147858696964367) (-0.8693467336683417,0.914568859928518) (-0.8592964824120602,0.9144675649199742) (-0.8492462311557789,0.9141087124083966) (-0.8391959798994976,0.9140338035281039) (-0.829145728643216,0.9147777116678704) (-0.8190954773869348,0.915244892694302) (-0.8090452261306533,0.9141140243663842) (-0.7989949748743718,0.9143862601289487) (-0.7889447236180904,0.91409721010713) (-0.7788944723618091,0.9141121118255148) (-0.7688442211055276,0.913457149437133) (-0.7587939698492462,0.9134901222526562) (-0.7487437185929648,0.9130407212093874) (-0.7386934673366834,0.9130570392782407) (-0.7286432160804021,0.9123559757095495) (-0.7185929648241205,0.9119963974113989) (-0.7085427135678393,0.9118404180145585) (-0.6984924623115578,0.9116793946005275) (-0.6884422110552764,0.9114768479727893) (-0.678391959798995,0.911717030066099) (-0.6683417085427136,0.9117127367407694) (-0.6582914572864321,0.911669766750643) (-0.6482412060301508,0.9116654453546919) (-0.6381909547738693,0.9121480823611665) (-0.628140703517588,0.9116218414583729) (-0.6180904522613065,0.9119868025822467) (-0.6080402010050251,0.9116728621951878) (-0.5979899497487438,0.9116017691647383) (-0.5879396984924624,0.9113891476780234) (-0.5778894472361809,0.9110182358468976) (-0.5678391959798995,0.9110590003549179) (-0.5577889447236181,0.911151760566226) (-0.5477386934673367,0.9109990296418409) (-0.5376884422110553,0.9106707829732426) (-0.5276381909547738,0.9102792703667437) (-0.5175879396984925,0.9103027655733799) (-0.507537688442211,0.910146643134284) (-0.49748743718592964,0.9092971768600238) (-0.48743718592964824,0.9093066700558016) (-0.47738693467336674,0.9091586156912779) (-0.46733668341708545,0.9092608694168939) (-0.45728643216080406,0.9093040363726468) (-0.4472361809045226,0.9094714159657273) (-0.4371859296482412,0.9091537204719651) (-0.4271356783919598,0.9088445167569422) (-0.4170854271356783,0.9091200178012566) (-0.40703517587939697,0.9088698511471069) (-0.3969849246231155,0.9086101308507467) (-0.38693467336683424,0.9087400924876051) (-0.37688442211055273,0.9078293067149567) (-0.36683417085427134,0.9080742020554633) (-0.35678391959798994,0.907958902338831) (-0.34673366834170855,0.9082485957884693) (-0.3366834170854271,0.9081446565462528) (-0.32663316582914564,0.9074579303269645) (-0.31658291457286436,0.9074052062613835) (-0.3065326633165829,0.907167408892206) (-0.29648241206030146,0.9076261360536746) (-0.2864321608040201,0.9076558558924306) (-0.2763819095477387,0.9070793653383413) (-0.26633165829145733,0.9064633973654477) (-0.25628140703517593,0.906420209580371) (-0.24623115577889443,0.9059423883240464) (-0.236180904522613,0.906127187473856) (-0.22613065326633172,0.9058330526357192) (-0.21608040201005027,0.9060546572081126) (-0.20603015075376882,0.9060652518796192) (-0.19597989949748745,0.9057547817621913) (-0.18592964824120603,0.9058207026016059) (-0.1758793969849246,0.9059940726798787) (-0.16582914572864324,0.9060869406350356) (-0.15577889447236182,0.9058314167306238) (-0.14572864321608042,0.9057599879789217) (-0.135678391959799,0.9062601956978074) (-0.12562814070351758,0.906626122911158) (-0.11557788944723618,0.9064504479561553) (-0.10552763819095479,0.9059447600130752) (-0.09547738693467334,0.905929377773117) (-0.08542713567839195,0.9054030323499654) (-0.07537688442211053,0.9050111562168579) (-0.06532663316582918,0.9046742755636961) (-0.0552763819095477,0.9046149620266568) (-0.04522613065326632,0.904514426886113) (-0.0351758793969849,0.9048189997195295) (-0.025125628140703474,0.9047562538698531) (-0.015075376884422112,0.9045092613516896) (-0.005025125628140694,0.9043200296772111) (0.005025125628140765,0.903694002281506) (0.015075376884422073,0.9041228315829913) (0.02512562814070352,0.9045420750888021) (0.035175879396984966,0.9046304999172976) (0.045226130653266354,0.9049348566675107) (0.05527638190954779,0.9048099335383503) (0.06532663316582915,0.9053319360863035) (0.07537688442211059,0.9054692175460418) (0.08542713567839198,0.9052533573646957) (0.09547738693467332,0.9058257683234731) (0.10552763819095473,0.9054491855086315) (0.11557788944723618,0.9052982335826378) (0.1256281407035176,0.9053502165517535) (0.13567839195979903,0.9053844431437861) (0.1457286432160804,0.9053264686983272) (0.15577889447236185,0.9050505016731958) (0.16582914572864324,0.9053069514880943) (0.17587939698492464,0.905808600359406) (0.18592964824120609,0.9055165454431455) (0.19597989949748737,0.9056919553526172) (0.20603015075376882,0.9054622686783862) (0.21608040201005027,0.9056463988628342) (0.2261306532663317,0.9055545237174483) (0.23618090452261306,0.9056586950105532) (0.24623115577889448,0.9057680117434064) (0.2562814070351759,0.9062630558658902) (0.26633165829145733,0.906485033725546) (0.27638190954773867,0.9066246354717683) (0.2864321608040201,0.9076268332487537) (0.2964824120603015,0.9078601957539664) (0.30653266331658285,0.907833022263065) (0.31658291457286436,0.9079906877565379) (0.3266331658291457,0.9082010960751629) (0.33668341708542715,0.9084120022829049) (0.34673366834170855,0.908233232424553) (0.35678391959798994,0.9082193039449036) (0.36683417085427134,0.908067970818421) (0.3768844221105528,0.9079005496335661) (0.3869346733668342,0.9082599265928107) (0.39698492462311563,0.9084054913730284) (0.407035175879397,0.9085387700302007) (0.4170854271356784,0.9083390804722734) (0.42713567839195976,0.908377536500334) (0.4371859296482412,0.9088053208293124) (0.44723618090452266,0.9088732823655276) (0.45728643216080406,0.9090275616369131) (0.4673366834170854,0.9088977180146312) (0.47738693467336685,0.9091874447863142) (0.4874371859296483,0.909439586157637) (0.49748743718592964,0.9100820014047228) (0.5075376884422111,0.9100005105375678) (0.5175879396984925,0.9100867494880465) (0.5276381909547738,0.909992543560245) (0.5376884422110553,0.9101063738152779) (0.5477386934673366,0.9104648786457664) (0.5577889447236181,0.9108635960151817) (0.5678391959798995,0.9114763867455451) (0.5778894472361809,0.9111662497177454) (0.5879396984924624,0.911389012702131) (0.5979899497487438,0.9116207579236794) (0.6080402010050252,0.911885273358287) (0.6180904522613065,0.9121204587918336) (0.6281407035175879,0.9122461800498021) (0.6381909547738693,0.9124170340941583) (0.6482412060301508,0.9118673312379985) (0.6582914572864322,0.9118488776993113) (0.6683417085427137,0.9120326269749139) (0.678391959798995,0.912660951065263) (0.6884422110552764,0.9130292148082925) (0.6984924623115579,0.9129349371017252) (0.7085427135678392,0.9132673844237519) (0.7185929648241206,0.9135971596365091) (0.728643216080402,0.9141891159563141) (0.7386934673366834,0.914151624236777) (0.7487437185929648,0.9144939508439345) (0.7587939698492462,0.9144930843376647) (0.7688442211055276,0.9146216922027046) (0.7788944723618091,0.9147364211960105) (0.7889447236180904,0.914835573923349) (0.7989949748743719,0.9152469402724631) (0.8090452261306533,0.9161240066316417) (0.8190954773869347,0.9162796844130146) (0.8291457286432161,0.9163983753050995) (0.8391959798994975,0.9168147457481071) (0.8492462311557789,0.9173450749610096) (0.8592964824120604,0.9170922568446334) (0.8693467336683417,0.9169602427642527) (0.879396984924623,0.9170001037991241) (0.8894472361809046,0.9170188071110927) (0.8994974874371859,0.9168945953980411) (0.9095477386934674,0.9172532204321489) (0.9195979899497488,0.9170815357630908) (0.9296482412060301,0.9170851074267455) (0.9396984924623116,0.9174735370948367) (0.9497487437185931,0.9179002429368494) (0.9597989949748743,0.9178926887066632) (0.9698492462311558,0.9179468274554102) (0.9798994974874373,0.9185225586158251) (0.9899497487437185,0.9186454106345868) (1.0,0.9185828812018723) };
  \addplot[name path=L_jeffreys_grp, draw=none, mark=none, forget plot] coordinates {(-1.0,0.8897119928315192) (-0.9899497487437185,0.8894003884236986) (-0.9798994974874372,0.8896599884743095) (-0.9698492462311558,0.8898255904265354) (-0.9597989949748743,0.8898667471139942) (-0.9497487437185931,0.8897469229367259) (-0.9396984924623115,0.8896456481069017) (-0.9296482412060301,0.8894946636510911) (-0.9195979899497487,0.8897907098191236) (-0.9095477386934674,0.8894478196942968) (-0.8994974874371859,0.8885025981344156) (-0.8894472361809044,0.888654003627838) (-0.8793969849246231,0.8876187823630376) (-0.8693467336683417,0.8875410602802807) (-0.8592964824120602,0.8874338033221848) (-0.8492462311557789,0.8871117759571554) (-0.8391959798994976,0.8871354696159521) (-0.829145728643216,0.8877689430849278) (-0.8190954773869348,0.8881827946163667) (-0.8090452261306533,0.8872363143776408) (-0.7989949748743718,0.8874918102316882) (-0.7889447236180904,0.8872189582138067) (-0.7788944723618091,0.8871208009151103) (-0.7688442211055276,0.8865204827089778) (-0.7587939698492462,0.8867746222648586) (-0.7487437185929648,0.8862171176192093) (-0.7386934673366834,0.8863035583349637) (-0.7286432160804021,0.8857753214634726) (-0.7185929648241205,0.8855624922206209) (-0.7085427135678393,0.8854218604981894) (-0.6984924623115578,0.8852047642174385) (-0.6884422110552764,0.8850233449519449) (-0.678391959798995,0.8853027398725113) (-0.6683417085427136,0.8854962879070654) (-0.6582914572864321,0.8853605221883105) (-0.6482412060301508,0.8853970767773801) (-0.6381909547738693,0.8858428049015767) (-0.628140703517588,0.8853437446129705) (-0.6180904522613065,0.885654293191657) (-0.6080402010050251,0.8853466705035774) (-0.5979899497487438,0.8853040292883895) (-0.5879396984924624,0.8850136257365517) (-0.5778894472361809,0.8847876453584381) (-0.5678391959798995,0.8848747796594298) (-0.5577889447236181,0.8848745743446911) (-0.5477386934673367,0.8846260185838677) (-0.5376884422110553,0.8843690441151166) (-0.5276381909547738,0.8839974127504582) (-0.5175879396984925,0.8839629732206901) (-0.507537688442211,0.8839712828526436) (-0.49748743718592964,0.8830694399556109) (-0.48743718592964824,0.8830842144928782) (-0.47738693467336674,0.883040678972846) (-0.46733668341708545,0.8830005519446097) (-0.45728643216080406,0.882932994639188) (-0.4472361809045226,0.8831312441833772) (-0.4371859296482412,0.8828399376713238) (-0.4271356783919598,0.8825283995738419) (-0.4170854271356783,0.8827639117744623) (-0.40703517587939697,0.8825361465940449) (-0.3969849246231155,0.882169244102912) (-0.38693467336683424,0.8823215568862425) (-0.37688442211055273,0.8815024522093696) (-0.36683417085427134,0.8817607304369972) (-0.35678391959798994,0.8817669268670745) (-0.34673366834170855,0.8822075923065366) (-0.3366834170854271,0.8822320371362489) (-0.32663316582914564,0.8816808678738184) (-0.31658291457286436,0.8815918491697545) (-0.3065326633165829,0.8810877533732632) (-0.29648241206030146,0.8814760435508802) (-0.2864321608040201,0.8816387913281454) (-0.2763819095477387,0.8812832430143449) (-0.26633165829145733,0.8806581821675994) (-0.25628140703517593,0.8804919242220933) (-0.24623115577889443,0.8799761595417209) (-0.236180904522613,0.8800926089071327) (-0.22613065326633172,0.8798329924866501) (-0.21608040201005027,0.8800883756189429) (-0.20603015075376882,0.8800829404970366) (-0.19597989949748745,0.879790314384045) (-0.18592964824120603,0.8798811317163107) (-0.1758793969849246,0.8800909234708907) (-0.16582914572864324,0.8801199399203237) (-0.15577889447236182,0.8798863632428078) (-0.14572864321608042,0.8798515388719649) (-0.135678391959799,0.8803456263301208) (-0.12562814070351758,0.8807324125279824) (-0.11557788944723618,0.8806001809415724) (-0.10552763819095479,0.8800211959436166) (-0.09547738693467334,0.8799077238723674) (-0.08542713567839195,0.879306607119592) (-0.07537688442211053,0.8788767054986926) (-0.06532663316582918,0.8786253802292847) (-0.0552763819095477,0.8785685917050405) (-0.04522613065326632,0.878337226655577) (-0.0351758793969849,0.8786573069531249) (-0.025125628140703474,0.8785733629044702) (-0.015075376884422112,0.8782192790622683) (-0.005025125628140694,0.8780540141155437) (0.005025125628140765,0.8775423031809283) (0.015075376884422073,0.8779806326644134) (0.02512562814070352,0.8782680670625428) (0.035175879396984966,0.8781986727477151) (0.045226130653266354,0.8783479745890216) (0.05527638190954779,0.8780926442045053) (0.06532663316582915,0.8785617935475765) (0.07537688442211059,0.8787012644864803) (0.08542713567839198,0.8785300078596021) (0.09547738693467332,0.8790649130328928) (0.10552763819095473,0.878761068065862) (0.11557788944723618,0.8786754921223847) (0.1256281407035176,0.8789268276586579) (0.13567839195979903,0.8789061750886935) (0.1457286432160804,0.8788412501322146) (0.15577889447236185,0.8785796370375087) (0.16582914572864324,0.878774723805847) (0.17587939698492464,0.8790898085754174) (0.18592964824120609,0.8787896845776493) (0.19597989949748737,0.8791485139881862) (0.20603015075376882,0.878977171702723) (0.21608040201005027,0.8793568145914862) (0.2261306532663317,0.8793427745007831) (0.23618090452261306,0.8794628868532568) (0.24623115577889448,0.8796378002753358) (0.2562814070351759,0.8801202792950238) (0.26633165829145733,0.8802424918941391) (0.27638190954773867,0.8802225624798136) (0.2864321608040201,0.8811220550613051) (0.2964824120603015,0.8813188338514487) (0.30653266331658285,0.8813673283262021) (0.31658291457286436,0.8814371674130491) (0.3266331658291457,0.8814845851265839) (0.33668341708542715,0.8816651193615895) (0.34673366834170855,0.8814550702080175) (0.35678391959798994,0.8814569260848224) (0.36683417085427134,0.8813657477792375) (0.3768844221105528,0.8812513256300847) (0.3869346733668342,0.881577068261778) (0.39698492462311563,0.8817166040304323) (0.407035175879397,0.88183050728341) (0.4170854271356784,0.8815985917907166) (0.42713567839195976,0.8815119348981674) (0.4371859296482412,0.8819591546535118) (0.44723618090452266,0.8820182719974711) (0.45728643216080406,0.882303876434063) (0.4673366834170854,0.882314281756422) (0.47738693467336685,0.8824958067055828) (0.4874371859296483,0.88281003770044) (0.49748743718592964,0.8833126307272297) (0.5075376884422111,0.8833256205793115) (0.5175879396984925,0.8833166207463228) (0.5276381909547738,0.8831649384371436) (0.5376884422110553,0.8832141095479644) (0.5477386934673366,0.8835936616053975) (0.5577889447236181,0.8838768657700268) (0.5678391959798995,0.8844787398104177) (0.5778894472361809,0.8842464008592811) (0.5879396984924624,0.884398036293384) (0.5979899497487438,0.8845688774877772) (0.6080402010050252,0.884825142095512) (0.6180904522613065,0.8849689210004094) (0.6281407035175879,0.8850850115750533) (0.6381909547738693,0.8851924120756834) (0.6482412060301508,0.884658226634038) (0.6582914572864322,0.8845715021843704) (0.6683417085427137,0.8847925271564987) (0.678391959798995,0.8853759244453097) (0.6884422110552764,0.8856745969677263) (0.6984924623115579,0.8856659948348105) (0.7085427135678392,0.8860288496560335) (0.7185929648241206,0.8863354663763545) (0.728643216080402,0.8868904204962039) (0.7386934673366834,0.8868270549686491) (0.7487437185929648,0.8871567142675101) (0.7587939698492462,0.8871463437789958) (0.7688442211055276,0.8871497019488748) (0.7788944723618091,0.887302471199167) (0.7889447236180904,0.8874376237440448) (0.7989949748743719,0.8877992952722696) (0.8090452261306533,0.8887129424790462) (0.8190954773869347,0.8887038357269372) (0.8291457286432161,0.888814060848989) (0.8391959798994975,0.8892071993005313) (0.8492462311557789,0.8897912924586512) (0.8592964824120604,0.8895434465824581) (0.8693467336683417,0.8895167314588124) (0.879396984924623,0.8895007670546017) (0.8894472361809046,0.8894602341910953) (0.8994974874371859,0.8893154791954986) (0.9095477386934674,0.8895985575287526) (0.9195979899497488,0.8890727492584536) (0.9296482412060301,0.8890588304779438) (0.9396984924623116,0.889391303019338) (0.9497487437185931,0.8897979814367191) (0.9597989949748743,0.8898324930422752) (0.9698492462311558,0.8898584537681512) (0.9798994974874373,0.8903899089402508) (0.9899497487437185,0.8905224548423273) (1.0,0.8903862082502721) };
  \addplot[ours_jeffreys, fill, fill opacity=0.18, draw=none, mark=none, forget plot] fill between [of=U_jeffreys_grp and L_jeffreys_grp];
  \addplot+[ours_jeffreys, mark=none] coordinates {(-1.0,0.9035094882633431) (-0.9899497487437185,0.9030924398530901) (-0.9798994974874372,0.9033753921458031) (-0.9698492462311558,0.9035874325030548) (-0.9597989949748743,0.9037254868733743) (-0.9497487437185931,0.9036361583689788) (-0.9396984924623115,0.9034975021144835) (-0.9296482412060301,0.9031829223872655) (-0.9195979899497487,0.9034906127073393) (-0.9095477386934674,0.9030949054007665) (-0.8994974874371859,0.9021293359520551) (-0.8894472361809044,0.9021511268242945) (-0.8793969849246231,0.9012023260297372) (-0.8693467336683417,0.9010549601043993) (-0.8592964824120602,0.9009506841210795) (-0.8492462311557789,0.900610244182776) (-0.8391959798994976,0.900584636572028) (-0.829145728643216,0.9012733273763991) (-0.8190954773869348,0.9017138436553344) (-0.8090452261306533,0.9006751693720125) (-0.7989949748743718,0.9009390351803185) (-0.7889447236180904,0.9006580841604683) (-0.7788944723618091,0.9006164563703125) (-0.7688442211055276,0.8999888160730554) (-0.7587939698492462,0.9001323722587574) (-0.7487437185929648,0.8996289194142983) (-0.7386934673366834,0.8996802988066022) (-0.7286432160804021,0.8990656485865111) (-0.7185929648241205,0.8987794448160099) (-0.7085427135678393,0.898631139256374) (-0.6984924623115578,0.898442079408983) (-0.6884422110552764,0.8982500964623671) (-0.678391959798995,0.8985098849693052) (-0.6683417085427136,0.8986045123239174) (-0.6582914572864321,0.8985151444694768) (-0.6482412060301508,0.898531261066036) (-0.6381909547738693,0.8989954436313716) (-0.628140703517588,0.8984827930356717) (-0.6180904522613065,0.8988205478869519) (-0.6080402010050251,0.8985097663493826) (-0.5979899497487438,0.8984528992265639) (-0.5879396984924624,0.8982013867072876) (-0.5778894472361809,0.8979029406026678) (-0.5678391959798995,0.8979668900071739) (-0.5577889447236181,0.8980131674554586) (-0.5477386934673367,0.8978125241128543) (-0.5376884422110553,0.8975199135441796) (-0.5276381909547738,0.897138341558601) (-0.5175879396984925,0.897132869397035) (-0.507537688442211,0.8970589629934638) (-0.49748743718592964,0.8961833084078173) (-0.48743718592964824,0.8961954422743399) (-0.47738693467336674,0.8960996473320619) (-0.46733668341708545,0.8961307106807518) (-0.45728643216080406,0.8961185155059174) (-0.4472361809045226,0.8963013300745523) (-0.4371859296482412,0.8959968290716445) (-0.4271356783919598,0.895686458165392) (-0.4170854271356783,0.8959419647878595) (-0.40703517587939697,0.8957029988705759) (-0.3969849246231155,0.8953896874768293) (-0.38693467336683424,0.8955308246869238) (-0.37688442211055273,0.8946658794621631) (-0.36683417085427134,0.8949174662462303) (-0.35678391959798994,0.8948629146029528) (-0.34673366834170855,0.8952280940475029) (-0.3366834170854271,0.8951883468412508) (-0.32663316582914564,0.8945693991003915) (-0.31658291457286436,0.894498527715569) (-0.3065326633165829,0.8941275811327346) (-0.29648241206030146,0.8945510898022774) (-0.2864321608040201,0.894647323610288) (-0.2763819095477387,0.8941813041763431) (-0.26633165829145733,0.8935607897665235) (-0.25628140703517593,0.8934560669012321) (-0.24623115577889443,0.8929592739328837) (-0.236180904522613,0.8931098981904944) (-0.22613065326633172,0.8928330225611847) (-0.21608040201005027,0.8930715164135278) (-0.20603015075376882,0.8930740961883279) (-0.19597989949748745,0.8927725480731181) (-0.18592964824120603,0.8928509171589583) (-0.1758793969849246,0.8930424980753847) (-0.16582914572864324,0.8931034402776796) (-0.15577889447236182,0.8928588899867158) (-0.14572864321608042,0.8928057634254433) (-0.135678391959799,0.8933029110139641) (-0.12562814070351758,0.8936792677195702) (-0.11557788944723618,0.8935253144488638) (-0.10552763819095479,0.8929829779783459) (-0.09547738693467334,0.8929185508227422) (-0.08542713567839195,0.8923548197347787) (-0.07537688442211053,0.8919439308577752) (-0.06532663316582918,0.8916498278964904) (-0.0552763819095477,0.8915917768658487) (-0.04522613065326632,0.891425826770845) (-0.0351758793969849,0.8917381533363272) (-0.025125628140703474,0.8916648083871617) (-0.015075376884422112,0.891364270206979) (-0.005025125628140694,0.8911870218963774) (0.005025125628140765,0.8906181527312171) (0.015075376884422073,0.8910517321237024) (0.02512562814070352,0.8914050710756725) (0.035175879396984966,0.8914145863325064) (0.045226130653266354,0.8916414156282662) (0.05527638190954779,0.8914512888714278) (0.06532663316582915,0.89194686481694) (0.07537688442211059,0.8920852410162611) (0.08542713567839198,0.8918916826121489) (0.09547738693467332,0.8924453406781829) (0.10552763819095473,0.8921051267872467) (0.11557788944723618,0.8919868628525113) (0.1256281407035176,0.8921385221052057) (0.13567839195979903,0.8921453091162398) (0.1457286432160804,0.8920838594152709) (0.15577889447236185,0.8918150693553523) (0.16582914572864324,0.8920408376469706) (0.17587939698492464,0.8924492044674117) (0.18592964824120609,0.8921531150103974) (0.19597989949748737,0.8924202346704017) (0.20603015075376882,0.8922197201905546) (0.21608040201005027,0.8925016067271602) (0.2261306532663317,0.8924486491091157) (0.23618090452261306,0.892560790931905) (0.24623115577889448,0.8927029060093711) (0.2562814070351759,0.893191667580457) (0.26633165829145733,0.8933637628098425) (0.27638190954773867,0.8934235989757909) (0.2864321608040201,0.8943744441550294) (0.2964824120603015,0.8945895148027075) (0.30653266331658285,0.8946001752946335) (0.31658291457286436,0.8947139275847935) (0.3266331658291457,0.8948428406008734) (0.33668341708542715,0.8950385608222472) (0.34673366834170855,0.8948441513162853) (0.35678391959798994,0.894838115014863) (0.36683417085427134,0.8947168592988293) (0.3768844221105528,0.8945759376318254) (0.3869346733668342,0.8949184974272943) (0.39698492462311563,0.8950610477017303) (0.407035175879397,0.8951846386568053) (0.4170854271356784,0.894968836131495) (0.42713567839195976,0.8949447356992507) (0.4371859296482412,0.8953822377414121) (0.44723618090452266,0.8954457771814993) (0.45728643216080406,0.8956657190354881) (0.4673366834170854,0.8956059998855266) (0.47738693467336685,0.8958416257459485) (0.4874371859296483,0.8961248119290385) (0.49748743718592964,0.8966973160659762) (0.5075376884422111,0.8966630655584397) (0.5175879396984925,0.8967016851171846) (0.5276381909547738,0.8965787409986943) (0.5376884422110553,0.8966602416816212) (0.5477386934673366,0.897029270125582) (0.5577889447236181,0.8973702308926043) (0.5678391959798995,0.8979775632779814) (0.5778894472361809,0.8977063252885132) (0.5879396984924624,0.8978935244977575) (0.5979899497487438,0.8980948177057283) (0.6080402010050252,0.8983552077268995) (0.6180904522613065,0.8985446898961215) (0.6281407035175879,0.8986655958124277) (0.6381909547738693,0.8988047230849209) (0.6482412060301508,0.8982627789360182) (0.6582914572864322,0.8982101899418409) (0.6683417085427137,0.8984125770657063) (0.678391959798995,0.8990184377552863) (0.6884422110552764,0.8993519058880094) (0.6984924623115579,0.8993004659682678) (0.7085427135678392,0.8996481170398927) (0.7185929648241206,0.8999663130064318) (0.728643216080402,0.900539768226259) (0.7386934673366834,0.900489339602713) (0.7487437185929648,0.9008253325557223) (0.7587939698492462,0.9008197140583303) (0.7688442211055276,0.9008856970757897) (0.7788944723618091,0.9010194461975888) (0.7889447236180904,0.9011365988336969) (0.7989949748743719,0.9015231177723664) (0.8090452261306533,0.902418474555344) (0.8190954773869347,0.9024917600699759) (0.8291457286432161,0.9026062180770442) (0.8391959798994975,0.9030109725243192) (0.8492462311557789,0.9035681837098304) (0.8592964824120604,0.9033178517135457) (0.8693467336683417,0.9032384871115325) (0.879396984924623,0.9032504354268629) (0.8894472361809046,0.903239520651094) (0.8994974874371859,0.9031050372967698) (0.9095477386934674,0.9034258889804507) (0.9195979899497488,0.9030771425107722) (0.9296482412060301,0.9030719689523447) (0.9396984924623116,0.9034324200570873) (0.9497487437185931,0.9038491121867842) (0.9597989949748743,0.9038625908744692) (0.9698492462311558,0.9039026406117807) (0.9798994974874373,0.904456233778038) (0.9899497487437185,0.904583932738457) (1.0,0.9044845447260722) };
  \addplot+[ours_jeffreys, mark=none, line width=0.4pt, dashed, forget plot] coordinates {(-1.0,0.7956015120124357) (1.0,0.7956015120124357)};
  \addplot+[ours_jeffreys, mark=*, only marks, mark size=\PanelArgminSize, forget plot] coordinates {(0.005025125628140765,0.8906181527312171)};
  \addplot[name path=U_ig_robust_grp, draw=none, mark=none, forget plot] coordinates {(-1.0,0.9247521382951152) (-0.9899497487437184,0.9242645754891394) (-0.9798994974874372,0.9231927133915784) (-0.9698492462311558,0.9225263479221348) (-0.9597989949748743,0.9220532830782239) (-0.9497487437185931,0.9220024244888965) (-0.9396984924623115,0.9225725964140654) (-0.9296482412060301,0.9216727150676445) (-0.9195979899497488,0.9209632009074973) (-0.9095477386934673,0.9200147622468243) (-0.8994974874371859,0.9189334938708739) (-0.8894472361809044,0.918823287349763) (-0.8793969849246231,0.9179984120295224) (-0.8693467336683416,0.9176113598250749) (-0.8592964824120601,0.9179504150256383) (-0.8492462311557789,0.9179747976075021) (-0.8391959798994975,0.9179200228841563) (-0.8291457286432159,0.9180089892398998) (-0.8190954773869347,0.9175308478253176) (-0.8090452261306532,0.9171266216944371) (-0.7989949748743718,0.9165530532188138) (-0.7889447236180904,0.9167016723261054) (-0.7788944723618091,0.9164126586321072) (-0.7688442211055275,0.9170151306543679) (-0.7587939698492463,0.9166446069000818) (-0.7487437185929648,0.9166438949882325) (-0.7386934673366834,0.9162291580185831) (-0.728643216080402,0.916268064335703) (-0.7185929648241206,0.9158324814521608) (-0.7085427135678392,0.9159798545195316) (-0.6984924623115576,0.9156503960718423) (-0.6884422110552765,0.915508065351022) (-0.6783919597989949,0.9154413263608958) (-0.6683417085427135,0.9148129565164955) (-0.6582914572864322,0.9149561247713219) (-0.6482412060301508,0.9147268932598139) (-0.6381909547738692,0.9147383936811726) (-0.628140703517588,0.914971546719409) (-0.6180904522613065,0.9149639549492479) (-0.6080402010050251,0.9147283272527295) (-0.5979899497487438,0.9147581061237645) (-0.5879396984924623,0.9139913216227044) (-0.5778894472361809,0.913378582722888) (-0.5678391959798996,0.913209433063796) (-0.5577889447236181,0.9130172618663148) (-0.5477386934673366,0.9130232130783547) (-0.5376884422110553,0.9133054179411136) (-0.5276381909547738,0.9131041720257812) (-0.5175879396984925,0.9130673295255805) (-0.507537688442211,0.9125826963303204) (-0.49748743718592964,0.9122867686655124) (-0.48743718592964824,0.9126263485986953) (-0.47738693467336685,0.912129865959764) (-0.4673366834170854,0.9120378308661694) (-0.457286432160804,0.9117321694793284) (-0.4472361809045226,0.912075053605102) (-0.43718592964824116,0.9122295011291163) (-0.4271356783919598,0.9116778777190374) (-0.4170854271356783,0.9117118073508559) (-0.40703517587939697,0.9115631779645575) (-0.3969849246231156,0.9114198120539799) (-0.38693467336683407,0.911691000050776) (-0.37688442211055273,0.9115518603847448) (-0.3668341708542714,0.9113340293679328) (-0.35678391959798994,0.9111635054853857) (-0.3467336683417085,0.9107218700984758) (-0.33668341708542704,0.91062969971328) (-0.3266331658291457,0.9104871303070899) (-0.31658291457286436,0.9105289998064693) (-0.30653266331658285,0.9104447212475091) (-0.29648241206030146,0.9103049932793738) (-0.2864321608040202,0.9104498879254921) (-0.2763819095477386,0.9108334948296412) (-0.26633165829145733,0.910760066244642) (-0.2562814070351758,0.9105718732174862) (-0.24623115577889443,0.9106154336121135) (-0.2361809045226131,0.9104960301306594) (-0.22613065326633164,0.9105294407012764) (-0.2160804020100502,0.9105905509398339) (-0.20603015075376882,0.9100352676922758) (-0.19597989949748737,0.9099354224241032) (-0.18592964824120603,0.9095669059611601) (-0.1758793969849246,0.9093753907996874) (-0.16582914572864324,0.9088482038498175) (-0.15577889447236182,0.9079377618823695) (-0.14572864321608042,0.908090430137645) (-0.135678391959799,0.9083293707590466) (-0.12562814070351758,0.9085216000185119) (-0.11557788944723618,0.9083713248012721) (-0.10552763819095479,0.908188289660923) (-0.0954773869346734,0.908290836445151) (-0.0854271356783919,0.9080064463912073) (-0.0753768844221106,0.9078663645977884) (-0.06532663316582912,0.9073898412644492) (-0.055276381909547756,0.9077474766081367) (-0.04522613065326632,0.9075701934209629) (-0.03517587939698496,0.9074077609177621) (-0.025125628140703526,0.9073827729361305) (-0.015075376884422112,0.9072869077454288) (-0.005025125628140694,0.9071044471953588) (0.00502512562814067,0.906830282276812) (0.015075376884422073,0.9066836174743822) (0.02512562814070352,0.9067232273096841) (0.035175879396984966,0.9074036049914295) (0.045226130653266264,0.9071242585879266) (0.0552763819095477,0.907279000183932) (0.06532663316582915,0.907276058405874) (0.07537688442211059,0.9074746951793524) (0.08542713567839191,0.907550259099674) (0.09547738693467332,0.9071090769242738) (0.10552763819095481,0.906679310139401) (0.11557788944723618,0.9065707716967668) (0.12562814070351755,0.9064423066303938) (0.13567839195979894,0.9059543239295111) (0.1457286432160804,0.9061893511166552) (0.15577889447236185,0.9058859939743761) (0.16582914572864324,0.9062503673237484) (0.17587939698492464,0.9059751900475939) (0.18592964824120609,0.9059633600705103) (0.19597989949748745,0.9054560275100584) (0.20603015075376882,0.9052718364706038) (0.21608040201005022,0.9050784513589871) (0.22613065326633164,0.9051726459849831) (0.23618090452261306,0.9053411447747797) (0.24623115577889443,0.905275623675715) (0.2562814070351758,0.9053358003870122) (0.2663316582914573,0.9054743570779948) (0.2763819095477386,0.9053976121945476) (0.28643216080402006,0.9056031825024495) (0.2964824120603015,0.9059374339453447) (0.30653266331658285,0.9060737912339885) (0.31658291457286436,0.9057341234865909) (0.3266331658291457,0.9055899770674035) (0.33668341708542715,0.9054820134927344) (0.3467336683417085,0.9058842003976826) (0.35678391959798994,0.9057974060371896) (0.36683417085427134,0.9060418416150832) (0.37688442211055273,0.9063061565420995) (0.3869346733668341,0.9063769785482182) (0.3969849246231155,0.9062313189834783) (0.40703517587939697,0.9063539422105061) (0.4170854271356784,0.9065297900586332) (0.4271356783919598,0.9066768678061679) (0.4371859296482412,0.9066451430830449) (0.44723618090452255,0.9068593510921656) (0.457286432160804,0.906970185794498) (0.4673366834170854,0.9073469109569082) (0.4773869346733668,0.9076372722961067) (0.48743718592964824,0.9077269533270972) (0.49748743718592964,0.9077923441455064) (0.507537688442211,0.9082120467694571) (0.5175879396984924,0.9082124889040791) (0.5276381909547738,0.9081563364745231) (0.5376884422110553,0.9083467578455026) (0.5477386934673366,0.9084461007683322) (0.5577889447236181,0.9084585239495944) (0.5678391959798995,0.9086201703667376) (0.5778894472361809,0.9088648353505712) (0.5879396984924623,0.9092819645604165) (0.5979899497487438,0.9098856909306481) (0.6080402010050251,0.9097685013303892) (0.6180904522613065,0.9095899076436522) (0.6281407035175879,0.9094540828137453) (0.6381909547738693,0.9096653358513694) (0.6482412060301508,0.9092669110030799) (0.6582914572864321,0.9095874157283463) (0.6683417085427136,0.9098191838356398) (0.678391959798995,0.9100508982272262) (0.6884422110552764,0.9104074751176874) (0.6984924623115577,0.9104721423186987) (0.7085427135678392,0.9102673470237262) (0.7185929648241206,0.9105447426393423) (0.728643216080402,0.9106798125630791) (0.7386934673366834,0.9108450309762239) (0.7487437185929647,0.9109791645122894) (0.7587939698492462,0.9116101437339511) (0.7688442211055275,0.9119068600705105) (0.778894472361809,0.912127034026508) (0.7889447236180904,0.9121492522370606) (0.7989949748743719,0.9120877933686202) (0.8090452261306532,0.9119934712095379) (0.8190954773869347,0.9120879648907938) (0.829145728643216,0.9118791423745094) (0.8391959798994975,0.9118725280395311) (0.8492462311557788,0.9118970272130449) (0.8592964824120602,0.9124772651045353) (0.8693467336683417,0.912687978657281) (0.8793969849246231,0.9129974266464669) (0.8894472361809045,0.9133080935447632) (0.8994974874371858,0.9132171028055072) (0.9095477386934673,0.9132142945020056) (0.9195979899497487,0.9132033831678197) (0.9296482412060301,0.9134978792871656) (0.9396984924623115,0.9135177987781741) (0.949748743718593,0.9138587177452788) (0.9597989949748744,0.9140026444823347) (0.9698492462311558,0.9140248658377953) (0.9798994974874371,0.9146349922281877) (0.9899497487437185,0.9146327471049732) (0.9999999999999999,0.914433638478665) };
  \addplot[name path=L_ig_robust_grp, draw=none, mark=none, forget plot] coordinates {(-1.0,0.897591435361127) (-0.9899497487437184,0.8971725926944067) (-0.9798994974874372,0.8964455904763154) (-0.9698492462311558,0.8956933422630677) (-0.9597989949748743,0.8951016790876177) (-0.9497487437185931,0.89507377469627) (-0.9396984924623115,0.8953731917455163) (-0.9296482412060301,0.8945690919536118) (-0.9195979899497488,0.8941153520641163) (-0.9095477386934673,0.8932340122979502) (-0.8994974874371859,0.8922628208606164) (-0.8894472361809044,0.8921773805893394) (-0.8793969849246231,0.8914940249354648) (-0.8693467336683416,0.8911074939980507) (-0.8592964824120601,0.8915202919095999) (-0.8492462311557789,0.8917170833029087) (-0.8391959798994975,0.8916459598866419) (-0.8291457286432159,0.8916330718939668) (-0.8190954773869347,0.8914256879053232) (-0.8090452261306532,0.8908835767677832) (-0.7989949748743718,0.8903765965429826) (-0.7889447236180904,0.8903898102818986) (-0.7788944723618091,0.8901664230445031) (-0.7688442211055275,0.8907345003307814) (-0.7587939698492463,0.8904560305746534) (-0.7487437185929648,0.8905733857176534) (-0.7386934673366834,0.8901653878308962) (-0.728643216080402,0.890049927691222) (-0.7185929648241206,0.8898877922710371) (-0.7085427135678392,0.8902700967776321) (-0.6984924623115576,0.8898878391742906) (-0.6884422110552765,0.8897448360350865) (-0.6783919597989949,0.8896875077430872) (-0.6683417085427135,0.8889808492736209) (-0.6582914572864322,0.8890944037348926) (-0.6482412060301508,0.8888610786267624) (-0.6381909547738692,0.8888335058348513) (-0.628140703517588,0.8890425967694876) (-0.6180904522613065,0.8890213312379027) (-0.6080402010050251,0.8887445053689077) (-0.5979899497487438,0.8886823955017863) (-0.5879396984924623,0.8880191847173474) (-0.5778894472361809,0.8874424004020571) (-0.5678391959798996,0.8874514842243945) (-0.5577889447236181,0.8872965752277571) (-0.5477386934673366,0.8871692411734231) (-0.5376884422110553,0.8872015550915401) (-0.5276381909547738,0.8871276240163001) (-0.5175879396984925,0.8870512745236512) (-0.507537688442211,0.8867246635391369) (-0.49748743718592964,0.8865124863756954) (-0.48743718592964824,0.8868013286162383) (-0.47738693467336685,0.8864567432896621) (-0.4673366834170854,0.8863407299793722) (-0.457286432160804,0.8860251530322599) (-0.4472361809045226,0.8862577440215228) (-0.43718592964824116,0.8862170342613717) (-0.4271356783919598,0.8858516840465903) (-0.4170854271356783,0.885982707823264) (-0.40703517587939697,0.8858308302396104) (-0.3969849246231156,0.8857313748710477) (-0.38693467336683407,0.8859518451767957) (-0.37688442211055273,0.8856700797772838) (-0.3668341708542714,0.8854120196861643) (-0.35678391959798994,0.8853279838425949) (-0.3467336683417085,0.8848967254185018) (-0.33668341708542704,0.8846105484571826) (-0.3266331658291457,0.8844288707201283) (-0.31658291457286436,0.884368680820525) (-0.30653266331658285,0.884316383975851) (-0.29648241206030146,0.8841740088382956) (-0.2864321608040202,0.884270120825233) (-0.2763819095477386,0.8847198859247879) (-0.26633165829145733,0.884627116191467) (-0.2562814070351758,0.8843582020768176) (-0.24623115577889443,0.8844743719230082) (-0.2361809045226131,0.8843215521088171) (-0.22613065326633164,0.8844373531922863) (-0.2160804020100502,0.8844215905276053) (-0.20603015075376882,0.883947136696354) (-0.19597989949748737,0.8840558813624277) (-0.18592964824120603,0.8835569883432689) (-0.1758793969849246,0.8833973275813464) (-0.16582914572864324,0.8828975719264031) (-0.15577889447236182,0.8821997694682885) (-0.14572864321608042,0.8824568783851771) (-0.135678391959799,0.8825794403547119) (-0.12562814070351758,0.882721907003901) (-0.11557788944723618,0.8825568191850752) (-0.10552763819095479,0.8825361657607873) (-0.0954773869346734,0.8826737370800146) (-0.0854271356783919,0.8822808559145962) (-0.0753768844221106,0.882144354082456) (-0.06532663316582912,0.8817792209955615) (-0.055276381909547756,0.8820866019957166) (-0.04522613065326632,0.8816588792027765) (-0.03517587939698496,0.8816540233051805) (-0.025125628140703526,0.8816882863143404) (-0.015075376884422112,0.8816480143448832) (-0.005025125628140694,0.8813445198002477) (0.00502512562814067,0.8811289183953024) (0.015075376884422073,0.8810732325843014) (0.02512562814070352,0.8809759326441404) (0.035175879396984966,0.8818101075967991) (0.045226130653266264,0.8815656043283605) (0.0552763819095477,0.8816779682741119) (0.06532663316582915,0.8816499953659352) (0.07537688442211059,0.881829295957197) (0.08542713567839191,0.8818829024867978) (0.09547738693467332,0.8815458365837696) (0.10552763819095481,0.8810665355277055) (0.11557788944723618,0.8810558773352539) (0.12562814070351755,0.8807218399495144) (0.13567839195979894,0.8801717677647525) (0.1457286432160804,0.8803100202772287) (0.15577889447236185,0.8801267076518932) (0.16582914572864324,0.8804810975380283) (0.17587939698492464,0.8802405446882386) (0.18592964824120609,0.8801600141263467) (0.19597989949748745,0.8795690351860349) (0.20603015075376882,0.8793995570519257) (0.21608040201005022,0.8792032222691115) (0.22613065326633164,0.879140758052364) (0.23618090452261306,0.8792379293799624) (0.24623115577889443,0.8791778072466242) (0.2562814070351758,0.8792154318993548) (0.2663316582914573,0.8795950830941301) (0.2763819095477386,0.8794856975729549) (0.28643216080402006,0.8795176850495601) (0.2964824120603015,0.8798438110730972) (0.30653266331658285,0.8800146419105599) (0.31658291457286436,0.8797458582354879) (0.3266331658291457,0.8796285805008802) (0.33668341708542715,0.8795566682878859) (0.3467336683417085,0.8800171761383748) (0.35678391959798994,0.8798733884345269) (0.36683417085427134,0.8800683463220381) (0.37688442211055273,0.8803912616802282) (0.3869346733668341,0.8804692078665942) (0.3969849246231155,0.8803572746856365) (0.40703517587939697,0.8803323445851993) (0.4170854271356784,0.880331317026175) (0.4271356783919598,0.8804837810977446) (0.4371859296482412,0.8804827469367829) (0.44723618090452255,0.8806657674584992) (0.457286432160804,0.8808421245209564) (0.4673366834170854,0.8811505830591305) (0.4773869346733668,0.8814760678864632) (0.48743718592964824,0.8818080579096287) (0.49748743718592964,0.8818668785490821) (0.507537688442211,0.8821981593920479) (0.5175879396984924,0.8821826174584751) (0.5276381909547738,0.8821337161532474) (0.5376884422110553,0.8822344203412406) (0.5477386934673366,0.8822759952616324) (0.5577889447236181,0.8821725718536972) (0.5678391959798995,0.8823017516691365) (0.5778894472361809,0.882594220848473) (0.5879396984924623,0.8828965000452799) (0.5979899497487438,0.8834265616102701) (0.6080402010050251,0.8832865274225273) (0.6180904522613065,0.8830776338139228) (0.6281407035175879,0.8829427364563479) (0.6381909547738693,0.8832606632382317) (0.6482412060301508,0.8829709499446176) (0.6582914572864321,0.8831172170388668) (0.6683417085427136,0.8833862233255021) (0.678391959798995,0.8836157525789023) (0.6884422110552764,0.8839304558614185) (0.6984924623115577,0.8840024431489074) (0.7085427135678392,0.8838549381034243) (0.7185929648241206,0.8841316051752217) (0.728643216080402,0.884235618238101) (0.7386934673366834,0.8842993885353054) (0.7487437185929647,0.8844126036055929) (0.7587939698492462,0.8850539550627298) (0.7688442211055275,0.8853192382787916) (0.778894472361809,0.8854621435589225) (0.7889447236180904,0.8853208408052774) (0.7989949748743719,0.8852367391893687) (0.8090452261306532,0.8852679658308376) (0.8190954773869347,0.8853154472955682) (0.829145728643216,0.8851259080690732) (0.8391959798994975,0.8851143948774386) (0.8492462311557788,0.8850307454944004) (0.8592964824120602,0.8855947229466118) (0.8693467336683417,0.885681034313696) (0.8793969849246231,0.8860258319435805) (0.8894472361809045,0.8861612632999939) (0.8994974874371858,0.8861311284222084) (0.9095477386934673,0.8861232670268211) (0.9195979899497487,0.8861576837527896) (0.9296482412060301,0.8864105961887389) (0.9396984924623115,0.8863958604102701) (0.949748743718593,0.8867151514225275) (0.9597989949748744,0.8868338444055542) (0.9698492462311558,0.8868175456447366) (0.9798994974874371,0.8873877689837341) (0.9899497487437185,0.8873665267293446) (0.9999999999999999,0.887266076832696) };
  \addplot[ours_ig_robust, fill, fill opacity=0.18, draw=none, mark=none, forget plot] fill between [of=U_ig_robust_grp and L_ig_robust_grp];
  \addplot+[ours_ig_robust, mark=none] coordinates {(-1.0,0.911171786828121) (-0.9899497487437184,0.9107185840917731) (-0.9798994974874372,0.9098191519339469) (-0.9698492462311558,0.9091098450926013) (-0.9597989949748743,0.9085774810829208) (-0.9497487437185931,0.9085380995925832) (-0.9396984924623115,0.9089728940797909) (-0.9296482412060301,0.9081209035106281) (-0.9195979899497488,0.9075392764858068) (-0.9095477386934673,0.9066243872723873) (-0.8994974874371859,0.9055981573657451) (-0.8894472361809044,0.9055003339695512) (-0.8793969849246231,0.9047462184824936) (-0.8693467336683416,0.9043594269115628) (-0.8592964824120601,0.9047353534676191) (-0.8492462311557789,0.9048459404552054) (-0.8391959798994975,0.9047829913853991) (-0.8291457286432159,0.9048210305669333) (-0.8190954773869347,0.9044782678653204) (-0.8090452261306532,0.9040050992311102) (-0.7989949748743718,0.9034648248808982) (-0.7889447236180904,0.903545741304002) (-0.7788944723618091,0.9032895408383051) (-0.7688442211055275,0.9038748154925746) (-0.7587939698492463,0.9035503187373676) (-0.7487437185929648,0.903608640352943) (-0.7386934673366834,0.9031972729247396) (-0.728643216080402,0.9031589960134625) (-0.7185929648241206,0.9028601368615989) (-0.7085427135678392,0.9031249756485819) (-0.6984924623115576,0.9027691176230664) (-0.6884422110552765,0.9026264506930542) (-0.6783919597989949,0.9025644170519915) (-0.6683417085427135,0.9018969028950582) (-0.6582914572864322,0.9020252642531073) (-0.6482412060301508,0.9017939859432882) (-0.6381909547738692,0.901785949758012) (-0.628140703517588,0.9020070717444483) (-0.6180904522613065,0.9019926430935753) (-0.6080402010050251,0.9017364163108186) (-0.5979899497487438,0.9017202508127754) (-0.5879396984924623,0.9010052531700259) (-0.5778894472361809,0.9004104915624725) (-0.5678391959798996,0.9003304586440952) (-0.5577889447236181,0.900156918547036) (-0.5477386934673366,0.9000962271258889) (-0.5376884422110553,0.9002534865163269) (-0.5276381909547738,0.9001158980210406) (-0.5175879396984925,0.9000593020246158) (-0.507537688442211,0.8996536799347287) (-0.49748743718592964,0.8993996275206039) (-0.48743718592964824,0.8997138386074668) (-0.47738693467336685,0.899293304624713) (-0.4673366834170854,0.8991892804227708) (-0.457286432160804,0.8988786612557942) (-0.4472361809045226,0.8991663988133124) (-0.43718592964824116,0.899223267695244) (-0.4271356783919598,0.8987647808828139) (-0.4170854271356783,0.89884725758706) (-0.40703517587939697,0.898697004102084) (-0.3969849246231156,0.8985755934625138) (-0.38693467336683407,0.8988214226137858) (-0.37688442211055273,0.8986109700810143) (-0.3668341708542714,0.8983730245270486) (-0.35678391959798994,0.8982457446639903) (-0.3467336683417085,0.8978092977584888) (-0.33668341708542704,0.8976201240852313) (-0.3266331658291457,0.8974580005136091) (-0.31658291457286436,0.8974488403134971) (-0.30653266331658285,0.89738055261168) (-0.29648241206030146,0.8972395010588347) (-0.2864321608040202,0.8973600043753626) (-0.2763819095477386,0.8977766903772145) (-0.26633165829145733,0.8976935912180545) (-0.2562814070351758,0.8974650376471519) (-0.24623115577889443,0.8975449027675608) (-0.2361809045226131,0.8974087911197383) (-0.22613065326633164,0.8974833969467814) (-0.2160804020100502,0.8975060707337196) (-0.20603015075376882,0.8969912021943149) (-0.19597989949748737,0.8969956518932655) (-0.18592964824120603,0.8965619471522145) (-0.1758793969849246,0.8963863591905169) (-0.16582914572864324,0.8958728878881103) (-0.15577889447236182,0.895068765675329) (-0.14572864321608042,0.895273654261411) (-0.135678391959799,0.8954544055568793) (-0.12562814070351758,0.8956217535112064) (-0.11557788944723618,0.8954640719931737) (-0.10552763819095479,0.8953622277108552) (-0.0954773869346734,0.8954822867625828) (-0.0854271356783919,0.8951436511529017) (-0.0753768844221106,0.8950053593401222) (-0.06532663316582912,0.8945845311300054) (-0.055276381909547756,0.8949170393019267) (-0.04522613065326632,0.8946145363118697) (-0.03517587939698496,0.8945308921114713) (-0.025125628140703526,0.8945355296252354) (-0.015075376884422112,0.894467461045156) (-0.005025125628140694,0.8942244834978033) (0.00502512562814067,0.8939796003360572) (0.015075376884422073,0.8938784250293418) (0.02512562814070352,0.8938495799769123) (0.035175879396984966,0.8946068562941143) (0.045226130653266264,0.8943449314581435) (0.0552763819095477,0.894478484229022) (0.06532663316582915,0.8944630268859046) (0.07537688442211059,0.8946519955682747) (0.08542713567839191,0.8947165807932359) (0.09547738693467332,0.8943274567540217) (0.10552763819095481,0.8938729228335532) (0.11557788944723618,0.8938133245160104) (0.12562814070351755,0.8935820732899541) (0.13567839195979894,0.8930630458471318) (0.1457286432160804,0.893249685696942) (0.15577889447236185,0.8930063508131346) (0.16582914572864324,0.8933657324308883) (0.17587939698492464,0.8931078673679163) (0.18592964824120609,0.8930616870984285) (0.19597989949748745,0.8925125313480466) (0.20603015075376882,0.8923356967612648) (0.21608040201005022,0.8921408368140493) (0.22613065326633164,0.8921567020186736) (0.23618090452261306,0.892289537077371) (0.24623115577889443,0.8922267154611696) (0.2562814070351758,0.8922756161431835) (0.2663316582914573,0.8925347200860625) (0.2763819095477386,0.8924416548837513) (0.28643216080402006,0.8925604337760048) (0.2964824120603015,0.8928906225092209) (0.30653266331658285,0.8930442165722742) (0.31658291457286436,0.8927399908610394) (0.3266331658291457,0.8926092787841419) (0.33668341708542715,0.8925193408903102) (0.3467336683417085,0.8929506882680287) (0.35678391959798994,0.8928353972358583) (0.36683417085427134,0.8930550939685606) (0.37688442211055273,0.8933487091111638) (0.3869346733668341,0.8934230932074062) (0.3969849246231155,0.8932942968345574) (0.40703517587939697,0.8933431433978527) (0.4170854271356784,0.8934305535424041) (0.4271356783919598,0.8935803244519562) (0.4371859296482412,0.8935639450099139) (0.44723618090452255,0.8937625592753324) (0.457286432160804,0.8939061551577272) (0.4673366834170854,0.8942487470080194) (0.4773869346733668,0.8945566700912849) (0.48743718592964824,0.8947675056183629) (0.49748743718592964,0.8948296113472942) (0.507537688442211,0.8952051030807525) (0.5175879396984924,0.8951975531812771) (0.5276381909547738,0.8951450263138853) (0.5376884422110553,0.8952905890933716) (0.5477386934673366,0.8953610480149823) (0.5577889447236181,0.8953155479016458) (0.5678391959798995,0.895460961017937) (0.5778894472361809,0.8957295280995221) (0.5879396984924623,0.8960892323028482) (0.5979899497487438,0.8966561262704591) (0.6080402010050251,0.8965275143764583) (0.6180904522613065,0.8963337707287875) (0.6281407035175879,0.8961984096350466) (0.6381909547738693,0.8964629995448006) (0.6482412060301508,0.8961189304738487) (0.6582914572864321,0.8963523163836066) (0.6683417085427136,0.8966027035805709) (0.678391959798995,0.8968333254030643) (0.6884422110552764,0.8971689654895529) (0.6984924623115577,0.897237292733803) (0.7085427135678392,0.8970611425635753) (0.7185929648241206,0.897338173907282) (0.728643216080402,0.89745771540059) (0.7386934673366834,0.8975722097557647) (0.7487437185929647,0.8976958840589412) (0.7587939698492462,0.8983320493983404) (0.7688442211055275,0.898613049174651) (0.778894472361809,0.8987945887927152) (0.7889447236180904,0.898735046521169) (0.7989949748743719,0.8986622662789945) (0.8090452261306532,0.8986307185201877) (0.8190954773869347,0.898701706093181) (0.829145728643216,0.8985025252217913) (0.8391959798994975,0.8984934614584849) (0.8492462311557788,0.8984638863537227) (0.8592964824120602,0.8990359940255735) (0.8693467336683417,0.8991845064854885) (0.8793969849246231,0.8995116292950237) (0.8894472361809045,0.8997346784223785) (0.8994974874371858,0.8996741156138578) (0.9095477386934673,0.8996687807644134) (0.9195979899497487,0.8996805334603046) (0.9296482412060301,0.8999542377379522) (0.9396984924623115,0.8999568295942221) (0.949748743718593,0.9002869345839032) (0.9597989949748744,0.9004182444439445) (0.9698492462311558,0.9004212057412659) (0.9798994974874371,0.9010113806059609) (0.9899497487437185,0.9009996369171589) (0.9999999999999999,0.9008498576556805) };
  \addplot+[ours_ig_robust, mark=none, line width=0.4pt, dashed, forget plot] coordinates {(-1.0,0.7895182280060294) (1.0,0.7895182280060294)};
  \addplot+[ours_ig_robust, mark=*, only marks, mark size=\PanelArgminSize, forget plot] coordinates {(0.21608040201005022,0.8921408368140493)};
  \addplot[name path=U_ig_mle_grp, draw=none, mark=none, forget plot] coordinates {(-1.0,0.9195376900982916) (-0.9899497487437184,0.9194338874410998) (-0.9798994974874372,0.9187057254572446) (-0.9698492462311558,0.9188430966755113) (-0.9597989949748743,0.9192196320295414) (-0.9497487437185931,0.9192282513372301) (-0.9396984924623115,0.9194246485187453) (-0.9296482412060301,0.9192663084363227) (-0.9195979899497488,0.9193759016116636) (-0.9095477386934673,0.9191799261344434) (-0.8994974874371859,0.9189889168933297) (-0.8894472361809044,0.9188314793731605) (-0.8793969849246231,0.9187049934664516) (-0.8693467336683416,0.9187169277151906) (-0.8592964824120601,0.9186004305394807) (-0.8492462311557789,0.9185475680417519) (-0.8391959798994975,0.9180991249922081) (-0.8291457286432159,0.9177189855611712) (-0.8190954773869347,0.916816171345438) (-0.8090452261306532,0.9164523021851174) (-0.7989949748743718,0.9160514105458137) (-0.7889447236180904,0.9160424953352291) (-0.7788944723618091,0.9160322025770833) (-0.7688442211055275,0.9159473420217021) (-0.7587939698492463,0.9158322352716772) (-0.7487437185929648,0.9158227570164212) (-0.7386934673366834,0.915513449705655) (-0.728643216080402,0.9157037011261774) (-0.7185929648241206,0.9156542364783264) (-0.7085427135678392,0.915414533209255) (-0.6984924623115576,0.9145730268450776) (-0.6884422110552765,0.9143360894644814) (-0.6783919597989949,0.9140236255288717) (-0.6683417085427135,0.91413952219923) (-0.6582914572864322,0.9137464709113358) (-0.6482412060301508,0.9141130379756182) (-0.6381909547738692,0.9141119128669845) (-0.628140703517588,0.9139107015989281) (-0.6180904522613065,0.9140805564985958) (-0.6080402010050251,0.9138865998114243) (-0.5979899497487438,0.9141170003147654) (-0.5879396984924623,0.9144924335722256) (-0.5778894472361809,0.9144650167819505) (-0.5678391959798996,0.9144576691806292) (-0.5577889447236181,0.914205003160632) (-0.5477386934673366,0.9138753942683002) (-0.5376884422110553,0.9142247569928513) (-0.5276381909547738,0.9135358392651509) (-0.5175879396984925,0.9133478759005651) (-0.507537688442211,0.9130612833646428) (-0.49748743718592964,0.9132504515698213) (-0.48743718592964824,0.9136666938607052) (-0.47738693467336685,0.9135713787793035) (-0.4673366834170854,0.9129428140324148) (-0.457286432160804,0.9127600170358073) (-0.4472361809045226,0.9124782781765495) (-0.43718592964824116,0.9122554127464618) (-0.4271356783919598,0.9121056214388775) (-0.4170854271356783,0.911945766455809) (-0.40703517587939697,0.9117038717195007) (-0.3969849246231156,0.911623922319714) (-0.38693467336683407,0.9111972156183723) (-0.37688442211055273,0.9110457618304582) (-0.3668341708542714,0.9104866726066467) (-0.35678391959798994,0.9105212098193348) (-0.3467336683417085,0.9104359168964635) (-0.33668341708542704,0.9100513807320468) (-0.3266331658291457,0.9105471736118677) (-0.31658291457286436,0.9103664438281898) (-0.30653266331658285,0.9095635974628555) (-0.29648241206030146,0.9099814255518142) (-0.2864321608040202,0.9096528061013669) (-0.2763819095477386,0.9094030885663531) (-0.26633165829145733,0.9093911679283693) (-0.2562814070351758,0.9088606814946819) (-0.24623115577889443,0.9082346550204496) (-0.2361809045226131,0.9083325223412749) (-0.22613065326633164,0.9082010134963211) (-0.2160804020100502,0.9085239003209331) (-0.20603015075376882,0.9082506306389345) (-0.19597989949748737,0.9081245190461631) (-0.18592964824120603,0.907788682294105) (-0.1758793969849246,0.9077927363677862) (-0.16582914572864324,0.9077757772822802) (-0.15577889447236182,0.9074694598066206) (-0.14572864321608042,0.9071472701445906) (-0.135678391959799,0.9070340334039314) (-0.12562814070351758,0.9068879940683122) (-0.11557788944723618,0.9067844125508543) (-0.10552763819095479,0.9069172079360441) (-0.0954773869346734,0.90675428515504) (-0.0854271356783919,0.90691490120215) (-0.0753768844221106,0.9067791503004742) (-0.06532663316582912,0.9067201974871045) (-0.055276381909547756,0.9062338658161422) (-0.04522613065326632,0.9058925099630863) (-0.03517587939698496,0.9056482465748746) (-0.025125628140703526,0.9054311057756539) (-0.015075376884422112,0.9053644625181301) (-0.005025125628140694,0.9053834564313616) (0.00502512562814067,0.9051841613537701) (0.015075376884422073,0.9054052510691576) (0.02512562814070352,0.9057143032821799) (0.035175879396984966,0.905646698876416) (0.045226130653266264,0.9057988398620266) (0.0552763819095477,0.905959943783178) (0.06532663316582915,0.9057449366712963) (0.07537688442211059,0.9059042063177957) (0.08542713567839191,0.9064250391730659) (0.09547738693467332,0.9066105094383563) (0.10552763819095481,0.9064219994770911) (0.11557788944723618,0.9065393672908213) (0.12562814070351755,0.9065316649393403) (0.13567839195979894,0.9064342335074943) (0.1457286432160804,0.9060762530368369) (0.15577889447236185,0.9065094863700601) (0.16582914572864324,0.906493790379711) (0.17587939698492464,0.907243339420323) (0.18592964824120609,0.9074344154019375) (0.19597989949748745,0.9077929854863384) (0.20603015075376882,0.9075919990282834) (0.21608040201005022,0.9080081273545147) (0.22613065326633164,0.9079756034227996) (0.23618090452261306,0.9084141491502751) (0.24623115577889443,0.9082954284715414) (0.2562814070351758,0.9086188665882499) (0.2663316582914573,0.9086921744789472) (0.2763819095477386,0.9086307839023167) (0.28643216080402006,0.9082621861569977) (0.2964824120603015,0.9083986105686159) (0.30653266331658285,0.9082667642434328) (0.31658291457286436,0.9081614541932941) (0.3266331658291457,0.9078052187349932) (0.33668341708542715,0.9079787438774967) (0.3467336683417085,0.9084512395287935) (0.35678391959798994,0.908395835869238) (0.36683417085427134,0.9086212390282842) (0.37688442211055273,0.9088821514008065) (0.3869346733668341,0.9085846355495524) (0.3969849246231155,0.9085712924321301) (0.40703517587939697,0.9088237615076518) (0.4170854271356784,0.9086783919604339) (0.4271356783919598,0.9085708140241475) (0.4371859296482412,0.9088802309932442) (0.44723618090452255,0.9085228745260412) (0.457286432160804,0.9087076558135392) (0.4673366834170854,0.908790675949754) (0.4773869346733668,0.9092029951065219) (0.48743718592964824,0.9091951048426138) (0.49748743718592964,0.9091646303641372) (0.507537688442211,0.9088617809044874) (0.5175879396984924,0.9092165747691832) (0.5276381909547738,0.9094961721182648) (0.5376884422110553,0.9093448763914767) (0.5477386934673366,0.9097361821500186) (0.5577889447236181,0.9097851915741008) (0.5678391959798995,0.9097522916124454) (0.5778894472361809,0.9101629278039722) (0.5879396984924623,0.9100120881097026) (0.5979899497487438,0.9102540417660538) (0.6080402010050251,0.9106111752961162) (0.6180904522613065,0.9108141538272303) (0.6281407035175879,0.9106689100363665) (0.6381909547738693,0.9107321835526799) (0.6482412060301508,0.9110751663160146) (0.6582914572864321,0.9106421865614679) (0.6683417085427136,0.9102874217101424) (0.678391959798995,0.9106298492184531) (0.6884422110552764,0.9109747363310273) (0.6984924623115577,0.9108916680516446) (0.7085427135678392,0.9109001325363597) (0.7185929648241206,0.9107914385050571) (0.728643216080402,0.9107919456983157) (0.7386934673366834,0.9107410589921534) (0.7487437185929647,0.9106402193456528) (0.7587939698492462,0.9105371976253214) (0.7688442211055275,0.9108502990777363) (0.778894472361809,0.9110004896144356) (0.7889447236180904,0.9112622824308599) (0.7989949748743719,0.9111592618544452) (0.8090452261306532,0.9112354247670963) (0.8190954773869347,0.9109075842824329) (0.829145728643216,0.9107376259221281) (0.8391959798994975,0.9106940813481723) (0.8492462311557788,0.9105649800127068) (0.8592964824120602,0.9109502727874319) (0.8693467336683417,0.9110341133089165) (0.8793969849246231,0.9113306712720286) (0.8894472361809045,0.9119469955002655) (0.8994974874371858,0.9119113406569586) (0.9095477386934673,0.9122183767856201) (0.9195979899497487,0.9124318229912406) (0.9296482412060301,0.9123096254830744) (0.9396984924623115,0.9125152024974043) (0.949748743718593,0.9124737114171746) (0.9597989949748744,0.9132554979981227) (0.9698492462311558,0.9132749790612931) (0.9798994974874371,0.9134861718183612) (0.9899497487437185,0.9140842209025899) (0.9999999999999999,0.9143620161820948) };
  \addplot[name path=L_ig_mle_grp, draw=none, mark=none, forget plot] coordinates {(-1.0,0.8923580416026781) (-0.9899497487437184,0.8922104309826558) (-0.9798994974874372,0.8916437525095595) (-0.9698492462311558,0.8917188281036118) (-0.9597989949748743,0.8919037267240755) (-0.9497487437185931,0.8918743585134932) (-0.9396984924623115,0.8920023176145856) (-0.9296482412060301,0.8916259959286618) (-0.9195979899497488,0.8916585643373909) (-0.9095477386934673,0.8913743928189656) (-0.8994974874371859,0.8914500055316289) (-0.8894472361809044,0.8912310640886241) (-0.8793969849246231,0.891066424375378) (-0.8693467336683416,0.8911600585454601) (-0.8592964824120601,0.8909880418907657) (-0.8492462311557789,0.8909763978731562) (-0.8391959798994975,0.8905718317598819) (-0.8291457286432159,0.8900739374021374) (-0.8190954773869347,0.8891893375363701) (-0.8090452261306532,0.8888636551907307) (-0.7989949748743718,0.8887854386424365) (-0.7889447236180904,0.8889740134835646) (-0.7788944723618091,0.8888405299562898) (-0.7688442211055275,0.8885647432845947) (-0.7587939698492463,0.8884415160553059) (-0.7487437185929648,0.8881460214284338) (-0.7386934673366834,0.8879816094842916) (-0.728643216080402,0.8882551368901833) (-0.7185929648241206,0.8880145349206172) (-0.7085427135678392,0.8876869485972377) (-0.6984924623115576,0.8868463124095967) (-0.6884422110552765,0.886398405204555) (-0.6783919597989949,0.8860700988087499) (-0.6683417085427135,0.8862327402706546) (-0.6582914572864322,0.8858142642410999) (-0.6482412060301508,0.8861117393335263) (-0.6381909547738692,0.8862174159840114) (-0.628140703517588,0.8859321537463256) (-0.6180904522613065,0.8861854272167894) (-0.6080402010050251,0.8859164857698816) (-0.5979899497487438,0.8859548380554576) (-0.5879396984924623,0.8863483156815491) (-0.5778894472361809,0.8865060712193064) (-0.5678391959798996,0.8864383227926447) (-0.5577889447236181,0.8860704578957818) (-0.5477386934673366,0.8855722396890395) (-0.5376884422110553,0.885846870353645) (-0.5276381909547738,0.8852315967336395) (-0.5175879396984925,0.8850188161031993) (-0.507537688442211,0.8846685267928822) (-0.49748743718592964,0.8851511481472208) (-0.48743718592964824,0.8855632833632368) (-0.47738693467336685,0.8853589525030973) (-0.4673366834170854,0.8847040193453043) (-0.457286432160804,0.884606811873797) (-0.4472361809045226,0.884670288191989) (-0.43718592964824116,0.8843708852719766) (-0.4271356783919598,0.8843543765808418) (-0.4170854271356783,0.8840439765815915) (-0.40703517587939697,0.8837050369534996) (-0.3969849246231156,0.883624731360831) (-0.38693467336683407,0.8829858546174674) (-0.37688442211055273,0.88300912421638) (-0.3668341708542714,0.8824224670662628) (-0.35678391959798994,0.8826009601136308) (-0.3467336683417085,0.8826755572068348) (-0.33668341708542704,0.8822665505198702) (-0.3266331658291457,0.8829661426246284) (-0.31658291457286436,0.8826524251231667) (-0.30653266331658285,0.881811048546252) (-0.29648241206030146,0.8822478757420655) (-0.2864321608040202,0.8818321345075255) (-0.2763819095477386,0.8815104832559626) (-0.26633165829145733,0.8816415545106739) (-0.2562814070351758,0.8812204096376562) (-0.24623115577889443,0.8806367567692883) (-0.2361809045226131,0.880773407262062) (-0.22613065326633164,0.8806589699360733) (-0.2160804020100502,0.8809825388714183) (-0.20603015075376882,0.8807914987120654) (-0.19597989949748737,0.8805369849260241) (-0.18592964824120603,0.8802167025189946) (-0.1758793969849246,0.8802140673863972) (-0.16582914572864324,0.8804137469837775) (-0.15577889447236182,0.8802881438386787) (-0.14572864321608042,0.8799885554322278) (-0.135678391959799,0.879864936333577) (-0.12562814070351758,0.8797523468496781) (-0.11557788944723618,0.8797041245015148) (-0.10552763819095479,0.8798378000690723) (-0.0954773869346734,0.8798408292448837) (-0.0854271356783919,0.8800779575513411) (-0.0753768844221106,0.8799210597966748) (-0.06532663316582912,0.8795435369804453) (-0.055276381909547756,0.8791951742238997) (-0.04522613065326632,0.8791028443481228) (-0.03517587939698496,0.878750332374267) (-0.025125628140703526,0.8784840196393677) (-0.015075376884422112,0.8784904323439775) (-0.005025125628140694,0.8785336384814131) (0.00502512562814067,0.8783784510796138) (0.015075376884422073,0.8786269713685441) (0.02512562814070352,0.8788649221452345) (0.035175879396984966,0.8788100478201428) (0.045226130653266264,0.8788583940351938) (0.0552763819095477,0.8789723834067911) (0.06532663316582915,0.8788984843624211) (0.07537688442211059,0.8789938632901865) (0.08542713567839191,0.8796208828195842) (0.09547738693467332,0.8797084411643241) (0.10552763819095481,0.8795458346226873) (0.11557788944723618,0.8796148538546775) (0.12562814070351755,0.8794652855450288) (0.13567839195979894,0.8792962180037812) (0.1457286432160804,0.8789734103662148) (0.15577889447236185,0.8793914264749559) (0.16582914572864324,0.8794034691455612) (0.17587939698492464,0.8800440882462899) (0.18592964824120609,0.8801525369955953) (0.19597989949748745,0.8805214645149518) (0.20603015075376882,0.8803427134024783) (0.21608040201005022,0.8806126942291251) (0.22613065326633164,0.8806514733313119) (0.23618090452261306,0.881046917315403) (0.24623115577889443,0.8809528615146247) (0.2562814070351758,0.8813087029849153) (0.2663316582914573,0.881334905189166) (0.2763819095477386,0.8811675453889742) (0.28643216080402006,0.8808251655678476) (0.2964824120603015,0.8810067331997868) (0.30653266331658285,0.8808725643301913) (0.31658291457286436,0.8809072734891172) (0.3266331658291457,0.8806468099115001) (0.33668341708542715,0.8808435479282228) (0.3467336683417085,0.8813371603522767) (0.35678391959798994,0.8812537206655849) (0.36683417085427134,0.8813005923530861) (0.37688442211055273,0.8816457064641197) (0.3869346733668341,0.8815283543596291) (0.3969849246231155,0.881553378631224) (0.40703517587939697,0.881813215565006) (0.4170854271356784,0.8817529167810246) (0.4271356783919598,0.8816345179131226) (0.4371859296482412,0.8817823708705161) (0.44723618090452255,0.8815519807885177) (0.457286432160804,0.881815264892996) (0.4673366834170854,0.8818757345218172) (0.4773869346733668,0.8820949930848312) (0.48743718592964824,0.8821087050109898) (0.49748743718592964,0.8820754230331698) (0.507537688442211,0.8817323473782341) (0.5175879396984924,0.882035123093424) (0.5276381909547738,0.8823969810619268) (0.5376884422110553,0.8823179667934491) (0.5477386934673366,0.8827026369271781) (0.5577889447236181,0.8827351707996743) (0.5678391959798995,0.8827572367751289) (0.5778894472361809,0.8831419536452148) (0.5879396984924623,0.883000621247552) (0.5979899497487438,0.8833201256033929) (0.6080402010050251,0.883610134271721) (0.6180904522613065,0.8836521596596614) (0.6281407035175879,0.8834854346502639) (0.6381909547738693,0.8835447398917856) (0.6482412060301508,0.8837291674366423) (0.6582914572864321,0.8833646311510074) (0.6683417085427136,0.8830803349902614) (0.678391959798995,0.883357525094658) (0.6884422110552764,0.8835827653883671) (0.6984924623115577,0.8834342872348611) (0.7085427135678392,0.8833729674787572) (0.7185929648241206,0.8831657313695338) (0.728643216080402,0.8831764505850533) (0.7386934673366834,0.8831931930789851) (0.7487437185929647,0.8829509586089638) (0.7587939698492462,0.8827932014940756) (0.7688442211055275,0.8831198212903548) (0.778894472361809,0.8833445750261365) (0.7889447236180904,0.8835855619770415) (0.7989949748743719,0.8834873361776058) (0.8090452261306532,0.8835391283920441) (0.8190954773869347,0.8832454545612979) (0.829145728643216,0.8832223700199625) (0.8391959798994975,0.8831664090470615) (0.8492462311557788,0.8830255343552383) (0.8592964824120602,0.8832183243253195) (0.8693467336683417,0.8832642909908036) (0.8793969849246231,0.8834332583379856) (0.8894472361809045,0.8839815719573445) (0.8994974874371858,0.8840553479883408) (0.9095477386934673,0.8842857893395627) (0.9195979899497487,0.8843282237312869) (0.9296482412060301,0.884319196845129) (0.9396984924623115,0.8844966197076706) (0.949748743718593,0.8844116484604021) (0.9597989949748744,0.8850848046278386) (0.9698492462311558,0.8850417333789322) (0.9798994974874371,0.885203788775128) (0.9899497487437185,0.8857529951793939) (0.9999999999999999,0.8860074764729372) };
  \addplot[ours_ig_mle, fill, fill opacity=0.18, draw=none, mark=none, forget plot] fill between [of=U_ig_mle_grp and L_ig_mle_grp];
  \addplot+[ours_ig_mle, mark=none] coordinates {(-1.0,0.9059478658504848) (-0.9899497487437184,0.9058221592118778) (-0.9798994974874372,0.9051747389834021) (-0.9698492462311558,0.9052809623895616) (-0.9597989949748743,0.9055616793768084) (-0.9497487437185931,0.9055513049253616) (-0.9396984924623115,0.9057134830666654) (-0.9296482412060301,0.9054461521824922) (-0.9195979899497488,0.9055172329745272) (-0.9095477386934673,0.9052771594767045) (-0.8994974874371859,0.9052194612124793) (-0.8894472361809044,0.9050312717308923) (-0.8793969849246231,0.9048857089209148) (-0.8693467336683416,0.9049384931303254) (-0.8592964824120601,0.9047942362151232) (-0.8492462311557789,0.904761982957454) (-0.8391959798994975,0.904335478376045) (-0.8291457286432159,0.9038964614816543) (-0.8190954773869347,0.903002754440904) (-0.8090452261306532,0.9026579786879241) (-0.7989949748743718,0.9024184245941251) (-0.7889447236180904,0.9025082544093969) (-0.7788944723618091,0.9024363662666866) (-0.7688442211055275,0.9022560426531484) (-0.7587939698492463,0.9021368756634915) (-0.7487437185929648,0.9019843892224275) (-0.7386934673366834,0.9017475295949733) (-0.728643216080402,0.9019794190081803) (-0.7185929648241206,0.9018343856994718) (-0.7085427135678392,0.9015507409032464) (-0.6984924623115576,0.9007096696273371) (-0.6884422110552765,0.9003672473345182) (-0.6783919597989949,0.9000468621688108) (-0.6683417085427135,0.9001861312349423) (-0.6582914572864322,0.8997803675762178) (-0.6482412060301508,0.9001123886545722) (-0.6381909547738692,0.9001646644254979) (-0.628140703517588,0.8999214276726268) (-0.6180904522613065,0.9001329918576926) (-0.6080402010050251,0.8999015427906529) (-0.5979899497487438,0.9000359191851115) (-0.5879396984924623,0.9004203746268874) (-0.5778894472361809,0.9004855440006284) (-0.5678391959798996,0.9004479959866369) (-0.5577889447236181,0.9001377305282069) (-0.5477386934673366,0.8997238169786699) (-0.5376884422110553,0.9000358136732481) (-0.5276381909547738,0.8993837179993952) (-0.5175879396984925,0.8991833460018822) (-0.507537688442211,0.8988649050787625) (-0.49748743718592964,0.8992007998585211) (-0.48743718592964824,0.899614988611971) (-0.47738693467336685,0.8994651656412004) (-0.4673366834170854,0.8988234166888596) (-0.457286432160804,0.8986834144548022) (-0.4472361809045226,0.8985742831842692) (-0.43718592964824116,0.8983131490092192) (-0.4271356783919598,0.8982299990098597) (-0.4170854271356783,0.8979948715187003) (-0.40703517587939697,0.8977044543365001) (-0.3969849246231156,0.8976243268402725) (-0.38693467336683407,0.8970915351179198) (-0.37688442211055273,0.8970274430234191) (-0.3668341708542714,0.8964545698364548) (-0.35678391959798994,0.8965610849664828) (-0.3467336683417085,0.8965557370516491) (-0.33668341708542704,0.8961589656259585) (-0.3266331658291457,0.896756658118248) (-0.31658291457286436,0.8965094344756782) (-0.30653266331658285,0.8956873230045538) (-0.29648241206030146,0.8961146506469398) (-0.2864321608040202,0.8957424703044462) (-0.2763819095477386,0.8954567859111578) (-0.26633165829145733,0.8955163612195216) (-0.2562814070351758,0.8950405455661691) (-0.24623115577889443,0.894435705894869) (-0.2361809045226131,0.8945529648016685) (-0.22613065326633164,0.8944299917161972) (-0.2160804020100502,0.8947532195961757) (-0.20603015075376882,0.8945210646755) (-0.19597989949748737,0.8943307519860936) (-0.18592964824120603,0.8940026924065498) (-0.1758793969849246,0.8940034018770917) (-0.16582914572864324,0.8940947621330289) (-0.15577889447236182,0.8938788018226497) (-0.14572864321608042,0.8935679127884092) (-0.135678391959799,0.8934494848687542) (-0.12562814070351758,0.8933201704589951) (-0.11557788944723618,0.8932442685261845) (-0.10552763819095479,0.8933775040025582) (-0.0954773869346734,0.8932975571999618) (-0.0854271356783919,0.8934964293767456) (-0.0753768844221106,0.8933501050485745) (-0.06532663316582912,0.8931318672337749) (-0.055276381909547756,0.8927145200200209) (-0.04522613065326632,0.8924976771556046) (-0.03517587939698496,0.8921992894745708) (-0.025125628140703526,0.8919575627075108) (-0.015075376884422112,0.8919274474310538) (-0.005025125628140694,0.8919585474563874) (0.00502512562814067,0.8917813062166919) (0.015075376884422073,0.8920161112188508) (0.02512562814070352,0.8922896127137072) (0.035175879396984966,0.8922283733482794) (0.045226130653266264,0.8923286169486102) (0.0552763819095477,0.8924661635949845) (0.06532663316582915,0.8923217105168587) (0.07537688442211059,0.8924490348039911) (0.08542713567839191,0.893022960996325) (0.09547738693467332,0.8931594753013402) (0.10552763819095481,0.8929839170498892) (0.11557788944723618,0.8930771105727494) (0.12562814070351755,0.8929984752421846) (0.13567839195979894,0.8928652257556378) (0.1457286432160804,0.8925248317015259) (0.15577889447236185,0.892950456422508) (0.16582914572864324,0.8929486297626361) (0.17587939698492464,0.8936437138333064) (0.18592964824120609,0.8937934761987664) (0.19597989949748745,0.8941572250006451) (0.20603015075376882,0.8939673562153808) (0.21608040201005022,0.8943104107918199) (0.22613065326633164,0.8943135383770557) (0.23618090452261306,0.894730533232839) (0.24623115577889443,0.8946241449930831) (0.2562814070351758,0.8949637847865826) (0.2663316582914573,0.8950135398340566) (0.2763819095477386,0.8948991646456455) (0.28643216080402006,0.8945436758624227) (0.2964824120603015,0.8947026718842014) (0.30653266331658285,0.8945696642868121) (0.31658291457286436,0.8945343638412057) (0.3266331658291457,0.8942260143232467) (0.33668341708542715,0.8944111459028597) (0.3467336683417085,0.8948941999405351) (0.35678391959798994,0.8948247782674115) (0.36683417085427134,0.8949609156906851) (0.37688442211055273,0.8952639289324631) (0.3869346733668341,0.8950564949545907) (0.3969849246231155,0.895062335531677) (0.40703517587939697,0.8953184885363289) (0.4170854271356784,0.8952156543707293) (0.4271356783919598,0.8951026659686351) (0.4371859296482412,0.8953313009318802) (0.44723618090452255,0.8950374276572794) (0.457286432160804,0.8952614603532676) (0.4673366834170854,0.8953332052357856) (0.4773869346733668,0.8956489940956766) (0.48743718592964824,0.8956519049268018) (0.49748743718592964,0.8956200266986535) (0.507537688442211,0.8952970641413608) (0.5175879396984924,0.8956258489313036) (0.5276381909547738,0.8959465765900958) (0.5376884422110553,0.8958314215924629) (0.5477386934673366,0.8962194095385984) (0.5577889447236181,0.8962601811868876) (0.5678391959798995,0.8962547641937871) (0.5778894472361809,0.8966524407245935) (0.5879396984924623,0.8965063546786273) (0.5979899497487438,0.8967870836847234) (0.6080402010050251,0.8971106547839186) (0.6180904522613065,0.8972331567434458) (0.6281407035175879,0.8970771723433152) (0.6381909547738693,0.8971384617222328) (0.6482412060301508,0.8974021668763285) (0.6582914572864321,0.8970034088562376) (0.6683417085427136,0.8966838783502019) (0.678391959798995,0.8969936871565556) (0.6884422110552764,0.8972787508596972) (0.6984924623115577,0.8971629776432528) (0.7085427135678392,0.8971365500075584) (0.7185929648241206,0.8969785849372954) (0.728643216080402,0.8969841981416845) (0.7386934673366834,0.8969671260355693) (0.7487437185929647,0.8967955889773083) (0.7587939698492462,0.8966651995596985) (0.7688442211055275,0.8969850601840456) (0.778894472361809,0.897172532320286) (0.7889447236180904,0.8974239222039507) (0.7989949748743719,0.8973232990160255) (0.8090452261306532,0.8973872765795702) (0.8190954773869347,0.8970765194218654) (0.829145728643216,0.8969799979710453) (0.8391959798994975,0.8969302451976169) (0.8492462311557788,0.8967952571839726) (0.8592964824120602,0.8970842985563757) (0.8693467336683417,0.89714920214986) (0.8793969849246231,0.8973819648050071) (0.8894472361809045,0.897964283728805) (0.8994974874371858,0.8979833443226497) (0.9095477386934673,0.8982520830625914) (0.9195979899497487,0.8983800233612638) (0.9296482412060301,0.8983144111641017) (0.9396984924623115,0.8985059111025374) (0.949748743718593,0.8984426799387883) (0.9597989949748744,0.8991701513129806) (0.9698492462311558,0.8991583562201126) (0.9798994974874371,0.8993449802967446) (0.9899497487437185,0.8999186080409919) (0.9999999999999999,0.900184746327516) };
  \addplot+[ours_ig_mle, mark=none, line width=0.4pt, dashed, forget plot] coordinates {(-1.0,0.7963180085564017) (1.0,0.7963180085564017)};
  \addplot+[ours_ig_mle, mark=*, only marks, mark size=\PanelArgminSize, forget plot] coordinates {(0.00502512562814067,0.8917813062166919)};
  \addplot+[ours_oracle, mark=none, forget plot] coordinates {(-1.0,0.31670797994473593) (1.0,0.31670797994473593)};
%
			\renewcommand\nextgroupplot[1][]{\origNGP[#1, ymin=0.10, ymax=1.5,
					ytick={0.25,0.5,0.75,1.0,1.25,1.5}]}%
			% panel_profile_group.tex — \nextgroupplot fragment.
% Host preamble MUST load: \usepgfplotslibrary{groupplots},
%                          \usepgfplotslibrary{fillbetween}.
% Must be \input inside a \begin{groupplot}[...] with these
% options set at the GROUPPLOT level (not per-cell):
%   panel_axis, xmin=-1, xmax=1, xtick={\PanelXTicks},
%   xlabel={\PanelXLabel}, ylabel={\PanelYLabel},
%   ylabels at=edge left, yticklabels at=edge left,
%   group style={group size=N by 1, horizontal sep=...}
% That way pgfplots' native `ylabels at=edge left` shows the y-
% annotations only on the leftmost cell.  Per-cell options here
% are LIMITED to `title=` so the row's wrapper retains all axis
% control.  See `profile_row_template` for the canonical wrapper.

\nextgroupplot[title={ETH3D / SuperPoint+LightGlue / \(\bF\)}]
  \addplot[name path=U_ransac_grp, draw=none, mark=none, forget plot] coordinates {(-1.0,0.8232370098819639) (-0.9899497487437187,0.8232370098819639) (-0.9798994974874371,0.8232370098819639) (-0.9698492462311558,0.8232370098819639) (-0.9597989949748745,0.8232370098819639) (-0.9497487437185929,0.8232370098819639) (-0.9396984924623116,0.8232370098819639) (-0.9296482412060303,0.8089087135344567) (-0.9195979899497487,0.8089087135344567) (-0.9095477386934674,0.8089087135344567) (-0.899497487437186,0.8089087135344567) (-0.8894472361809044,0.8089087135344567) (-0.8793969849246231,0.8089087135344567) (-0.8693467336683418,0.8089087135344567) (-0.8592964824120601,0.781496574756295) (-0.8492462311557789,0.781496574756295) (-0.8391959798994972,0.781496574756295) (-0.829145728643216,0.781496574756295) (-0.8190954773869348,0.781496574756295) (-0.8090452261306531,0.781496574756295) (-0.7989949748743718,0.7325962309929377) (-0.7889447236180905,0.7325962309929377) (-0.7788944723618089,0.7325962309929377) (-0.7688442211055275,0.7325962309929377) (-0.7587939698492463,0.7325962309929377) (-0.7487437185929647,0.7325962309929377) (-0.7386934673366835,0.7461146741898321) (-0.7286432160804021,0.7461146741898321) (-0.7185929648241205,0.7461146741898321) (-0.7085427135678393,0.7461146741898321) (-0.6984924623115579,0.7254042734788345) (-0.6884422110552763,0.7254042734788345) (-0.678391959798995,0.7254042734788345) (-0.6683417085427137,0.7254042734788345) (-0.658291457286432,0.7254042734788345) (-0.6482412060301508,0.7105199308051688) (-0.6381909547738696,0.7105199308051688) (-0.6281407035175878,0.7105199308051688) (-0.6180904522613065,0.7105199308051688) (-0.6080402010050253,0.6828567706918001) (-0.5979899497487436,0.6828567706918001) (-0.5879396984924623,0.6828567706918001) (-0.5778894472361811,0.667746784935975) (-0.5678391959798994,0.667746784935975) (-0.5577889447236181,0.667746784935975) (-0.5477386934673368,0.667746784935975) (-0.5376884422110553,0.6512004312297981) (-0.5276381909547739,0.6512004312297981) (-0.5175879396984923,0.6512004312297981) (-0.507537688442211,0.6355168351082783) (-0.4974874371859297,0.6355168351082783) (-0.4874371859296481,0.6340518073688401) (-0.47738693467336674,0.6340518073688401) (-0.4673366834170855,0.6340518073688401) (-0.45728643216080395,0.6332778461261583) (-0.44723618090452255,0.6332778461261583) (-0.4371859296482413,0.6332778461261583) (-0.42713567839195965,0.6273026741296078) (-0.4170854271356784,0.6273026741296078) (-0.40703517587939714,0.6183767757446552) (-0.39698492462311546,0.6183767757446552) (-0.3869346733668342,0.6220311257679798) (-0.3768844221105529,0.6220311257679798) (-0.3668341708542712,0.6220311257679798) (-0.35678391959798994,0.6111042170630894) (-0.34673366834170866,0.6111042170630894) (-0.33668341708542704,0.6014864404957818) (-0.32663316582914576,0.6014864404957818) (-0.3165829145728644,0.601154380055979) (-0.30653266331658285,0.595245610879787) (-0.29648241206030157,0.595245610879787) (-0.2864321608040203,0.5921607186008478) (-0.27638190954773856,0.5921607186008478) (-0.26633165829145733,0.5910727846673867) (-0.2562814070351761,0.5910727846673867) (-0.24623115577889443,0.5974501730274944) (-0.23618090452261317,0.5948329018381996) (-0.2261306532663315,0.5948329018381996) (-0.2160804020100502,0.5719479793903128) (-0.206030150753769,0.5786632343429731) (-0.19597989949748723,0.5786632343429731) (-0.18592964824120595,0.5774194182307253) (-0.1758793969849247,0.5745921642668111) (-0.1658291457286431,0.578483364500349) (-0.15577889447236182,0.578483364500349) (-0.14572864321608053,0.5698184999023949) (-0.13567839195979886,0.5790685062047065) (-0.12562814070351758,0.577590699823595) (-0.1155778894472363,0.577590699823595) (-0.10552763819095466,0.5803263938928722) (-0.0954773869346734,0.5747667589748503) (-0.08542713567839208,0.5672776670457363) (-0.07537688442211042,0.5673212408067351) (-0.06532663316582912,0.5710562839850577) (-0.05527638190954787,0.5691601189796762) (-0.045226130653266215,0.572944278969549) (-0.0351758793969849,0.5629919396613396) (-0.02512562814070368,0.5797256855726206) (-0.015075376884422013,0.5763694085390726) (-0.005025125628140743,0.5742071246824061) (0.005025125628140574,0.574597326380043) (0.015075376884422167,0.583438845812802) (0.02512562814070352,0.5821217467949972) (0.03517587939698479,0.5808045488138783) (0.04522613065326644,0.5873149193115578) (0.0552763819095477,0.5906223654711972) (0.06532663316582898,0.6040084750988178) (0.07537688442211059,0.6130687839200816) (0.08542713567839183,0.6205590597621761) (0.09547738693467356,0.6128629590129266) (0.10552763819095481,0.6291401610861387) (0.11557788944723604,0.6259697764062196) (0.12562814070351774,0.6151074631777215) (0.13567839195979903,0.6212914962147776) (0.1457286432160803,0.6275422026774082) (0.15577889447236196,0.6278503040237926) (0.16582914572864324,0.652645246633279) (0.17587939698492452,0.6530028956226311) (0.1859296482412062,0.6651267747926087) (0.19597989949748745,0.6762428196391912) (0.2060301507537687,0.6742989256439234) (0.2160804020100504,0.6750597791820351) (0.22613065326633164,0.6834237610480864) (0.2361809045226129,0.6801676847360789) (0.2462311557788946,0.6938500486444379) (0.2562814070351759,0.6823314264607439) (0.26633165829145716,0.6853176681328269) (0.2763819095477388,0.6901757367440013) (0.28643216080402006,0.692427255925996) (0.29648241206030135,0.7144345735453914) (0.30653266331658297,0.7031971569990428) (0.31658291457286425,0.7212795877588151) (0.3266331658291456,0.7390041171125254) (0.33668341708542715,0.7477405128509464) (0.3467336683417085,0.7646481927900806) (0.3567839195979898,0.7658342837305361) (0.36683417085427145,0.7954812334391297) (0.37688442211055273,0.8110984221636725) (0.38693467336683435,0.8297414491688556) (0.3969849246231157,0.8322949649596052) (0.40703517587939697,0.8493521044078345) (0.41708542713567853,0.8658685934067871) (0.42713567839195987,0.8761482559293255) (0.4371859296482411,0.8911730341717176) (0.44723618090452283,0.9043984123241192) (0.45728643216080406,0.917863423026684) (0.46733668341708534,0.9371319461024684) (0.47738693467336696,0.96594048867562) (0.4874371859296483,0.9608461826438419) (0.4974874371859295,0.9905067816709078) (0.5075376884422113,1.0133389009205542) (0.5175879396984925,1.0208123164889271) (0.5276381909547737,1.0423912929226569) (0.5376884422110554,1.070170164344422) (0.5477386934673367,1.1014488965326803) (0.557788944723618,1.140468395797587) (0.5678391959798996,1.1593006859639927) (0.5778894472361809,1.1807591664029509) (0.5879396984924622,1.2028525530558458) (0.5979899497487439,1.2273721309258134) (0.6080402010050251,1.2419018175569376) (0.6180904522613063,1.2346246146525313) (0.628140703517588,1.2596755446606136) (0.6381909547738693,1.2961092319925434) (0.6482412060301505,1.333171972960459) (0.6582914572864322,1.3682723492724274) (0.6683417085427136,1.4016704601447336) (0.6783919597989948,1.423020946974164) (0.6884422110552765,1.4742009255246151) (0.6984924623115577,1.5544711337327108) (0.7085427135678392,1.5636689040306504) (0.7185929648241206,1.5871847726165214) (0.728643216080402,1.6035195420952015) (0.7386934673366834,1.6231158248164768) (0.7487437185929648,1.657552253438005) (0.7587939698492461,1.7061478246510167) (0.7688442211055275,1.7446217247263474) (0.7788944723618091,1.7441372348070214) (0.7889447236180905,1.8011064527137823) (0.7989949748743719,1.8615599199925001) (0.8090452261306533,1.9039355052205988) (0.8190954773869348,1.9364511621195746) (0.829145728643216,1.9910476884152484) (0.8391959798994976,2.0299088902986657) (0.849246231155779,2.058717402419585) (0.8592964824120602,2.08385951073976) (0.8693467336683417,2.088993589561906) (0.8793969849246231,2.1301840587671164) (0.8894472361809044,2.174976510137009) (0.8994974874371859,2.1854299276972595) (0.9095477386934674,2.216488852484115) (0.9195979899497486,2.2604517126101697) (0.9296482412060301,2.3103020192911594) (0.9396984924623116,2.360174755893755) (0.949748743718593,2.3856689432048737) (0.9597989949748743,2.4472173870918374) (0.9698492462311559,2.5495973419493656) (0.9798994974874373,2.6011699804551207) (0.9899497487437185,2.668566631313701) (1.0,2.7001898699792095) };
  \addplot[name path=L_ransac_grp, draw=none, mark=none, forget plot] coordinates {(-1.0,0.7424997368126725) (-0.9899497487437187,0.7424997368126725) (-0.9798994974874371,0.7424997368126725) (-0.9698492462311558,0.7424997368126725) (-0.9597989949748745,0.7424997368126725) (-0.9497487437185929,0.7424997368126725) (-0.9396984924623116,0.7424997368126725) (-0.9296482412060303,0.7085649895850149) (-0.9195979899497487,0.7085649895850149) (-0.9095477386934674,0.7085649895850149) (-0.899497487437186,0.7085649895850149) (-0.8894472361809044,0.7085649895850149) (-0.8793969849246231,0.7085649895850149) (-0.8693467336683418,0.7085649895850149) (-0.8592964824120601,0.6904342928532594) (-0.8492462311557789,0.6904342928532594) (-0.8391959798994972,0.6904342928532594) (-0.829145728643216,0.6904342928532594) (-0.8190954773869348,0.6904342928532594) (-0.8090452261306531,0.6904342928532594) (-0.7989949748743718,0.644378308416143) (-0.7889447236180905,0.644378308416143) (-0.7788944723618089,0.644378308416143) (-0.7688442211055275,0.644378308416143) (-0.7587939698492463,0.644378308416143) (-0.7487437185929647,0.644378308416143) (-0.7386934673366835,0.6466226108682245) (-0.7286432160804021,0.6466226108682245) (-0.7185929648241205,0.6466226108682245) (-0.7085427135678393,0.6466226108682245) (-0.6984924623115579,0.6360560025676946) (-0.6884422110552763,0.6360560025676946) (-0.678391959798995,0.6360560025676946) (-0.6683417085427137,0.6360560025676946) (-0.658291457286432,0.6360560025676946) (-0.6482412060301508,0.6156630173242672) (-0.6381909547738696,0.6156630173242672) (-0.6281407035175878,0.6156630173242672) (-0.6180904522613065,0.6156630173242672) (-0.6080402010050253,0.5952485035372663) (-0.5979899497487436,0.5952485035372663) (-0.5879396984924623,0.5952485035372663) (-0.5778894472361811,0.5883091720137227) (-0.5678391959798994,0.5883091720137227) (-0.5577889447236181,0.5883091720137227) (-0.5477386934673368,0.5883091720137227) (-0.5376884422110553,0.5808453665079896) (-0.5276381909547739,0.5808453665079896) (-0.5175879396984923,0.5808453665079896) (-0.507537688442211,0.5680476587854151) (-0.4974874371859297,0.5680476587854151) (-0.4874371859296481,0.5702680858466713) (-0.47738693467336674,0.5702680858466713) (-0.4673366834170855,0.5702680858466713) (-0.45728643216080395,0.5614879214948134) (-0.44723618090452255,0.5614879214948134) (-0.4371859296482413,0.5614879214948134) (-0.42713567839195965,0.5539932897259597) (-0.4170854271356784,0.5539932897259597) (-0.40703517587939714,0.5410521810167941) (-0.39698492462311546,0.5410521810167941) (-0.3869346733668342,0.5416965576142507) (-0.3768844221105529,0.5416965576142507) (-0.3668341708542712,0.5416965576142507) (-0.35678391959798994,0.5356311168949939) (-0.34673366834170866,0.5356311168949939) (-0.33668341708542704,0.5280169655635365) (-0.32663316582914576,0.5280169655635365) (-0.3165829145728644,0.5254723462894665) (-0.30653266331658285,0.5184606258113693) (-0.29648241206030157,0.5184606258113693) (-0.2864321608040203,0.5197045546807136) (-0.27638190954773856,0.5197045546807136) (-0.26633165829145733,0.5272508402348404) (-0.2562814070351761,0.5272508402348404) (-0.24623115577889443,0.5270630327169641) (-0.23618090452261317,0.5267789553245376) (-0.2261306532663315,0.5267789553245376) (-0.2160804020100502,0.5116052621741138) (-0.206030150753769,0.5174011229651675) (-0.19597989949748723,0.5174011229651675) (-0.18592964824120595,0.5114986248143276) (-0.1758793969849247,0.5046286235239735) (-0.1658291457286431,0.50748280439259) (-0.15577889447236182,0.50748280439259) (-0.14572864321608053,0.49833366003760277) (-0.13567839195979886,0.5069306956656803) (-0.12562814070351758,0.5087885855590714) (-0.1155778894472363,0.5087885855590714) (-0.10552763819095466,0.516792367877932) (-0.0954773869346734,0.5128823532833312) (-0.08542713567839208,0.5025776727666259) (-0.07537688442211042,0.504860851919606) (-0.06532663316582912,0.5120065081398857) (-0.05527638190954787,0.5134272274134974) (-0.045226130653266215,0.5136005675219616) (-0.0351758793969849,0.49901449964128136) (-0.02512562814070368,0.5229782183267252) (-0.015075376884422013,0.5129247051456448) (-0.005025125628140743,0.5115276921487781) (0.005025125628140574,0.5134551773910012) (0.015075376884422167,0.5245669562052427) (0.02512562814070352,0.5240523038418473) (0.03517587939698479,0.5176954977464371) (0.04522613065326644,0.5312299941914927) (0.0552763819095477,0.5349282169338089) (0.06532663316582898,0.5448855050170315) (0.07537688442211059,0.5554410398400705) (0.08542713567839183,0.5557839431864219) (0.09547738693467356,0.5519811235515946) (0.10552763819095481,0.5612449066311016) (0.11557788944723604,0.5592356386912906) (0.12562814070351774,0.5496243389130636) (0.13567839195979903,0.5567112210075921) (0.1457286432160803,0.5621231256914863) (0.15577889447236196,0.5654549342753823) (0.16582914572864324,0.5840278748420826) (0.17587939698492452,0.5827859492640809) (0.1859296482412062,0.5942554131300651) (0.19597989949748745,0.6077331830397275) (0.2060301507537687,0.6014580706085741) (0.2160804020100504,0.607614504654591) (0.22613065326633164,0.6054674117220455) (0.2361809045226129,0.610939136217042) (0.2462311557788946,0.6150256334947779) (0.2562814070351759,0.6082520669644225) (0.26633165829145716,0.6150757567158754) (0.2763819095477388,0.6218281807123229) (0.28643216080402006,0.6224789673467125) (0.29648241206030135,0.6365437800037752) (0.30653266331658297,0.6365629431752365) (0.31658291457286425,0.6428388437300626) (0.3266331658291456,0.6638205328096272) (0.33668341708542715,0.6803879012920115) (0.3467336683417085,0.6823765909327106) (0.3567839195979898,0.6942318799315952) (0.36683417085427145,0.7161581976037956) (0.37688442211055273,0.7264999983603985) (0.38693467336683435,0.7392458147372706) (0.3969849246231157,0.7412373072143941) (0.40703517587939697,0.760934631463333) (0.41708542713567853,0.778591082621491) (0.42713567839195987,0.8002969497848346) (0.4371859296482411,0.8176794272960489) (0.44723618090452283,0.8292406413302676) (0.45728643216080406,0.8486637315958843) (0.46733668341708534,0.861847575973505) (0.47738693467336696,0.8800378252242751) (0.4874371859296483,0.8769939278843646) (0.4974874371859295,0.9089980219734776) (0.5075376884422113,0.9177570725315592) (0.5175879396984925,0.9338559172252734) (0.5276381909547737,0.9490403356737152) (0.5376884422110554,0.9664082810415973) (0.5477386934673367,0.9982526571835785) (0.557788944723618,1.0323421617714361) (0.5678391959798996,1.0497470450740325) (0.5778894472361809,1.0756823325475529) (0.5879396984924622,1.104317161098179) (0.5979899497487439,1.1389658430123848) (0.6080402010050251,1.1431886581014363) (0.6180904522613063,1.1262342365308502) (0.628140703517588,1.1414056457107185) (0.6381909547738693,1.1774396215841367) (0.6482412060301505,1.207137930849934) (0.6582914572864322,1.2367899801626316) (0.6683417085427136,1.2749752348869512) (0.6783919597989948,1.3117821003892642) (0.6884422110552765,1.3338343193404054) (0.6984924623115577,1.3853844807911309) (0.7085427135678392,1.4031001732209625) (0.7185929648241206,1.4437000297519373) (0.728643216080402,1.46548971295186) (0.7386934673366834,1.4838766392942135) (0.7487437185929648,1.513747155606617) (0.7587939698492461,1.5370069932283081) (0.7688442211055275,1.5702478696779871) (0.7788944723618091,1.5745807427243068) (0.7889447236180905,1.6259433811429589) (0.7989949748743719,1.686881490832456) (0.8090452261306533,1.7450102972015165) (0.8190954773869348,1.7662696861341494) (0.829145728643216,1.8092912407911759) (0.8391959798994976,1.8370180174046624) (0.849246231155779,1.8547428822811884) (0.8592964824120602,1.88180560882354) (0.8693467336683417,1.888473731021894) (0.8793969849246231,1.9145764376785408) (0.8894472361809044,1.9433510618719234) (0.8994974874371859,1.9429552807606507) (0.9095477386934674,1.9708424582270834) (0.9195979899497486,2.0308224847387013) (0.9296482412060301,2.0955618629735873) (0.9396984924623116,2.15002956846083) (0.949748743718593,2.174437015392303) (0.9597989949748743,2.195516660334927) (0.9698492462311559,2.2576914259203718) (0.9798994974874373,2.290824142582228) (0.9899497487437185,2.355240524256195) (1.0,2.419999733772909) };
  \addplot[ours_ransac, fill, fill opacity=0.18, draw=none, mark=none, forget plot] fill between [of=U_ransac_grp and L_ransac_grp];
  \addplot+[ours_ransac, mark=none] coordinates {(-1.0,0.7828683733473182) (-0.9899497487437187,0.7828683733473182) (-0.9798994974874371,0.7828683733473182) (-0.9698492462311558,0.7828683733473182) (-0.9597989949748745,0.7828683733473182) (-0.9497487437185929,0.7828683733473182) (-0.9396984924623116,0.7828683733473182) (-0.9296482412060303,0.7587368515597358) (-0.9195979899497487,0.7587368515597358) (-0.9095477386934674,0.7587368515597358) (-0.899497487437186,0.7587368515597358) (-0.8894472361809044,0.7587368515597358) (-0.8793969849246231,0.7587368515597358) (-0.8693467336683418,0.7587368515597358) (-0.8592964824120601,0.7359654338047772) (-0.8492462311557789,0.7359654338047772) (-0.8391959798994972,0.7359654338047772) (-0.829145728643216,0.7359654338047772) (-0.8190954773869348,0.7359654338047772) (-0.8090452261306531,0.7359654338047772) (-0.7989949748743718,0.6884872697045403) (-0.7889447236180905,0.6884872697045403) (-0.7788944723618089,0.6884872697045403) (-0.7688442211055275,0.6884872697045403) (-0.7587939698492463,0.6884872697045403) (-0.7487437185929647,0.6884872697045403) (-0.7386934673366835,0.6963686425290283) (-0.7286432160804021,0.6963686425290283) (-0.7185929648241205,0.6963686425290283) (-0.7085427135678393,0.6963686425290283) (-0.6984924623115579,0.6807301380232645) (-0.6884422110552763,0.6807301380232645) (-0.678391959798995,0.6807301380232645) (-0.6683417085427137,0.6807301380232645) (-0.658291457286432,0.6807301380232645) (-0.6482412060301508,0.663091474064718) (-0.6381909547738696,0.663091474064718) (-0.6281407035175878,0.663091474064718) (-0.6180904522613065,0.663091474064718) (-0.6080402010050253,0.6390526371145332) (-0.5979899497487436,0.6390526371145332) (-0.5879396984924623,0.6390526371145332) (-0.5778894472361811,0.6280279784748488) (-0.5678391959798994,0.6280279784748488) (-0.5577889447236181,0.6280279784748488) (-0.5477386934673368,0.6280279784748488) (-0.5376884422110553,0.6160228988688938) (-0.5276381909547739,0.6160228988688938) (-0.5175879396984923,0.6160228988688938) (-0.507537688442211,0.6017822469468467) (-0.4974874371859297,0.6017822469468467) (-0.4874371859296481,0.6021599466077557) (-0.47738693467336674,0.6021599466077557) (-0.4673366834170855,0.6021599466077557) (-0.45728643216080395,0.5973828838104859) (-0.44723618090452255,0.5973828838104859) (-0.4371859296482413,0.5973828838104859) (-0.42713567839195965,0.5906479819277838) (-0.4170854271356784,0.5906479819277838) (-0.40703517587939714,0.5797144783807247) (-0.39698492462311546,0.5797144783807247) (-0.3869346733668342,0.5818638416911153) (-0.3768844221105529,0.5818638416911153) (-0.3668341708542712,0.5818638416911153) (-0.35678391959798994,0.5733676669790416) (-0.34673366834170866,0.5733676669790416) (-0.33668341708542704,0.5647517030296592) (-0.32663316582914576,0.5647517030296592) (-0.3165829145728644,0.5633133631727227) (-0.30653266331658285,0.5568531183455782) (-0.29648241206030157,0.5568531183455782) (-0.2864321608040203,0.5559326366407807) (-0.27638190954773856,0.5559326366407807) (-0.26633165829145733,0.5591618124511135) (-0.2562814070351761,0.5591618124511135) (-0.24623115577889443,0.5622566028722292) (-0.23618090452261317,0.5608059285813686) (-0.2261306532663315,0.5608059285813686) (-0.2160804020100502,0.5417766207822133) (-0.206030150753769,0.5480321786540703) (-0.19597989949748723,0.5480321786540703) (-0.18592964824120595,0.5444590215225265) (-0.1758793969849247,0.5396103938953923) (-0.1658291457286431,0.5429830844464695) (-0.15577889447236182,0.5429830844464695) (-0.14572864321608053,0.5340760799699988) (-0.13567839195979886,0.5429996009351934) (-0.12562814070351758,0.5431896426913332) (-0.1155778894472363,0.5431896426913332) (-0.10552763819095466,0.5485593808854021) (-0.0954773869346734,0.5438245561290908) (-0.08542713567839208,0.5349276699061811) (-0.07537688442211042,0.5360910463631705) (-0.06532663316582912,0.5415313960624717) (-0.05527638190954787,0.5412936731965868) (-0.045226130653266215,0.5432724232457553) (-0.0351758793969849,0.5310032196513105) (-0.02512562814070368,0.5513519519496729) (-0.015075376884422013,0.5446470568423587) (-0.005025125628140743,0.5428674084155921) (0.005025125628140574,0.5440262518855221) (0.015075376884422167,0.5540029010090224) (0.02512562814070352,0.5530870253184222) (0.03517587939698479,0.5492500232801577) (0.04522613065326644,0.5592724567515253) (0.0552763819095477,0.562775291202503) (0.06532663316582898,0.5744469900579247) (0.07537688442211059,0.584254911880076) (0.08542713567839183,0.588171501474299) (0.09547738693467356,0.5824220412822606) (0.10552763819095481,0.5951925338586201) (0.11557788944723604,0.5926027075487551) (0.12562814070351774,0.5823659010453925) (0.13567839195979903,0.5890013586111849) (0.1457286432160803,0.5948326641844472) (0.15577889447236196,0.5966526191495874) (0.16582914572864324,0.6183365607376808) (0.17587939698492452,0.617894422443356) (0.1859296482412062,0.6296910939613369) (0.19597989949748745,0.6419880013394593) (0.2060301507537687,0.6378784981262487) (0.2160804020100504,0.641337141918313) (0.22613065326633164,0.644445586385066) (0.2361809045226129,0.6455534104765605) (0.2462311557788946,0.6544378410696079) (0.2562814070351759,0.6452917467125832) (0.26633165829145716,0.6501967124243512) (0.2763819095477388,0.6560019587281621) (0.28643216080402006,0.6574531116363542) (0.29648241206030135,0.6754891767745833) (0.30653266331658297,0.6698800500871397) (0.31658291457286425,0.6820592157444388) (0.3266331658291456,0.7014123249610763) (0.33668341708542715,0.714064207071479) (0.3467336683417085,0.7235123918613956) (0.3567839195979898,0.7300330818310656) (0.36683417085427145,0.7558197155214627) (0.37688442211055273,0.7687992102620355) (0.38693467336683435,0.7844936319530631) (0.3969849246231157,0.7867661360869996) (0.40703517587939697,0.8051433679355837) (0.41708542713567853,0.8222298380141391) (0.42713567839195987,0.8382226028570801) (0.4371859296482411,0.8544262307338832) (0.44723618090452283,0.8668195268271934) (0.45728643216080406,0.8832635773112841) (0.46733668341708534,0.8994897610379867) (0.47738693467336696,0.9229891569499475) (0.4874371859296483,0.9189200552641033) (0.4974874371859295,0.9497524018221927) (0.5075376884422113,0.9655479867260567) (0.5175879396984925,0.9773341168571003) (0.5276381909547737,0.995715814298186) (0.5376884422110554,1.0182892226930096) (0.5477386934673367,1.0498507768581293) (0.557788944723618,1.0864052787845115) (0.5678391959798996,1.1045238655190126) (0.5778894472361809,1.1282207494752519) (0.5879396984924622,1.1535848570770124) (0.5979899497487439,1.1831689869690991) (0.6080402010050251,1.192545237829187) (0.6180904522613063,1.1804294255916907) (0.628140703517588,1.200540595185666) (0.6381909547738693,1.23677442678834) (0.6482412060301505,1.2701549519051964) (0.6582914572864322,1.3025311647175295) (0.6683417085427136,1.3383228475158424) (0.6783919597989948,1.367401523681714) (0.6884422110552765,1.4040176224325103) (0.6984924623115577,1.4699278072619208) (0.7085427135678392,1.4833845386258064) (0.7185929648241206,1.5154424011842293) (0.728643216080402,1.5345046275235308) (0.7386934673366834,1.5534962320553451) (0.7487437185929648,1.585649704522311) (0.7587939698492461,1.6215774089396624) (0.7688442211055275,1.6574347972021672) (0.7788944723618091,1.6593589887656641) (0.7889447236180905,1.7135249169283706) (0.7989949748743719,1.7742207054124781) (0.8090452261306533,1.8244729012110577) (0.8190954773869348,1.851360424126862) (0.829145728643216,1.9001694646032121) (0.8391959798994976,1.9334634538516642) (0.849246231155779,1.9567301423503867) (0.8592964824120602,1.98283255978165) (0.8693467336683417,1.9887336602919) (0.8793969849246231,2.0223802482228286) (0.8894472361809044,2.0591637860044663) (0.8994974874371859,2.064192604228955) (0.9095477386934674,2.0936656553555992) (0.9195979899497486,2.1456370986744355) (0.9296482412060301,2.2029319411323733) (0.9396984924623116,2.2551021621772924) (0.949748743718593,2.2800529792985884) (0.9597989949748743,2.3213670237133823) (0.9698492462311559,2.4036443839348687) (0.9798994974874373,2.4459970615186744) (0.9899497487437185,2.511903577784948) (1.0,2.560094801876059) };
  \addplot+[ours_ransac, mark=none, line width=0.4pt, dashed, forget plot] coordinates {(-1.0,0.28308814087642103) (1.0,0.28308814087642103)};
  \addplot+[ours_ransac, mark=*, only marks, mark size=\PanelArgminSize, forget plot] coordinates {(-0.0351758793969849,0.5310032196513105)};
  \addplot[name path=U_msac_grp, draw=none, mark=none, forget plot] coordinates {(-1.0,0.7739547712338593) (-0.9899497487437187,0.7703318386160816) (-0.9798994974874371,0.7722077430036105) (-0.9698492462311558,0.7683035844288992) (-0.9597989949748745,0.7668280220252738) (-0.9497487437185929,0.7713406226628302) (-0.9396984924623116,0.7671243491808012) (-0.9296482412060303,0.7635571657691472) (-0.9195979899497487,0.7566106388011744) (-0.9095477386934674,0.7506192244861268) (-0.899497487437186,0.7465879550852452) (-0.8894472361809045,0.7488161988362028) (-0.8793969849246231,0.7491390876456626) (-0.8693467336683419,0.7458207949260053) (-0.8592964824120602,0.7427530241266033) (-0.8492462311557789,0.737213846142271) (-0.8391959798994973,0.7296993938716865) (-0.8291457286432161,0.7233910346748652) (-0.8190954773869348,0.7220167044879955) (-0.8090452261306531,0.7173351720781304) (-0.7989949748743719,0.7154571379967769) (-0.7889447236180905,0.7096975692324071) (-0.7788944723618089,0.7052356837890541) (-0.7688442211055276,0.7045892538489088) (-0.7587939698492464,0.7003341683133716) (-0.7487437185929646,0.6981526952901977) (-0.7386934673366834,0.6842694604733781) (-0.7286432160804022,0.6757410269866022) (-0.7185929648241205,0.6700906231305801) (-0.7085427135678393,0.6681145035298913) (-0.6984924623115579,0.665647526502279) (-0.6884422110552763,0.6605703253771118) (-0.678391959798995,0.6544142382743502) (-0.6683417085427137,0.6531431892995818) (-0.6582914572864321,0.6521287781056748) (-0.6482412060301508,0.646237425815116) (-0.6381909547738696,0.6442336173663402) (-0.6281407035175878,0.642634293803071) (-0.6180904522613067,0.6409952562404286) (-0.6080402010050253,0.6406364763399838) (-0.5979899497487436,0.6372800227069784) (-0.5879396984924624,0.6351835525659169) (-0.5778894472361812,0.6349060479654685) (-0.5678391959798994,0.6264057994931456) (-0.5577889447236181,0.6248575294210604) (-0.5477386934673368,0.6223555125790364) (-0.5376884422110553,0.6163972587562346) (-0.527638190954774,0.6144173824862771) (-0.5175879396984923,0.6112383561887134) (-0.507537688442211,0.612051608455891) (-0.4974874371859297,0.6105950138986664) (-0.487437185929648,0.607423625617263) (-0.47738693467336685,0.6031034078306361) (-0.4673366834170856,0.6017368623749615) (-0.45728643216080384,0.6014778696374335) (-0.4472361809045226,0.597150324546859) (-0.4371859296482413,0.5953490332230142) (-0.42713567839195965,0.5935705047942443) (-0.4170854271356784,0.5919282325098028) (-0.40703517587939714,0.5879245071274085) (-0.39698492462311546,0.5822161844655733) (-0.3869346733668342,0.580951726767746) (-0.37688442211055295,0.578620176369458) (-0.3668341708542713,0.5793430452542929) (-0.35678391959798994,0.5748389353084659) (-0.34673366834170866,0.5750153429345977) (-0.33668341708542704,0.5728023073682156) (-0.32663316582914576,0.572623631412998) (-0.31658291457286447,0.5704814571237471) (-0.30653266331658285,0.5728689958435677) (-0.29648241206030157,0.5669424411756674) (-0.2864321608040203,0.567165760299256) (-0.2763819095477386,0.5654164686568556) (-0.26633165829145733,0.5669701898485119) (-0.2562814070351761,0.569336943412637) (-0.24623115577889443,0.5696559037883289) (-0.2361809045226131,0.5692056197320504) (-0.2261306532663315,0.5689248465924867) (-0.2160804020100502,0.5680382150425095) (-0.2060301507537689,0.5622430793074518) (-0.19597989949748723,0.5622480473117116) (-0.18592964824120603,0.5621478058414611) (-0.17587939698492477,0.5609688188558525) (-0.1658291457286431,0.5619643963967667) (-0.15577889447236173,0.5628405250700398) (-0.14572864321608048,0.5599423981806833) (-0.13567839195979886,0.5608827608736314) (-0.12562814070351758,0.5587211947473808) (-0.11557788944723638,0.5535853240454041) (-0.10552763819095459,0.5560690573947001) (-0.0954773869346734,0.5572734450047847) (-0.08542713567839214,0.562090947316174) (-0.07537688442211042,0.5631314369889034) (-0.06532663316582912,0.561837504687493) (-0.05527638190954792,0.562853971705147) (-0.045226130653266215,0.5618048707881526) (-0.03517587939698501,0.5592592559783962) (-0.02512562814070373,0.5596964878074849) (-0.015075376884422061,0.5649959810998244) (-0.005025125628140743,0.5669142028141576) (0.005025125628140479,0.5679291410695957) (0.015075376884422167,0.5697592730237626) (0.025125628140703425,0.5699607233170132) (0.0351758793969847,0.5697449247381563) (0.045226130653266354,0.5738690700120167) (0.05527638190954762,0.5743655324443865) (0.0653266331658289,0.5738113188251847) (0.07537688442211059,0.5723503210823977) (0.08542713567839191,0.5728179276302033) (0.09547738693467347,0.5759529798672103) (0.10552763819095481,0.5748054987468557) (0.11557788944723604,0.5732032637313831) (0.12562814070351774,0.5718442440104234) (0.13567839195979903,0.5671079733114952) (0.14572864321608026,0.567769183851011) (0.15577889447236196,0.5698290050011131) (0.16582914572864324,0.5718022290629113) (0.17587939698492452,0.5743356567601359) (0.18592964824120614,0.574625510076866) (0.19597989949748745,0.573441362168982) (0.2060301507537687,0.5731827087355775) (0.21608040201005033,0.572662130241716) (0.2261306532663317,0.5707199062517683) (0.23618090452261295,0.5733858149988094) (0.2462311557788946,0.5749778516296187) (0.2562814070351759,0.5774506222642584) (0.2663316582914571,0.5780418012588837) (0.2763819095477387,0.5783242846476413) (0.28643216080402006,0.5786193867308431) (0.2964824120603013,0.5810295695846286) (0.30653266331658297,0.585609815440483) (0.31658291457286425,0.588173966185519) (0.3266331658291455,0.5877967838674473) (0.33668341708542715,0.5872488959956195) (0.3467336683417085,0.5888962709665758) (0.3567839195979898,0.5903313660427669) (0.36683417085427145,0.5921825156558989) (0.37688442211055273,0.5921679518642391) (0.38693467336683435,0.5921335340576541) (0.39698492462311563,0.5922734999064685) (0.40703517587939686,0.5938336958982889) (0.41708542713567853,0.5956270473516964) (0.4271356783919598,0.5968339886416253) (0.4371859296482411,0.6008028193874341) (0.44723618090452255,0.60438044434352) (0.457286432160804,0.6058155611963121) (0.4673366834170855,0.6086801385240744) (0.4773869346733668,0.6118539169078872) (0.48743718592964824,0.6124566243289115) (0.4974874371859297,0.6135207304006693) (0.5075376884422109,0.6148966863235188) (0.5175879396984924,0.6178024695700713) (0.5276381909547739,0.6201348183590034) (0.5376884422110553,0.6206712617853408) (0.5477386934673366,0.622062289747128) (0.5577889447236181,0.6217302492494337) (0.5678391959798994,0.6231357809721492) (0.5778894472361809,0.6251371406236994) (0.5879396984924623,0.6257435291552877) (0.5979899497487435,0.6267505560477149) (0.6080402010050251,0.629999568256941) (0.6180904522613065,0.6308874352263275) (0.628140703517588,0.6311063990717726) (0.6381909547738693,0.6328349570294993) (0.6482412060301508,0.6346949570477612) (0.6582914572864322,0.6365572387395066) (0.6683417085427136,0.6362856666574304) (0.678391959798995,0.6366907295466017) (0.6884422110552765,0.6398218998993703) (0.6984924623115577,0.6417230506664375) (0.7085427135678392,0.6381765507519692) (0.7185929648241206,0.6402288004331445) (0.728643216080402,0.6418761671624037) (0.7386934673366834,0.643427466873953) (0.7487437185929648,0.6487514708443423) (0.7587939698492461,0.6480451448427811) (0.7688442211055276,0.650318953487945) (0.7788944723618091,0.6503741035115319) (0.7889447236180905,0.6510782831574028) (0.7989949748743718,0.65271052436538) (0.8090452261306533,0.6538837555598208) (0.8190954773869348,0.654581975919349) (0.829145728643216,0.6568089284602017) (0.8391959798994975,0.659099081389254) (0.8492462311557789,0.6613529557438277) (0.8592964824120601,0.6613529557438277) (0.8693467336683417,0.6628321007671217) (0.8793969849246231,0.6635620806092547) (0.8894472361809044,0.6643870327324866) (0.8994974874371859,0.6656097167900622) (0.9095477386934674,0.667233148782257) (0.9195979899497486,0.6672975282571473) (0.9296482412060301,0.6672975282571473) (0.9396984924623115,0.6706792062834687) (0.949748743718593,0.6707350134283802) (0.9597989949748743,0.6722775343443724) (0.9698492462311558,0.6730754239364888) (0.9798994974874372,0.6744700733772986) (0.9899497487437185,0.67380005726126) (1.0,0.6747137993004588) };
  \addplot[name path=L_msac_grp, draw=none, mark=none, forget plot] coordinates {(-1.0,0.6731900632850836) (-0.9899497487437187,0.6710541614021628) (-0.9798994974874371,0.6723189917449628) (-0.9698492462311558,0.6699459560661394) (-0.9597989949748745,0.6696585912798175) (-0.9497487437185929,0.6735080242930865) (-0.9396984924623116,0.6716752651896084) (-0.9296482412060303,0.6627840480142672) (-0.9195979899497487,0.6555711978705376) (-0.9095477386934674,0.6532460702140558) (-0.899497487437186,0.6500719810152507) (-0.8894472361809045,0.6537939264752266) (-0.8793969849246231,0.65369242202147) (-0.8693467336683419,0.6527648110788172) (-0.8592964824120602,0.6511033929469173) (-0.8492462311557789,0.6417444205857816) (-0.8391959798994973,0.63495761156517) (-0.8291457286432161,0.629711945323956) (-0.8190954773869348,0.6306469724564769) (-0.8090452261306531,0.6254268193285183) (-0.7989949748743719,0.6229608327089589) (-0.7889447236180905,0.61655473084273) (-0.7788944723618089,0.6124597463583215) (-0.7688442211055276,0.6078366385155255) (-0.7587939698492464,0.6052555428148519) (-0.7487437185929646,0.6033639906188273) (-0.7386934673366834,0.5938091477381924) (-0.7286432160804022,0.5890832874059513) (-0.7185929648241205,0.5854894964351747) (-0.7085427135678393,0.5844297744481398) (-0.6984924623115579,0.5847499053499932) (-0.6884422110552763,0.5814845632255481) (-0.678391959798995,0.5795096022412182) (-0.6683417085427137,0.5780350604826503) (-0.6582914572864321,0.5754390171863717) (-0.6482412060301508,0.5699835048915349) (-0.6381909547738696,0.5688982359099309) (-0.6281407035175878,0.5667831649272198) (-0.6180904522613067,0.5656126080729695) (-0.6080402010050253,0.5642116998195704) (-0.5979899497487436,0.5617151861459124) (-0.5879396984924624,0.5597960865071046) (-0.5778894472361812,0.5586116772236553) (-0.5678391959798994,0.5516251363966165) (-0.5577889447236181,0.5505506940639712) (-0.5477386934673368,0.5495922197299482) (-0.5376884422110553,0.5446018702905205) (-0.527638190954774,0.5427528101137677) (-0.5175879396984923,0.5383505122597334) (-0.507537688442211,0.5392686061582422) (-0.4974874371859297,0.5374207456915596) (-0.487437185929648,0.5348543898898538) (-0.47738693467336685,0.53087457385501) (-0.4673366834170856,0.5330253834379612) (-0.45728643216080384,0.5322858292697724) (-0.4472361809045226,0.5270214784475884) (-0.4371859296482413,0.527153241103095) (-0.42713567839195965,0.5255616753036287) (-0.4170854271356784,0.5244817439252558) (-0.40703517587939714,0.5192298103809293) (-0.39698492462311546,0.5135585926470891) (-0.3869346733668342,0.5125379375696146) (-0.37688442211055295,0.5099649489587855) (-0.3668341708542713,0.5087236649112967) (-0.35678391959798994,0.503852825678282) (-0.34673366834170866,0.5053161619838771) (-0.33668341708542704,0.5032587868512037) (-0.32663316582914576,0.5007667139120482) (-0.31658291457286447,0.4993190196178914) (-0.30653266331658285,0.5009032191010433) (-0.29648241206030157,0.4982921349462692) (-0.2864321608040203,0.4975254978172502) (-0.2763819095477386,0.4953271240763288) (-0.26633165829145733,0.4964744497818738) (-0.2562814070351761,0.49793251606329103) (-0.24623115577889443,0.49797284996721575) (-0.2361809045226131,0.49789007551657904) (-0.2261306532663315,0.4998266679917295) (-0.2160804020100502,0.4989531298639941) (-0.2060301507537689,0.4945790625386162) (-0.19597989949748723,0.4956182354744668) (-0.18592964824120603,0.4966145080673412) (-0.17587939698492477,0.49364416767260605) (-0.1658291457286431,0.49610762051894436) (-0.15577889447236173,0.49631412726473945) (-0.14572864321608048,0.4927978105022898) (-0.13567839195979886,0.49206417167376737) (-0.12562814070351758,0.4889401963534442) (-0.11557788944723638,0.48531095961216947) (-0.10552763819095459,0.4878508214204659) (-0.0954773869346734,0.48773353364970573) (-0.08542713567839214,0.4940160511460129) (-0.07537688442211042,0.49728365099838656) (-0.06532663316582912,0.496634491617373) (-0.05527638190954792,0.497019696976679) (-0.045226130653266215,0.494136348476966) (-0.03517587939698501,0.49373195045762097) (-0.02512562814070373,0.4947672106099389) (-0.015075376884422061,0.497735228435316) (-0.005025125628140743,0.4985695963573145) (0.005025125628140479,0.5011877094125523) (0.015075376884422167,0.5022466198720761) (0.025125628140703425,0.5001183646261533) (0.0351758793969847,0.5003799709848251) (0.045226130653266354,0.5049264675147743) (0.05527638190954762,0.5055046158883681) (0.0653266331658289,0.5065194219456338) (0.07537688442211059,0.5015886913284051) (0.08542713567839191,0.5034937135600045) (0.09547738693467347,0.506724839402839) (0.10552763819095481,0.5057205650382145) (0.11557788944723604,0.5040547718366784) (0.12562814070351774,0.5004871956066399) (0.13567839195979903,0.4957444697598346) (0.14572864321608026,0.4962869740740466) (0.15577889447236196,0.4983914945301747) (0.16582914572864324,0.5003213094060018) (0.17587939698492452,0.5018439007939888) (0.18592964824120614,0.5038131422929144) (0.19597989949748745,0.5022374688445588) (0.2060301507537687,0.5027678270698056) (0.21608040201005033,0.5028058854950578) (0.2261306532663317,0.5010129345478107) (0.23618090452261295,0.5032910404803548) (0.2462311557788946,0.5065967032542791) (0.2562814070351759,0.5079397136305739) (0.2663316582914571,0.51006943465094) (0.2763819095477387,0.5117750524210345) (0.28643216080402006,0.5145329693425948) (0.2964824120603013,0.5168824873955595) (0.30653266331658297,0.5202156909703406) (0.31658291457286425,0.5236914842114841) (0.3266331658291455,0.5244968689730555) (0.33668341708542715,0.5237340567103224) (0.3467336683417085,0.5267360370928015) (0.3567839195979898,0.5273275912357669) (0.36683417085427145,0.5288518903584558) (0.37688442211055273,0.5283107212596838) (0.38693467336683435,0.5283010097005862) (0.39698492462311563,0.5299012421739238) (0.40703517587939686,0.532245603560476) (0.41708542713567853,0.5344825803287384) (0.4271356783919598,0.5363494029164243) (0.4371859296482411,0.5394651046555682) (0.44723618090452255,0.5424513894494817) (0.457286432160804,0.5438999082936449) (0.4673366834170855,0.545698807091468) (0.4773869346733668,0.5481039350636818) (0.48743718592964824,0.5491390248335525) (0.4974874371859297,0.5490714256386777) (0.5075376884422109,0.5499391229621566) (0.5175879396984924,0.5526218456082165) (0.5276381909547739,0.5549546879051079) (0.5376884422110553,0.5550538987282104) (0.5477386934673366,0.5553822639603849) (0.5577889447236181,0.5540939725297842) (0.5678391959798994,0.5552618373366653) (0.5778894472361809,0.5572019828800843) (0.5879396984924623,0.558600754963257) (0.5979899497487435,0.5600794232300415) (0.6080402010050251,0.5629947204827996) (0.6180904522613065,0.5634878739432455) (0.628140703517588,0.5625273345282257) (0.6381909547738693,0.5643153932181215) (0.6482412060301508,0.5660321338456217) (0.6582914572864322,0.5679372371498232) (0.6683417085427136,0.5683876178923684) (0.678391959798995,0.5681600104230399) (0.6884422110552765,0.5701859889307412) (0.6984924623115577,0.571974732402061) (0.7085427135678392,0.5699753192884052) (0.7185929648241206,0.5708959500363862) (0.728643216080402,0.5713386469461558) (0.7386934673366834,0.5731565928111997) (0.7487437185929648,0.5772962879087571) (0.7587939698492461,0.5758608839354353) (0.7688442211055276,0.577155121672102) (0.7788944723618091,0.5771149090083395) (0.7889447236180905,0.577809430036106) (0.7989949748743718,0.5785204782846339) (0.8090452261306533,0.5804455855593301) (0.8190954773869348,0.5811705533945699) (0.829145728643216,0.5814441241649848) (0.8391959798994975,0.583040766129585) (0.8492462311557789,0.5855092742922338) (0.8592964824120601,0.5855092742922338) (0.8693467336683417,0.5867760249698833) (0.8793969849246231,0.5873171078788176) (0.8894472361809044,0.5878401606647983) (0.8994974874371859,0.5893588036107328) (0.9095477386934674,0.5904389456336595) (0.9195979899497486,0.5903389674739993) (0.9296482412060301,0.5903389674739993) (0.9396984924623115,0.5926237744303814) (0.949748743718593,0.5926790409109873) (0.9597989949748743,0.5941072847243266) (0.9698492462311558,0.5948206297274515) (0.9798994974874372,0.5953263537529305) (0.9899497487437185,0.594436262195717) (1.0,0.5950511135434411) };
  \addplot[ours_msac, fill, fill opacity=0.18, draw=none, mark=none, forget plot] fill between [of=U_msac_grp and L_msac_grp];
  \addplot+[ours_msac, mark=none] coordinates {(-1.0,0.7235724172594714) (-0.9899497487437187,0.7206930000091222) (-0.9798994974874371,0.7222633673742866) (-0.9698492462311558,0.7191247702475193) (-0.9597989949748745,0.7182433066525457) (-0.9497487437185929,0.7224243234779584) (-0.9396984924623116,0.7193998071852048) (-0.9296482412060303,0.7131706068917072) (-0.9195979899497487,0.706090918335856) (-0.9095477386934674,0.7019326473500913) (-0.899497487437186,0.6983299680502479) (-0.8894472361809045,0.7013050626557147) (-0.8793969849246231,0.7014157548335663) (-0.8693467336683419,0.6992928030024113) (-0.8592964824120602,0.6969282085367603) (-0.8492462311557789,0.6894791333640263) (-0.8391959798994973,0.6823285027184283) (-0.8291457286432161,0.6765514899994106) (-0.8190954773869348,0.6763318384722362) (-0.8090452261306531,0.6713809957033243) (-0.7989949748743719,0.6692089853528679) (-0.7889447236180905,0.6631261500375686) (-0.7788944723618089,0.6588477150736878) (-0.7688442211055276,0.6562129461822171) (-0.7587939698492464,0.6527948555641118) (-0.7487437185929646,0.6507583429545125) (-0.7386934673366834,0.6390393041057852) (-0.7286432160804022,0.6324121571962767) (-0.7185929648241205,0.6277900597828774) (-0.7085427135678393,0.6262721389890156) (-0.6984924623115579,0.6251987159261361) (-0.6884422110552763,0.62102744430133) (-0.678391959798995,0.6169619202577842) (-0.6683417085427137,0.6155891248911161) (-0.6582914572864321,0.6137838976460233) (-0.6482412060301508,0.6081104653533255) (-0.6381909547738696,0.6065659266381356) (-0.6281407035175878,0.6047087293651454) (-0.6180904522613067,0.6033039321566991) (-0.6080402010050253,0.6024240880797771) (-0.5979899497487436,0.5994976044264454) (-0.5879396984924624,0.5974898195365107) (-0.5778894472361812,0.5967588625945619) (-0.5678391959798994,0.589015467944881) (-0.5577889447236181,0.5877041117425158) (-0.5477386934673368,0.5859738661544923) (-0.5376884422110553,0.5804995645233776) (-0.527638190954774,0.5785850963000224) (-0.5175879396984923,0.5747944342242234) (-0.507537688442211,0.5756601073070666) (-0.4974874371859297,0.574007879795113) (-0.487437185929648,0.5711390077535584) (-0.47738693467336685,0.5669889908428231) (-0.4673366834170856,0.5673811229064614) (-0.45728643216080384,0.566881849453603) (-0.4472361809045226,0.5620859014972237) (-0.4371859296482413,0.5612511371630546) (-0.42713567839195965,0.5595660900489365) (-0.4170854271356784,0.5582049882175293) (-0.40703517587939714,0.5535771587541689) (-0.39698492462311546,0.5478873885563312) (-0.3869346733668342,0.5467448321686803) (-0.37688442211055295,0.5442925626641217) (-0.3668341708542713,0.5440333550827948) (-0.35678391959798994,0.539345880493374) (-0.34673366834170866,0.5401657524592374) (-0.33668341708542704,0.5380305471097097) (-0.32663316582914576,0.5366951726625231) (-0.31658291457286447,0.5349002383708192) (-0.30653266331658285,0.5368861074723055) (-0.29648241206030157,0.5326172880609683) (-0.2864321608040203,0.5323456290582531) (-0.2763819095477386,0.5303717963665922) (-0.26633165829145733,0.5317223198151928) (-0.2562814070351761,0.533634729737964) (-0.24623115577889443,0.5338143768777723) (-0.2361809045226131,0.5335478476243147) (-0.2261306532663315,0.5343757572921081) (-0.2160804020100502,0.5334956724532518) (-0.2060301507537689,0.528411070923034) (-0.19597989949748723,0.5289331413930892) (-0.18592964824120603,0.5293811569544011) (-0.17587939698492477,0.5273064932642293) (-0.1658291457286431,0.5290360084578555) (-0.15577889447236173,0.5295773261673896) (-0.14572864321608048,0.5263701043414866) (-0.13567839195979886,0.5264734662736994) (-0.12562814070351758,0.5238306955504125) (-0.11557788944723638,0.5194481418287867) (-0.10552763819095459,0.521959939407583) (-0.0954773869346734,0.5225034893272452) (-0.08542713567839214,0.5280534992310935) (-0.07537688442211042,0.530207543993645) (-0.06532663316582912,0.529235998152433) (-0.05527638190954792,0.529936834340913) (-0.045226130653266215,0.5279706096325593) (-0.03517587939698501,0.5264956032180086) (-0.02512562814070373,0.5272318492087119) (-0.015075376884422061,0.5313656047675702) (-0.005025125628140743,0.532741899585736) (0.005025125628140479,0.534558425241074) (0.015075376884422167,0.5360029464479193) (0.025125628140703425,0.5350395439715833) (0.0351758793969847,0.5350624478614907) (0.045226130653266354,0.5393977687633955) (0.05527638190954762,0.5399350741663773) (0.0653266331658289,0.5401653703854092) (0.07537688442211059,0.5369695062054014) (0.08542713567839191,0.5381558205951039) (0.09547738693467347,0.5413389096350246) (0.10552763819095481,0.5402630318925351) (0.11557788944723604,0.5386290177840307) (0.12562814070351774,0.5361657198085317) (0.13567839195979903,0.5314262215356649) (0.14572864321608026,0.5320280789625288) (0.15577889447236196,0.5341102497656439) (0.16582914572864324,0.5360617692344566) (0.17587939698492452,0.5380897787770623) (0.18592964824120614,0.5392193261848902) (0.19597989949748745,0.5378394155067704) (0.2060301507537687,0.5379752679026916) (0.21608040201005033,0.5377340078683869) (0.2261306532663317,0.5358664203997895) (0.23618090452261295,0.5383384277395821) (0.2462311557788946,0.5407872774419489) (0.2562814070351759,0.5426951679474161) (0.2663316582914571,0.5440556179549119) (0.2763819095477387,0.5450496685343379) (0.28643216080402006,0.546576178036719) (0.2964824120603013,0.5489560284900941) (0.30653266331658297,0.5529127532054118) (0.31658291457286425,0.5559327251985016) (0.3266331658291455,0.5561468264202514) (0.33668341708542715,0.5554914763529709) (0.3467336683417085,0.5578161540296886) (0.3567839195979898,0.5588294786392669) (0.36683417085427145,0.5605172030071773) (0.37688442211055273,0.5602393365619615) (0.38693467336683435,0.5602172718791202) (0.39698492462311563,0.5610873710401961) (0.40703517587939686,0.5630396497293825) (0.41708542713567853,0.5650548138402174) (0.4271356783919598,0.5665916957790248) (0.4371859296482411,0.5701339620215011) (0.44723618090452255,0.5734159168965008) (0.457286432160804,0.5748577347449785) (0.4673366834170855,0.5771894728077712) (0.4773869346733668,0.5799789259857845) (0.48743718592964824,0.580797824581232) (0.4974874371859297,0.5812960780196735) (0.5075376884422109,0.5824179046428377) (0.5175879396984924,0.5852121575891439) (0.5276381909547739,0.5875447531320557) (0.5376884422110553,0.5878625802567756) (0.5477386934673366,0.5887222768537564) (0.5577889447236181,0.5879121108896089) (0.5678391959798994,0.5891988091544073) (0.5778894472361809,0.5911695617518918) (0.5879396984924623,0.5921721420592724) (0.5979899497487435,0.5934149896388782) (0.6080402010050251,0.5964971443698703) (0.6180904522613065,0.5971876545847865) (0.628140703517588,0.5968168667999991) (0.6381909547738693,0.5985751751238104) (0.6482412060301508,0.6003635454466915) (0.6582914572864322,0.602247237944665) (0.6683417085427136,0.6023366422748994) (0.678391959798995,0.6024253699848208) (0.6884422110552765,0.6050039444150558) (0.6984924623115577,0.6068488915342493) (0.7085427135678392,0.6040759350201872) (0.7185929648241206,0.6055623752347653) (0.728643216080402,0.6066074070542797) (0.7386934673366834,0.6082920298425764) (0.7487437185929648,0.6130238793765497) (0.7587939698492461,0.6119530143891082) (0.7688442211055276,0.6137370375800235) (0.7788944723618091,0.6137445062599357) (0.7889447236180905,0.6144438565967544) (0.7989949748743718,0.615615501325007) (0.8090452261306533,0.6171646705595755) (0.8190954773869348,0.6178762646569594) (0.829145728643216,0.6191265263125932) (0.8391959798994975,0.6210699237594195) (0.8492462311557789,0.6234311150180307) (0.8592964824120601,0.6234311150180307) (0.8693467336683417,0.6248040628685025) (0.8793969849246231,0.6254395942440362) (0.8894472361809044,0.6261135966986424) (0.8994974874371859,0.6274842602003975) (0.9095477386934674,0.6288360472079583) (0.9195979899497486,0.6288182478655733) (0.9296482412060301,0.6288182478655733) (0.9396984924623115,0.631651490356925) (0.949748743718593,0.6317070271696837) (0.9597989949748743,0.6331924095343495) (0.9698492462311558,0.6339480268319702) (0.9798994974874372,0.6348982135651146) (0.9899497487437185,0.6341181597284885) (1.0,0.63488245642195) };
  \addplot+[ours_msac, mark=none, line width=0.4pt, dashed, forget plot] coordinates {(-1.0,0.39966080400959836) (1.0,0.39966080400959836)};
  \addplot+[ours_msac, mark=*, only marks, mark size=\PanelArgminSize, forget plot] coordinates {(-0.11557788944723638,0.5194481418287867)};
  \addplot[name path=U_gau_grp, draw=none, mark=none, forget plot] coordinates {(-1.0,0.744966015318872) (-0.9899497487437187,0.7403781425440071) (-0.9798994974874371,0.7355663015154795) (-0.9698492462311558,0.7274873335587951) (-0.9597989949748745,0.7225514806649164) (-0.9497487437185929,0.7211919001678742) (-0.9396984924623116,0.7191691374799892) (-0.9296482412060303,0.7181322809592946) (-0.9195979899497487,0.7152038042104262) (-0.9095477386934674,0.7112059385513592) (-0.899497487437186,0.7069631028410187) (-0.8894472361809044,0.7007457178259004) (-0.8793969849246231,0.6987623125727493) (-0.8693467336683418,0.6958884193572534) (-0.8592964824120601,0.6890765903158574) (-0.8492462311557789,0.6783766983302407) (-0.8391959798994972,0.673007379502002) (-0.829145728643216,0.6669107830617828) (-0.8190954773869348,0.6650979763941482) (-0.8090452261306531,0.6605427853536946) (-0.7989949748743718,0.6637783102090872) (-0.7889447236180905,0.6598936204437018) (-0.7788944723618089,0.6571174710976964) (-0.7688442211055275,0.6546069967201662) (-0.7587939698492463,0.6455605805789429) (-0.7487437185929647,0.6474275525680135) (-0.7386934673366835,0.6447776260146932) (-0.7286432160804021,0.6392144783349195) (-0.7185929648241205,0.6378797313134486) (-0.7085427135678393,0.6332777349601167) (-0.6984924623115579,0.6340617200822704) (-0.6884422110552763,0.6324011674574413) (-0.678391959798995,0.6320128488058453) (-0.6683417085427137,0.6294342015188026) (-0.658291457286432,0.6259463333677685) (-0.6482412060301508,0.6247338262237122) (-0.6381909547738696,0.6209376557991828) (-0.6281407035175878,0.6176925553825089) (-0.6180904522613065,0.6123220842439834) (-0.6080402010050253,0.6085104000216374) (-0.5979899497487436,0.6091290305558875) (-0.5879396984924623,0.60704069845034) (-0.5778894472361811,0.6060952047913096) (-0.5678391959798994,0.6039658661091455) (-0.5577889447236181,0.6008147976037986) (-0.5477386934673368,0.5996609774704684) (-0.5376884422110553,0.5976543815707726) (-0.5276381909547739,0.5985916905650268) (-0.5175879396984923,0.596663439692456) (-0.507537688442211,0.592821410271956) (-0.4974874371859297,0.5907868153698209) (-0.4874371859296481,0.5872946088236887) (-0.47738693467336674,0.5864692076102047) (-0.4673366834170855,0.5863145627425419) (-0.45728643216080395,0.586449704237737) (-0.44723618090452255,0.5839968435859767) (-0.4371859296482413,0.5837296619991834) (-0.42713567839195965,0.5818118966730852) (-0.4170854271356784,0.5758727969914423) (-0.40703517587939714,0.5757280549932237) (-0.39698492462311546,0.5731580451714325) (-0.3869346733668342,0.5711802472688676) (-0.3768844221105529,0.5717150059001532) (-0.3668341708542712,0.5721400597544538) (-0.35678391959798994,0.5721243607512181) (-0.34673366834170866,0.571163097946739) (-0.33668341708542704,0.5708381160310069) (-0.32663316582914576,0.5711102479357852) (-0.3165829145728644,0.5720702105185166) (-0.30653266331658285,0.5701529429108031) (-0.29648241206030157,0.570659225423431) (-0.2864321608040203,0.5668466983790081) (-0.27638190954773856,0.5654301328331653) (-0.26633165829145733,0.564465507891647) (-0.2562814070351761,0.5674766481615902) (-0.24623115577889443,0.5672439206127075) (-0.23618090452261317,0.5676746984860878) (-0.2261306532663315,0.5668294456227515) (-0.2160804020100502,0.5654708727453898) (-0.206030150753769,0.5631522413278235) (-0.19597989949748723,0.5627321349715435) (-0.18592964824120595,0.561979923811534) (-0.1758793969849247,0.5623752894139553) (-0.1658291457286431,0.5612993864970047) (-0.15577889447236182,0.5595461167780502) (-0.14572864321608053,0.5600485401811797) (-0.13567839195979886,0.5624950137225913) (-0.12562814070351758,0.5617768173863497) (-0.1155778894472363,0.5615804002692899) (-0.10552763819095466,0.560340756758593) (-0.0954773869346734,0.5607979682697068) (-0.08542713567839208,0.5607619445059634) (-0.07537688442211042,0.5631652444589335) (-0.06532663316582912,0.5632903364253617) (-0.05527638190954787,0.5615932967578761) (-0.045226130653266215,0.5641275407788072) (-0.0351758793969849,0.5644633611489944) (-0.02512562814070368,0.5654789274106863) (-0.015075376884422013,0.5657285413816711) (-0.005025125628140743,0.5664233975106443) (0.005025125628140574,0.566925750841663) (0.015075376884422167,0.5711462506141991) (0.02512562814070352,0.5727060173574012) (0.03517587939698479,0.5691179633319763) (0.04522613065326644,0.568342098464484) (0.0552763819095477,0.5686481019342826) (0.06532663316582898,0.5688705523654463) (0.07537688442211059,0.569813396125058) (0.08542713567839183,0.5705829951758888) (0.09547738693467356,0.5706270075759486) (0.10552763819095481,0.5692334555320103) (0.11557788944723604,0.5692002007086727) (0.12562814070351774,0.5679880964410751) (0.13567839195979903,0.56933047659616) (0.1457286432160803,0.5722957314905824) (0.15577889447236196,0.5727499400267824) (0.16582914572864324,0.5716668433587846) (0.17587939698492452,0.5711718868008595) (0.1859296482412062,0.569790520574724) (0.19597989949748745,0.5718595969686406) (0.2060301507537687,0.5748193968306881) (0.2160804020100504,0.5751974252813643) (0.22613065326633164,0.57738658205236) (0.2361809045226129,0.5778257780247261) (0.2462311557788946,0.579562770964405) (0.2562814070351759,0.5801747796876907) (0.26633165829145716,0.5806308561180296) (0.2763819095477388,0.5806848847753578) (0.28643216080402006,0.5798050916644764) (0.29648241206030135,0.5813695558936618) (0.30653266331658297,0.5834208637291977) (0.31658291457286425,0.5830100876295283) (0.3266331658291456,0.5848975642510539) (0.33668341708542715,0.5847472345680602) (0.3467336683417085,0.5859247049767248) (0.3567839195979898,0.5867890759544302) (0.36683417085427145,0.588422756824671) (0.37688442211055273,0.5887232751901106) (0.38693467336683435,0.591041976579431) (0.3969849246231157,0.5906173467503593) (0.40703517587939697,0.5962187345960837) (0.41708542713567853,0.5977479557022322) (0.42713567839195987,0.5977479557022322) (0.4371859296482411,0.5975732847899666) (0.44723618090452283,0.5982558135113456) (0.45728643216080406,0.5996836546308922) (0.46733668341708534,0.5997662895151218) (0.47738693467336696,0.6014819124105492) (0.4874371859296483,0.6025316287611497) (0.4974874371859295,0.6048358448060758) (0.5075376884422113,0.6051458202421743) (0.5175879396984925,0.6072269168458818) (0.5276381909547737,0.6077977276587752) (0.5376884422110554,0.6095741890080113) (0.5477386934673367,0.6087040835337164) (0.557788944723618,0.6128731515918776) (0.5678391959798996,0.6158575052099411) (0.5778894472361809,0.6164101209598858) (0.5879396984924622,0.614681823815033) (0.5979899497487439,0.6166691758203356) (0.6080402010050251,0.617315583462563) (0.6180904522613063,0.6173656605154526) (0.628140703517588,0.6172179116256196) (0.6381909547738693,0.617891367742603) (0.6482412060301505,0.6200590237464743) (0.6582914572864322,0.6206204020635753) (0.6683417085427136,0.6206543365464048) (0.6783919597989948,0.6230166688702018) (0.6884422110552765,0.6235210915418663) (0.6984924623115577,0.6247226962556752) (0.7085427135678392,0.6245188995280192) (0.7185929648241206,0.6271402388712697) (0.728643216080402,0.6285027690343696) (0.7386934673366834,0.6280134239680397) (0.7487437185929648,0.628297940028236) (0.7587939698492461,0.6281841031931212) (0.7688442211055275,0.6281841031931212) (0.7788944723618091,0.6285107314162158) (0.7889447236180905,0.6305789318819298) (0.7989949748743719,0.6305789318819298) (0.8090452261306533,0.6335898318089075) (0.8190954773869348,0.6338797431427464) (0.829145728643216,0.6367080912794757) (0.8391959798994976,0.638139279713304) (0.849246231155779,0.6389348593023552) (0.8592964824120602,0.6394839970080936) (0.8693467336683417,0.639532841692535) (0.8793969849246231,0.6403506839508643) (0.8894472361809044,0.6403506839508643) (0.8994974874371859,0.6427176020317383) (0.9095477386934674,0.6434314086283429) (0.9195979899497486,0.6434314086283429) (0.9296482412060301,0.6434314086283429) (0.9396984924623116,0.6441720576237084) (0.949748743718593,0.6441720576237084) (0.9597989949748743,0.6447783210422564) (0.9698492462311559,0.6451164719308693) (0.9798994974874373,0.6452709498343513) (0.9899497487437185,0.6452709498343513) (1.0,0.6452709498343513) };
  \addplot[name path=L_gau_grp, draw=none, mark=none, forget plot] coordinates {(-1.0,0.6464452073553417) (-0.9899497487437187,0.6421701431852317) (-0.9798994974874371,0.6355454641806025) (-0.9698492462311558,0.6312800700633264) (-0.9597989949748745,0.6290396037486484) (-0.9497487437185929,0.6293402390664964) (-0.9396984924623116,0.6262691664975778) (-0.9296482412060303,0.6252579632854344) (-0.9195979899497487,0.622253745565032) (-0.9095477386934674,0.6187022875419177) (-0.899497487437186,0.615110718781655) (-0.8894472361809044,0.6091234457669525) (-0.8793969849246231,0.607849861066678) (-0.8693467336683418,0.601203700341667) (-0.8592964824120601,0.5942948995914547) (-0.8492462311557789,0.5886291588717855) (-0.8391959798994972,0.585707928767015) (-0.829145728643216,0.5839597049176258) (-0.8190954773869348,0.5830613138975889) (-0.8090452261306531,0.5795105791330138) (-0.7989949748743718,0.5823310891535626) (-0.7889447236180905,0.5802606270252979) (-0.7788944723618089,0.5777982884332833) (-0.7688442211055275,0.5776374914993654) (-0.7587939698492463,0.5725305226724843) (-0.7487437185929647,0.5726495597439891) (-0.7386934673366835,0.5705513652072726) (-0.7286432160804021,0.5657598292558194) (-0.7185929648241205,0.5636273173712216) (-0.7085427135678393,0.5607711274135768) (-0.6984924623115579,0.5624178435593724) (-0.6884422110552763,0.5614810334178583) (-0.678391959798995,0.5596930247443778) (-0.6683417085427137,0.5565505081205973) (-0.658291457286432,0.5534013690826846) (-0.6482412060301508,0.5517373762129131) (-0.6381909547738696,0.5482243994871694) (-0.6281407035175878,0.5472498591606428) (-0.6180904522613065,0.5427304472162013) (-0.6080402010050253,0.5396304209909478) (-0.5979899497487436,0.53990462063821) (-0.5879396984924623,0.5381845877974241) (-0.5778894472361811,0.5377219860572918) (-0.5678391959798994,0.5345253514856613) (-0.5577889447236181,0.5326840970431971) (-0.5477386934673368,0.5309428653624548) (-0.5376884422110553,0.5290922646757702) (-0.5276381909547739,0.5294888382482446) (-0.5175879396984923,0.5286279033938989) (-0.507537688442211,0.5242377029591283) (-0.4974874371859297,0.5225559955455173) (-0.4874371859296481,0.5178210980911023) (-0.47738693467336674,0.5174461892475386) (-0.4673366834170855,0.5177275732327353) (-0.45728643216080395,0.5171421138108687) (-0.44723618090452255,0.5142404545527207) (-0.4371859296482413,0.5128510696736595) (-0.42713567839195965,0.5091707414633895) (-0.4170854271356784,0.5054530648011377) (-0.40703517587939714,0.502863729367791) (-0.39698492462311546,0.5019012716185877) (-0.3869346733668342,0.49886501271196537) (-0.3768844221105529,0.49946744931809567) (-0.3668341708542712,0.4990844296432236) (-0.35678391959798994,0.4999882384519287) (-0.34673366834170866,0.4996822127526385) (-0.33668341708542704,0.5012405169432502) (-0.32663316582914576,0.5011289937954329) (-0.3165829145728644,0.5015003638459706) (-0.30653266331658285,0.49955230853201316) (-0.29648241206030157,0.5001867270896226) (-0.2864321608040203,0.4986260774120814) (-0.27638190954773856,0.4964026992427145) (-0.26633165829145733,0.49438543290867276) (-0.2562814070351761,0.4971563553256625) (-0.24623115577889443,0.49740655302247017) (-0.23618090452261317,0.4967545718068921) (-0.2261306532663315,0.4975450283809791) (-0.2160804020100502,0.4962648593758863) (-0.206030150753769,0.49486624225298387) (-0.19597989949748723,0.49420086379248085) (-0.18592964824120595,0.4922447225069225) (-0.1758793969849247,0.49181360253617434) (-0.1658291457286431,0.492116720621761) (-0.15577889447236182,0.49033932682783976) (-0.14572864321608053,0.49097467563275804) (-0.13567839195979886,0.4917037337443293) (-0.12562814070351758,0.4896115543673164) (-0.1155778894472363,0.4894857778271706) (-0.10552763819095466,0.4905466514900562) (-0.0954773869346734,0.49281669946907747) (-0.08542713567839208,0.4929256930727283) (-0.07537688442211042,0.49506624380871944) (-0.06532663316582912,0.49462597954747134) (-0.05527638190954787,0.4943008194306345) (-0.045226130653266215,0.4957150415657117) (-0.0351758793969849,0.49479754924128055) (-0.02512562814070368,0.4956492016695574) (-0.015075376884422013,0.4951545619552284) (-0.005025125628140743,0.495733883378839) (0.005025125628140574,0.4965266724493871) (0.015075376884422167,0.4990830242039413) (0.02512562814070352,0.4996429243398715) (0.03517587939698479,0.49743784264658) (0.04522613065326644,0.49874483717524676) (0.0552763819095477,0.4991378164261097) (0.06532663316582898,0.498771595195707) (0.07537688442211059,0.49928942481819144) (0.08542713567839183,0.4998373225149978) (0.09547738693467356,0.49986250076971495) (0.10552763819095481,0.4988085043440621) (0.11557788944723604,0.49925109921675453) (0.12562814070351774,0.4982497773570602) (0.13567839195979903,0.49811560623020584) (0.1457286432160803,0.5002942308736732) (0.15577889447236196,0.5006169705668428) (0.16582914572864324,0.5010575578133254) (0.17587939698492452,0.5003348422078742) (0.1859296482412062,0.49920426191705225) (0.19597989949748745,0.49977751002521376) (0.2060301507537687,0.5015246146600434) (0.2160804020100504,0.5024836933238264) (0.22613065326633164,0.5044893467874472) (0.2361809045226129,0.5055345631817548) (0.2462311557788946,0.5078848840831199) (0.2562814070351759,0.5101471082054211) (0.26633165829145716,0.5099169120114364) (0.2763819095477388,0.5112355189978062) (0.28643216080402006,0.5121590715849893) (0.29648241206030135,0.5143331897816321) (0.30653266331658297,0.5171506813482628) (0.31658291457286425,0.5178693946923814) (0.3266331658291456,0.5182425558619805) (0.33668341708542715,0.5184083650119669) (0.3467336683417085,0.519654592190199) (0.3567839195979898,0.5217482267810208) (0.36683417085427145,0.5237071537390418) (0.37688442211055273,0.524913687079346) (0.38693467336683435,0.5288388285044867) (0.3969849246231157,0.5282063314575269) (0.40703517587939697,0.5336360412746022) (0.41708542713567853,0.5362882427056496) (0.42713567839195987,0.5362882427056496) (0.4371859296482411,0.5365104554111264) (0.44723618090452283,0.5368822860801636) (0.45728643216080406,0.5388096899183304) (0.46733668341708534,0.5402748442966038) (0.47738693467336696,0.5418141834028715) (0.4874371859296483,0.5431107264099049) (0.4974874371859295,0.5436805319264406) (0.5075376884422113,0.5435128872007852) (0.5175879396984925,0.5444512792938617) (0.5276381909547737,0.5459033522420869) (0.5376884422110554,0.5474958616351882) (0.5477386934673367,0.5476468056943627) (0.557788944723618,0.5512606527720877) (0.5678391959798996,0.5532547039237967) (0.5778894472361809,0.553878969009528) (0.5879396984924622,0.5524096811195917) (0.5979899497487439,0.5537807709016306) (0.6080402010050251,0.5543151159887045) (0.6180904522613063,0.5541735878524178) (0.628140703517588,0.5548637066320146) (0.6381909547738693,0.5550581442942941) (0.6482412060301505,0.5572606009749697) (0.6582914572864322,0.5578473890455371) (0.6683417085427136,0.5578366727749455) (0.6783919597989948,0.5599807217862244) (0.6884422110552765,0.5593808585835323) (0.6984924623115577,0.5596623092992364) (0.7085427135678392,0.559772915955984) (0.7185929648241206,0.5618709578808954) (0.728643216080402,0.5629452078339907) (0.7386934673366834,0.5621257129932407) (0.7487437185929648,0.562046941229146) (0.7587939698492461,0.5621435700266979) (0.7688442211055275,0.5621435700266979) (0.7788944723618091,0.5619631184218269) (0.7889447236180905,0.5635627331863947) (0.7989949748743719,0.5635627331863947) (0.8090452261306533,0.5657643438336266) (0.8190954773869348,0.5651818764931159) (0.829145728643216,0.567393971364817) (0.8391959798994976,0.5684671085625986) (0.849246231155779,0.5690129470368985) (0.8592964824120602,0.5696622782854394) (0.8693467336683417,0.5695984561296008) (0.8793969849246231,0.5702829170449301) (0.8894472361809044,0.5702829170449301) (0.8994974874371859,0.5716590235941509) (0.9095477386934674,0.5722748136206549) (0.9195979899497486,0.5722748136206549) (0.9296482412060301,0.5722748136206549) (0.9396984924623116,0.5727504095832605) (0.949748743718593,0.5727504095832605) (0.9597989949748743,0.5734634195015806) (0.9698492462311559,0.5727095334859977) (0.9798994974874373,0.5735112562952546) (0.9899497487437185,0.5735112562952546) (1.0,0.5735112562952546) };
  \addplot[ours_gau, fill, fill opacity=0.18, draw=none, mark=none, forget plot] fill between [of=U_gau_grp and L_gau_grp];
  \addplot+[ours_gau, mark=none] coordinates {(-1.0,0.6957056113371068) (-0.9899497487437187,0.6912741428646194) (-0.9798994974874371,0.685555882848041) (-0.9698492462311558,0.6793837018110608) (-0.9597989949748745,0.6757955422067824) (-0.9497487437185929,0.6752660696171853) (-0.9396984924623116,0.6727191519887835) (-0.9296482412060303,0.6716951221223645) (-0.9195979899497487,0.6687287748877291) (-0.9095477386934674,0.6649541130466384) (-0.899497487437186,0.6610369108113369) (-0.8894472361809044,0.6549345817964265) (-0.8793969849246231,0.6533060868197137) (-0.8693467336683418,0.6485460598494602) (-0.8592964824120601,0.641685744953656) (-0.8492462311557789,0.6335029286010131) (-0.8391959798994972,0.6293576541345085) (-0.829145728643216,0.6254352439897043) (-0.8190954773869348,0.6240796451458686) (-0.8090452261306531,0.6200266822433542) (-0.7989949748743718,0.6230546996813249) (-0.7889447236180905,0.6200771237344999) (-0.7788944723618089,0.6174578797654898) (-0.7688442211055275,0.6161222441097658) (-0.7587939698492463,0.6090455516257136) (-0.7487437185929647,0.6100385561560013) (-0.7386934673366835,0.6076644956109829) (-0.7286432160804021,0.6024871537953694) (-0.7185929648241205,0.6007535243423351) (-0.7085427135678393,0.5970244311868468) (-0.6984924623115579,0.5982397818208214) (-0.6884422110552763,0.5969411004376498) (-0.678391959798995,0.5958529367751115) (-0.6683417085427137,0.5929923548196999) (-0.658291457286432,0.5896738512252265) (-0.6482412060301508,0.5882356012183126) (-0.6381909547738696,0.5845810276431761) (-0.6281407035175878,0.5824712072715759) (-0.6180904522613065,0.5775262657300924) (-0.6080402010050253,0.5740704105062926) (-0.5979899497487436,0.5745168255970488) (-0.5879396984924623,0.572612643123882) (-0.5778894472361811,0.5719085954243007) (-0.5678391959798994,0.5692456087974034) (-0.5577889447236181,0.5667494473234979) (-0.5477386934673368,0.5653019214164616) (-0.5376884422110553,0.5633733231232714) (-0.5276381909547739,0.5640402644066357) (-0.5175879396984923,0.5626456715431775) (-0.507537688442211,0.5585295566155422) (-0.4974874371859297,0.5566714054576691) (-0.4874371859296481,0.5525578534573955) (-0.47738693467336674,0.5519576984288717) (-0.4673366834170855,0.5520210679876386) (-0.45728643216080395,0.5517959090243029) (-0.44723618090452255,0.5491186490693487) (-0.4371859296482413,0.5482903658364214) (-0.42713567839195965,0.5454913190682373) (-0.4170854271356784,0.54066293089629) (-0.40703517587939714,0.5392958921805073) (-0.39698492462311546,0.5375296583950101) (-0.3869346733668342,0.5350226299904165) (-0.3768844221105529,0.5355912276091245) (-0.3668341708542712,0.5356122446988387) (-0.35678391959798994,0.5360562996015734) (-0.34673366834170866,0.5354226553496888) (-0.33668341708542704,0.5360393164871285) (-0.32663316582914576,0.536119620865609) (-0.3165829145728644,0.5367852871822436) (-0.30653266331658285,0.5348526257214081) (-0.29648241206030157,0.5354229762565268) (-0.2864321608040203,0.5327363878955448) (-0.27638190954773856,0.5309164160379399) (-0.26633165829145733,0.5294254704001599) (-0.2562814070351761,0.5323165017436263) (-0.24623115577889443,0.5323252368175888) (-0.23618090452261317,0.5322146351464899) (-0.2261306532663315,0.5321872370018653) (-0.2160804020100502,0.5308678660606381) (-0.206030150753769,0.5290092417904037) (-0.19597989949748723,0.5284664993820122) (-0.18592964824120595,0.5271123231592283) (-0.1758793969849247,0.5270944459750648) (-0.1658291457286431,0.5267080535593829) (-0.15577889447236182,0.524942721802945) (-0.14572864321608053,0.5255116079069688) (-0.13567839195979886,0.5270993737334603) (-0.12562814070351758,0.525694185876833) (-0.1155778894472363,0.5255330890482303) (-0.10552763819095466,0.5254437041243246) (-0.0954773869346734,0.5268073338693922) (-0.08542713567839208,0.5268438187893458) (-0.07537688442211042,0.5291157441338264) (-0.06532663316582912,0.5289581579864165) (-0.05527638190954787,0.5279470580942552) (-0.045226130653266215,0.5299212911722595) (-0.0351758793969849,0.5296304551951375) (-0.02512562814070368,0.5305640645401218) (-0.015075376884422013,0.5304415516684498) (-0.005025125628140743,0.5310786404447416) (0.005025125628140574,0.531726211645525) (0.015075376884422167,0.5351146374090702) (0.02512562814070352,0.5361744708486363) (0.03517587939698479,0.5332779029892781) (0.04522613065326644,0.5335434678198654) (0.0552763819095477,0.5338929591801962) (0.06532663316582898,0.5338210737805766) (0.07537688442211059,0.5345514104716247) (0.08542713567839183,0.5352101588454433) (0.09547738693467356,0.5352447541728318) (0.10552763819095481,0.5340209799380362) (0.11557788944723604,0.5342256499627136) (0.12562814070351774,0.5331189368990676) (0.13567839195979903,0.5337230414131829) (0.1457286432160803,0.5362949811821278) (0.15577889447236196,0.5366834552968126) (0.16582914572864324,0.536362200586055) (0.17587939698492452,0.5357533645043668) (0.1859296482412062,0.5344973912458881) (0.19597989949748745,0.5358185534969272) (0.2060301507537687,0.5381720057453657) (0.2160804020100504,0.5388405593025953) (0.22613065326633164,0.5409379644199036) (0.2361809045226129,0.5416801706032405) (0.2462311557788946,0.5437238275237625) (0.2562814070351759,0.5451609439465559) (0.26633165829145716,0.545273884064733) (0.2763819095477388,0.545960201886582) (0.28643216080402006,0.5459820816247328) (0.29648241206030135,0.547851372837647) (0.30653266331658297,0.5502857725387302) (0.31658291457286425,0.5504397411609548) (0.3266331658291456,0.5515700600565172) (0.33668341708542715,0.5515777997900135) (0.3467336683417085,0.5527896485834619) (0.3567839195979898,0.5542686513677255) (0.36683417085427145,0.5560649552818564) (0.37688442211055273,0.5568184811347283) (0.38693467336683435,0.5599404025419589) (0.3969849246231157,0.5594118391039431) (0.40703517587939697,0.5649273879353429) (0.41708542713567853,0.5670180992039409) (0.42713567839195987,0.5670180992039409) (0.4371859296482411,0.5670418701005465) (0.44723618090452283,0.5675690497957546) (0.45728643216080406,0.5692466722746113) (0.46733668341708534,0.5700205669058628) (0.47738693467336696,0.5716480479067103) (0.4874371859296483,0.5728211775855273) (0.4974874371859295,0.5742581883662582) (0.5075376884422113,0.5743293537214798) (0.5175879396984925,0.5758390980698718) (0.5276381909547737,0.576850539950431) (0.5376884422110554,0.5785350253215997) (0.5477386934673367,0.5781754446140396) (0.557788944723618,0.5820669021819826) (0.5678391959798996,0.5845561045668689) (0.5778894472361809,0.5851445449847069) (0.5879396984924622,0.5835457524673123) (0.5979899497487439,0.5852249733609831) (0.6080402010050251,0.5858153497256338) (0.6180904522613063,0.5857696241839352) (0.628140703517588,0.5860408091288171) (0.6381909547738693,0.5864747560184486) (0.6482412060301505,0.588659812360722) (0.6582914572864322,0.5892338955545562) (0.6683417085427136,0.5892455046606752) (0.6783919597989948,0.5914986953282131) (0.6884422110552765,0.5914509750626993) (0.6984924623115577,0.5921925027774558) (0.7085427135678392,0.5921459077420016) (0.7185929648241206,0.5945055983760825) (0.728643216080402,0.5957239884341802) (0.7386934673366834,0.5950695684806402) (0.7487437185929648,0.595172440628691) (0.7587939698492461,0.5951638366099096) (0.7688442211055275,0.5951638366099096) (0.7788944723618091,0.5952369249190214) (0.7889447236180905,0.5970708325341623) (0.7989949748743719,0.5970708325341623) (0.8090452261306533,0.599677087821267) (0.8190954773869348,0.5995308098179312) (0.829145728643216,0.6020510313221463) (0.8391959798994976,0.6033031941379513) (0.849246231155779,0.6039739031696268) (0.8592964824120602,0.6045731376467665) (0.8693467336683417,0.6045656489110679) (0.8793969849246231,0.6053168004978972) (0.8894472361809044,0.6053168004978972) (0.8994974874371859,0.6071883128129446) (0.9095477386934674,0.6078531111244989) (0.9195979899497486,0.6078531111244989) (0.9296482412060301,0.6078531111244989) (0.9396984924623116,0.6084612336034845) (0.949748743718593,0.6084612336034845) (0.9597989949748743,0.6091208702719185) (0.9698492462311559,0.6089130027084335) (0.9798994974874373,0.6093911030648029) (0.9899497487437185,0.6093911030648029) (1.0,0.6093911030648029) };
  \addplot+[ours_gau, mark=none, line width=0.4pt, dashed, forget plot] coordinates {(-1.0,0.41695186124509553) (1.0,0.41695186124509553)};
  \addplot+[ours_gau, mark=*, only marks, mark size=\PanelArgminSize, forget plot] coordinates {(-0.15577889447236182,0.524942721802945)};
  \addplot[name path=U_magsac_grp, draw=none, mark=none, forget plot] coordinates {(-0.9999999999999999,0.761257826286303) (-0.9899497487437185,0.7563411887936724) (-0.979899497487437,0.7568018191647978) (-0.9698492462311556,0.754281389438763) (-0.9597989949748744,0.7549506266492753) (-0.9497487437185927,0.7512234030688909) (-0.9396984924623114,0.7443753399941762) (-0.9296482412060302,0.7377734817437179) (-0.9195979899497485,0.7386791312887757) (-0.9095477386934673,0.7358562412157994) (-0.8994974874371859,0.7329452615380869) (-0.8894472361809044,0.7323771814754604) (-0.879396984924623,0.7277815284312426) (-0.8693467336683418,0.7274693618675353) (-0.85929648241206,0.7218190565900423) (-0.8492462311557788,0.7197024838543162) (-0.8391959798994976,0.7158623755706137) (-0.8291457286432159,0.7107595534517435) (-0.8190954773869346,0.7091426199201212) (-0.8090452261306533,0.7055641213772481) (-0.7989949748743717,0.6986516473439036) (-0.7889447236180904,0.6885694521088886) (-0.7788944723618088,0.6864067703227608) (-0.7688442211055275,0.6829348202063851) (-0.7587939698492462,0.6787778928221526) (-0.7487437185929645,0.6787953265660884) (-0.7386934673366833,0.6718971247358919) (-0.728643216080402,0.6692428310601252) (-0.7185929648241203,0.6580650536662671) (-0.708542713567839,0.6562302487312174) (-0.6984924623115578,0.652554412902214) (-0.688442211055276,0.6554355984625401) (-0.6783919597989948,0.6491882498948672) (-0.6683417085427136,0.6419436105075969) (-0.6582914572864319,0.6325796220215727) (-0.6482412060301506,0.6345527766233516) (-0.6381909547738694,0.6330545151818551) (-0.6281407035175877,0.6309956199656881) (-0.6180904522613064,0.6292666829034157) (-0.6080402010050252,0.6262741165661169) (-0.5979899497487434,0.6237149615337374) (-0.5879396984924622,0.6214809021547961) (-0.577889447236181,0.6203638650819352) (-0.5678391959798993,0.6164834692463821) (-0.5577889447236181,0.6126193829517212) (-0.5477386934673367,0.6072563917468288) (-0.537688442211055,0.6030462526740238) (-0.5276381909547738,0.6038143562148396) (-0.5175879396984925,0.601321223398953) (-0.5075376884422109,0.6024608636934503) (-0.4974874371859296,0.5987591883589547) (-0.4874371859296483,0.5979973542163396) (-0.4773869346733667,0.5980703556531147) (-0.4673366834170854,0.595653841626467) (-0.4572864321608041,0.5960806050334388) (-0.4472361809045225,0.5955745742185762) (-0.43718592964824116,0.5939478561125046) (-0.42713567839195987,0.5911041497103402) (-0.4170854271356782,0.5901171229661172) (-0.4070351758793969,0.5854292264670391) (-0.3969849246231153,0.588880116612731) (-0.386934673366834,0.5861000181251247) (-0.3768844221105528,0.5837235015254288) (-0.36683417085427106,0.5774776229520165) (-0.3567839195979898,0.5773719613866222) (-0.34673366834170855,0.5768365376353923) (-0.3366834170854269,0.5753321921267006) (-0.32663316582914564,0.5734052432954684) (-0.31658291457286436,0.5716808223059929) (-0.3065326633165827,0.5707115101977054) (-0.29648241206030135,0.5709591397269271) (-0.2864321608040202,0.5688909178284253) (-0.27638190954773845,0.5682000476775507) (-0.26633165829145716,0.5706069023940099) (-0.25628140703517593,0.5720402350605713) (-0.24623115577889426,0.5711753366764664) (-0.236180904522613,0.5695539099737659) (-0.22613065326633172,0.5680108377021306) (-0.21608040201005002,0.5685513549527202) (-0.20603015075376882,0.5662283949917816) (-0.19597989949748745,0.5632361761383887) (-0.1859296482412059,0.5646202172137194) (-0.1758793969849246,0.5645480861981351) (-0.16582914572864324,0.5643932231604019) (-0.15577889447236168,0.5634320710943445) (-0.14572864321608042,0.5646403711295448) (-0.13567839195979906,0.561491863867236) (-0.12562814070351747,0.561376279487703) (-0.11557788944723613,0.561174527493654) (-0.10552763819095484,0.5596194970092482) (-0.09547738693467316,0.5616371810809621) (-0.08542713567839195,0.562032475475334) (-0.07537688442211071,0.5595242261447065) (-0.06532663316582901,0.5601184058674175) (-0.0552763819095477,0.5640617966478331) (-0.04522613065326643,0.5633120280714075) (-0.03517587939698485,0.5627787316076248) (-0.025125628140703526,0.5613604038646761) (-0.015075376884422212,0.5609676032701937) (-0.005025125628140596,0.5589375871186385) (0.00502512562814067,0.56100077434251) (0.015075376884422353,0.562966051278758) (0.02512562814070361,0.5634435214896936) (0.035175879396984876,0.5625561071172293) (0.04522613065326653,0.5640819775352629) (0.05527638190954779,0.5624871318674319) (0.06532663316582915,0.5619186549819741) (0.07537688442211074,0.5632250986476989) (0.08542713567839207,0.5657275199978472) (0.09547738693467332,0.5668478633200955) (0.10552763819095497,0.5667738132548319) (0.11557788944723627,0.5634183936168569) (0.12562814070351755,0.5631454195175913) (0.13567839195979922,0.5657764310257025) (0.14572864321608045,0.5641493455771641) (0.1557788944723617,0.5681890538583841) (0.16582914572864335,0.5704446434616325) (0.17587939698492464,0.5709233999949791) (0.1859296482412059,0.5695152520270719) (0.19597989949748745,0.5680706758095596) (0.2060301507537689,0.5668149957589458) (0.21608040201005027,0.5666317370540941) (0.22613065326633164,0.5671866180383811) (0.236180904522613,0.5673525246645001) (0.24623115577889454,0.5694184163125708) (0.2562814070351758,0.5706293476165921) (0.26633165829145733,0.5704368590522529) (0.2763819095477387,0.5717865312378324) (0.28643216080402,0.5713064562960182) (0.2964824120603015,0.5717992746010001) (0.30653266331658297,0.5736906519100566) (0.3165829145728644,0.5738996633219708) (0.32663316582914576,0.5737582299857797) (0.33668341708542715,0.5744942743884656) (0.34673366834170866,0.5750484214188828) (0.35678391959798994,0.5757962289678734) (0.36683417085427145,0.5755479527722617) (0.3768844221105529,0.5766987805999217) (0.3869346733668341,0.5782601201201031) (0.3969849246231155,0.5790265871460842) (0.407035175879397,0.5791995789261091) (0.4170854271356783,0.5797757489192454) (0.4271356783919598,0.5812250878379547) (0.4371859296482413,0.5819981177917415) (0.44723618090452255,0.5850476507106761) (0.457286432160804,0.587118768926336) (0.4673366834170855,0.5859604785859386) (0.4773869346733668,0.5877085385343955) (0.48743718592964824,0.5920534184079232) (0.4974874371859297,0.5942045616875372) (0.5075376884422111,0.5933779754321685) (0.5175879396984924,0.5946233534372786) (0.5276381909547738,0.5970956908388009) (0.5376884422110554,0.5977495574383747) (0.5477386934673366,0.5977535501559756) (0.5577889447236181,0.5992889686125575) (0.5678391959798995,0.6007977259176189) (0.5778894472361809,0.6012708540244879) (0.5879396984924623,0.601481340618799) (0.5979899497487438,0.6058029681763918) (0.6080402010050251,0.6075054746668789) (0.6180904522613065,0.6072674217059699) (0.6281407035175879,0.6084081705572892) (0.6381909547738693,0.6104545518096391) (0.6482412060301507,0.6142257929829559) (0.6582914572864322,0.6144711976012521) (0.6683417085427135,0.6148687606281698) (0.678391959798995,0.6157627162320937) (0.6884422110552765,0.6158474089549401) (0.6984924623115577,0.6150530209572707) (0.7085427135678392,0.6175357946336144) (0.7185929648241206,0.6189224645067971) (0.7286432160804021,0.6196574410934239) (0.7386934673366834,0.6218634223626942) (0.7487437185929647,0.6225387531371241) (0.7587939698492463,0.6269481841222565) (0.7688442211055275,0.628253549118047) (0.778894472361809,0.6282796143990952) (0.7889447236180905,0.630353531892799) (0.7989949748743718,0.6285043809341523) (0.8090452261306532,0.6293000398294045) (0.8190954773869347,0.629830816389569) (0.829145728643216,0.631779490019091) (0.8391959798994975,0.6341523965815112) (0.8492462311557789,0.6348010647441525) (0.8592964824120601,0.6401335783011252) (0.8693467336683416,0.6423076306875088) (0.8793969849246231,0.6433080076504246) (0.8894472361809044,0.6433212850031409) (0.8994974874371858,0.6465965135442081) (0.9095477386934674,0.6492428091171714) (0.9195979899497488,0.6493210529869277) (0.9296482412060301,0.6522113261724375) (0.9396984924623115,0.6536988519443439) (0.949748743718593,0.6543112956738226) (0.9597989949748743,0.6550884114700617) (0.9698492462311558,0.655843884125347) (0.9798994974874372,0.6579573014667863) (0.9899497487437185,0.6588031144903365) (0.9999999999999999,0.6627892851658024) };
  \addplot[name path=L_magsac_grp, draw=none, mark=none, forget plot] coordinates {(-0.9999999999999999,0.6682387698930343) (-0.9899497487437185,0.6635033970870992) (-0.979899497487437,0.663298926290155) (-0.9698492462311556,0.6616501479134616) (-0.9597989949748744,0.6624110045754595) (-0.9497487437185927,0.6561632894028204) (-0.9396984924623114,0.6506423180465173) (-0.9296482412060302,0.6458867533240377) (-0.9195979899497485,0.6479567707293961) (-0.9095477386934673,0.6475237738417119) (-0.8994974874371859,0.6449044158641013) (-0.8894472361809044,0.6394427367317616) (-0.879396984924623,0.6344366879486718) (-0.8693467336683418,0.6349236625263909) (-0.85929648241206,0.628519441709979) (-0.8492462311557788,0.6235751839946986) (-0.8391959798994976,0.6198833643461195) (-0.8291457286432159,0.6172348699899908) (-0.8190954773869346,0.6151527822122624) (-0.8090452261306533,0.6122421890164189) (-0.7989949748743717,0.6045975258832165) (-0.7889447236180904,0.5980759205215208) (-0.7788944723618088,0.5956727182994621) (-0.7688442211055275,0.5938052291706238) (-0.7587939698492462,0.59147330679888) (-0.7487437185929645,0.5892998862563641) (-0.7386934673366833,0.5866569297426759) (-0.728643216080402,0.5850032181901871) (-0.7185929648241203,0.5784272397842509) (-0.708542713567839,0.5778928702603267) (-0.6984924623115578,0.5763137414617225) (-0.688442211055276,0.5758247153210979) (-0.6783919597989948,0.573891372854974) (-0.6683417085427136,0.5688256806783016) (-0.6582914572864319,0.5637695069514153) (-0.6482412060301506,0.5658948089903018) (-0.6381909547738694,0.5661123287538578) (-0.6281407035175877,0.5662106530140705) (-0.6180904522613064,0.5644454954103073) (-0.6080402010050252,0.5611687618423252) (-0.5979899497487434,0.5560646023325526) (-0.5879396984924622,0.5534519853976311) (-0.577889447236181,0.5511188011815258) (-0.5678391959798993,0.5462000455149869) (-0.5577889447236181,0.5439845956319957) (-0.5477386934673367,0.5396512162982916) (-0.537688442211055,0.5360097654211037) (-0.5276381909547738,0.5385161129795916) (-0.5175879396984925,0.5370404963588907) (-0.5075376884422109,0.5374098179173099) (-0.4974874371859296,0.5340801632350072) (-0.4874371859296483,0.5335143813661208) (-0.4773869346733667,0.5345617068477225) (-0.4673366834170854,0.5320306871588099) (-0.4572864321608041,0.5336302331669724) (-0.4472361809045225,0.5315919726065353) (-0.43718592964824116,0.5295592050312977) (-0.42713567839195987,0.5278554337686803) (-0.4170854271356782,0.5258009320124922) (-0.4070351758793969,0.5220969398426927) (-0.3969849246231153,0.5234338562747473) (-0.386934673366834,0.5218119625257888) (-0.3768844221105528,0.5210644343664875) (-0.36683417085427106,0.5143425284478732) (-0.3567839195979898,0.5138815596981696) (-0.34673366834170855,0.5127376018638812) (-0.3366834170854269,0.5117197608879768) (-0.32663316582914564,0.5085542799141962) (-0.31658291457286436,0.5057928346729627) (-0.3065326633165827,0.5036717456233832) (-0.29648241206030135,0.5033566266280802) (-0.2864321608040202,0.5017957514389163) (-0.27638190954773845,0.5022510546735713) (-0.26633165829145716,0.5059297195368497) (-0.25628140703517593,0.5070670980920462) (-0.24623115577889426,0.5067114898913144) (-0.236180904522613,0.5046118956949819) (-0.22613065326633172,0.5021968637585114) (-0.21608040201005002,0.5023277363806723) (-0.20603015075376882,0.5008944370116631) (-0.19597989949748745,0.498496432243264) (-0.1859296482412059,0.4989814905052496) (-0.1758793969849246,0.500091320476671) (-0.16582914572864324,0.5000301595797849) (-0.15577889447236168,0.49923080084707244) (-0.14572864321608042,0.500516270003572) (-0.13567839195979906,0.4981826139076434) (-0.12562814070351747,0.498449908677493) (-0.11557788944723613,0.49865069096366266) (-0.10552763819095484,0.49719954403128497) (-0.09547738693467316,0.49784739176570264) (-0.08542713567839195,0.49864446799623335) (-0.07537688442211071,0.4960277176103927) (-0.06532663316582901,0.49723318124811333) (-0.0552763819095477,0.4995072203891991) (-0.04522613065326643,0.49929187734855063) (-0.03517587939698485,0.4992958310093648) (-0.025125628140703526,0.49889037358940697) (-0.015075376884422212,0.4994729291265062) (-0.005025125628140596,0.4982786382005006) (0.00502512562814067,0.49985985818970946) (0.015075376884422353,0.5005377783909017) (0.02512562814070361,0.5004025288145559) (0.035175879396984876,0.4994778731564461) (0.04522613065326653,0.49984592680443124) (0.05527638190954779,0.4977837549408052) (0.06532663316582915,0.49712217893705896) (0.07537688442211074,0.4987139870552267) (0.08542713567839207,0.5003874568961396) (0.09547738693467332,0.5017889209796365) (0.10552763819095497,0.5003221442820374) (0.11557788944723627,0.49775943293720326) (0.12562814070351755,0.4978617398077854) (0.13567839195979922,0.4995985606504525) (0.14572864321608045,0.4996748613008565) (0.1557788944723617,0.5033753330751743) (0.16582914572864335,0.5060443277193901) (0.17587939698492464,0.5066222354138238) (0.1859296482412059,0.505638131885217) (0.19597989949748745,0.5038230800676365) (0.2060301507537689,0.5028645615859427) (0.21608040201005027,0.503133210871173) (0.22613065326633164,0.5040145055563579) (0.236180904522613,0.5038714776394941) (0.24623115577889454,0.505500112735463) (0.2562814070351758,0.5063280239243756) (0.26633165829145733,0.5056695591839199) (0.2763819095477387,0.5063002697231133) (0.28643216080402,0.5057028892593101) (0.2964824120603015,0.506444770011077) (0.30653266331658297,0.5086524816207587) (0.3165829145728644,0.5079109764444358) (0.32663316582914576,0.5079591556303906) (0.33668341708542715,0.5103720596507482) (0.34673366834170866,0.5124151630134571) (0.35678391959798994,0.5134579018792086) (0.36683417085427145,0.5134837217491505) (0.3768844221105529,0.5152831304727745) (0.3869346733668341,0.5175673917830043) (0.3969849246231155,0.5181400553937845) (0.407035175879397,0.5176052573430447) (0.4170854271356783,0.5182111067929591) (0.4271356783919598,0.520833011250364) (0.4371859296482413,0.5222227157771941) (0.44723618090452255,0.5260097735466116) (0.457286432160804,0.5288513214005992) (0.4673366834170855,0.5281934110480465) (0.4773869346733668,0.5306248276992177) (0.48743718592964824,0.5340637515303764) (0.4974874371859297,0.5360428500838964) (0.5075376884422111,0.5359354726147446) (0.5175879396984924,0.5368638439971386) (0.5276381909547738,0.5388529338354707) (0.5376884422110554,0.5394166955623236) (0.5477386934673366,0.5395085608868773) (0.5577889447236181,0.5415513818400923) (0.5678391959798995,0.5436561682299623) (0.5778894472361809,0.5432896501485239) (0.5879396984924623,0.5435789582440257) (0.5979899497487438,0.5452173429910682) (0.6080402010050251,0.5468812152329228) (0.6180904522613065,0.5470101887997225) (0.6281407035175879,0.5487906951263475) (0.6381909547738693,0.5504051170579154) (0.6482412060301507,0.5535053389708134) (0.6582914572864322,0.5537149411893062) (0.6683417085427135,0.5540382010796396) (0.678391959798995,0.5544458256711967) (0.6884422110552765,0.5543369517433899) (0.6984924623115577,0.5527656325326658) (0.7085427135678392,0.5549744393400002) (0.7185929648241206,0.5559993117745022) (0.7286432160804021,0.5563683050338013) (0.7386934673366834,0.5587272183279928) (0.7487437185929647,0.5592394819923162) (0.7587939698492463,0.5630650781222192) (0.7688442211055275,0.5639158324955149) (0.778894472361809,0.5639695786071732) (0.7889447236180905,0.5651010272126044) (0.7989949748743718,0.5632115110048387) (0.8090452261306532,0.5640629539877694) (0.8190954773869347,0.5638422410805912) (0.829145728643216,0.565672110769178) (0.8391959798994975,0.5680766181023742) (0.8492462311557789,0.5670717028951974) (0.8592964824120601,0.5713089409978253) (0.8693467336683416,0.573112513116562) (0.8793969849246231,0.5735435108307533) (0.8894472361809044,0.573593253585347) (0.8994974874371858,0.5757946557896773) (0.9095477386934674,0.5779888314126502) (0.9195979899497488,0.5778845768159612) (0.9296482412060301,0.5802623745643999) (0.9396984924623115,0.5814055421769322) (0.949748743718593,0.5821707341360391) (0.9597989949748743,0.5827281284771495) (0.9698492462311558,0.5833736564836315) (0.9798994974874372,0.585324161092335) (0.9899497487437185,0.5859667187978089) (0.9999999999999999,0.5888304094305248) };
  \addplot[ours_magsac, fill, fill opacity=0.18, draw=none, mark=none, forget plot] fill between [of=U_magsac_grp and L_magsac_grp];
  \addplot+[ours_magsac, mark=none] coordinates {(-0.9999999999999999,0.7147482980896687) (-0.9899497487437185,0.7099222929403858) (-0.979899497487437,0.7100503727274764) (-0.9698492462311556,0.7079657686761123) (-0.9597989949748744,0.7086808156123674) (-0.9497487437185927,0.7036933462358557) (-0.9396984924623114,0.6975088290203467) (-0.9296482412060302,0.6918301175338778) (-0.9195979899497485,0.6933179510090859) (-0.9095477386934673,0.6916900075287556) (-0.8994974874371859,0.6889248387010941) (-0.8894472361809044,0.685909959103611) (-0.879396984924623,0.6811091081899572) (-0.8693467336683418,0.6811965121969631) (-0.85929648241206,0.6751692491500106) (-0.8492462311557788,0.6716388339245074) (-0.8391959798994976,0.6678728699583666) (-0.8291457286432159,0.6639972117208671) (-0.8190954773869346,0.6621477010661918) (-0.8090452261306533,0.6589031551968335) (-0.7989949748743717,0.65162458661356) (-0.7889447236180904,0.6433226863152047) (-0.7788944723618088,0.6410397443111114) (-0.7688442211055275,0.6383700246885045) (-0.7587939698492462,0.6351255998105163) (-0.7487437185929645,0.6340476064112263) (-0.7386934673366833,0.6292770272392839) (-0.728643216080402,0.6271230246251561) (-0.7185929648241203,0.618246146725259) (-0.708542713567839,0.617061559495772) (-0.6984924623115578,0.6144340771819683) (-0.688442211055276,0.615630156891819) (-0.6783919597989948,0.6115398113749206) (-0.6683417085427136,0.6053846455929492) (-0.6582914572864319,0.598174564486494) (-0.6482412060301506,0.6002237928068267) (-0.6381909547738694,0.5995834219678564) (-0.6281407035175877,0.5986031364898793) (-0.6180904522613064,0.5968560891568615) (-0.6080402010050252,0.593721439204221) (-0.5979899497487434,0.589889781933145) (-0.5879396984924622,0.5874664437762136) (-0.577889447236181,0.5857413331317305) (-0.5678391959798993,0.5813417573806845) (-0.5577889447236181,0.5783019892918585) (-0.5477386934673367,0.5734538040225602) (-0.537688442211055,0.5695280090475637) (-0.5276381909547738,0.5711652345972156) (-0.5175879396984925,0.5691808598789219) (-0.5075376884422109,0.5699353408053801) (-0.4974874371859296,0.5664196757969809) (-0.4874371859296483,0.5657558677912302) (-0.4773869346733667,0.5663160312504186) (-0.4673366834170854,0.5638422643926384) (-0.4572864321608041,0.5648554191002056) (-0.4472361809045225,0.5635832734125558) (-0.43718592964824116,0.5617535305719011) (-0.42713567839195987,0.5594797917395102) (-0.4170854271356782,0.5579590274893047) (-0.4070351758793969,0.5537630831548659) (-0.3969849246231153,0.5561569864437391) (-0.386934673366834,0.5539559903254567) (-0.3768844221105528,0.5523939679459582) (-0.36683417085427106,0.5459100756999449) (-0.3567839195979898,0.5456267605423959) (-0.34673366834170855,0.5447870697496368) (-0.3366834170854269,0.5435259765073387) (-0.32663316582914564,0.5409797616048323) (-0.31658291457286436,0.5387368284894778) (-0.3065326633165827,0.5371916279105443) (-0.29648241206030135,0.5371578831775037) (-0.2864321608040202,0.5353433346336708) (-0.27638190954773845,0.535225551175561) (-0.26633165829145716,0.5382683109654298) (-0.25628140703517593,0.5395536665763088) (-0.24623115577889426,0.5389434132838904) (-0.236180904522613,0.5370829028343739) (-0.22613065326633172,0.535103850730321) (-0.21608040201005002,0.5354395456666963) (-0.20603015075376882,0.5335614160017224) (-0.19597989949748745,0.5308663041908264) (-0.1859296482412059,0.5318008538594845) (-0.1758793969849246,0.5323197033374031) (-0.16582914572864324,0.5322116913700934) (-0.15577889447236168,0.5313314359707084) (-0.14572864321608042,0.5325783205665584) (-0.13567839195979906,0.5298372388874397) (-0.12562814070351747,0.529913094082598) (-0.11557788944723613,0.5299126092286583) (-0.10552763819095484,0.5284095205202666) (-0.09547738693467316,0.5297422864233323) (-0.08542713567839195,0.5303384717357836) (-0.07537688442211071,0.5277759718775497) (-0.06532663316582901,0.5286757935577654) (-0.0552763819095477,0.5317845085185161) (-0.04522613065326643,0.531301952709979) (-0.03517587939698485,0.5310372813084948) (-0.025125628140703526,0.5301253887270415) (-0.015075376884422212,0.5302202661983499) (-0.005025125628140596,0.5286081126595695) (0.00502512562814067,0.5304303162661097) (0.015075376884422353,0.5317519148348299) (0.02512562814070361,0.5319230251521248) (0.035175879396984876,0.5310169901368377) (0.04522613065326653,0.531963952169847) (0.05527638190954779,0.5301354434041186) (0.06532663316582915,0.5295204169595166) (0.07537688442211074,0.5309695428514628) (0.08542713567839207,0.5330574884469934) (0.09547738693467332,0.534318392149866) (0.10552763819095497,0.5335479787684346) (0.11557788944723627,0.5305889132770301) (0.12562814070351755,0.5305035796626884) (0.13567839195979922,0.5326874958380775) (0.14572864321608045,0.5319121034390103) (0.1557788944723617,0.5357821934667792) (0.16582914572864335,0.5382444855905113) (0.17587939698492464,0.5387728177044014) (0.1859296482412059,0.5375766919561444) (0.19597989949748745,0.5359468779385981) (0.2060301507537689,0.5348397786724443) (0.21608040201005027,0.5348824739626336) (0.22613065326633164,0.5356005617973695) (0.236180904522613,0.5356120011519971) (0.24623115577889454,0.5374592645240169) (0.2562814070351758,0.5384786857704839) (0.26633165829145733,0.5380532091180864) (0.2763819095477387,0.5390434004804728) (0.28643216080402,0.5385046727776641) (0.2964824120603015,0.5391220223060386) (0.30653266331658297,0.5411715667654077) (0.3165829145728644,0.5409053198832033) (0.32663316582914576,0.5408586928080852) (0.33668341708542715,0.5424331670196069) (0.34673366834170866,0.54373179221617) (0.35678391959798994,0.544627065423541) (0.36683417085427145,0.5445158372607061) (0.3768844221105529,0.5459909555363481) (0.3869346733668341,0.5479137559515537) (0.3969849246231155,0.5485833212699344) (0.407035175879397,0.5484024181345769) (0.4170854271356783,0.5489934278561023) (0.4271356783919598,0.5510290495441593) (0.4371859296482413,0.5521104167844678) (0.44723618090452255,0.5555287121286439) (0.457286432160804,0.5579850451634676) (0.4673366834170855,0.5570769448169925) (0.4773869346733668,0.5591666831168066) (0.48743718592964824,0.5630585849691498) (0.4974874371859297,0.5651237058857168) (0.5075376884422111,0.5646567240234566) (0.5175879396984924,0.5657435987172086) (0.5276381909547738,0.5679743123371358) (0.5376884422110554,0.5685831265003491) (0.5477386934673366,0.5686310555214265) (0.5577889447236181,0.5704201752263249) (0.5678391959798995,0.5722269470737906) (0.5778894472361809,0.5722802520865059) (0.5879396984924623,0.5725301494314123) (0.5979899497487438,0.57551015558373) (0.6080402010050251,0.5771933449499008) (0.6180904522613065,0.5771388052528462) (0.6281407035175879,0.5785994328418184) (0.6381909547738693,0.5804298344337773) (0.6482412060301507,0.5838655659768847) (0.6582914572864322,0.5840930693952792) (0.6683417085427135,0.5844534808539047) (0.678391959798995,0.5851042709516452) (0.6884422110552765,0.585092180349165) (0.6984924623115577,0.5839093267449682) (0.7085427135678392,0.5862551169868073) (0.7185929648241206,0.5874608881406497) (0.7286432160804021,0.5880128730636126) (0.7386934673366834,0.5902953203453435) (0.7487437185929647,0.5908891175647202) (0.7587939698492463,0.5950066311222378) (0.7688442211055275,0.596084690806781) (0.778894472361809,0.5961245965031342) (0.7889447236180905,0.5977272795527017) (0.7989949748743718,0.5958579459694955) (0.8090452261306532,0.596681496908587) (0.8190954773869347,0.5968365287350801) (0.829145728643216,0.5987258003941345) (0.8391959798994975,0.6011145073419427) (0.8492462311557789,0.600936383819675) (0.8592964824120601,0.6057212596494752) (0.8693467336683416,0.6077100719020354) (0.8793969849246231,0.608425759240589) (0.8894472361809044,0.6084572692942439) (0.8994974874371858,0.6111955846669427) (0.9095477386934674,0.6136158202649108) (0.9195979899497488,0.6136028149014444) (0.9296482412060301,0.6162368503684187) (0.9396984924623115,0.617552197060638) (0.949748743718593,0.6182410149049309) (0.9597989949748743,0.6189082699736056) (0.9698492462311558,0.6196087703044892) (0.9798994974874372,0.6216407312795607) (0.9899497487437185,0.6223849166440727) (0.9999999999999999,0.6258098472981636) };
  \addplot+[ours_magsac, mark=none, line width=0.4pt, dashed, forget plot] coordinates {(-1.0,0.41472775796327965) (1.0,0.41472775796327965)};
  \addplot+[ours_magsac, mark=*, only marks, mark size=\PanelArgminSize, forget plot] coordinates {(-0.07537688442211071,0.5277759718775497)};
  \addplot[name path=U_jeffreys_grp, draw=none, mark=none, forget plot] coordinates {(-0.9999999999999999,0.6388047407207913) (-0.9899497487437185,0.6356178969669165) (-0.979899497487437,0.6311985157211463) (-0.9698492462311556,0.6316271245701034) (-0.9597989949748744,0.6271110200658) (-0.9497487437185927,0.6250989932948078) (-0.9396984924623114,0.6244454908217567) (-0.9296482412060302,0.6240233060365188) (-0.9195979899497485,0.6218324110703094) (-0.9095477386934673,0.6209685610253781) (-0.8994974874371859,0.6217826951967742) (-0.8894472361809044,0.6216397605153932) (-0.879396984924623,0.6203743777670787) (-0.8693467336683418,0.6192270865197905) (-0.85929648241206,0.6164777129093527) (-0.8492462311557789,0.6149245204193161) (-0.8391959798994976,0.6136658286559332) (-0.8291457286432159,0.6124658609941827) (-0.8190954773869346,0.6118137653596242) (-0.809045226130653,0.6112869300332393) (-0.7989949748743717,0.6091395130860906) (-0.7889447236180904,0.6071016006407496) (-0.7788944723618088,0.6062080787206003) (-0.7688442211055275,0.6054984065880313) (-0.7587939698492462,0.6045469915038191) (-0.7487437185929646,0.6001618710721791) (-0.7386934673366833,0.5960215537438438) (-0.728643216080402,0.591361681485341) (-0.7185929648241204,0.5918801499910962) (-0.708542713567839,0.592665816917426) (-0.6984924623115578,0.5889133588055894) (-0.6884422110552761,0.5877629367183982) (-0.6783919597989948,0.5856580709188266) (-0.6683417085427136,0.5843786509389166) (-0.658291457286432,0.5862286188238209) (-0.6482412060301507,0.5860620695937047) (-0.6381909547738694,0.5847516219770414) (-0.6281407035175878,0.5838477233719376) (-0.6180904522613064,0.5827947069684588) (-0.6080402010050252,0.5822021073107437) (-0.5979899497487435,0.5801053536328309) (-0.5879396984924623,0.5771663455632007) (-0.577889447236181,0.5750244994310749) (-0.5678391959798993,0.5748537150434788) (-0.5577889447236181,0.5757364233251323) (-0.5477386934673366,0.5764596803509822) (-0.5376884422110553,0.5757374762751536) (-0.5276381909547738,0.57416542914915) (-0.5175879396984924,0.5741618508813234) (-0.5075376884422109,0.5739651987217632) (-0.49748743718592964,0.5736061217713658) (-0.4874371859296482,0.5729305100238705) (-0.4773869346733667,0.5738885300455692) (-0.46733668341708545,0.5726718362800904) (-0.45728643216080395,0.5738067299130043) (-0.44723618090452255,0.5764272472213414) (-0.4371859296482412,0.5772555273271838) (-0.42713567839195976,0.5768498166903068) (-0.41708542713567825,0.5772721241195826) (-0.4070351758793971,0.5767632831014615) (-0.3969849246231155,0.5786006155573237) (-0.38693467336683407,0.5778563499203087) (-0.3768844221105528,0.5763418789097768) (-0.36683417085427134,0.5737917230775752) (-0.35678391959798983,0.5735002545840189) (-0.34673366834170843,0.5733687707917954) (-0.3366834170854271,0.5741198463301044) (-0.32663316582914564,0.5751586771789143) (-0.31658291457286425,0.5731738790181193) (-0.30653266331658297,0.5718805603658109) (-0.29648241206030146,0.568391112816898) (-0.28643216080402,0.5675392828114172) (-0.2763819095477387,0.5684941657205281) (-0.2663316582914572,0.5700029534108866) (-0.2562814070351758,0.5683776775859074) (-0.2462311557788945,0.5669564313581765) (-0.236180904522613,0.5643216287139894) (-0.22613065326633158,0.5654827682799104) (-0.21608040201005027,0.5656840549963107) (-0.20603015075376882,0.5668531577624237) (-0.19597989949748737,0.5661050395926898) (-0.18592964824120603,0.5694176217421923) (-0.1758793969849246,0.5715968093275645) (-0.1658291457286431,0.5726469504649044) (-0.15577889447236168,0.5717906736545919) (-0.14572864321608042,0.5697870267573475) (-0.13567839195979894,0.5691819606853684) (-0.12562814070351747,0.5683783373819454) (-0.11557788944723618,0.5703098747417302) (-0.10552763819095472,0.5685147650252085) (-0.09547738693467328,0.5670129924101428) (-0.08542713567839201,0.568839957994731) (-0.07537688442211048,0.5675798104855181) (-0.06532663316582907,0.5668450118297941) (-0.055276381909547756,0.567705744668852) (-0.04522613065326632,0.5677950261301118) (-0.03517587939698485,0.5689096779701922) (-0.025125628140703578,0.570182690619524) (-0.015075376884422061,0.5690241156957331) (-0.005025125628140596,0.5690569996965021) (0.00502512562814067,0.5721246544801913) (0.015075376884422073,0.5726437958279481) (0.02512562814070361,0.5743074969457298) (0.035175879396985056,0.5757496786897848) (0.045226130653266264,0.5764398352487915) (0.05527638190954779,0.5762434212257093) (0.06532663316582923,0.5750072309557286) (0.0753768844221105,0.5737769171191431) (0.08542713567839198,0.5745408403542978) (0.09547738693467347,0.5753615991322684) (0.10552763819095473,0.5765575274981507) (0.11557788944723618,0.5761519114117168) (0.1256281407035177,0.5762570272541339) (0.13567839195979894,0.5762005335782536) (0.1457286432160804,0.5764824765459958) (0.1557788944723619,0.5776293947156447) (0.16582914572864316,0.577988295709343) (0.17587939698492458,0.577446001343426) (0.18592964824120614,0.5779404109563524) (0.19597989949748731,0.5773925541792072) (0.20603015075376882,0.5758767960796297) (0.21608040201005027,0.5764994909615477) (0.2261306532663318,0.5758541877126626) (0.236180904522613,0.5758232606929897) (0.24623115577889448,0.5764859148660169) (0.25628140703517593,0.5788573135050216) (0.2663316582914572,0.5775215261687973) (0.27638190954773867,0.5769466407472204) (0.2864321608040202,0.5771502553968314) (0.2964824120603014,0.5808919505408507) (0.30653266331658285,0.580225749661608) (0.31658291457286425,0.5808158236551728) (0.32663316582914576,0.5801419459683709) (0.33668341708542715,0.5792336873226284) (0.3467336683417085,0.5780370916303875) (0.35678391959799,0.5780370916303875) (0.36683417085427134,0.5785934088318642) (0.3768844221105528,0.5784520824381176) (0.3869346733668342,0.5780659609273192) (0.3969849246231155,0.5775752456964719) (0.407035175879397,0.5775443163140404) (0.4170854271356784,0.5769985707001328) (0.42713567839195976,0.5761570129808427) (0.43718592964824116,0.5761570129808427) (0.4472361809045226,0.5760243021791556) (0.45728643216080406,0.5766780147571715) (0.46733668341708545,0.5754947291030027) (0.4773869346733668,0.5761245459547) (0.4874371859296483,0.5771047211266347) (0.49748743718592964,0.5764021794081283) (0.5075376884422109,0.5770234208391618) (0.5175879396984924,0.578064274233594) (0.5276381909547738,0.5766114878660967) (0.5376884422110553,0.5758019212997989) (0.5477386934673366,0.5766825750795973) (0.5577889447236181,0.5763798054217262) (0.5678391959798995,0.5772821751628957) (0.5778894472361809,0.5766571800219794) (0.5879396984924622,0.5776444469129731) (0.5979899497487436,0.5779062360473747) (0.6080402010050252,0.5785926758045199) (0.6180904522613064,0.5794756320977809) (0.6281407035175879,0.579441554086949) (0.6381909547738693,0.5788210221840582) (0.6482412060301508,0.5788210221840582) (0.6582914572864322,0.5793312350083756) (0.6683417085427136,0.5800540313354745) (0.678391959798995,0.578914236254282) (0.6884422110552764,0.5805455682184447) (0.6984924623115577,0.5801535733766726) (0.7085427135678392,0.5788607113272345) (0.7185929648241205,0.5783765792772355) (0.728643216080402,0.5783692984171475) (0.7386934673366834,0.5795498827625509) (0.7487437185929648,0.5787074820338232) (0.7587939698492462,0.579278263818261) (0.7688442211055276,0.5805145926246529) (0.778894472361809,0.5798577839209893) (0.7889447236180904,0.5799465992666005) (0.7989949748743718,0.5798487099815666) (0.8090452261306532,0.5803761577807368) (0.8190954773869347,0.5805351419637161) (0.829145728643216,0.5799270501119562) (0.8391959798994975,0.5792878172171745) (0.8492462311557789,0.5799467605813612) (0.8592964824120602,0.5799467605813612) (0.8693467336683416,0.5794309742979398) (0.8793969849246231,0.5794028266426899) (0.8894472361809045,0.5788868184712717) (0.8994974874371859,0.5776054329908377) (0.9095477386934674,0.5777352204947428) (0.9195979899497487,0.5768311839703237) (0.9296482412060301,0.5772970102953802) (0.9396984924623116,0.577856045358913) (0.949748743718593,0.577856045358913) (0.9597989949748743,0.5768716091870477) (0.9698492462311558,0.5768716091870477) (0.9798994974874372,0.5763774021912546) (0.9899497487437185,0.5763774021912546) (0.9999999999999999,0.5763774021912546) };
  \addplot[name path=L_jeffreys_grp, draw=none, mark=none, forget plot] coordinates {(-0.9999999999999999,0.5687536200008314) (-0.9899497487437185,0.5648857992525339) (-0.979899497487437,0.5618940100426096) (-0.9698492462311556,0.5633650583866842) (-0.9597989949748744,0.5591997780251433) (-0.9497487437185927,0.5549733264492772) (-0.9396984924623114,0.5550506449469537) (-0.9296482412060302,0.555600696197265) (-0.9195979899497485,0.5533330511761761) (-0.9095477386934673,0.5533065248855831) (-0.8994974874371859,0.5527897177985104) (-0.8894472361809044,0.55364002728137) (-0.879396984924623,0.552193284618353) (-0.8693467336683418,0.5504000127269432) (-0.85929648241206,0.5506828259909765) (-0.8492462311557789,0.5483238615610582) (-0.8391959798994976,0.5455302910117429) (-0.8291457286432159,0.544669103829069) (-0.8190954773869346,0.5440788763118899) (-0.809045226130653,0.5425169412827684) (-0.7989949748743717,0.539784238538122) (-0.7889447236180904,0.5376055323510368) (-0.7788944723618088,0.5369221278923966) (-0.7688442211055275,0.5373090758928669) (-0.7587939698492462,0.5367603252521931) (-0.7487437185929646,0.5330080687219484) (-0.7386934673366833,0.5285810674573426) (-0.728643216080402,0.5256605391811852) (-0.7185929648241204,0.5247070665968472) (-0.708542713567839,0.5248873126991275) (-0.6984924623115578,0.5200711775057348) (-0.6884422110552761,0.5186149446112452) (-0.6783919597989948,0.5168588243927792) (-0.6683417085427136,0.5166024330087824) (-0.658291457286432,0.5177123084959124) (-0.6482412060301507,0.5153601265119725) (-0.6381909547738694,0.5142583333315456) (-0.6281407035175878,0.5146769754818213) (-0.6180904522613064,0.5131783244134779) (-0.6080402010050252,0.5122768555095263) (-0.5979899497487435,0.5106765380171538) (-0.5879396984924623,0.5085671916523676) (-0.577889447236181,0.5066619544567632) (-0.5678391959798993,0.506923801726256) (-0.5577889447236181,0.5082910536046641) (-0.5477386934673366,0.5089184577650622) (-0.5376884422110553,0.5082841784145532) (-0.5276381909547738,0.507071038751894) (-0.5175879396984924,0.5068589015690198) (-0.5075376884422109,0.5070519747869299) (-0.49748743718592964,0.5055515762849346) (-0.4874371859296482,0.5042177407489404) (-0.4773869346733667,0.5062016095885076) (-0.46733668341708545,0.5050012855350821) (-0.45728643216080395,0.5062934283257747) (-0.44723618090452255,0.5084730960828407) (-0.4371859296482412,0.5108637150142208) (-0.42713567839195976,0.5081210286908716) (-0.41708542713567825,0.5090453474180598) (-0.4070351758793971,0.5083314420851277) (-0.3969849246231155,0.5098746487074916) (-0.38693467336683407,0.5093090117219041) (-0.3768844221105528,0.5057384412344413) (-0.36683417085427134,0.5038566577024367) (-0.35678391959798983,0.5050993094758267) (-0.34673366834170843,0.504923518165715) (-0.3366834170854271,0.5068940672522474) (-0.32663316582914564,0.5082637973391251) (-0.31658291457286425,0.5055198484988401) (-0.30653266331658297,0.5042189597061567) (-0.29648241206030146,0.5008226387292968) (-0.28643216080402,0.49934533499225336) (-0.2763819095477387,0.4994330583510197) (-0.2663316582914572,0.49981613976472467) (-0.2562814070351758,0.49924355502927786) (-0.2462311557788945,0.49679737365581017) (-0.236180904522613,0.49314473479398835) (-0.22613065326633158,0.4944360238326677) (-0.21608040201005027,0.49372542201048697) (-0.20603015075376882,0.49516691688059095) (-0.19597989949748737,0.4946786923472916) (-0.18592964824120603,0.49722780902808333) (-0.1758793969849246,0.5003838232675986) (-0.1658291457286431,0.5015662711972263) (-0.15577889447236168,0.4999191073040317) (-0.14572864321608042,0.4980488284776496) (-0.13567839195979894,0.49739969563515424) (-0.12562814070351747,0.4956469952802309) (-0.11557788944723618,0.4977906737996035) (-0.10552763819095472,0.4956102243348324) (-0.09547738693467328,0.4950659169557707) (-0.08542713567839201,0.497106273857526) (-0.07537688442211048,0.4955690518013421) (-0.06532663316582907,0.49532822352573413) (-0.055276381909547756,0.49633902053497075) (-0.04522613065326632,0.49635112323175246) (-0.03517587939698485,0.49781807662991506) (-0.025125628140703578,0.5020996128493171) (-0.015075376884422061,0.5015569014540123) (-0.005025125628140596,0.5014378715972722) (0.00502512562814067,0.5050712558563292) (0.015075376884422073,0.5062033525452309) (0.02512562814070361,0.508004342655735) (0.035175879396985056,0.5081469432665832) (0.045226130653266264,0.5086740392264523) (0.05527638190954779,0.5080387809875908) (0.06532663316582923,0.5071372738337214) (0.0753768844221105,0.5067674928108759) (0.08542713567839198,0.5066047117231349) (0.09547738693467347,0.5063078522974213) (0.10552763819095473,0.5075390926455126) (0.11557788944723618,0.506846657502271) (0.1256281407035177,0.5055164266871455) (0.13567839195979894,0.5058743632694434) (0.1457286432160804,0.5054646006216098) (0.1557788944723619,0.5068523137103659) (0.16582914572864316,0.5064123373711369) (0.17587939698492458,0.5056863782225701) (0.18592964824120614,0.5053883734326541) (0.19597989949748731,0.5047312793991431) (0.20603015075376882,0.5034016105169029) (0.21608040201005027,0.5032989975827142) (0.2261306532663318,0.5025749084670192) (0.236180904522613,0.5025603320758013) (0.24623115577889448,0.5028653103154741) (0.25628140703517593,0.5054392790794698) (0.2663316582914572,0.5048354389920412) (0.27638190954773867,0.5042347608275193) (0.2864321608040202,0.5036449168972) (0.2964824120603014,0.5074919269329161) (0.30653266331658285,0.5068291616097526) (0.31658291457286425,0.5080977896333136) (0.32663316582914576,0.5076740461806984) (0.33668341708542715,0.5081189459379482) (0.3467336683417085,0.5068040640034076) (0.35678391959799,0.5068040640034076) (0.36683417085427134,0.5076092618737457) (0.3768844221105528,0.5076169452061752) (0.3869346733668342,0.5067494119091266) (0.3969849246231155,0.5059888513266696) (0.407035175879397,0.5060517153678097) (0.4170854271356784,0.5053368896817205) (0.42713567839195976,0.5048256043007355) (0.43718592964824116,0.5048256043007355) (0.4472361809045226,0.504975437179219) (0.45728643216080406,0.505495531138479) (0.46733668341708545,0.5041656387974388) (0.4773869346733668,0.5047117456690968) (0.4874371859296483,0.5059772160942138) (0.49748743718592964,0.5056259892004008) (0.5075376884422109,0.5061785852899647) (0.5175879396984924,0.5064831323554122) (0.5276381909547738,0.506216211618329) (0.5376884422110553,0.5056379720701032) (0.5477386934673366,0.5061955806208723) (0.5577889447236181,0.5063219864108175) (0.5678391959798995,0.5068147986157007) (0.5778894472361809,0.5061866476073674) (0.5879396984924622,0.5076974804552739) (0.5979899497487436,0.5075767553505169) (0.6080402010050252,0.508088405434143) (0.6180904522613064,0.5086319642349338) (0.6281407035175879,0.5078139786940924) (0.6381909547738693,0.5070103420211609) (0.6482412060301508,0.5070103420211609) (0.6582914572864322,0.5078242372588063) (0.6683417085427136,0.508291054033206) (0.678391959798995,0.506874237315318) (0.6884422110552764,0.509910823310222) (0.6984924623115577,0.5089386483655628) (0.7085427135678392,0.5077027572341598) (0.7185929648241205,0.5067532526804515) (0.728643216080402,0.5069186153532287) (0.7386934673366834,0.5073783571088324) (0.7487437185929648,0.5085504055548291) (0.7587939698492462,0.5092661369503166) (0.7688442211055276,0.5105160968241627) (0.778894472361809,0.5099723718451379) (0.7889447236180904,0.5099796703108991) (0.7989949748743718,0.5101014781360804) (0.8090452261306532,0.5104047522907037) (0.8190954773869347,0.5104118065467418) (0.829145728643216,0.5098234754353708) (0.8391959798994975,0.5091347808879585) (0.8492462311557789,0.5094343129812323) (0.8592964824120602,0.5094343129812323) (0.8693467336683416,0.5087492482204632) (0.8793969849246231,0.5098160430181841) (0.8894472361809045,0.5090588289306672) (0.8994974874371859,0.508133925546827) (0.9095477386934674,0.5079019525899706) (0.9195979899497487,0.50852737593919) (0.9296482412060301,0.5092320977656105) (0.9396984924623116,0.510845888654905) (0.949748743718593,0.510845888654905) (0.9597989949748743,0.5104214255281431) (0.9698492462311558,0.5104214255281431) (0.9798994974874372,0.5096666072163144) (0.9899497487437185,0.5096666072163144) (0.9999999999999999,0.5096666072163144) };
  \addplot[ours_jeffreys, fill, fill opacity=0.18, draw=none, mark=none, forget plot] fill between [of=U_jeffreys_grp and L_jeffreys_grp];
  \addplot+[ours_jeffreys, mark=none] coordinates {(-0.9999999999999999,0.6037791803608114) (-0.9899497487437185,0.6002518481097252) (-0.979899497487437,0.596546262881878) (-0.9698492462311556,0.5974960914783938) (-0.9597989949748744,0.5931553990454717) (-0.9497487437185927,0.5900361598720425) (-0.9396984924623114,0.5897480678843552) (-0.9296482412060302,0.5898120011168919) (-0.9195979899497485,0.5875827311232428) (-0.9095477386934673,0.5871375429554806) (-0.8994974874371859,0.5872862064976423) (-0.8894472361809044,0.5876398938983816) (-0.879396984924623,0.5862838311927159) (-0.8693467336683418,0.5848135496233668) (-0.85929648241206,0.5835802694501646) (-0.8492462311557789,0.5816241909901871) (-0.8391959798994976,0.579598059833838) (-0.8291457286432159,0.5785674824116258) (-0.8190954773869346,0.577946320835757) (-0.809045226130653,0.5769019356580039) (-0.7989949748743717,0.5744618758121063) (-0.7889447236180904,0.5723535664958932) (-0.7788944723618088,0.5715651033064985) (-0.7688442211055275,0.5714037412404491) (-0.7587939698492462,0.5706536583780061) (-0.7487437185929646,0.5665849698970638) (-0.7386934673366833,0.5623013106005932) (-0.728643216080402,0.5585111103332631) (-0.7185929648241204,0.5582936082939717) (-0.708542713567839,0.5587765648082768) (-0.6984924623115578,0.5544922681556621) (-0.6884422110552761,0.5531889406648217) (-0.6783919597989948,0.5512584476558029) (-0.6683417085427136,0.5504905419738495) (-0.658291457286432,0.5519704636598667) (-0.6482412060301507,0.5507110980528386) (-0.6381909547738694,0.5495049776542935) (-0.6281407035175878,0.5492623494268795) (-0.6180904522613064,0.5479865156909683) (-0.6080402010050252,0.547239481410135) (-0.5979899497487435,0.5453909458249924) (-0.5879396984924623,0.5428667686077842) (-0.577889447236181,0.5408432269439191) (-0.5678391959798993,0.5408887583848674) (-0.5577889447236181,0.5420137384648982) (-0.5477386934673366,0.5426890690580222) (-0.5376884422110553,0.5420108273448534) (-0.5276381909547738,0.540618233950522) (-0.5175879396984924,0.5405103762251716) (-0.5075376884422109,0.5405085867543465) (-0.49748743718592964,0.5395788490281502) (-0.4874371859296482,0.5385741253864055) (-0.4773869346733667,0.5400450698170384) (-0.46733668341708545,0.5388365609075862) (-0.45728643216080395,0.5400500791193895) (-0.44723618090452255,0.5424501716520911) (-0.4371859296482412,0.5440596211707023) (-0.42713567839195976,0.5424854226905892) (-0.41708542713567825,0.5431587357688212) (-0.4070351758793971,0.5425473625932946) (-0.3969849246231155,0.5442376321324076) (-0.38693467336683407,0.5435826808211064) (-0.3768844221105528,0.5410401600721091) (-0.36683417085427134,0.538824190390006) (-0.35678391959798983,0.5392997820299228) (-0.34673366834170843,0.5391461444787552) (-0.3366834170854271,0.5405069567911759) (-0.32663316582914564,0.5417112372590197) (-0.31658291457286425,0.5393468637584797) (-0.30653266331658297,0.5380497600359838) (-0.29648241206030146,0.5346068757730974) (-0.28643216080402,0.5334423089018353) (-0.2763819095477387,0.5339636120357739) (-0.2663316582914572,0.5349095465878057) (-0.2562814070351758,0.5338106163075926) (-0.2462311557788945,0.5318769025069934) (-0.236180904522613,0.5287331817539889) (-0.22613065326633158,0.5299593960562891) (-0.21608040201005027,0.5297047385033988) (-0.20603015075376882,0.5310100373215073) (-0.19597989949748737,0.5303918659699907) (-0.18592964824120603,0.5333227153851379) (-0.1758793969849246,0.5359903162975815) (-0.1658291457286431,0.5371066108310654) (-0.15577889447236168,0.5358548904793118) (-0.14572864321608042,0.5339179276174986) (-0.13567839195979894,0.5332908281602613) (-0.12562814070351747,0.5320126663310881) (-0.11557788944723618,0.5340502742706669) (-0.10552763819095472,0.5320624946800204) (-0.09547738693467328,0.5310394546829568) (-0.08542713567839201,0.5329731159261285) (-0.07537688442211048,0.5315744311434301) (-0.06532663316582907,0.5310866176777641) (-0.055276381909547756,0.5320223826019114) (-0.04522613065326632,0.5320730746809321) (-0.03517587939698485,0.5333638773000536) (-0.025125628140703578,0.5361411517344206) (-0.015075376884422061,0.5352905085748727) (-0.005025125628140596,0.5352474356468871) (0.00502512562814067,0.5385979551682603) (0.015075376884422073,0.5394235741865895) (0.02512562814070361,0.5411559198007324) (0.035175879396985056,0.541948310978184) (0.045226130653266264,0.5425569372376219) (0.05527638190954779,0.54214110110665) (0.06532663316582923,0.541072252394725) (0.0753768844221105,0.5402722049650095) (0.08542713567839198,0.5405727760387163) (0.09547738693467347,0.5408347257148448) (0.10552763819095473,0.5420483100718316) (0.11557788944723618,0.5414992844569939) (0.1256281407035177,0.5408867269706397) (0.13567839195979894,0.5410374484238485) (0.1457286432160804,0.5409735385838028) (0.1557788944723619,0.5422408542130053) (0.16582914572864316,0.5422003165402399) (0.17587939698492458,0.5415661897829981) (0.18592964824120614,0.5416643921945032) (0.19597989949748731,0.5410619167891751) (0.20603015075376882,0.5396392032982663) (0.21608040201005027,0.539899244272131) (0.2261306532663318,0.5392145480898409) (0.236180904522613,0.5391917963843955) (0.24623115577889448,0.5396756125907455) (0.25628140703517593,0.5421482962922457) (0.2663316582914572,0.5411784825804192) (0.27638190954773867,0.5405907007873698) (0.2864321608040202,0.5403975861470157) (0.2964824120603014,0.5441919387368834) (0.30653266331658285,0.5435274556356803) (0.31658291457286425,0.5444568066442432) (0.32663316582914576,0.5439079960745347) (0.33668341708542715,0.5436763166302883) (0.3467336683417085,0.5424205778168976) (0.35678391959799,0.5424205778168976) (0.36683417085427134,0.5431013353528049) (0.3768844221105528,0.5430345138221464) (0.3869346733668342,0.5424076864182229) (0.3969849246231155,0.5417820485115707) (0.407035175879397,0.541798015840925) (0.4170854271356784,0.5411677301909267) (0.42713567839195976,0.5404913086407891) (0.43718592964824116,0.5404913086407891) (0.4472361809045226,0.5404998696791873) (0.45728643216080406,0.5410867729478253) (0.46733668341708545,0.5398301839502208) (0.4773869346733668,0.5404181458118984) (0.4874371859296483,0.5415409686104242) (0.49748743718592964,0.5410140843042646) (0.5075376884422109,0.5416010030645633) (0.5175879396984924,0.5422737032945031) (0.5276381909547738,0.5414138497422128) (0.5376884422110553,0.540719946684951) (0.5477386934673366,0.5414390778502348) (0.5577889447236181,0.5413508959162718) (0.5678391959798995,0.5420484868892982) (0.5778894472361809,0.5414219138146734) (0.5879396984924622,0.5426709636841235) (0.5979899497487436,0.5427414956989458) (0.6080402010050252,0.5433405406193315) (0.6180904522613064,0.5440537981663573) (0.6281407035175879,0.5436277663905207) (0.6381909547738693,0.5429156821026095) (0.6482412060301508,0.5429156821026095) (0.6582914572864322,0.543577736133591) (0.6683417085427136,0.5441725426843402) (0.678391959798995,0.5428942367848) (0.6884422110552764,0.5452281957643333) (0.6984924623115577,0.5445461108711177) (0.7085427135678392,0.5432817342806971) (0.7185929648241205,0.5425649159788435) (0.728643216080402,0.5426439568851881) (0.7386934673366834,0.5434641199356917) (0.7487437185929648,0.5436289437943261) (0.7587939698492462,0.5442722003842888) (0.7688442211055276,0.5455153447244078) (0.778894472361809,0.5449150778830636) (0.7889447236180904,0.5449631347887498) (0.7989949748743718,0.5449750940588235) (0.8090452261306532,0.5453904550357203) (0.8190954773869347,0.545473474255229) (0.829145728643216,0.5448752627736635) (0.8391959798994975,0.5442112990525665) (0.8492462311557789,0.5446905367812968) (0.8592964824120602,0.5446905367812968) (0.8693467336683416,0.5440901112592015) (0.8793969849246231,0.544609434830437) (0.8894472361809045,0.5439728237009694) (0.8994974874371859,0.5428696792688323) (0.9095477386934674,0.5428185865423567) (0.9195979899497487,0.5426792799547568) (0.9296482412060301,0.5432645540304953) (0.9396984924623116,0.544350967006909) (0.949748743718593,0.544350967006909) (0.9597989949748743,0.5436465173575954) (0.9698492462311558,0.5436465173575954) (0.9798994974874372,0.5430220047037845) (0.9899497487437185,0.5430220047037845) (0.9999999999999999,0.5430220047037845) };
  \addplot+[ours_jeffreys, mark=none, line width=0.4pt, dashed, forget plot] coordinates {(-1.0,0.4430012466399693) (1.0,0.4430012466399693)};
  \addplot+[ours_jeffreys, mark=*, only marks, mark size=\PanelArgminSize, forget plot] coordinates {(-0.236180904522613,0.5287331817539889)};
  \addplot[name path=U_ig_robust_grp, draw=none, mark=none, forget plot] coordinates {(-0.9999999999999999,0.6305544769177985) (-0.9899497487437185,0.628331677292106) (-0.979899497487437,0.6278356817132358) (-0.9698492462311556,0.6222211560454427) (-0.9597989949748744,0.6215933548990352) (-0.9497487437185927,0.620559064008493) (-0.9396984924623114,0.6208095045469095) (-0.9296482412060302,0.6208697862247472) (-0.9195979899497485,0.6190124376164007) (-0.9095477386934673,0.6193064181884795) (-0.8994974874371859,0.6173397579546963) (-0.8894472361809044,0.6165961863791184) (-0.879396984924623,0.6161231305270244) (-0.8693467336683418,0.6146093139127051) (-0.85929648241206,0.6120960585275439) (-0.8492462311557789,0.6110945659591838) (-0.8391959798994976,0.6094986557804735) (-0.8291457286432159,0.6073077724075082) (-0.8190954773869346,0.6074201185374385) (-0.809045226130653,0.6079689627499933) (-0.7989949748743717,0.6043748266772961) (-0.7889447236180904,0.602941798673867) (-0.7788944723618088,0.6029105610055441) (-0.7688442211055275,0.601137054216043) (-0.7587939698492462,0.5978167259803848) (-0.7487437185929646,0.593228517624782) (-0.7386934673366833,0.5890162252918874) (-0.728643216080402,0.5876076144385626) (-0.7185929648241204,0.5868481143482668) (-0.708542713567839,0.5857574511772321) (-0.6984924623115578,0.582076480205822) (-0.6884422110552761,0.5812442778380271) (-0.6783919597989948,0.5799804237239883) (-0.6683417085427136,0.5793574961574579) (-0.658291457286432,0.5814453822683275) (-0.6482412060301507,0.5802951293374748) (-0.6381909547738694,0.5805128217386781) (-0.6281407035175878,0.578341704510829) (-0.6180904522613064,0.5774049634731577) (-0.6080402010050252,0.5775065298126273) (-0.5979899497487435,0.5736606763250787) (-0.5879396984924623,0.5729619759781915) (-0.577889447236181,0.57094507659843) (-0.5678391959798993,0.5705377704178789) (-0.5577889447236181,0.5702422248063473) (-0.5477386934673366,0.5720876578920836) (-0.5376884422110553,0.5702177402296666) (-0.5276381909547738,0.5704091320484143) (-0.5175879396984924,0.5716027226386537) (-0.5075376884422109,0.5708572475895813) (-0.49748743718592964,0.5704613434559231) (-0.4874371859296482,0.5697710870928795) (-0.4773869346733667,0.5696642569797058) (-0.46733668341708545,0.5711368829129791) (-0.45728643216080395,0.5710455509680764) (-0.44723618090452255,0.5728579995827144) (-0.4371859296482412,0.5744238549310481) (-0.42713567839195976,0.5744834092903057) (-0.41708542713567825,0.5736230176400073) (-0.4070351758793971,0.5730258002735922) (-0.3969849246231155,0.5744183751579265) (-0.38693467336683407,0.5737384918327785) (-0.3768844221105528,0.5711827950490188) (-0.36683417085427134,0.5684761836669272) (-0.35678391959798983,0.5684907652902474) (-0.34673366834170843,0.5685626037861006) (-0.3366834170854271,0.5696220831782142) (-0.32663316582914564,0.5702126468396793) (-0.31658291457286425,0.5691842496847362) (-0.30653266331658297,0.56690578837133) (-0.29648241206030146,0.5643219270486981) (-0.28643216080402,0.5633558247664666) (-0.2763819095477387,0.5642456107555145) (-0.2663316582914572,0.565502733946664) (-0.2562814070351758,0.5633627603796509) (-0.2462311557788945,0.5620722132922192) (-0.236180904522613,0.5602807517325208) (-0.22613065326633158,0.559591726341772) (-0.21608040201005027,0.5607172634304742) (-0.20603015075376882,0.5624523445484999) (-0.19597989949748737,0.5611270489272764) (-0.18592964824120603,0.5637449341607611) (-0.1758793969849246,0.5668842280387666) (-0.1658291457286431,0.568172150912113) (-0.15577889447236168,0.5660485894711922) (-0.14572864321608042,0.5655956516109983) (-0.13567839195979894,0.5655282784528214) (-0.12562814070351747,0.5639155611259657) (-0.11557788944723618,0.565865090555627) (-0.10552763819095472,0.5639845372956145) (-0.09547738693467328,0.564033327684327) (-0.08542713567839201,0.5637376553410132) (-0.07537688442211048,0.5631409339666913) (-0.06532663316582907,0.5633653554261555) (-0.055276381909547756,0.5623422732666602) (-0.04522613065326632,0.5631046031589783) (-0.03517587939698485,0.5635104589553911) (-0.025125628140703578,0.5646165546433374) (-0.015075376884422061,0.5647868959742468) (-0.005025125628140596,0.5648044264105434) (0.00502512562814067,0.5661965217298855) (0.015075376884422073,0.5687699582977512) (0.02512562814070361,0.5700128154209189) (0.035175879396985056,0.5715918538271212) (0.045226130653266264,0.5729283428117178) (0.05527638190954779,0.5722622215253119) (0.06532663316582923,0.5713753336379678) (0.0753768844221105,0.5704932614664886) (0.08542713567839198,0.5723851699485909) (0.09547738693467347,0.5731436236020873) (0.10552763819095473,0.5729614916796538) (0.11557788944723618,0.5729394401568102) (0.1256281407035177,0.5729293231675958) (0.13567839195979894,0.5729335961787744) (0.1457286432160804,0.5738091890528446) (0.1557788944723619,0.5741024234685644) (0.16582914572864316,0.5747204601337712) (0.17587939698492458,0.5741487197868027) (0.18592964824120614,0.5744034425116427) (0.19597989949748731,0.5739494743236936) (0.20603015075376882,0.5724269086760498) (0.21608040201005027,0.5737344454829438) (0.2261306532663318,0.5723036831884988) (0.236180904522613,0.5730097224004924) (0.24623115577889448,0.5728094260992141) (0.25628140703517593,0.574735737923311) (0.2663316582914572,0.5744345094461736) (0.27638190954773867,0.5738546273878385) (0.2864321608040202,0.5734330854572253) (0.2964824120603014,0.5766099896380064) (0.30653266331658285,0.5766764793480023) (0.31658291457286425,0.577084851482561) (0.32663316582914576,0.5757440580019499) (0.33668341708542715,0.5757440580019499) (0.3467336683417085,0.5746976104722317) (0.35678391959799,0.5741573414121391) (0.36683417085427134,0.5748013043564912) (0.3768844221105528,0.5748563830726006) (0.3869346733668342,0.5751789458735148) (0.3969849246231155,0.573881756292047) (0.407035175879397,0.5732072848954394) (0.4170854271356784,0.5731622328978433) (0.42713567839195976,0.5723317991476096) (0.43718592964824116,0.5723317991476096) (0.4472361809045226,0.5715236584784488) (0.45728643216080406,0.5722211855386096) (0.46733668341708545,0.5728620123651624) (0.4773869346733668,0.5721107798581436) (0.4874371859296483,0.572691878762573) (0.49748743718592964,0.5718021177781983) (0.5075376884422109,0.5730029744861651) (0.5175879396984924,0.5744330302435994) (0.5276381909547738,0.5740623235076714) (0.5376884422110553,0.5727758755092531) (0.5477386934673366,0.5722298610815119) (0.5577889447236181,0.5714843115845772) (0.5678391959798995,0.5719400647171562) (0.5778894472361809,0.5723534124682943) (0.5879396984924622,0.5727571974307029) (0.5979899497487436,0.5725844713383674) (0.6080402010050252,0.5725844713383674) (0.6180904522613064,0.5739901342621927) (0.6281407035175879,0.5743088883825037) (0.6381909547738693,0.5741295058527071) (0.6482412060301508,0.5734882035129256) (0.6582914572864322,0.5741000925997588) (0.6683417085427136,0.5749523383302115) (0.678391959798995,0.5743293516067779) (0.6884422110552764,0.5757321033105565) (0.6984924623115577,0.575935478838672) (0.7085427135678392,0.575935478838672) (0.7185929648241205,0.5753984564394141) (0.728643216080402,0.5748404795499782) (0.7386934673366834,0.5747917326506147) (0.7487437185929648,0.5756298510480529) (0.7587939698492462,0.5753176552329924) (0.7688442211055276,0.5769896267589848) (0.778894472361809,0.5763156403008158) (0.7889447236180904,0.5762267872203868) (0.7989949748743718,0.5756512862097658) (0.8090452261306532,0.5769545487424841) (0.8190954773869347,0.5769040542939057) (0.829145728643216,0.5762701587811985) (0.8391959798994975,0.5756983417765261) (0.8492462311557789,0.5759196340881648) (0.8592964824120602,0.576501296049769) (0.8693467336683416,0.5764832704916453) (0.8793969849246231,0.5764756776889478) (0.8894472361809045,0.5754607172417692) (0.8994974874371859,0.5739461365478643) (0.9095477386934674,0.5740960450534576) (0.9195979899497487,0.5732915490649826) (0.9296482412060301,0.5739618942985484) (0.9396984924623116,0.5745014013668849) (0.949748743718593,0.5745014013668849) (0.9597989949748743,0.5736247429441103) (0.9698492462311558,0.5736247429441103) (0.9798994974874372,0.5731250884822972) (0.9899497487437185,0.5731250884822972) (0.9999999999999999,0.5731250884822972) };
  \addplot[name path=L_ig_robust_grp, draw=none, mark=none, forget plot] coordinates {(-0.9999999999999999,0.563182343083897) (-0.9899497487437185,0.5603191053350661) (-0.979899497487437,0.5602744403396254) (-0.9698492462311556,0.5533823774012301) (-0.9597989949748744,0.5514935700374163) (-0.9497487437185927,0.550762871381302) (-0.9396984924623114,0.5528222296565668) (-0.9296482412060302,0.5505519535939234) (-0.9195979899497485,0.5508069318729406) (-0.9095477386934673,0.5490246549074893) (-0.8994974874371859,0.5481898925359495) (-0.8894472361809044,0.5479280687985367) (-0.879396984924623,0.5460261048880971) (-0.8693467336683418,0.547875513168758) (-0.85929648241206,0.5449043875420694) (-0.8492462311557789,0.5446310917631907) (-0.8391959798994976,0.542056374805942) (-0.8291457286432159,0.5407208579171328) (-0.8190954773869346,0.54051344289656) (-0.809045226130653,0.539400306383585) (-0.7989949748743717,0.5364216962934767) (-0.7889447236180904,0.5347967333440218) (-0.7788944723618088,0.5346474443216322) (-0.7688442211055275,0.5345323407941989) (-0.7587939698492462,0.5311223137649597) (-0.7487437185929646,0.5272412308127565) (-0.7386934673366833,0.5249039012557999) (-0.728643216080402,0.5225561551997099) (-0.7185929648241204,0.5219658334175227) (-0.708542713567839,0.5193603311403672) (-0.6984924623115578,0.515301236438287) (-0.6884422110552761,0.5150114391618044) (-0.6783919597989948,0.513668154467056) (-0.6683417085427136,0.5137740805694457) (-0.658291457286432,0.5145578456433594) (-0.6482412060301507,0.5119610777506548) (-0.6381909547738694,0.5116906106185519) (-0.6281407035175878,0.510063576381832) (-0.6180904522613064,0.5091819384504237) (-0.6080402010050252,0.5087689051008152) (-0.5979899497487435,0.5064422167112643) (-0.5879396984924623,0.5057368185906265) (-0.577889447236181,0.5039528399376348) (-0.5678391959798993,0.5031883206372489) (-0.5577889447236181,0.5034180980458091) (-0.5477386934673366,0.5052812731913626) (-0.5376884422110553,0.5034316865774751) (-0.5276381909547738,0.5029659354646226) (-0.5175879396984924,0.5040964795526268) (-0.5075376884422109,0.5035165687792086) (-0.49748743718592964,0.5020579392175817) (-0.4874371859296482,0.5007450061713116) (-0.4773869346733667,0.5013335235766362) (-0.46733668341708545,0.5022715665934273) (-0.45728643216080395,0.502302945536066) (-0.44723618090452255,0.5043339935933236) (-0.4371859296482412,0.5072611951608581) (-0.42713567839195976,0.5053139983145258) (-0.41708542713567825,0.5048400244968368) (-0.4070351758793971,0.5042193993530644) (-0.3969849246231155,0.5059938174146276) (-0.38693467336683407,0.5053988503844445) (-0.3768844221105528,0.5011147313375955) (-0.36683417085427134,0.499522503350278) (-0.35678391959798983,0.5011518311260804) (-0.34673366834170843,0.5001335648476123) (-0.3366834170854271,0.5020113502738501) (-0.32663316582914564,0.5026346068691228) (-0.31658291457286425,0.5010699301251474) (-0.30653266331658297,0.497216862134719) (-0.29648241206030146,0.4953903693235971) (-0.28643216080402,0.49405308434189676) (-0.2763819095477387,0.49410731958865317) (-0.2663316582914572,0.4944522029799864) (-0.2562814070351758,0.49296827844056335) (-0.2462311557788945,0.49082199634041357) (-0.236180904522613,0.48899499851706785) (-0.22613065326633158,0.4884232307283957) (-0.21608040201005027,0.48911613915860386) (-0.20603015075376882,0.4909015763134166) (-0.19597989949748737,0.48975128442866556) (-0.18592964824120603,0.49138846359509725) (-0.1758793969849246,0.4950881967235722) (-0.1658291457286431,0.49609038707198433) (-0.15577889447236168,0.49459660794459537) (-0.14572864321608042,0.49351433941341116) (-0.13567839195979894,0.4937504599342765) (-0.12562814070351747,0.49169687616421054) (-0.11557788944723618,0.4936309478641188) (-0.10552763819095472,0.49141084147169045) (-0.09547738693467328,0.4914895874449627) (-0.08542713567839201,0.4923762201625543) (-0.07537688442211048,0.49188635228775746) (-0.06532663316582907,0.49350630297208475) (-0.055276381909547756,0.4920788262700161) (-0.04522613065326632,0.4925865252296699) (-0.03517587939698485,0.49336664326526225) (-0.025125628140703578,0.49657995035729957) (-0.015075376884422061,0.4970272352619034) (-0.005025125628140596,0.49704955932915507) (0.00502512562814067,0.498969343855577) (0.015075376884422073,0.5013553660126813) (0.02512562814070361,0.5031547098154837) (0.035175879396985056,0.5029547233159692) (0.045226130653266264,0.5041714500261867) (0.05527638190954779,0.5027224142488) (0.06532663316582923,0.5024879203655898) (0.0753768844221105,0.5020851990507214) (0.08542713567839198,0.5032075702743137) (0.09547738693467347,0.5030614730654743) (0.10552763819095473,0.5032068905736155) (0.11557788944723618,0.5032919648527407) (0.1256281407035177,0.5018015730107652) (0.13567839195979894,0.5021035668706598) (0.1457286432160804,0.5022073187605528) (0.1557788944723619,0.502062429483403) (0.16582914572864316,0.5027129942706423) (0.17587939698492458,0.5021225774471993) (0.18592964824120614,0.5019556346882825) (0.19597989949748731,0.5012592790882319) (0.20603015075376882,0.5000310637596617) (0.21608040201005027,0.5003258311495558) (0.2261306532663318,0.4990852673472711) (0.236180904522613,0.49966760682407196) (0.24623115577889448,0.499349634949679) (0.25628140703517593,0.5014444118303459) (0.2663316582914572,0.5021629545237052) (0.27638190954773867,0.5014460597093066) (0.2864321608040202,0.499972667203891) (0.2964824120603014,0.5033796772440993) (0.30653266331658285,0.5032335247712333) (0.31658291457286425,0.5048646871414312) (0.32663316582914576,0.5046626366141065) (0.33668341708542715,0.5046626366141065) (0.3467336683417085,0.5041104993843843) (0.35678391959799,0.5035022066820624) (0.36683417085427134,0.5042381984515503) (0.3768844221105528,0.5042893983371617) (0.3869346733668342,0.5041242793297922) (0.3969849246231155,0.5029343462529298) (0.407035175879397,0.5022886821961631) (0.4170854271356784,0.502351070885105) (0.42713567839195976,0.5018280794577301) (0.43718592964824116,0.5018280794577301) (0.4472361809045226,0.5011966797632869) (0.45728643216080406,0.5018647390258512) (0.46733668341708545,0.5025030510818904) (0.4773869346733668,0.5018945334335388) (0.4874371859296483,0.5026037516334259) (0.49748743718592964,0.5023391579044102) (0.5075376884422109,0.5037777001669854) (0.5175879396984924,0.5040571484584575) (0.5276381909547738,0.5052143568322423) (0.5376884422110553,0.503964126324479) (0.5477386934673366,0.5031885056213322) (0.5577889447236181,0.5026678765260862) (0.5678391959798995,0.503408391515562) (0.5778894472361809,0.5042392866182703) (0.5879396984924622,0.5050501595417771) (0.5979899497487436,0.5050694847182258) (0.6080402010050252,0.5050694847182258) (0.6180904522613064,0.5062954382182641) (0.6281407035175879,0.5061232808244386) (0.6381909547738693,0.5051983373030676) (0.6482412060301508,0.5044960572430908) (0.6582914572864322,0.5052626868988328) (0.6683417085427136,0.5056106753898988) (0.678391959798995,0.5050033829651673) (0.6884422110552764,0.5072543009106154) (0.6984924623115577,0.5068453129987477) (0.7085427135678392,0.5068453129987477) (0.7185929648241205,0.506056860402422) (0.728643216080402,0.5054754553295648) (0.7386934673366834,0.5046298501399731) (0.7487437185929648,0.5059937624641642) (0.7587939698492462,0.5068919766966091) (0.7688442211055276,0.508898246730387) (0.778894472361809,0.5083708416835522) (0.7889447236180904,0.5082861972189556) (0.7989949748743718,0.5076510871533995) (0.8090452261306532,0.5083581828449875) (0.8190954773869347,0.5083557499024696) (0.829145728643216,0.50775420017633) (0.8391959798994975,0.5070451659769616) (0.8492462311557789,0.5066033200284145) (0.8592964824120602,0.5074486847410767) (0.8693467336683416,0.5074770582571776) (0.8793969849246231,0.5075459443223395) (0.8894472361809045,0.5069775397655936) (0.8994974874371859,0.5061666587281076) (0.9095477386934674,0.5058789612806109) (0.9195979899497487,0.5064920570503388) (0.9296482412060301,0.5070130573531398) (0.9396984924623116,0.5087200611353454) (0.949748743718593,0.5087200611353454) (0.9597989949748743,0.5081069268030692) (0.9698492462311558,0.5081069268030692) (0.9798994974874372,0.5073561801879634) (0.9899497487437185,0.5073561801879634) (0.9999999999999999,0.5073561801879634) };
  \addplot[ours_ig_robust, fill, fill opacity=0.18, draw=none, mark=none, forget plot] fill between [of=U_ig_robust_grp and L_ig_robust_grp];
  \addplot+[ours_ig_robust, mark=none] coordinates {(-0.9999999999999999,0.5968684100008478) (-0.9899497487437185,0.5943253913135861) (-0.979899497487437,0.5940550610264306) (-0.9698492462311556,0.5878017667233364) (-0.9597989949748744,0.5865434624682258) (-0.9497487437185927,0.5856609676948975) (-0.9396984924623114,0.5868158671017382) (-0.9296482412060302,0.5857108699093353) (-0.9195979899497485,0.5849096847446706) (-0.9095477386934673,0.5841655365479844) (-0.8994974874371859,0.5827648252453229) (-0.8894472361809044,0.5822621275888276) (-0.879396984924623,0.5810746177075607) (-0.8693467336683418,0.5812424135407316) (-0.85929648241206,0.5785002230348066) (-0.8492462311557789,0.5778628288611872) (-0.8391959798994976,0.5757775152932078) (-0.8291457286432159,0.5740143151623205) (-0.8190954773869346,0.5739667807169992) (-0.809045226130653,0.5736846345667892) (-0.7989949748743717,0.5703982614853864) (-0.7889447236180904,0.5688692660089444) (-0.7788944723618088,0.5687790026635882) (-0.7688442211055275,0.5678346975051209) (-0.7587939698492462,0.5644695198726722) (-0.7487437185929646,0.5602348742187693) (-0.7386934673366833,0.5569600632738436) (-0.728643216080402,0.5550818848191362) (-0.7185929648241204,0.5544069738828947) (-0.708542713567839,0.5525588911587996) (-0.6984924623115578,0.5486888583220545) (-0.6884422110552761,0.5481278584999157) (-0.6783919597989948,0.5468242890955222) (-0.6683417085427136,0.5465657883634518) (-0.658291457286432,0.5480016139558435) (-0.6482412060301507,0.5461281035440648) (-0.6381909547738694,0.546101716178615) (-0.6281407035175878,0.5442026404463305) (-0.6180904522613064,0.5432934509617907) (-0.6080402010050252,0.5431377174567212) (-0.5979899497487435,0.5400514465181715) (-0.5879396984924623,0.539349397284409) (-0.577889447236181,0.5374489582680324) (-0.5678391959798993,0.5368630455275639) (-0.5577889447236181,0.5368301614260782) (-0.5477386934673366,0.5386844655417231) (-0.5376884422110553,0.5368247134035709) (-0.5276381909547738,0.5366875337565185) (-0.5175879396984924,0.5378496010956403) (-0.5075376884422109,0.537186908184395) (-0.49748743718592964,0.5362596413367524) (-0.4874371859296482,0.5352580466320955) (-0.4773869346733667,0.535498890278171) (-0.46733668341708545,0.5367042247532032) (-0.45728643216080395,0.5366742482520712) (-0.44723618090452255,0.538595996588019) (-0.4371859296482412,0.5408425250459531) (-0.42713567839195976,0.5398987038024158) (-0.41708542713567825,0.5392315210684221) (-0.4070351758793971,0.5386225998133283) (-0.3969849246231155,0.540206096286277) (-0.38693467336683407,0.5395686711086115) (-0.3768844221105528,0.5361487631933072) (-0.36683417085427134,0.5339993435086026) (-0.35678391959798983,0.5348212982081639) (-0.34673366834170843,0.5343480843168564) (-0.3366834170854271,0.5358167167260321) (-0.32663316582914564,0.536423626854401) (-0.31658291457286425,0.5351270899049418) (-0.30653266331658297,0.5320613252530245) (-0.29648241206030146,0.5298561481861476) (-0.28643216080402,0.5287044545541817) (-0.2763819095477387,0.5291764651720838) (-0.2663316582914572,0.5299774684633252) (-0.2562814070351758,0.5281655194101071) (-0.2462311557788945,0.5264471048163164) (-0.236180904522613,0.5246378751247943) (-0.22613065326633158,0.5240074785350839) (-0.21608040201005027,0.524916701294539) (-0.20603015075376882,0.5266769604309582) (-0.19597989949748737,0.525439166677971) (-0.18592964824120603,0.5275666988779292) (-0.1758793969849246,0.5309862123811694) (-0.1658291457286431,0.5321312689920487) (-0.15577889447236168,0.5303225987078938) (-0.14572864321608042,0.5295549955122048) (-0.13567839195979894,0.5296393691935489) (-0.12562814070351747,0.5278062186450881) (-0.11557788944723618,0.5297480192098729) (-0.10552763819095472,0.5276976893836525) (-0.09547738693467328,0.5277614575646449) (-0.08542713567839201,0.5280569377517838) (-0.07537688442211048,0.5275136431272244) (-0.06532663316582907,0.5284358291991201) (-0.055276381909547756,0.5272105497683381) (-0.04522613065326632,0.527845564194324) (-0.03517587939698485,0.5284385511103267) (-0.025125628140703578,0.5305982525003184) (-0.015075376884422061,0.5309070656180751) (-0.005025125628140596,0.5309269928698492) (0.00502512562814067,0.5325829327927313) (0.015075376884422073,0.5350626621552163) (0.02512562814070361,0.5365837626182013) (0.035175879396985056,0.5372732885715452) (0.045226130653266264,0.5385498964189522) (0.05527638190954779,0.537492317887056) (0.06532663316582923,0.5369316270017788) (0.0753768844221105,0.536289230258605) (0.08542713567839198,0.5377963701114523) (0.09547738693467347,0.5381025483337808) (0.10552763819095473,0.5380841911266346) (0.11557788944723618,0.5381157025047755) (0.1256281407035177,0.5373654480891805) (0.13567839195979894,0.5375185815247171) (0.1457286432160804,0.5380082539066987) (0.1557788944723619,0.5380824264759837) (0.16582914572864316,0.5387167272022068) (0.17587939698492458,0.538135648617001) (0.18592964824120614,0.5381795385999626) (0.19597989949748731,0.5376043767059627) (0.20603015075376882,0.5362289862178558) (0.21608040201005027,0.5370301383162498) (0.2261306532663318,0.5356944752678849) (0.236180904522613,0.5363386646122822) (0.24623115577889448,0.5360795305244466) (0.25628140703517593,0.5380900748768285) (0.2663316582914572,0.5382987319849394) (0.27638190954773867,0.5376503435485726) (0.2864321608040202,0.5367028763305581) (0.2964824120603014,0.5399948334410528) (0.30653266331658285,0.5399550020596178) (0.31658291457286425,0.5409747693119961) (0.32663316582914576,0.5402033473080282) (0.33668341708542715,0.5402033473080282) (0.3467336683417085,0.539404054928308) (0.35678391959799,0.5388297740471008) (0.36683417085427134,0.5395197514040208) (0.3768844221105528,0.5395728907048811) (0.3869346733668342,0.5396516126016535) (0.3969849246231155,0.5384080512724884) (0.407035175879397,0.5377479835458012) (0.4170854271356784,0.5377566518914741) (0.42713567839195976,0.5370799393026698) (0.43718592964824116,0.5370799393026698) (0.4472361809045226,0.5363601691208678) (0.45728643216080406,0.5370429622822304) (0.46733668341708545,0.5376825317235264) (0.4773869346733668,0.5370026566458412) (0.4874371859296483,0.5376478151979994) (0.49748743718592964,0.5370706378413043) (0.5075376884422109,0.5383903373265753) (0.5175879396984924,0.5392450893510284) (0.5276381909547738,0.5396383401699568) (0.5376884422110553,0.538370000916866) (0.5477386934673366,0.537709183351422) (0.5577889447236181,0.5370760940553317) (0.5678391959798995,0.5376742281163591) (0.5778894472361809,0.5382963495432823) (0.5879396984924622,0.53890367848624) (0.5979899497487436,0.5388269780282966) (0.6080402010050252,0.5388269780282966) (0.6180904522613064,0.5401427862402284) (0.6281407035175879,0.5402160846034711) (0.6381909547738693,0.5396639215778873) (0.6482412060301508,0.5389921303780082) (0.6582914572864322,0.5396813897492958) (0.6683417085427136,0.5402815068600552) (0.678391959798995,0.5396663672859726) (0.6884422110552764,0.5414932021105859) (0.6984924623115577,0.5413903959187099) (0.7085427135678392,0.5413903959187099) (0.7185929648241205,0.5407276584209181) (0.728643216080402,0.5401579674397715) (0.7386934673366834,0.5397107913952939) (0.7487437185929648,0.5408118067561085) (0.7587939698492462,0.5411048159648008) (0.7688442211055276,0.5429439367446859) (0.778894472361809,0.542343240992184) (0.7889447236180904,0.5422564922196712) (0.7989949748743718,0.5416511866815826) (0.8090452261306532,0.5426563657937358) (0.8190954773869347,0.5426299020981876) (0.829145728643216,0.5420121794787642) (0.8391959798994975,0.5413717538767439) (0.8492462311557789,0.5412614770582896) (0.8592964824120602,0.5419749903954229) (0.8693467336683416,0.5419801643744114) (0.8793969849246231,0.5420108110056436) (0.8894472361809045,0.5412191285036814) (0.8994974874371859,0.540056397637986) (0.9095477386934674,0.5399875031670343) (0.9195979899497487,0.5398918030576607) (0.9296482412060301,0.5404874758258441) (0.9396984924623116,0.5416107312511151) (0.949748743718593,0.5416107312511151) (0.9597989949748743,0.5408658348735897) (0.9698492462311558,0.5408658348735897) (0.9798994974874372,0.5402406343351303) (0.9899497487437185,0.5402406343351303) (0.9999999999999999,0.5402406343351303) };
  \addplot+[ours_ig_robust, mark=none, line width=0.4pt, dashed, forget plot] coordinates {(-1.0,0.4413841475971843) (1.0,0.4413841475971843)};
  \addplot+[ours_ig_robust, mark=*, only marks, mark size=\PanelArgminSize, forget plot] coordinates {(-0.22613065326633158,0.5240074785350839)};
  \addplot[name path=U_ig_mle_grp, draw=none, mark=none, forget plot] coordinates {(-0.9999999999999999,0.6296331345510767) (-0.9899497487437185,0.627565034577261) (-0.979899497487437,0.6277403329060915) (-0.9698492462311556,0.6220640699390956) (-0.9597989949748744,0.6211729087130718) (-0.9497487437185927,0.6208814504270251) (-0.9396984924623114,0.6205220250638497) (-0.9296482412060302,0.6204752196326877) (-0.9195979899497485,0.618282103088367) (-0.9095477386934673,0.6185694186415284) (-0.8994974874371859,0.6173175799973053) (-0.8894472361809044,0.6166313704730012) (-0.879396984924623,0.6167294019031917) (-0.8693467336683418,0.6142214377724473) (-0.85929648241206,0.6129023144574358) (-0.8492462311557789,0.6112527231127634) (-0.8391959798994976,0.6105474893835071) (-0.8291457286432159,0.6082258258237019) (-0.8190954773869346,0.6069266935468279) (-0.809045226130653,0.6064033710919158) (-0.7989949748743717,0.6026201591898774) (-0.7889447236180904,0.601193030740351) (-0.7788944723618088,0.6022559392030506) (-0.7688442211055275,0.6014340452957864) (-0.7587939698492462,0.5972716616475428) (-0.7487437185929646,0.5931591170731163) (-0.7386934673366833,0.5887116840082209) (-0.728643216080402,0.5863828347638221) (-0.7185929648241204,0.5864342381113075) (-0.708542713567839,0.584531722708209) (-0.6984924623115578,0.5811634773761216) (-0.6884422110552761,0.580508310253438) (-0.6783919597989948,0.578736358544137) (-0.6683417085427136,0.5788451476431209) (-0.658291457286432,0.5807399236636143) (-0.6482412060301507,0.5802286863252621) (-0.6381909547738694,0.5794536162988281) (-0.6281407035175878,0.5782507370120831) (-0.6180904522613064,0.5767082693360448) (-0.6080402010050252,0.577640310817446) (-0.5979899497487435,0.5738584358154346) (-0.5879396984924623,0.5717738653973575) (-0.577889447236181,0.5710980853573007) (-0.5678391959798993,0.5696311705087083) (-0.5577889447236181,0.5691710968525587) (-0.5477386934673366,0.569523990804177) (-0.5376884422110553,0.5693441149438394) (-0.5276381909547738,0.5687190893799049) (-0.5175879396984924,0.5686289483599984) (-0.5075376884422109,0.5692356887453907) (-0.49748743718592964,0.5690341619817432) (-0.4874371859296482,0.5677423951911185) (-0.4773869346733667,0.5680206388831075) (-0.46733668341708545,0.5687567872204533) (-0.45728643216080395,0.5686842891437659) (-0.44723618090452255,0.5704383094094354) (-0.4371859296482412,0.5720126801032877) (-0.42713567839195976,0.5719275628514586) (-0.41708542713567825,0.571719263698426) (-0.4070351758793971,0.5711535405477357) (-0.3969849246231155,0.5727995216319806) (-0.38693467336683407,0.5721105785063667) (-0.3768844221105528,0.570383140373175) (-0.36683417085427134,0.5678223847336611) (-0.35678391959798983,0.5673188235014105) (-0.34673366834170843,0.567161823842127) (-0.3366834170854271,0.568300741306287) (-0.32663316582914564,0.5682915194959413) (-0.31658291457286425,0.5678255907816785) (-0.30653266331658297,0.5654853770228406) (-0.29648241206030146,0.5637758056327833) (-0.28643216080402,0.56194961735413) (-0.2763819095477387,0.5628438586219061) (-0.2663316582914572,0.5640432251644291) (-0.2562814070351758,0.5635812442036513) (-0.2462311557788945,0.5603663558826525) (-0.236180904522613,0.5589741110946475) (-0.22613065326633158,0.5596553207494682) (-0.21608040201005027,0.5589768189573249) (-0.20603015075376882,0.5588276357687813) (-0.19597989949748737,0.5607936283186276) (-0.18592964824120603,0.5629457630663028) (-0.1758793969849246,0.5659361113980073) (-0.1658291457286431,0.5663740982405998) (-0.15577889447236168,0.5656568368313698) (-0.14572864321608042,0.5647421856748466) (-0.13567839195979894,0.5647677902529139) (-0.12562814070351747,0.5622773672475955) (-0.11557788944723618,0.5650190530995141) (-0.10552763819095472,0.5630206497253004) (-0.09547738693467328,0.5630386639096134) (-0.08542713567839201,0.5619681132302925) (-0.07537688442211048,0.5620168580334206) (-0.06532663316582907,0.5625580223291087) (-0.055276381909547756,0.5608060062017045) (-0.04522613065326632,0.5623970895575355) (-0.03517587939698485,0.5629318858404205) (-0.025125628140703578,0.5635514507439036) (-0.015075376884422061,0.5643598108122733) (-0.005025125628140596,0.5644683655450463) (0.00502512562814067,0.5642095011917552) (0.015075376884422073,0.5671911609672557) (0.02512562814070361,0.5681233513337569) (0.035175879396985056,0.5707543755623775) (0.045226130653266264,0.5720118972321077) (0.05527638190954779,0.572284168900674) (0.06532663316582923,0.5699638512682906) (0.0753768844221105,0.5692475971259654) (0.08542713567839198,0.5708202194508275) (0.09547738693467347,0.5723155334279714) (0.10552763819095473,0.5722606244048631) (0.11557788944723618,0.5723302968271649) (0.1256281407035177,0.5720502023129355) (0.13567839195979894,0.5725741203559332) (0.1457286432160804,0.5733748068282001) (0.1557788944723619,0.573501663100906) (0.16582914572864316,0.5742199243964297) (0.17587939698492458,0.5737018971003905) (0.18592964824120614,0.5739388275781654) (0.19597989949748731,0.5732565585243374) (0.20603015075376882,0.5714419803080608) (0.21608040201005027,0.5734697102683701) (0.2261306532663318,0.5720499514897086) (0.236180904522613,0.5720499514897086) (0.24623115577889448,0.572608628035626) (0.25628140703517593,0.5738788646562509) (0.2663316582914572,0.5748296202873902) (0.27638190954773867,0.573643287950939) (0.2864321608040202,0.5739279704930978) (0.2964824120603014,0.5738879087688873) (0.30653266331658285,0.5764153710026018) (0.31658291457286425,0.5768417572612348) (0.32663316582914576,0.5750217550797524) (0.33668341708542715,0.5755299632479113) (0.3467336683417085,0.5749373886582153) (0.35678391959799,0.5739989121289317) (0.36683417085427134,0.5745581648444646) (0.3768844221105528,0.5746167913357635) (0.3869346733668342,0.574102484896845) (0.3969849246231155,0.5734374853118225) (0.407035175879397,0.5728977458374819) (0.4170854271356784,0.5728971682684391) (0.42713567839195976,0.5722109344360452) (0.43718592964824116,0.5720260345069681) (0.4472361809045226,0.5719795092814676) (0.45728643216080406,0.5721186586694382) (0.46733668341708545,0.572737274781547) (0.4773869346733668,0.5721123085482662) (0.4874371859296483,0.5726585642469213) (0.49748743718592964,0.5733518944307477) (0.5075376884422109,0.5718368410474144) (0.5175879396984924,0.5735089903564228) (0.5276381909547738,0.5740599589979227) (0.5376884422110553,0.5730771721671721) (0.5477386934673366,0.57258290625026) (0.5577889447236181,0.5712953453688603) (0.5678391959798995,0.5722729271440007) (0.5778894472361809,0.571762425294479) (0.5879396984924622,0.5727595873502052) (0.5979899497487436,0.5727595873502052) (0.6080402010050252,0.5726416885235149) (0.6180904522613064,0.5740030397376142) (0.6281407035175879,0.5740030397376142) (0.6381909547738693,0.5741176954391728) (0.6482412060301508,0.5734535344493608) (0.6582914572864322,0.5734535344493608) (0.6683417085427136,0.5747717986370383) (0.678391959798995,0.5747717986370383) (0.6884422110552764,0.5748488874712885) (0.6984924623115577,0.5756966104434211) (0.7085427135678392,0.5756966104434211) (0.7185929648241205,0.5752415195224517) (0.728643216080402,0.5744694537019357) (0.7386934673366834,0.574385265272158) (0.7487437185929648,0.5750767728887174) (0.7587939698492462,0.5751618277222954) (0.7688442211055276,0.5761071844263672) (0.778894472361809,0.5759744897244469) (0.7889447236180904,0.5760241434177944) (0.7989949748743718,0.5753916428057574) (0.8090452261306532,0.5760009185631736) (0.8190954773869347,0.5758874726656975) (0.829145728643216,0.5765410435635314) (0.8391959798994975,0.5759544540806633) (0.8492462311557789,0.575315629223159) (0.8592964824120602,0.5755084224993343) (0.8693467336683416,0.5760779923274997) (0.8793969849246231,0.5753988248649367) (0.8894472361809045,0.575360102248534) (0.8994974874371859,0.5748939925935725) (0.9095477386934674,0.5737625868606282) (0.9195979899497487,0.572918133047136) (0.9296482412060301,0.5733765033978245) (0.9396984924623116,0.5740428380768176) (0.949748743718593,0.5740428380768176) (0.9597989949748743,0.5736454588407175) (0.9698492462311558,0.573181098644464) (0.9798994974874372,0.5727519497999215) (0.9899497487437185,0.5727519497999215) (0.9999999999999999,0.5727519497999215) };
  \addplot[name path=L_ig_mle_grp, draw=none, mark=none, forget plot] coordinates {(-0.9999999999999999,0.5614709534445189) (-0.9899497487437185,0.5582904185849298) (-0.979899497487437,0.559636244804242) (-0.9698492462311556,0.5535644086391962) (-0.9597989949748744,0.5509845228240856) (-0.9497487437185927,0.5525217006914903) (-0.9396984924623114,0.5535968963347835) (-0.9296482412060302,0.551812241580222) (-0.9195979899497485,0.5508471519708611) (-0.9095477386934673,0.5494536971490894) (-0.8994974874371859,0.5482150431576315) (-0.8894472361809044,0.5467194358241113) (-0.879396984924623,0.5465822271355647) (-0.8693467336683418,0.5472242762264916) (-0.85929648241206,0.5466155052974674) (-0.8492462311557789,0.5446690513554846) (-0.8391959798994976,0.5427017783648947) (-0.8291457286432159,0.5413658225172007) (-0.8190954773869346,0.5399455472565468) (-0.809045226130653,0.538486964032704) (-0.7989949748743717,0.5345219203666713) (-0.7889447236180904,0.5331934425232661) (-0.7788944723618088,0.5345082793277178) (-0.7688442211055275,0.5349487808442307) (-0.7587939698492462,0.5307750436713027) (-0.7487437185929646,0.5277652796932841) (-0.7386934673366833,0.5255232331335966) (-0.728643216080402,0.5227921810943928) (-0.7185929648241204,0.522637859748369) (-0.708542713567839,0.5194655432222335) (-0.6984924623115578,0.5163138029473273) (-0.6884422110552761,0.5156177291631386) (-0.6783919597989948,0.513890832965125) (-0.6683417085427136,0.5147100321063949) (-0.658291457286432,0.5155867375413684) (-0.6482412060301507,0.5133205270037164) (-0.6381909547738694,0.5128636153418488) (-0.6281407035175878,0.5117099161976295) (-0.6180904522613064,0.5099533057649328) (-0.6080402010050252,0.510155535588959) (-0.5979899497487435,0.5077970427903482) (-0.5879396984924623,0.5062788115186435) (-0.577889447236181,0.5056393856172587) (-0.5678391959798993,0.5042300572898248) (-0.5577889447236181,0.5045857960576161) (-0.5477386934673366,0.5053906077741229) (-0.5376884422110553,0.5055698860950939) (-0.5276381909547738,0.5046275131881969) (-0.5175879396984924,0.5044381403130742) (-0.5075376884422109,0.5051455152362903) (-0.49748743718592964,0.5034039360107512) (-0.4874371859296482,0.5014614348630109) (-0.4773869346733667,0.5029898054788182) (-0.46733668341708545,0.5032444870526511) (-0.45728643216080395,0.5033395581244823) (-0.44723618090452255,0.5057334604995656) (-0.4371859296482412,0.5085671031435246) (-0.42713567839195976,0.5066270995122737) (-0.41708542713567825,0.5066903774302011) (-0.4070351758793971,0.5059993325017005) (-0.3969849246231155,0.5075501662450164) (-0.38693467336683407,0.5070068903854028) (-0.3768844221105528,0.5034491962193817) (-0.36683417085427134,0.5016033852231155) (-0.35678391959798983,0.5024131345081652) (-0.34673366834170843,0.5016481414085086) (-0.3366834170854271,0.50343842966541) (-0.32663316582914564,0.5034857737044189) (-0.31658291457286425,0.5027014968036785) (-0.30653266331658297,0.4987270298752889) (-0.29648241206030146,0.49731276305444716) (-0.28643216080402,0.495279398626116) (-0.2763819095477387,0.49538644569921) (-0.2663316582914572,0.4968924145929748) (-0.2562814070351758,0.4955213964006763) (-0.2462311557788945,0.4909871784520815) (-0.236180904522613,0.48970396228135094) (-0.22613065326633158,0.4902481182874901) (-0.21608040201005027,0.4889431387020161) (-0.20603015075376882,0.4891992843624763) (-0.19597989949748737,0.49116482088873503) (-0.18592964824120603,0.4922526372093294) (-0.1758793969849246,0.49599778954072393) (-0.1658291457286431,0.496614818868827) (-0.15577889447236168,0.49620269093366093) (-0.14572864321608042,0.4945210726758894) (-0.13567839195979894,0.49469844386342793) (-0.12562814070351747,0.4918154412486105) (-0.11557788944723618,0.4945914160393259) (-0.10552763819095472,0.49252396795517994) (-0.09547738693467328,0.4925026280931226) (-0.08542713567839201,0.492681832075448) (-0.07537688442211048,0.4927272350061976) (-0.06532663316582907,0.49380170071370477) (-0.055276381909547756,0.49253641879050886) (-0.04522613065326632,0.49366528984204233) (-0.03517587939698485,0.49441754643608393) (-0.025125628140703578,0.49706831581489686) (-0.015075376884422061,0.4983353368200222) (-0.005025125628140596,0.49829894040369827) (0.00502512562814067,0.4996584258842043) (0.015075376884422073,0.5014369259233207) (0.02512562814070361,0.5034337834777971) (0.035175879396985056,0.5046090364968704) (0.045226130653266264,0.5057206595545521) (0.05527638190954779,0.5058413042904361) (0.06532663316582923,0.5032462179464271) (0.0753768844221105,0.5031423219601541) (0.08542713567839198,0.504548110249606) (0.09547738693467347,0.5047618918311306) (0.10552763819095473,0.504854710263207) (0.11557788944723618,0.5046763162567198) (0.1256281407035177,0.5032355735470899) (0.13567839195979894,0.5032702348881054) (0.1457286432160804,0.5035403955357027) (0.1557788944723619,0.5034522843811109) (0.16582914572864316,0.5041658019547428) (0.17587939698492458,0.5035273005267199) (0.18592964824120614,0.503397266820562) (0.19597989949748731,0.502784701862768) (0.20603015075376882,0.5008417119639965) (0.21608040201005027,0.5026216242992427) (0.2261306532663318,0.5004083866327604) (0.236180904522613,0.5004083866327604) (0.24623115577889448,0.5007466883059337) (0.25628140703517593,0.5021811759266561) (0.2663316582914572,0.5042265210780634) (0.27638190954773867,0.5030770682716235) (0.2864321608040202,0.5022177231138322) (0.2964824120603014,0.5021123306376363) (0.30653266331658285,0.5047262948271612) (0.31658291457286425,0.5062188675411023) (0.32663316582914576,0.5053684665651695) (0.33668341708542715,0.5059176936124065) (0.3467336683417085,0.5053173906620884) (0.35678391959799,0.5049413743969953) (0.36683417085427134,0.5055354506509586) (0.3768844221105528,0.5055526264252513) (0.3869346733668342,0.5048866922787147) (0.3969849246231155,0.5042493144920837) (0.407035175879397,0.5034448791555785) (0.4170854271356784,0.5033209896895733) (0.42713567839195976,0.5026722157213187) (0.43718592964824116,0.5027692594954678) (0.4472361809045226,0.5029166446114023) (0.45728643216080406,0.5027908287085914) (0.46733668341708545,0.5034584508072436) (0.4773869346733668,0.5026160691492787) (0.4874371859296483,0.5033792488022937) (0.49748743718592964,0.5039190316233545) (0.5075376884422109,0.5030548222584992) (0.5175879396984924,0.5041659205414994) (0.5276381909547738,0.5046752735647049) (0.5376884422110553,0.5051441464018473) (0.5477386934673366,0.5044602979417846) (0.5577889447236181,0.5032142463540168) (0.5678391959798995,0.5044556006003486) (0.5778894472361809,0.5038804191479495) (0.5879396984924622,0.5054577173979216) (0.5979899497487436,0.5054577173979216) (0.6080402010050252,0.5054183661176376) (0.6180904522613064,0.5064883967700448) (0.6281407035175879,0.5064883967700448) (0.6381909547738693,0.5055435709484383) (0.6482412060301508,0.5049821603372908) (0.6582914572864322,0.5049821603372908) (0.6683417085427136,0.5059700556588486) (0.678391959798995,0.5059700556588486) (0.6884422110552764,0.5067968696701837) (0.6984924623115577,0.507017924069735) (0.7085427135678392,0.507017924069735) (0.7185929648241205,0.5062958790014808) (0.728643216080402,0.5056671853086991) (0.7386934673366834,0.5050156106581493) (0.7487437185929648,0.5055144193248249) (0.7587939698492462,0.5063127150360314) (0.7688442211055276,0.5084991511480164) (0.778894472361809,0.5086620013898694) (0.7889447236180904,0.5085421979770552) (0.7989949748743718,0.5079827165920171) (0.8090452261306532,0.508670273859682) (0.8190954773869347,0.5080579673310244) (0.829145728643216,0.5088595783948382) (0.8391959798994975,0.5081801073945945) (0.8492462311557789,0.5075052291264147) (0.8592964824120602,0.5071139425593452) (0.8693467336683416,0.5078108179957386) (0.8793969849246231,0.5073908169929786) (0.8894472361809045,0.5075572910432253) (0.8994974874371859,0.5077568582497873) (0.9095477386934674,0.5063296271154722) (0.9195979899497487,0.5069481911203604) (0.9296482412060301,0.5077232903450463) (0.9396984924623116,0.5093143451887469) (0.949748743718593,0.5093143451887469) (0.9597989949748743,0.5095541521589327) (0.9698492462311558,0.5089022382791902) (0.9798994974874372,0.5081141622495388) (0.9899497487437185,0.5081141622495388) (0.9999999999999999,0.5081141622495388) };
  \addplot[ours_ig_mle, fill, fill opacity=0.18, draw=none, mark=none, forget plot] fill between [of=U_ig_mle_grp and L_ig_mle_grp];
  \addplot+[ours_ig_mle, mark=none] coordinates {(-0.9999999999999999,0.5955520439977978) (-0.9899497487437185,0.5929277265810954) (-0.979899497487437,0.5936882888551668) (-0.9698492462311556,0.5878142392891459) (-0.9597989949748744,0.5860787157685787) (-0.9497487437185927,0.5867015755592577) (-0.9396984924623114,0.5870594606993166) (-0.9296482412060302,0.5861437306064549) (-0.9195979899497485,0.584564627529614) (-0.9095477386934673,0.5840115578953089) (-0.8994974874371859,0.5827663115774684) (-0.8894472361809044,0.5816754031485563) (-0.879396984924623,0.5816558145193782) (-0.8693467336683418,0.5807228569994695) (-0.85929648241206,0.5797589098774516) (-0.8492462311557789,0.577960887234124) (-0.8391959798994976,0.5766246338742009) (-0.8291457286432159,0.5747958241704513) (-0.8190954773869346,0.5734361204016873) (-0.809045226130653,0.5724451675623099) (-0.7989949748743717,0.5685710397782744) (-0.7889447236180904,0.5671932366318085) (-0.7788944723618088,0.5683821092653842) (-0.7688442211055275,0.5681914130700085) (-0.7587939698492462,0.5640233526594227) (-0.7487437185929646,0.5604621983832002) (-0.7386934673366833,0.5571174585709088) (-0.728643216080402,0.5545875079291075) (-0.7185929648241204,0.5545360489298382) (-0.708542713567839,0.5519986329652212) (-0.6984924623115578,0.5487386401617245) (-0.6884422110552761,0.5480630197082883) (-0.6783919597989948,0.546313595754631) (-0.6683417085427136,0.5467775898747579) (-0.658291457286432,0.5481633306024913) (-0.6482412060301507,0.5467746066644893) (-0.6381909547738694,0.5461586158203384) (-0.6281407035175878,0.5449803266048563) (-0.6180904522613064,0.5433307875504888) (-0.6080402010050252,0.5438979232032025) (-0.5979899497487435,0.5408277393028914) (-0.5879396984924623,0.5390263384580005) (-0.577889447236181,0.5383687354872797) (-0.5678391959798993,0.5369306138992666) (-0.5577889447236181,0.5368784464550874) (-0.5477386934673366,0.53745729928915) (-0.5376884422110553,0.5374570005194667) (-0.5276381909547738,0.5366733012840509) (-0.5175879396984924,0.5365335443365363) (-0.5075376884422109,0.5371906019908405) (-0.49748743718592964,0.5362190489962472) (-0.4874371859296482,0.5346019150270647) (-0.4773869346733667,0.5355052221809629) (-0.46733668341708545,0.5360006371365522) (-0.45728643216080395,0.5360119236341241) (-0.44723618090452255,0.5380858849545005) (-0.4371859296482412,0.5402898916234061) (-0.42713567839195976,0.5392773311818662) (-0.41708542713567825,0.5392048205643135) (-0.4070351758793971,0.5385764365247181) (-0.3969849246231155,0.5401748439384985) (-0.38693467336683407,0.5395587344458848) (-0.3768844221105528,0.5369161682962783) (-0.36683417085427134,0.5347128849783883) (-0.35678391959798983,0.5348659790047878) (-0.34673366834170843,0.5344049826253178) (-0.3366834170854271,0.5358695854858485) (-0.32663316582914564,0.5358886466001801) (-0.31658291457286425,0.5352635437926785) (-0.30653266331658297,0.5321062034490648) (-0.29648241206030146,0.5305442843436152) (-0.28643216080402,0.528614507990123) (-0.2763819095477387,0.529115152160558) (-0.2663316582914572,0.530467819878702) (-0.2562814070351758,0.5295513203021638) (-0.2462311557788945,0.525676767167367) (-0.236180904522613,0.5243390366879992) (-0.22613065326633158,0.5249517195184792) (-0.21608040201005027,0.5239599788296705) (-0.20603015075376882,0.5240134600656288) (-0.19597989949748737,0.5259792246036813) (-0.18592964824120603,0.5275992001378161) (-0.1758793969849246,0.5309669504693656) (-0.1658291457286431,0.5314944585547134) (-0.15577889447236168,0.5309297638825153) (-0.14572864321608042,0.529631629175368) (-0.13567839195979894,0.5297331170581709) (-0.12562814070351747,0.527046404248103) (-0.11557788944723618,0.52980523456942) (-0.10552763819095472,0.5277723088402402) (-0.09547738693467328,0.527770646001368) (-0.08542713567839201,0.5273249726528703) (-0.07537688442211048,0.5273720465198091) (-0.06532663316582907,0.5281798615214067) (-0.055276381909547756,0.5266712124961067) (-0.04522613065326632,0.5280311896997889) (-0.03517587939698485,0.5286747161382522) (-0.025125628140703578,0.5303098832794002) (-0.015075376884422061,0.5313475738161477) (-0.005025125628140596,0.5313836529743723) (0.00502512562814067,0.5319339635379797) (0.015075376884422073,0.5343140434452882) (0.02512562814070361,0.535778567405777) (0.035175879396985056,0.537681706029624) (0.045226130653266264,0.5388662783933299) (0.05527638190954779,0.539062736595555) (0.06532663316582923,0.5366050346073589) (0.0753768844221105,0.5361949595430597) (0.08542713567839198,0.5376841648502168) (0.09547738693467347,0.538538712629551) (0.10552763819095473,0.5385576673340351) (0.11557788944723618,0.5385033065419423) (0.1256281407035177,0.5376428879300127) (0.13567839195979894,0.5379221776220193) (0.1457286432160804,0.5384576011819514) (0.1557788944723619,0.5384769737410084) (0.16582914572864316,0.5391928631755862) (0.17587939698492458,0.5386145988135552) (0.18592964824120614,0.5386680471993637) (0.19597989949748731,0.5380206301935527) (0.20603015075376882,0.5361418461360287) (0.21608040201005027,0.5380456672838064) (0.2261306532663318,0.5362291690612345) (0.236180904522613,0.5362291690612345) (0.24623115577889448,0.5366776581707798) (0.25628140703517593,0.5380300202914535) (0.2663316582914572,0.5395280706827268) (0.27638190954773867,0.5383601781112812) (0.2864321608040202,0.538072846803465) (0.2964824120603014,0.5380001197032618) (0.30653266331658285,0.5405708329148815) (0.31658291457286425,0.5415303124011686) (0.32663316582914576,0.540195110822461) (0.33668341708542715,0.5407238284301589) (0.3467336683417085,0.5401273896601518) (0.35678391959799,0.5394701432629635) (0.36683417085427134,0.5400468077477116) (0.3768844221105528,0.5400847088805074) (0.3869346733668342,0.5394945885877799) (0.3969849246231155,0.5388433999019531) (0.407035175879397,0.5381713124965302) (0.4170854271356784,0.5381090789790062) (0.42713567839195976,0.537441575078682) (0.43718592964824116,0.537397647001218) (0.4472361809045226,0.5374480769464349) (0.45728643216080406,0.5374547436890148) (0.46733668341708545,0.5380978627943953) (0.4773869346733668,0.5373641888487725) (0.4874371859296483,0.5380189065246075) (0.49748743718592964,0.5386354630270511) (0.5075376884422109,0.5374458316529568) (0.5175879396984924,0.5388374554489611) (0.5276381909547738,0.5393676162813138) (0.5376884422110553,0.5391106592845097) (0.5477386934673366,0.5385216020960223) (0.5577889447236181,0.5372547958614385) (0.5678391959798995,0.5383642638721746) (0.5778894472361809,0.5378214222212142) (0.5879396984924622,0.5391086523740634) (0.5979899497487436,0.5391086523740634) (0.6080402010050252,0.5390300273205763) (0.6180904522613064,0.5402457182538295) (0.6281407035175879,0.5402457182538295) (0.6381909547738693,0.5398306331938055) (0.6482412060301508,0.5392178473933258) (0.6582914572864322,0.5392178473933258) (0.6683417085427136,0.5403709271479434) (0.678391959798995,0.5403709271479434) (0.6884422110552764,0.5408228785707361) (0.6984924623115577,0.5413572672565781) (0.7085427135678392,0.5413572672565781) (0.7185929648241205,0.5407686992619662) (0.728643216080402,0.5400683195053174) (0.7386934673366834,0.5397004379651537) (0.7487437185929648,0.5402955961067711) (0.7587939698492462,0.5407372713791634) (0.7688442211055276,0.5423031677871918) (0.778894472361809,0.5423182455571581) (0.7889447236180904,0.5422831706974248) (0.7989949748743718,0.5416871796988872) (0.8090452261306532,0.5423355962114278) (0.8190954773869347,0.5419727199983609) (0.829145728643216,0.5427003109791848) (0.8391959798994975,0.5420672807376289) (0.8492462311557789,0.5414104291747869) (0.8592964824120602,0.5413111825293397) (0.8693467336683416,0.5419444051616191) (0.8793969849246231,0.5413948209289576) (0.8894472361809045,0.5414586966458796) (0.8994974874371859,0.5413254254216799) (0.9095477386934674,0.5400461069880502) (0.9195979899497487,0.5399331620837482) (0.9296482412060301,0.5405498968714354) (0.9396984924623116,0.5416785916327822) (0.949748743718593,0.5416785916327822) (0.9597989949748743,0.5415998054998251) (0.9698492462311558,0.5410416684618271) (0.9798994974874372,0.5404330560247301) (0.9899497487437185,0.5404330560247301) (0.9999999999999999,0.5404330560247301) };
  \addplot+[ours_ig_mle, mark=none, line width=0.4pt, dashed, forget plot] coordinates {(-1.0,0.4414050958348127) (1.0,0.4414050958348127)};
  \addplot+[ours_ig_mle, mark=*, only marks, mark size=\PanelArgminSize, forget plot] coordinates {(-0.21608040201005027,0.5239599788296705)};
  \addplot+[ours_oracle, mark=none, forget plot] coordinates {(-1.0,0.19700210717107242) (1.0,0.19700210717107242)};
%
			\renewcommand\nextgroupplot[1][]{\origNGP[#1, ymin=1.2, ymax=12,
					ytick={2,4,6,8,10,12}]}%
			% panel_profile_group.tex — \nextgroupplot fragment.
% Host preamble MUST load: \usepgfplotslibrary{groupplots},
%                          \usepgfplotslibrary{fillbetween}.
% Must be \input inside a \begin{groupplot}[...] with these
% options set at the GROUPPLOT level (not per-cell):
%   panel_axis, xmin=-1, xmax=1, xtick={\PanelXTicks},
%   xlabel={\PanelXLabel}, ylabel={\PanelYLabel},
%   ylabels at=edge left, yticklabels at=edge left,
%   group style={group size=N by 1, horizontal sep=...}
% That way pgfplots' native `ylabels at=edge left` shows the y-
% annotations only on the leftmost cell.  Per-cell options here
% are LIMITED to `title=` so the row's wrapper retains all axis
% control.  See `profile_row_template` for the canonical wrapper.

\nextgroupplot[title={HEB / SIFT / \(\bH\)}]
  \addplot[name path=U_ransac_grp, draw=none, mark=none, forget plot] coordinates {(-1.0,11.398714531298182) (-0.9899497487437188,11.398714531298182) (-0.979899497487437,11.398714531298182) (-0.9698492462311559,11.398714531298182) (-0.9597989949748745,11.398714531298182) (-0.9497487437185929,11.398714531298182) (-0.9396984924623116,11.398714531298182) (-0.9296482412060304,10.807347229622188) (-0.9195979899497487,10.807347229622188) (-0.9095477386934674,10.807347229622188) (-0.899497487437186,10.807347229622188) (-0.8894472361809044,10.807347229622188) (-0.8793969849246231,10.807347229622188) (-0.8693467336683419,10.807347229622188) (-0.8592964824120602,10.394512737033637) (-0.8492462311557789,10.394512737033637) (-0.8391959798994972,10.394512737033637) (-0.829145728643216,10.394512737033637) (-0.8190954773869348,10.394512737033637) (-0.8090452261306531,10.394512737033637) (-0.7989949748743719,10.054009380268077) (-0.7889447236180905,10.054009380268077) (-0.7788944723618089,10.054009380268077) (-0.7688442211055276,10.054009380268077) (-0.7587939698492464,10.054009380268077) (-0.7487437185929647,10.054009380268077) (-0.7386934673366834,9.732856410208713) (-0.7286432160804021,9.732856410208713) (-0.7185929648241205,9.732856410208713) (-0.7085427135678393,9.732856410208713) (-0.6984924623115579,9.50603565390598) (-0.6884422110552763,9.50603565390598) (-0.678391959798995,9.50603565390598) (-0.6683417085427137,9.50603565390598) (-0.6582914572864321,9.50603565390598) (-0.6482412060301508,9.2805773371328) (-0.6381909547738696,9.2805773371328) (-0.6281407035175879,9.2805773371328) (-0.6180904522613067,9.2805773371328) (-0.6080402010050253,9.15060017171945) (-0.5979899497487436,9.15060017171945) (-0.5879396984924624,9.15060017171945) (-0.5778894472361811,8.968676509848963) (-0.5678391959798994,8.968676509848963) (-0.5577889447236181,8.968676509848963) (-0.5477386934673368,8.968676509848963) (-0.5376884422110553,8.89149082612151) (-0.5276381909547739,8.89149082612151) (-0.5175879396984923,8.89149082612151) (-0.507537688442211,8.73633396982174) (-0.4974874371859298,8.73633396982174) (-0.4874371859296481,8.704925992070391) (-0.47738693467336674,8.704925992070391) (-0.4673366834170855,8.704925992070391) (-0.45728643216080384,8.590176625112234) (-0.4472361809045226,8.590176625112234) (-0.4371859296482414,8.590176625112234) (-0.4271356783919596,8.550907361802551) (-0.4170854271356784,8.550907361802551) (-0.40703517587939714,8.508281455760248) (-0.39698492462311546,8.508281455760248) (-0.3869346733668342,8.431697518757597) (-0.3768844221105529,8.431697518757597) (-0.3668341708542712,8.431697518757597) (-0.35678391959798994,8.368635588781244) (-0.3467336683417087,8.368635588781244) (-0.33668341708542704,8.296644139039664) (-0.3266331658291457,8.296644139039664) (-0.3165829145728644,8.31775339993833) (-0.30653266331658285,8.304732595159841) (-0.29648241206030157,8.304732595159841) (-0.2864321608040203,8.299245123793078) (-0.2763819095477386,8.299245123793078) (-0.26633165829145733,8.309718630850176) (-0.2562814070351761,8.309718630850176) (-0.24623115577889443,8.333738714585902) (-0.23618090452261317,8.242942017672636) (-0.2261306532663315,8.242942017672636) (-0.2160804020100502,8.22976429053788) (-0.206030150753769,8.229733113871006) (-0.19597989949748731,8.229733113871006) (-0.18592964824120595,8.215600516173103) (-0.1758793969849247,8.207023233531967) (-0.1658291457286431,8.260589216673308) (-0.15577889447236173,8.260589216673308) (-0.14572864321608053,8.23906951339002) (-0.13567839195979886,8.276788143560353) (-0.12562814070351758,8.26292545780994) (-0.1155778894472363,8.26292545780994) (-0.10552763819095466,8.249182430372562) (-0.09547738693467334,8.197557266479485) (-0.08542713567839214,8.231768649193517) (-0.07537688442211042,8.19392170014302) (-0.06532663316582912,8.210949436836183) (-0.05527638190954792,8.192027105690661) (-0.045226130653266264,8.235753461362377) (-0.03517587939698496,8.23742536029767) (-0.025125628140703626,8.289428178553141) (-0.015075376884422013,8.279072421853707) (-0.005025125628140791,8.27513717269639) (0.005025125628140479,8.28801489497626) (0.015075376884422259,8.291614879673942) (0.025125628140703425,8.300405062482888) (0.0351758793969847,8.30637302694954) (0.045226130653266354,8.326848422181799) (0.0552763819095477,8.381198023207082) (0.0653266331658289,8.387616088497683) (0.07537688442211066,8.349529884498315) (0.08542713567839183,8.324079993755898) (0.09547738693467356,8.36611140284091) (0.10552763819095481,8.426730487937808) (0.11557788944723611,8.439237278772667) (0.12562814070351774,8.439502351718685) (0.13567839195979903,8.446790741937406) (0.1457286432160803,8.468598934066042) (0.15577889447236196,8.48606631252187) (0.16582914572864324,8.525088156117203) (0.17587939698492452,8.522873460310294) (0.1859296482412062,8.57990167675485) (0.19597989949748745,8.615667143394363) (0.2060301507537687,8.678183137857499) (0.2160804020100504,8.685833547827556) (0.2261306532663317,8.700192875381866) (0.2361809045226129,8.710327162671186) (0.24623115577889454,8.756178220680614) (0.2562814070351758,8.765988833321684) (0.2663316582914571,8.769662056628393) (0.2763819095477388,8.787434100826257) (0.28643216080402006,8.764584551486099) (0.29648241206030135,8.789128121631215) (0.30653266331658297,8.818745869526184) (0.31658291457286436,8.849426931287741) (0.3266331658291455,8.85658458834356) (0.33668341708542726,8.906665312122685) (0.3467336683417085,8.918118795551536) (0.3567839195979898,8.947540050195133) (0.36683417085427145,8.958338444970643) (0.37688442211055273,8.98124300348846) (0.38693467336683435,9.032432692275805) (0.39698492462311563,9.057408932953221) (0.40703517587939686,9.089922533381804) (0.41708542713567864,9.117217036636898) (0.42713567839195987,9.151454080338018) (0.4371859296482411,9.208884429986165) (0.4472361809045228,9.245803072484328) (0.457286432160804,9.315908887327396) (0.4673366834170853,9.356514160303126) (0.47738693467336696,9.375941065015084) (0.48743718592964824,9.420152782526344) (0.4974874371859295,9.467467901641024) (0.5075376884422113,9.499187707861006) (0.5175879396984925,9.553869178404867) (0.5276381909547737,9.510148973984814) (0.5376884422110554,9.555995398855755) (0.5477386934673366,9.665190907503321) (0.557788944723618,9.628390089454989) (0.5678391959798996,9.721355039303612) (0.5778894472361809,9.75381803848228) (0.5879396984924622,9.74379030679584) (0.5979899497487436,9.771873527301937) (0.6080402010050252,9.798063331140192) (0.6180904522613065,9.90073188304022) (0.628140703517588,9.905277141228433) (0.6381909547738693,9.974492593601894) (0.6482412060301508,9.984827986889936) (0.6582914572864322,10.074885921174076) (0.6683417085427136,10.113029946566357) (0.678391959798995,10.092695967031814) (0.6884422110552765,10.176833688326548) (0.6984924623115577,10.246503424156636) (0.7085427135678392,10.280168131925235) (0.7185929648241206,10.30447203166404) (0.728643216080402,10.35101142514936) (0.7386934673366834,10.388429963814144) (0.7487437185929648,10.512290314446279) (0.7587939698492461,10.553588384842651) (0.7688442211055276,10.668042208900028) (0.7788944723618091,10.771238330243756) (0.7889447236180905,10.75777263280279) (0.7989949748743718,10.743561197908802) (0.8090452261306533,10.732147207655636) (0.8190954773869347,10.84173758843619) (0.829145728643216,10.885435932248008) (0.8391959798994975,10.9568640305368) (0.8492462311557789,11.025695981530149) (0.8592964824120601,11.010307171996807) (0.8693467336683417,11.057479057320682) (0.8793969849246231,11.156593944484255) (0.8894472361809044,11.249019444960757) (0.8994974874371858,11.333976459725026) (0.9095477386934674,11.290314716675569) (0.9195979899497487,11.439275970599635) (0.9296482412060301,11.594264170144404) (0.9396984924623116,11.60035405868004) (0.949748743718593,11.65772239920057) (0.9597989949748743,11.731798568558396) (0.9698492462311559,11.721605618117698) (0.9798994974874372,11.885084225082693) (0.9899497487437185,11.944302657749896) (1.0,12.036204873990522) };
  \addplot[name path=L_ransac_grp, draw=none, mark=none, forget plot] coordinates {(-1.0,10.73247753789645) (-0.9899497487437188,10.73247753789645) (-0.979899497487437,10.73247753789645) (-0.9698492462311559,10.73247753789645) (-0.9597989949748745,10.73247753789645) (-0.9497487437185929,10.73247753789645) (-0.9396984924623116,10.73247753789645) (-0.9296482412060304,10.211034176128233) (-0.9195979899497487,10.211034176128233) (-0.9095477386934674,10.211034176128233) (-0.899497487437186,10.211034176128233) (-0.8894472361809044,10.211034176128233) (-0.8793969849246231,10.211034176128233) (-0.8693467336683419,10.211034176128233) (-0.8592964824120602,9.81466359924503) (-0.8492462311557789,9.81466359924503) (-0.8391959798994972,9.81466359924503) (-0.829145728643216,9.81466359924503) (-0.8190954773869348,9.81466359924503) (-0.8090452261306531,9.81466359924503) (-0.7989949748743719,9.492957212169195) (-0.7889447236180905,9.492957212169195) (-0.7788944723618089,9.492957212169195) (-0.7688442211055276,9.492957212169195) (-0.7587939698492464,9.492957212169195) (-0.7487437185929647,9.492957212169195) (-0.7386934673366834,9.211298692730841) (-0.7286432160804021,9.211298692730841) (-0.7185929648241205,9.211298692730841) (-0.7085427135678393,9.211298692730841) (-0.6984924623115579,9.015385969132671) (-0.6884422110552763,9.015385969132671) (-0.678391959798995,9.015385969132671) (-0.6683417085427137,9.015385969132671) (-0.6582914572864321,9.015385969132671) (-0.6482412060301508,8.801412521313814) (-0.6381909547738696,8.801412521313814) (-0.6281407035175879,8.801412521313814) (-0.6180904522613067,8.801412521313814) (-0.6080402010050253,8.703190037703088) (-0.5979899497487436,8.703190037703088) (-0.5879396984924624,8.703190037703088) (-0.5778894472361811,8.48816876984173) (-0.5678391959798994,8.48816876984173) (-0.5577889447236181,8.48816876984173) (-0.5477386934673368,8.48816876984173) (-0.5376884422110553,8.379246022556133) (-0.5276381909547739,8.379246022556133) (-0.5175879396984923,8.379246022556133) (-0.507537688442211,8.265330027187598) (-0.4974874371859298,8.265330027187598) (-0.4874371859296481,8.249421317983034) (-0.47738693467336674,8.249421317983034) (-0.4673366834170855,8.249421317983034) (-0.45728643216080384,8.16330000377353) (-0.4472361809045226,8.16330000377353) (-0.4371859296482414,8.16330000377353) (-0.4271356783919596,8.134045937086224) (-0.4170854271356784,8.134045937086224) (-0.40703517587939714,8.120080773420575) (-0.39698492462311546,8.120080773420575) (-0.3869346733668342,8.0141952358951) (-0.3768844221105529,8.0141952358951) (-0.3668341708542712,8.0141952358951) (-0.35678391959798994,7.9580899264600005) (-0.3467336683417087,7.9580899264600005) (-0.33668341708542704,7.928126855504369) (-0.3266331658291457,7.928126855504369) (-0.3165829145728644,7.926804315571794) (-0.30653266331658285,7.9352172589686445) (-0.29648241206030157,7.9352172589686445) (-0.2864321608040203,7.942832876061448) (-0.2763819095477386,7.942832876061448) (-0.26633165829145733,7.936322962097057) (-0.2562814070351761,7.936322962097057) (-0.24623115577889443,7.941393726948942) (-0.23618090452261317,7.87224344040693) (-0.2261306532663315,7.87224344040693) (-0.2160804020100502,7.879493631891975) (-0.206030150753769,7.869413192129173) (-0.19597989949748731,7.869413192129173) (-0.18592964824120595,7.871635617783533) (-0.1758793969849247,7.846949767992829) (-0.1658291457286431,7.896285463817881) (-0.15577889447236173,7.896285463817881) (-0.14572864321608053,7.8800367797617366) (-0.13567839195979886,7.929564771376438) (-0.12562814070351758,7.915678913382778) (-0.1155778894472363,7.915678913382778) (-0.10552763819095466,7.92199766933105) (-0.09547738693467334,7.8649758793736035) (-0.08542713567839214,7.8851725551210885) (-0.07537688442211042,7.856632505480962) (-0.06532663316582912,7.8709230512179955) (-0.05527638190954792,7.875790767101296) (-0.045226130653266264,7.921399728357435) (-0.03517587939698496,7.919252362605182) (-0.025125628140703626,7.940719749576341) (-0.015075376884422013,7.933262454574988) (-0.005025125628140791,7.916402695402433) (0.005025125628140479,7.932472512182286) (0.015075376884422259,7.922550653450882) (0.025125628140703425,7.950761206153833) (0.0351758793969847,7.954465876531409) (0.045226130653266354,7.96859431796138) (0.0552763819095477,8.006395542837412) (0.0653266331658289,8.003331654930925) (0.07537688442211066,7.982110033634916) (0.08542713567839183,7.960527111772997) (0.09547738693467356,7.980364054369017) (0.10552763819095481,8.050530371059862) (0.11557788944723611,8.066017531267013) (0.12562814070351774,8.078541040915889) (0.13567839195979903,8.079226333404932) (0.1457286432160803,8.093867109147007) (0.15577889447236196,8.104846657809432) (0.16582914572864324,8.144984822259914) (0.17587939698492452,8.13439161683878) (0.1859296482412062,8.190035843647092) (0.19597989949748745,8.199228853291812) (0.2060301507537687,8.285004949182639) (0.2160804020100504,8.28934893649063) (0.2261306532663317,8.320564571175606) (0.2361809045226129,8.342675151659884) (0.24623115577889454,8.370065161726746) (0.2562814070351758,8.399699922544455) (0.2663316582914571,8.407265954260689) (0.2763819095477388,8.420083446469228) (0.28643216080402006,8.389845213201781) (0.29648241206030135,8.410470206632382) (0.30653266331658297,8.425586954706175) (0.31658291457286436,8.43609386028697) (0.3266331658291455,8.445236778559012) (0.33668341708542726,8.49631600558041) (0.3467336683417085,8.503355662974055) (0.3567839195979898,8.540498319144074) (0.36683417085427145,8.568347956044398) (0.37688442211055273,8.57707131299331) (0.38693467336683435,8.60355782197525) (0.39698492462311563,8.625205317644161) (0.40703517587939686,8.65007303449333) (0.41708542713567864,8.690692560268552) (0.42713567839195987,8.720294118710873) (0.4371859296482411,8.784738419887509) (0.4472361809045228,8.798186699738835) (0.457286432160804,8.854360500233373) (0.4673366834170853,8.906030111404322) (0.47738693467336696,8.935937021159972) (0.48743718592964824,9.003026161770382) (0.4974874371859295,9.024730418432503) (0.5075376884422113,9.051932554848536) (0.5175879396984925,9.125619850445325) (0.5276381909547737,9.073623376988737) (0.5376884422110554,9.122072967766934) (0.5477386934673366,9.217635937548552) (0.557788944723618,9.18137910589989) (0.5678391959798996,9.29591510682337) (0.5778894472361809,9.36215036593222) (0.5879396984924622,9.383083827758886) (0.5979899497487436,9.371162355186284) (0.6080402010050252,9.428953750218323) (0.6180904522613065,9.524693558689147) (0.628140703517588,9.538957875125364) (0.6381909547738693,9.580060595918422) (0.6482412060301508,9.6020549733238) (0.6582914572864322,9.661177029267684) (0.6683417085427136,9.690384042119401) (0.678391959798995,9.676462951189913) (0.6884422110552765,9.742404218852835) (0.6984924623115577,9.801520415808639) (0.7085427135678392,9.853156540884394) (0.7185929648241206,9.873245657368397) (0.728643216080402,9.896692958262244) (0.7386934673366834,9.943240964041706) (0.7487437185929648,10.026007124729073) (0.7587939698492461,10.058315160279339) (0.7688442211055276,10.194497071357896) (0.7788944723618091,10.289398688605084) (0.7889447236180905,10.310204317082015) (0.7989949748743718,10.295086222839126) (0.8090452261306533,10.293721288571543) (0.8190954773869347,10.43829959758081) (0.829145728643216,10.472946697788249) (0.8391959798994975,10.55397853418311) (0.8492462311557789,10.58149870171819) (0.8592964824120601,10.557013091449619) (0.8693467336683417,10.63336467400264) (0.8793969849246231,10.725693145184083) (0.8894472361809044,10.839476246468424) (0.8994974874371858,10.90551096529132) (0.9095477386934674,10.885539785242463) (0.9195979899497487,10.958677091243501) (0.9296482412060301,11.081796698947281) (0.9396984924623116,11.086433314106795) (0.949748743718593,11.136137155254609) (0.9597989949748743,11.259275758294248) (0.9698492462311559,11.228012942143136) (0.9798994974874372,11.405548022502273) (0.9899497487437185,11.455868706633408) (1.0,11.527058786080515) };
  \addplot[ours_ransac, fill, fill opacity=0.18, draw=none, mark=none, forget plot] fill between [of=U_ransac_grp and L_ransac_grp];
  \addplot+[ours_ransac, mark=none] coordinates {(-1.0,11.065596034597316) (-0.9899497487437188,11.065596034597316) (-0.979899497487437,11.065596034597316) (-0.9698492462311559,11.065596034597316) (-0.9597989949748745,11.065596034597316) (-0.9497487437185929,11.065596034597316) (-0.9396984924623116,11.065596034597316) (-0.9296482412060304,10.50919070287521) (-0.9195979899497487,10.50919070287521) (-0.9095477386934674,10.50919070287521) (-0.899497487437186,10.50919070287521) (-0.8894472361809044,10.50919070287521) (-0.8793969849246231,10.50919070287521) (-0.8693467336683419,10.50919070287521) (-0.8592964824120602,10.104588168139333) (-0.8492462311557789,10.104588168139333) (-0.8391959798994972,10.104588168139333) (-0.829145728643216,10.104588168139333) (-0.8190954773869348,10.104588168139333) (-0.8090452261306531,10.104588168139333) (-0.7989949748743719,9.773483296218636) (-0.7889447236180905,9.773483296218636) (-0.7788944723618089,9.773483296218636) (-0.7688442211055276,9.773483296218636) (-0.7587939698492464,9.773483296218636) (-0.7487437185929647,9.773483296218636) (-0.7386934673366834,9.472077551469777) (-0.7286432160804021,9.472077551469777) (-0.7185929648241205,9.472077551469777) (-0.7085427135678393,9.472077551469777) (-0.6984924623115579,9.260710811519326) (-0.6884422110552763,9.260710811519326) (-0.678391959798995,9.260710811519326) (-0.6683417085427137,9.260710811519326) (-0.6582914572864321,9.260710811519326) (-0.6482412060301508,9.040994929223308) (-0.6381909547738696,9.040994929223308) (-0.6281407035175879,9.040994929223308) (-0.6180904522613067,9.040994929223308) (-0.6080402010050253,8.926895104711269) (-0.5979899497487436,8.926895104711269) (-0.5879396984924624,8.926895104711269) (-0.5778894472361811,8.728422639845347) (-0.5678391959798994,8.728422639845347) (-0.5577889447236181,8.728422639845347) (-0.5477386934673368,8.728422639845347) (-0.5376884422110553,8.635368424338822) (-0.5276381909547739,8.635368424338822) (-0.5175879396984923,8.635368424338822) (-0.507537688442211,8.500831998504669) (-0.4974874371859298,8.500831998504669) (-0.4874371859296481,8.477173655026712) (-0.47738693467336674,8.477173655026712) (-0.4673366834170855,8.477173655026712) (-0.45728643216080384,8.376738314442882) (-0.4472361809045226,8.376738314442882) (-0.4371859296482414,8.376738314442882) (-0.4271356783919596,8.342476649444388) (-0.4170854271356784,8.342476649444388) (-0.40703517587939714,8.314181114590411) (-0.39698492462311546,8.314181114590411) (-0.3869346733668342,8.222946377326348) (-0.3768844221105529,8.222946377326348) (-0.3668341708542712,8.222946377326348) (-0.35678391959798994,8.163362757620622) (-0.3467336683417087,8.163362757620622) (-0.33668341708542704,8.112385497272017) (-0.3266331658291457,8.112385497272017) (-0.3165829145728644,8.122278857755061) (-0.30653266331658285,8.119974927064243) (-0.29648241206030157,8.119974927064243) (-0.2864321608040203,8.121038999927263) (-0.2763819095477386,8.121038999927263) (-0.26633165829145733,8.123020796473616) (-0.2562814070351761,8.123020796473616) (-0.24623115577889443,8.137566220767422) (-0.23618090452261317,8.057592729039783) (-0.2261306532663315,8.057592729039783) (-0.2160804020100502,8.054628961214927) (-0.206030150753769,8.04957315300009) (-0.19597989949748731,8.04957315300009) (-0.18592964824120595,8.043618066978318) (-0.1758793969849247,8.026986500762398) (-0.1658291457286431,8.078437340245594) (-0.15577889447236173,8.078437340245594) (-0.14572864321608053,8.059553146575878) (-0.13567839195979886,8.103176457468395) (-0.12562814070351758,8.08930218559636) (-0.1155778894472363,8.08930218559636) (-0.10552763819095466,8.085590049851806) (-0.09547738693467334,8.031266572926544) (-0.08542713567839214,8.058470602157303) (-0.07537688442211042,8.025277102811991) (-0.06532663316582912,8.040936244027089) (-0.05527638190954792,8.033908936395978) (-0.045226130653266264,8.078576594859905) (-0.03517587939698496,8.078338861451426) (-0.025125628140703626,8.115073964064742) (-0.015075376884422013,8.106167438214348) (-0.005025125628140791,8.095769934049411) (0.005025125628140479,8.110243703579274) (0.015075376884422259,8.107082766562412) (0.025125628140703425,8.12558313431836) (0.0351758793969847,8.130419451740474) (0.045226130653266354,8.14772137007159) (0.0552763819095477,8.193796783022247) (0.0653266331658289,8.195473871714304) (0.07537688442211066,8.165819959066615) (0.08542713567839183,8.142303552764448) (0.09547738693467356,8.173237728604963) (0.10552763819095481,8.238630429498835) (0.11557788944723611,8.25262740501984) (0.12562814070351774,8.259021696317287) (0.13567839195979903,8.263008537671169) (0.1457286432160803,8.281233021606525) (0.15577889447236196,8.295456485165651) (0.16582914572864324,8.335036489188559) (0.17587939698492452,8.328632538574537) (0.1859296482412062,8.384968760200971) (0.19597989949748745,8.407447998343088) (0.2060301507537687,8.481594043520069) (0.2160804020100504,8.487591242159093) (0.2261306532663317,8.510378723278736) (0.2361809045226129,8.526501157165535) (0.24623115577889454,8.56312169120368) (0.2562814070351758,8.58284437793307) (0.2663316582914571,8.588464005444541) (0.2763819095477388,8.603758773647742) (0.28643216080402006,8.57721488234394) (0.29648241206030135,8.599799164131799) (0.30653266331658297,8.62216641211618) (0.31658291457286436,8.642760395787356) (0.3266331658291455,8.650910683451286) (0.33668341708542726,8.701490658851547) (0.3467336683417085,8.710737229262795) (0.3567839195979898,8.744019184669604) (0.36683417085427145,8.76334320050752) (0.37688442211055273,8.779157158240885) (0.38693467336683435,8.817995257125528) (0.39698492462311563,8.841307125298691) (0.40703517587939686,8.869997783937567) (0.41708542713567864,8.903954798452725) (0.42713567839195987,8.935874099524446) (0.4371859296482411,8.996811424936837) (0.4472361809045228,9.021994886111582) (0.457286432160804,9.085134693780384) (0.4673366834170853,9.131272135853724) (0.47738693467336696,9.155939043087528) (0.48743718592964824,9.211589472148363) (0.4974874371859295,9.246099160036763) (0.5075376884422113,9.275560131354771) (0.5175879396984925,9.339744514425096) (0.5276381909547737,9.291886175486775) (0.5376884422110554,9.339034183311345) (0.5477386934673366,9.441413422525937) (0.557788944723618,9.40488459767744) (0.5678391959798996,9.508635073063491) (0.5778894472361809,9.55798420220725) (0.5879396984924622,9.563437067277363) (0.5979899497487436,9.57151794124411) (0.6080402010050252,9.613508540679257) (0.6180904522613065,9.712712720864683) (0.628140703517588,9.722117508176899) (0.6381909547738693,9.777276594760158) (0.6482412060301508,9.793441480106868) (0.6582914572864322,9.86803147522088) (0.6683417085427136,9.901706994342879) (0.678391959798995,9.884579459110864) (0.6884422110552765,9.959618953589692) (0.6984924623115577,10.024011919982637) (0.7085427135678392,10.066662336404814) (0.7185929648241206,10.088858844516219) (0.728643216080402,10.123852191705803) (0.7386934673366834,10.165835463927925) (0.7487437185929648,10.269148719587676) (0.7587939698492461,10.305951772560995) (0.7688442211055276,10.431269640128962) (0.7788944723618091,10.53031850942442) (0.7889447236180905,10.533988474942403) (0.7989949748743718,10.519323710373964) (0.8090452261306533,10.51293424811359) (0.8190954773869347,10.6400185930085) (0.829145728643216,10.679191315018128) (0.8391959798994975,10.755421282359954) (0.8492462311557789,10.803597341624169) (0.8592964824120601,10.783660131723213) (0.8693467336683417,10.845421865661661) (0.8793969849246231,10.94114354483417) (0.8894472361809044,11.04424784571459) (0.8994974874371858,11.119743712508173) (0.9095477386934674,11.087927250959016) (0.9195979899497487,11.198976530921568) (0.9296482412060301,11.338030434545843) (0.9396984924623116,11.343393686393417) (0.949748743718593,11.396929777227589) (0.9597989949748743,11.495537163426322) (0.9698492462311559,11.474809280130417) (0.9798994974874372,11.645316123792483) (0.9899497487437185,11.700085682191652) (1.0,11.781631830035519) };
  \addplot+[ours_ransac, mark=none, line width=0.4pt, dashed, forget plot] coordinates {(-1.0,3.5701666848288167) (1.0,3.5701666848288167)};
  \addplot+[ours_ransac, mark=*, only marks, mark size=\PanelArgminSize, forget plot] coordinates {(-0.07537688442211042,8.025277102811991)};
  \addplot[name path=U_msac_grp, draw=none, mark=none, forget plot] coordinates {(-0.9999999999999999,10.689628971966366) (-0.9899497487437187,10.66916432218583) (-0.979899497487437,10.604589859068867) (-0.9698492462311558,10.580718338328092) (-0.9597989949748745,10.522138930124166) (-0.9497487437185927,10.471056369810915) (-0.9396984924623115,10.459984118733583) (-0.9296482412060303,10.315018046070744) (-0.9195979899497485,10.19045259870216) (-0.9095477386934673,10.090811487009228) (-0.899497487437186,10.04195698089176) (-0.8894472361809044,9.945463808307956) (-0.8793969849246231,9.91018353813938) (-0.8693467336683418,9.861417351793277) (-0.8592964824120601,9.771497101675466) (-0.8492462311557789,9.695969485867048) (-0.8391959798994972,9.664884801539568) (-0.829145728643216,9.600856458404884) (-0.8190954773869347,9.5731171001615) (-0.809045226130653,9.536541895121974) (-0.7989949748743718,9.493214823334451) (-0.7889447236180904,9.443592148563297) (-0.7788944723618089,9.345028642621235) (-0.7688442211055275,9.292792803696276) (-0.7587939698492463,9.24544723593237) (-0.7487437185929646,9.210651346180274) (-0.7386934673366833,9.138077993797575) (-0.7286432160804021,9.13045255769128) (-0.7185929648241204,9.104638259124876) (-0.7085427135678392,9.079528131876673) (-0.6984924623115578,9.031980917938299) (-0.6884422110552761,8.94775008130786) (-0.6783919597989949,8.91670792332136) (-0.6683417085427137,8.862555894934983) (-0.658291457286432,8.82227546693058) (-0.6482412060301508,8.755744894764064) (-0.6381909547738694,8.726012774877692) (-0.6281407035175878,8.674423917268053) (-0.6180904522613064,8.626622616554247) (-0.6080402010050252,8.595388962595813) (-0.5979899497487435,8.562289662381641) (-0.5879396984924623,8.539821359378282) (-0.577889447236181,8.515507080478933) (-0.5678391959798994,8.482054474909775) (-0.5577889447236181,8.459121701564763) (-0.5477386934673368,8.41920386059693) (-0.5376884422110552,8.3625758819322) (-0.5276381909547738,8.330708207413219) (-0.5175879396984923,8.286272608291116) (-0.507537688442211,8.2693174959136) (-0.4974874371859297,8.245833361922958) (-0.487437185929648,8.230883087813165) (-0.47738693467336674,8.1887874176097) (-0.4673366834170855,8.16593240678617) (-0.4572864321608038,8.154320888720688) (-0.44723618090452255,8.130601430004457) (-0.4371859296482412,8.117725566771464) (-0.4271356783919596,8.089789557613054) (-0.4170854271356783,8.067625248909753) (-0.4070351758793971,8.043354565236433) (-0.3969849246231154,8.036232045060506) (-0.38693467336683407,8.022750036536095) (-0.37688442211055284,8.003408035685018) (-0.36683417085427117,7.976697596617621) (-0.3567839195979899,7.967623018596053) (-0.3467336683417086,7.962087007028678) (-0.33668341708542704,7.952050211508854) (-0.3266331658291457,7.961320987036304) (-0.3165829145728644,7.9525952481963165) (-0.30653266331658274,7.941801219314117) (-0.29648241206030146,7.918871737228276) (-0.2864321608040203,7.904791926311976) (-0.27638190954773856,7.893408549370246) (-0.2663316582914572,7.883870529727276) (-0.256281407035176,7.874655800586028) (-0.24623115577889435,7.871457946569195) (-0.2361809045226131,7.862158419562881) (-0.22613065326633142,7.861409077555866) (-0.2160804020100501,7.865792058236025) (-0.2060301507537689,7.857011447963052) (-0.19597989949748723,7.859646640751442) (-0.18592964824120595,7.8714695927290785) (-0.1758793969849247,7.863966535942603) (-0.16582914572864296,7.845520313655091) (-0.15577889447236168,7.850245694431656) (-0.14572864321608048,7.842322735256944) (-0.1356783919597988,7.822297259617935) (-0.12562814070351752,7.811351552808255) (-0.11557788944723625,7.815745163412901) (-0.10552763819095459,7.819347045985891) (-0.09547738693467328,7.822021034115372) (-0.08542713567839201,7.821713648589853) (-0.07537688442211037,7.80834145719029) (-0.06532663316582912,7.808428138803647) (-0.05527638190954787,7.809750084260959) (-0.04522613065326616,7.816723278226618) (-0.0351758793969849,7.842390190314557) (-0.025125628140703626,7.853549487755966) (-0.015075376884421912,7.845140065749931) (-0.005025125628140694,7.844528242380742) (0.005025125628140574,7.847523981492156) (0.015075376884422259,7.866596029552207) (0.02512562814070352,7.850824120471139) (0.03517587939698479,7.846343273244527) (0.04522613065326644,7.843177817778541) (0.0552763819095477,7.845243229351263) (0.06532663316582898,7.867994459917052) (0.07537688442211066,7.866346161503637) (0.08542713567839198,7.895598764362077) (0.09547738693467363,7.902521251838353) (0.10552763819095488,7.8739342353143025) (0.11557788944723618,7.874310482869834) (0.12562814070351783,7.892619702598433) (0.13567839195979908,7.899851526148981) (0.1457286432160804,7.912492771184316) (0.15577889447236204,7.924955304009299) (0.1658291457286433,7.958488822291155) (0.17587939698492458,7.96941746901181) (0.18592964824120625,7.97068368904872) (0.1959798994974875,7.9795596778101086) (0.20603015075376876,7.981558788883736) (0.21608040201005047,7.965903853887161) (0.22613065326633175,7.9757978635530655) (0.236180904522613,7.976624277412608) (0.24623115577889465,8.004720244994914) (0.25628140703517593,8.042369812584923) (0.2663316582914572,8.060244316458517) (0.27638190954773884,8.06427309944078) (0.2864321608040201,8.075872429491207) (0.2964824120603014,8.112960011971412) (0.3065326633165831,8.140750601714862) (0.31658291457286436,8.152508402356455) (0.3266331658291457,8.164685624930586) (0.33668341708542726,8.163457878182566) (0.34673366834170855,8.187627258118589) (0.35678391959799,8.200122798523726) (0.3668341708542715,8.213119148374588) (0.3768844221105528,8.23999637087947) (0.3869346733668343,8.252563049134235) (0.3969849246231157,8.271382920897127) (0.40703517587939697,8.29897675663988) (0.4170854271356784,8.318953622024774) (0.42713567839195987,8.337945628677872) (0.43718592964824116,8.358374581077964) (0.4472361809045226,8.4042399172043) (0.4572864321608041,8.419559966027425) (0.4673366834170855,8.43107330404483) (0.47738693467336685,8.448725525283683) (0.48743718592964835,8.491462626800757) (0.49748743718592975,8.51677739266343) (0.507537688442211,8.528100443647984) (0.5175879396984925,8.553924181201765) (0.527638190954774,8.575464014961787) (0.5376884422110553,8.593918480682474) (0.5477386934673367,8.60303590522207) (0.5577889447236182,8.62817705890757) (0.5678391959798995,8.656180772238047) (0.577889447236181,8.659811178112916) (0.5879396984924624,8.674302269420108) (0.5979899497487436,8.720508225689136) (0.6080402010050252,8.720680712732388) (0.6180904522613067,8.716349066048352) (0.6281407035175881,8.750535995255243) (0.6381909547738693,8.770252794507213) (0.6482412060301508,8.82902748677414) (0.6582914572864322,8.86228597383064) (0.6683417085427136,8.893803683785048) (0.678391959798995,8.921496856167519) (0.6884422110552765,8.951337316377407) (0.6984924623115578,9.015266786170812) (0.7085427135678393,9.04831655309136) (0.7185929648241206,9.087920012287547) (0.728643216080402,9.164051028437491) (0.7386934673366835,9.228830421114127) (0.748743718592965,9.260905344839928) (0.7587939698492462,9.283371359137158) (0.7688442211055276,9.296784040230646) (0.7788944723618092,9.3550349304149) (0.7889447236180905,9.391288144076286) (0.7989949748743719,9.438690043119621) (0.8090452261306533,9.477207531050452) (0.8190954773869348,9.51844503179064) (0.8291457286432161,9.57259513750566) (0.8391959798994976,9.581332948984912) (0.849246231155779,9.625329178670729) (0.8592964824120602,9.667092894923455) (0.8693467336683417,9.708151381456469) (0.8793969849246233,9.757387405230018) (0.8894472361809045,9.798534957445295) (0.899497487437186,9.86374193959956) (0.9095477386934674,9.922781547411528) (0.9195979899497487,9.97749205178669) (0.9296482412060302,10.003425504607836) (0.9396984924623116,10.05884427152296) (0.9497487437185931,10.139037801084523) (0.9597989949748744,10.246199280829703) (0.9698492462311559,10.333932193062099) (0.9798994974874372,10.404822507998825) (0.9899497487437185,10.45410383565538) (1.0,10.523519853732688) };
  \addplot[name path=L_msac_grp, draw=none, mark=none, forget plot] coordinates {(-0.9999999999999999,10.146832700017773) (-0.9899497487437187,10.115328344538552) (-0.979899497487437,10.039400923400667) (-0.9698492462311558,10.009350077125523) (-0.9597989949748745,9.949881335874322) (-0.9497487437185927,9.904647139029032) (-0.9396984924623115,9.896217332337642) (-0.9296482412060303,9.780437521851187) (-0.9195979899497485,9.663593800974978) (-0.9095477386934673,9.550869253895312) (-0.899497487437186,9.496887224281634) (-0.8894472361809044,9.423959783790195) (-0.8793969849246231,9.363306820257462) (-0.8693467336683418,9.307569922085209) (-0.8592964824120601,9.233100382278304) (-0.8492462311557789,9.173739365227904) (-0.8391959798994972,9.147736378956342) (-0.829145728643216,9.090788980278205) (-0.8190954773869347,9.0730572626811) (-0.809045226130653,9.021050728023507) (-0.7989949748743718,8.987358700915099) (-0.7889447236180904,8.932086567799018) (-0.7788944723618089,8.835982461948943) (-0.7688442211055275,8.79312483222444) (-0.7587939698492463,8.757046036390191) (-0.7487437185929646,8.728415815807626) (-0.7386934673366833,8.67270281305664) (-0.7286432160804021,8.664254763724854) (-0.7185929648241204,8.624981555713166) (-0.7085427135678392,8.590905251895892) (-0.6984924623115578,8.549969636626894) (-0.6884422110552761,8.471248219634965) (-0.6783919597989949,8.431083945364536) (-0.6683417085427137,8.388828950739567) (-0.658291457286432,8.35793548472498) (-0.6482412060301508,8.321336218206003) (-0.6381909547738694,8.309494968078635) (-0.6281407035175878,8.283033801971092) (-0.6180904522613064,8.263688994292401) (-0.6080402010050252,8.243448521843412) (-0.5979899497487435,8.200734611646649) (-0.5879396984924623,8.176933301301995) (-0.577889447236181,8.141649504629612) (-0.5678391959798994,8.09371953968403) (-0.5577889447236181,8.056769909633386) (-0.5477386934673368,8.020440239583653) (-0.5376884422110552,7.970162047470747) (-0.5276381909547738,7.946471422862414) (-0.5175879396984923,7.923077286384489) (-0.507537688442211,7.900177860021306) (-0.4974874371859297,7.877551963762432) (-0.487437185929648,7.865960770978575) (-0.47738693467336674,7.834896170009919) (-0.4673366834170855,7.821747132976576) (-0.4572864321608038,7.811212008190906) (-0.44723618090452255,7.7811774640559985) (-0.4371859296482412,7.769661605510319) (-0.4271356783919596,7.745476483720696) (-0.4170854271356783,7.723439012900701) (-0.4070351758793971,7.702606420899082) (-0.3969849246231154,7.696577073009058) (-0.38693467336683407,7.679670620612433) (-0.37688442211055284,7.661438439623512) (-0.36683417085427117,7.6369980073227826) (-0.3567839195979899,7.6218611633882665) (-0.3467336683417086,7.614486324079886) (-0.33668341708542704,7.6061116825626245) (-0.3266331658291457,7.607760757494278) (-0.3165829145728644,7.591981304015321) (-0.30653266331658274,7.578682755333382) (-0.29648241206030146,7.555777300786373) (-0.2864321608040203,7.546193555441581) (-0.27638190954773856,7.538701420744413) (-0.2663316582914572,7.531137351316983) (-0.256281407035176,7.522580634484996) (-0.24623115577889435,7.516050580872334) (-0.2361809045226131,7.5101867105116416) (-0.22613065326633142,7.511102638849062) (-0.2160804020100501,7.51605031644535) (-0.2060301507537689,7.5153307520405646) (-0.19597989949748723,7.51928568038903) (-0.18592964824120595,7.532911027367484) (-0.1758793969849247,7.532293764719104) (-0.16582914572864296,7.51089308064009) (-0.15577889447236168,7.51380339147167) (-0.14572864321608048,7.504541402872223) (-0.1356783919597988,7.479773872896752) (-0.12562814070351752,7.473297543543089) (-0.11557788944723625,7.482385728903987) (-0.10552763819095459,7.488900953856968) (-0.09547738693467328,7.491886650531309) (-0.08542713567839201,7.489052946597118) (-0.07537688442211037,7.473866265573475) (-0.06532663316582912,7.482828920946581) (-0.05527638190954787,7.474292086346856) (-0.04522613065326616,7.482372580484169) (-0.0351758793969849,7.499373401949057) (-0.025125628140703626,7.502345059023467) (-0.015075376884421912,7.494788461172471) (-0.005025125628140694,7.4907281672902) (0.005025125628140574,7.497641539469358) (0.015075376884422259,7.517914422674352) (0.02512562814070352,7.502375776491609) (0.03517587939698479,7.500840784183728) (0.04522613065326644,7.501507200501452) (0.0552763819095477,7.508291542517789) (0.06532663316582898,7.522057245405388) (0.07537688442211066,7.528163512257036) (0.08542713567839198,7.541508508109143) (0.09547738693467363,7.5408290848546145) (0.10552763819095488,7.516764721236373) (0.11557788944723618,7.50596976476055) (0.12562814070351783,7.515757405518444) (0.13567839195979908,7.5331349630372) (0.1457286432160804,7.548134167931459) (0.15577889447236204,7.558079562893309) (0.1658291457286433,7.5890614933261435) (0.17587939698492458,7.594269456408793) (0.18592964824120625,7.59330342400027) (0.1959798994974875,7.610461926621593) (0.20603015075376876,7.612748924275329) (0.21608040201005047,7.604218012828259) (0.22613065326633175,7.612388078441368) (0.236180904522613,7.618266153262814) (0.24623115577889465,7.646022055157923) (0.25628140703517593,7.68366578980635) (0.2663316582914572,7.704482990200932) (0.27638190954773884,7.706397973654209) (0.2864321608040201,7.721099286787636) (0.2964824120603014,7.764973392708571) (0.3065326633165831,7.795969381687859) (0.31658291457286436,7.811015094557198) (0.3266331658291457,7.832496403844545) (0.33668341708542726,7.830667150348948) (0.34673366834170855,7.852552663349267) (0.35678391959799,7.864006252933765) (0.3668341708542715,7.872786743650676) (0.3768844221105528,7.8944948883961885) (0.3869346733668343,7.906314109826049) (0.3969849246231157,7.916819536942002) (0.40703517587939697,7.939509199627151) (0.4170854271356784,7.9610394332370555) (0.42713567839195987,7.974511183094598) (0.43718592964824116,7.989092218352705) (0.4472361809045226,8.024950636963442) (0.4572864321608041,8.040300669287202) (0.4673366834170855,8.066240640311705) (0.47738693467336685,8.088085020652864) (0.48743718592964835,8.135159513683568) (0.49748743718592975,8.17037252744417) (0.507537688442211,8.180022823810743) (0.5175879396984925,8.20886821601012) (0.527638190954774,8.224461544454206) (0.5376884422110553,8.238206580600663) (0.5477386934673367,8.252223293718549) (0.5577889447236182,8.274739718358065) (0.5678391959798995,8.305959107088965) (0.577889447236181,8.312834034876767) (0.5879396984924624,8.32651102025074) (0.5979899497487436,8.364755518236644) (0.6080402010050252,8.370694850574132) (0.6180904522613067,8.372244142463853) (0.6281407035175881,8.400164674952146) (0.6381909547738693,8.420852659029636) (0.6482412060301508,8.44900788625736) (0.6582914572864322,8.479738896179553) (0.6683417085427136,8.502126443442037) (0.678391959798995,8.527634085784086) (0.6884422110552765,8.548780863594976) (0.6984924623115578,8.596811440887794) (0.7085427135678393,8.626382489958717) (0.7185929648241206,8.663687361697527) (0.728643216080402,8.73646194444019) (0.7386934673366835,8.79114661430773) (0.748743718592965,8.815384928211357) (0.7587939698492462,8.846311172856334) (0.7688442211055276,8.8650082783858) (0.7788944723618092,8.924905181913703) (0.7889447236180905,8.973287915299094) (0.7989949748743719,9.027137180605475) (0.8090452261306533,9.066135322049186) (0.8190954773869348,9.117471890563754) (0.8291457286432161,9.17695121118503) (0.8391959798994976,9.193057276963973) (0.849246231155779,9.249353490978045) (0.8592964824120602,9.288339913787942) (0.8693467336683417,9.314870541431011) (0.8793969849246233,9.364294388245446) (0.8894472361809045,9.414369077364277) (0.899497487437186,9.489646805289734) (0.9095477386934674,9.531746682995248) (0.9195979899497487,9.57490757344658) (0.9296482412060302,9.5926051991522) (0.9396984924623116,9.627841485296349) (0.9497487437185931,9.668049746739388) (0.9597989949748744,9.757492407130922) (0.9698492462311559,9.854669663680498) (0.9798994974874372,9.93378436908919) (0.9899497487437185,9.966275840011962) (1.0,10.033880829183508) };
  \addplot[ours_msac, fill, fill opacity=0.18, draw=none, mark=none, forget plot] fill between [of=U_msac_grp and L_msac_grp];
  \addplot+[ours_msac, mark=none] coordinates {(-0.9999999999999999,10.41823083599207) (-0.9899497487437187,10.392246333362191) (-0.979899497487437,10.321995391234767) (-0.9698492462311558,10.295034207726808) (-0.9597989949748745,10.236010132999244) (-0.9497487437185927,10.187851754419974) (-0.9396984924623115,10.178100725535613) (-0.9296482412060303,10.047727783960966) (-0.9195979899497485,9.927023199838569) (-0.9095477386934673,9.82084037045227) (-0.899497487437186,9.769422102586697) (-0.8894472361809044,9.684711796049076) (-0.8793969849246231,9.63674517919842) (-0.8693467336683418,9.584493636939243) (-0.8592964824120601,9.502298741976885) (-0.8492462311557789,9.434854425547476) (-0.8391959798994972,9.406310590247955) (-0.829145728643216,9.345822719341545) (-0.8190954773869347,9.3230871814213) (-0.809045226130653,9.27879631157274) (-0.7989949748743718,9.240286762124775) (-0.7889447236180904,9.187839358181158) (-0.7788944723618089,9.090505552285089) (-0.7688442211055275,9.042958817960358) (-0.7587939698492463,9.00124663616128) (-0.7487437185929646,8.96953358099395) (-0.7386934673366833,8.905390403427107) (-0.7286432160804021,8.897353660708067) (-0.7185929648241204,8.864809907419021) (-0.7085427135678392,8.835216691886282) (-0.6984924623115578,8.790975277282596) (-0.6884422110552761,8.709499150471412) (-0.6783919597989949,8.673895934342948) (-0.6683417085427137,8.625692422837275) (-0.658291457286432,8.59010547582778) (-0.6482412060301508,8.538540556485033) (-0.6381909547738694,8.517753871478163) (-0.6281407035175878,8.478728859619572) (-0.6180904522613064,8.445155805423324) (-0.6080402010050252,8.419418742219612) (-0.5979899497487435,8.381512137014145) (-0.5879396984924623,8.358377330340138) (-0.577889447236181,8.328578292554273) (-0.5678391959798994,8.287887007296902) (-0.5577889447236181,8.257945805599075) (-0.5477386934673368,8.219822050090292) (-0.5376884422110552,8.166368964701473) (-0.5276381909547738,8.138589815137816) (-0.5175879396984923,8.104674947337802) (-0.507537688442211,8.084747677967453) (-0.4974874371859297,8.061692662842695) (-0.487437185929648,8.04842192939587) (-0.47738693467336674,8.01184179380981) (-0.4673366834170855,7.993839769881373) (-0.4572864321608038,7.982766448455797) (-0.44723618090452255,7.955889447030228) (-0.4371859296482412,7.943693586140892) (-0.4271356783919596,7.917633020666875) (-0.4170854271356783,7.895532130905227) (-0.4070351758793971,7.872980493067757) (-0.3969849246231154,7.866404559034782) (-0.38693467336683407,7.851210328574265) (-0.37688442211055284,7.832423237654265) (-0.36683417085427117,7.806847801970202) (-0.3567839195979899,7.79474209099216) (-0.3467336683417086,7.788286665554282) (-0.33668341708542704,7.779080947035739) (-0.3266331658291457,7.784540872265291) (-0.3165829145728644,7.772288276105819) (-0.30653266331658274,7.76024198732375) (-0.29648241206030146,7.737324519007324) (-0.2864321608040203,7.725492740876779) (-0.27638190954773856,7.71605498505733) (-0.2663316582914572,7.7075039405221295) (-0.256281407035176,7.698618217535512) (-0.24623115577889435,7.6937542637207645) (-0.2361809045226131,7.686172565037261) (-0.22613065326633142,7.686255858202464) (-0.2160804020100501,7.690921187340687) (-0.2060301507537689,7.6861711000018085) (-0.19597989949748723,7.689466160570236) (-0.18592964824120595,7.702190310048281) (-0.1758793969849247,7.6981301503308535) (-0.16582914572864296,7.678206697147591) (-0.15577889447236168,7.682024542951663) (-0.14572864321608048,7.673432069064583) (-0.1356783919597988,7.6510355662573435) (-0.12562814070351752,7.642324548175672) (-0.11557788944723625,7.649065446158444) (-0.10552763819095459,7.6541239999214294) (-0.09547738693467328,7.6569538423233405) (-0.08542713567839201,7.6553832975934855) (-0.07537688442211037,7.6411038613818825) (-0.06532663316582912,7.645628529875114) (-0.05527638190954787,7.642021085303908) (-0.04522613065326616,7.649547929355394) (-0.0351758793969849,7.670881796131807) (-0.025125628140703626,7.677947273389717) (-0.015075376884421912,7.669964263461201) (-0.005025125628140694,7.667628204835471) (0.005025125628140574,7.672582760480757) (0.015075376884422259,7.6922552261132795) (0.02512562814070352,7.676599948481374) (0.03517587939698479,7.673592028714127) (0.04522613065326644,7.672342509139996) (0.0552763819095477,7.676767385934526) (0.06532663316582898,7.69502585266122) (0.07537688442211066,7.697254836880337) (0.08542713567839198,7.71855363623561) (0.09547738693467363,7.721675168346484) (0.10552763819095488,7.695349478275338) (0.11557788944723618,7.690140123815192) (0.12562814070351783,7.7041885540584385) (0.13567839195979908,7.7164932445930905) (0.1457286432160804,7.7303134695578875) (0.15577889447236204,7.741517433451304) (0.1658291457286433,7.773775157808649) (0.17587939698492458,7.7818434627103015) (0.18592964824120625,7.781993556524495) (0.1959798994974875,7.795010802215851) (0.20603015075376876,7.797153856579532) (0.21608040201005047,7.78506093335771) (0.22613065326633175,7.794092970997217) (0.236180904522613,7.797445215337711) (0.24623115577889465,7.8253711500764185) (0.25628140703517593,7.863017801195636) (0.2663316582914572,7.882363653329724) (0.27638190954773884,7.885335536547494) (0.2864321608040201,7.898485858139422) (0.2964824120603014,7.938966702339991) (0.3065326633165831,7.96835999170136) (0.31658291457286436,7.981761748456826) (0.3266331658291457,7.998591014387565) (0.33668341708542726,7.997062514265757) (0.34673366834170855,8.020089960733928) (0.35678391959799,8.032064525728746) (0.3668341708542715,8.042952946012631) (0.3768844221105528,8.06724562963783) (0.3869346733668343,8.079438579480142) (0.3969849246231157,8.094101228919564) (0.40703517587939697,8.119242978133515) (0.4170854271356784,8.139996527630915) (0.42713567839195987,8.156228405886235) (0.43718592964824116,8.173733399715335) (0.4472361809045226,8.214595277083871) (0.4572864321608041,8.229930317657313) (0.4673366834170855,8.248656972178267) (0.47738693467336685,8.268405272968273) (0.48743718592964835,8.313311070242163) (0.49748743718592975,8.3435749600538) (0.507537688442211,8.354061633729364) (0.5175879396984925,8.381396198605943) (0.527638190954774,8.399962779707996) (0.5376884422110553,8.416062530641568) (0.5477386934673367,8.42762959947031) (0.5577889447236182,8.451458388632817) (0.5678391959798995,8.481069939663506) (0.577889447236181,8.486322606494841) (0.5879396984924624,8.500406644835424) (0.5979899497487436,8.54263187196289) (0.6080402010050252,8.54568778165326) (0.6180904522613067,8.544296604256102) (0.6281407035175881,8.575350335103694) (0.6381909547738693,8.595552726768425) (0.6482412060301508,8.63901768651575) (0.6582914572864322,8.671012435005096) (0.6683417085427136,8.697965063613543) (0.678391959798995,8.724565470975802) (0.6884422110552765,8.750059089986191) (0.6984924623115578,8.806039113529303) (0.7085427135678393,8.837349521525038) (0.7185929648241206,8.875803686992537) (0.728643216080402,8.95025648643884) (0.7386934673366835,9.009988517710928) (0.748743718592965,9.038145136525642) (0.7587939698492462,9.064841265996746) (0.7688442211055276,9.080896159308223) (0.7788944723618092,9.139970056164302) (0.7889447236180905,9.18228802968769) (0.7989949748743719,9.232913611862548) (0.8090452261306533,9.271671426549819) (0.8190954773869348,9.317958461177197) (0.8291457286432161,9.374773174345345) (0.8391959798994976,9.387195112974442) (0.849246231155779,9.437341334824387) (0.8592964824120602,9.477716404355698) (0.8693467336683417,9.51151096144374) (0.8793969849246233,9.560840896737732) (0.8894472361809045,9.606452017404786) (0.899497487437186,9.676694372444647) (0.9095477386934674,9.727264115203388) (0.9195979899497487,9.776199812616635) (0.9296482412060302,9.798015351880018) (0.9396984924623116,9.843342878409654) (0.9497487437185931,9.903543773911956) (0.9597989949748744,10.001845843980313) (0.9698492462311559,10.094300928371299) (0.9798994974874372,10.169303438544008) (0.9899497487437185,10.21018983783367) (1.0,10.278700341458098) };
  \addplot+[ours_msac, mark=none, line width=0.4pt, dashed, forget plot] coordinates {(-1.0,4.4851597707191395) (1.0,4.4851597707191395)};
  \addplot+[ours_msac, mark=*, only marks, mark size=\PanelArgminSize, forget plot] coordinates {(-0.07537688442211037,7.6411038613818825)};
  \addplot[name path=U_gau_grp, draw=none, mark=none, forget plot] coordinates {(-1.0,10.078980431675117) (-0.9899497487437188,10.007611457236225) (-0.979899497487437,9.93805966579205) (-0.9698492462311559,9.888596189339765) (-0.9597989949748745,9.834526930758997) (-0.9497487437185929,9.782471309979234) (-0.9396984924623116,9.718809699900048) (-0.9296482412060304,9.681030753135602) (-0.9195979899497487,9.61316141854479) (-0.9095477386934674,9.56315198951624) (-0.899497487437186,9.530326508257517) (-0.8894472361809044,9.46996865931901) (-0.8793969849246231,9.401685890282936) (-0.8693467336683419,9.339653110456904) (-0.8592964824120602,9.275296202097238) (-0.8492462311557789,9.209027030446597) (-0.8391959798994972,9.166830844260573) (-0.829145728643216,9.11418525245808) (-0.8190954773869348,9.07062819847203) (-0.8090452261306531,9.01341041570132) (-0.7989949748743719,8.970203169313335) (-0.7889447236180905,8.887750206906578) (-0.7788944723618089,8.869699619140539) (-0.7688442211055276,8.842016728479907) (-0.7587939698492464,8.76258796803062) (-0.7487437185929647,8.743580658132634) (-0.7386934673366834,8.698089033327893) (-0.7286432160804021,8.658209069921451) (-0.7185929648241205,8.622986878343163) (-0.7085427135678393,8.58202082262307) (-0.6984924623115579,8.558158574639261) (-0.6884422110552763,8.532500257900377) (-0.678391959798995,8.50349987160429) (-0.6683417085427137,8.484232266617262) (-0.6582914572864321,8.459916529182077) (-0.6482412060301508,8.424141339924148) (-0.6381909547738696,8.362963489451904) (-0.6281407035175879,8.31756763382793) (-0.6180904522613067,8.298130902477592) (-0.6080402010050253,8.25791603048686) (-0.5979899497487436,8.237665515283947) (-0.5879396984924624,8.232469492276959) (-0.5778894472361811,8.2058440454487) (-0.5678391959798994,8.179070926279785) (-0.5577889447236181,8.160061404771854) (-0.5477386934673368,8.128886428231361) (-0.5376884422110553,8.116751672018193) (-0.5276381909547739,8.082974039112965) (-0.5175879396984923,8.055794102349306) (-0.507537688442211,8.043664776762803) (-0.4974874371859298,8.023697191370745) (-0.4874371859296481,7.994760199142989) (-0.47738693467336674,7.979959994424503) (-0.4673366834170855,7.95106872121326) (-0.45728643216080384,7.936649548969336) (-0.4472361809045226,7.9184783085526655) (-0.4371859296482414,7.901821851535795) (-0.4271356783919596,7.887206043868949) (-0.4170854271356784,7.890200898558601) (-0.40703517587939714,7.8815800102357985) (-0.39698492462311546,7.8680328112058735) (-0.3869346733668342,7.862413575814965) (-0.3768844221105529,7.851413915998379) (-0.3668341708542712,7.854045163477051) (-0.35678391959798994,7.839310136779917) (-0.3467336683417087,7.830681825738328) (-0.33668341708542704,7.826935456917458) (-0.3266331658291457,7.8193350734452896) (-0.3165829145728644,7.814077486572488) (-0.30653266331658285,7.814702849061843) (-0.29648241206030157,7.812468451005821) (-0.2864321608040203,7.821689737254761) (-0.2763819095477386,7.81653608814035) (-0.26633165829145733,7.812971276440559) (-0.2562814070351761,7.806982399434277) (-0.24623115577889443,7.805013597377028) (-0.23618090452261317,7.792396240036413) (-0.2261306532663315,7.792595814765097) (-0.2160804020100502,7.786576719937328) (-0.206030150753769,7.76879160699205) (-0.19597989949748731,7.774227319350856) (-0.18592964824120595,7.7914455961275) (-0.1758793969849247,7.793343722908775) (-0.1658291457286431,7.783652802437248) (-0.15577889447236173,7.776123519108319) (-0.14572864321608053,7.767184550898671) (-0.13567839195979886,7.76329951248007) (-0.12562814070351758,7.7604560965180776) (-0.1155778894472363,7.755837112182287) (-0.10552763819095466,7.760906872711132) (-0.09547738693467334,7.767827791946154) (-0.08542713567839214,7.761729805617986) (-0.07537688442211042,7.771135826697228) (-0.06532663316582912,7.7685987588404295) (-0.05527638190954792,7.764790052082205) (-0.045226130653266264,7.758521679363275) (-0.03517587939698496,7.752707298625155) (-0.025125628140703626,7.7561841279859465) (-0.015075376884422013,7.745911582948462) (-0.005025125628140791,7.743768995019149) (0.005025125628140479,7.750246313221355) (0.015075376884422259,7.745509188077244) (0.025125628140703425,7.752708551622851) (0.0351758793969847,7.74745976007992) (0.045226130653266354,7.746252024273002) (0.0552763819095477,7.742801705825564) (0.0653266331658289,7.7610235864010875) (0.07537688442211066,7.74914290582184) (0.08542713567839183,7.772709157463668) (0.09547738693467356,7.780598841954764) (0.10552763819095481,7.802764121828054) (0.11557788944723611,7.811313001288157) (0.12562814070351774,7.8419395087986015) (0.13567839195979903,7.871926137314175) (0.1457286432160803,7.88534673454533) (0.15577889447236196,7.905690623180978) (0.16582914572864324,7.888147892718649) (0.17587939698492452,7.902119477856798) (0.1859296482412062,7.923233896854873) (0.19597989949748745,7.930302190922827) (0.2060301507537687,7.932026191461742) (0.2160804020100504,7.943644757747573) (0.2261306532663317,7.958926910629705) (0.2361809045226129,7.978349621970854) (0.24623115577889454,7.998034097655531) (0.2562814070351758,8.02327706661744) (0.2663316582914571,8.017553761567342) (0.2763819095477388,8.021547524351684) (0.28643216080402006,8.051039623385302) (0.29648241206030135,8.07297742889751) (0.30653266331658297,8.106107743853999) (0.31658291457286436,8.117518417462158) (0.3266331658291455,8.124320072446766) (0.33668341708542726,8.137437032503918) (0.3467336683417085,8.161135462014718) (0.3567839195979898,8.19320991894131) (0.36683417085427145,8.220557725414341) (0.37688442211055273,8.239666998619876) (0.38693467336683435,8.265922466538276) (0.39698492462311563,8.263511848072909) (0.40703517587939686,8.286486969245326) (0.41708542713567864,8.334714132245491) (0.42713567839195987,8.363388278125276) (0.4371859296482411,8.373683212719152) (0.4472361809045228,8.39864131354666) (0.457286432160804,8.422002161701638) (0.4673366834170853,8.43551695719229) (0.47738693467336696,8.454876082334593) (0.48743718592964824,8.474118358666134) (0.4974874371859295,8.495307252887049) (0.5075376884422113,8.52668479908926) (0.5175879396984925,8.538423737563642) (0.5276381909547737,8.55164394285398) (0.5376884422110554,8.565281830276609) (0.5477386934673366,8.583821607258809) (0.557788944723618,8.622048476371742) (0.5678391959798996,8.639770679348791) (0.5778894472361809,8.659203062643982) (0.5879396984924622,8.68889196356039) (0.5979899497487436,8.690913118751437) (0.6080402010050252,8.721177563501886) (0.6180904522613065,8.744824450671643) (0.628140703517588,8.756212925581588) (0.6381909547738693,8.779229233062958) (0.6482412060301508,8.799743217873857) (0.6582914572864322,8.818146609252306) (0.6683417085427136,8.840062765301436) (0.678391959798995,8.851447700897522) (0.6884422110552765,8.867577451651263) (0.6984924623115577,8.889764068160742) (0.7085427135678392,8.938442077946602) (0.7185929648241206,8.952895171090187) (0.728643216080402,8.964564455163089) (0.7386934673366834,8.988231692494317) (0.7487437185929648,8.993771893396625) (0.7587939698492461,9.020232632187518) (0.7688442211055276,9.054808807642846) (0.7788944723618091,9.066070546302704) (0.7889447236180905,9.078638927755442) (0.7989949748743718,9.09999914553139) (0.8090452261306533,9.119552719240788) (0.8190954773869347,9.127808894951075) (0.829145728643216,9.13393648852657) (0.8391959798994975,9.13513303153291) (0.8492462311557789,9.148717251827371) (0.8592964824120601,9.162890748488262) (0.8693467336683417,9.169371667943055) (0.8793969849246231,9.174480200958497) (0.8894472361809044,9.192843346900391) (0.8994974874371858,9.20752909338725) (0.9095477386934674,9.223665969727618) (0.9195979899497487,9.233395594793697) (0.9296482412060301,9.238642249175273) (0.9396984924623116,9.232931201875687) (0.949748743718593,9.237826620314529) (0.9597989949748743,9.242152076079021) (0.9698492462311559,9.2531568455635) (0.9798994974874372,9.258638692620382) (0.9899497487437185,9.263020169888064) (1.0,9.269063637061299) };
  \addplot[name path=L_gau_grp, draw=none, mark=none, forget plot] coordinates {(-1.0,9.553637547459711) (-0.9899497487437188,9.498727371712002) (-0.979899497487437,9.440303843670826) (-0.9698492462311559,9.38774128254801) (-0.9597989949748745,9.327435666545217) (-0.9497487437185929,9.271569168102983) (-0.9396984924623116,9.195789285874449) (-0.9296482412060304,9.154337292152427) (-0.9195979899497487,9.095992903903257) (-0.9095477386934674,9.048809869903728) (-0.899497487437186,9.024586051652754) (-0.8894472361809044,8.955974861886121) (-0.8793969849246231,8.880855758765753) (-0.8693467336683419,8.824425683030846) (-0.8592964824120602,8.766428463240157) (-0.8492462311557789,8.712014767944614) (-0.8391959798994972,8.680994469156413) (-0.829145728643216,8.630515205586063) (-0.8190954773869348,8.583984519708679) (-0.8090452261306531,8.509342644297645) (-0.7989949748743719,8.466916323164023) (-0.7889447236180905,8.40786740985071) (-0.7788944723618089,8.39118790595826) (-0.7688442211055276,8.369783527506192) (-0.7587939698492464,8.324322011218772) (-0.7487437185929647,8.303701596633184) (-0.7386934673366834,8.278837482494463) (-0.7286432160804021,8.26457555569506) (-0.7185929648241205,8.23687238575307) (-0.7085427135678393,8.1974689058019) (-0.6984924623115579,8.172392978622828) (-0.6884422110552763,8.141586439210457) (-0.678391959798995,8.098991451861746) (-0.6683417085427137,8.07127737733888) (-0.6582914572864321,8.040771257098513) (-0.6482412060301508,8.009357311932021) (-0.6381909547738696,7.95446805309257) (-0.6281407035175879,7.923858140742565) (-0.6180904522613067,7.912421862786406) (-0.6080402010050253,7.89276957980762) (-0.5979899497487436,7.8745161192259365) (-0.5879396984924624,7.871035649653117) (-0.5778894472361811,7.84591178069525) (-0.5678391959798994,7.826057557300555) (-0.5577889447236181,7.8120164491039406) (-0.5477386934673368,7.780365262367994) (-0.5376884422110553,7.765236692161393) (-0.5276381909547739,7.736424474166613) (-0.5175879396984923,7.7178377526169575) (-0.507537688442211,7.707840140374192) (-0.4974874371859298,7.693937608433755) (-0.4874371859296481,7.672868764824522) (-0.47738693467336674,7.6624460915837105) (-0.4673366834170855,7.634822263857134) (-0.45728643216080384,7.619026455376761) (-0.4472361809045226,7.6007523192580395) (-0.4371859296482414,7.584676095125662) (-0.4271356783919596,7.566613264403581) (-0.4170854271356784,7.5675090878146465) (-0.40703517587939714,7.552240462080467) (-0.39698492462311546,7.539525623232022) (-0.3869346733668342,7.52877336078036) (-0.3768844221105529,7.519748696138893) (-0.3668341708542712,7.523108068241489) (-0.35678391959798994,7.513161489720695) (-0.3467336683417087,7.5088350732873295) (-0.33668341708542704,7.495958027990334) (-0.3266331658291457,7.494083044701679) (-0.3165829145728644,7.4843552276036185) (-0.30653266331658285,7.480771604575227) (-0.29648241206030157,7.471977211594695) (-0.2864321608040203,7.475305698540774) (-0.2763819095477386,7.471937567915256) (-0.26633165829145733,7.4661080081626325) (-0.2562814070351761,7.455744182138307) (-0.24623115577889443,7.459960851077697) (-0.23618090452261317,7.441559067565831) (-0.2261306532663315,7.4485167673941) (-0.2160804020100502,7.446674137646127) (-0.206030150753769,7.437097511324172) (-0.19597989949748731,7.44607528781881) (-0.18592964824120595,7.45381063933737) (-0.1758793969849247,7.452766659000524) (-0.1658291457286431,7.454843355995751) (-0.15577889447236173,7.448342172908747) (-0.14572864321608053,7.443009915717598) (-0.13567839195979886,7.434593776705678) (-0.12562814070351758,7.432890651043597) (-0.1155778894472363,7.421310813270531) (-0.10552763819095466,7.424824630818801) (-0.09547738693467334,7.435072702074249) (-0.08542713567839214,7.43288528785433) (-0.07537688442211042,7.45281836629217) (-0.06532663316582912,7.448155470152626) (-0.05527638190954792,7.445994310587109) (-0.045226130653266264,7.438569231368676) (-0.03517587939698496,7.437149858088133) (-0.025125628140703626,7.438392234360425) (-0.015075376884422013,7.4272865914279596) (-0.005025125628140791,7.418444427639127) (0.005025125628140479,7.4215258184640565) (0.015075376884422259,7.413895517765581) (0.025125628140703425,7.419620028073634) (0.0351758793969847,7.418894658711995) (0.045226130653266354,7.417589235568145) (0.0552763819095477,7.416654996006442) (0.0653266331658289,7.436975307940319) (0.07537688442211066,7.434214746933376) (0.08542713567839183,7.455187715435391) (0.09547738693467356,7.458223729682438) (0.10552763819095481,7.471368866856872) (0.11557788944723611,7.475310320246873) (0.12562814070351774,7.498535894983156) (0.13567839195979903,7.523806260295918) (0.1457286432160803,7.535764009506636) (0.15577889447236196,7.551128203775078) (0.16582914572864324,7.539004996660335) (0.17587939698492452,7.556767384919623) (0.1859296482412062,7.57517948648063) (0.19597989949748745,7.577760125961733) (0.2060301507537687,7.581143906206214) (0.2160804020100504,7.589145028993051) (0.2261306532663317,7.603504189309778) (0.2361809045226129,7.6306431973140905) (0.24623115577889454,7.655001026729978) (0.2562814070351758,7.675771004408529) (0.2663316582914571,7.673651711155615) (0.2763819095477388,7.678038745658319) (0.28643216080402006,7.704411005863857) (0.29648241206030135,7.724251463373086) (0.30653266331658297,7.769088458103624) (0.31658291457286436,7.782446867244213) (0.3266331658291455,7.79365391589999) (0.33668341708542726,7.800794402345743) (0.3467336683417085,7.8170358399258335) (0.3567839195979898,7.840527795606409) (0.36683417085427145,7.8742450601305585) (0.37688442211055273,7.892825861355528) (0.38693467336683435,7.909164533108755) (0.39698492462311563,7.912724976826804) (0.40703517587939686,7.926028399990213) (0.41708542713567864,7.967201167995415) (0.42713567839195987,7.991839659787287) (0.4371859296482411,8.001671906871293) (0.4472361809045228,8.030734505199772) (0.457286432160804,8.043237628856236) (0.4673366834170853,8.054263378449475) (0.47738693467336696,8.07058482458369) (0.48743718592964824,8.091256743032504) (0.4974874371859295,8.108904237089092) (0.5075376884422113,8.14709402684155) (0.5175879396984925,8.163863232229835) (0.5276381909547737,8.17589618814972) (0.5376884422110554,8.196290957491334) (0.5477386934673366,8.220128144455948) (0.557788944723618,8.258888926771977) (0.5678391959798996,8.27321282858905) (0.5778894472361809,8.286236739246885) (0.5879396984924622,8.312334719270094) (0.5979899497487436,8.319480150117558) (0.6080402010050252,8.352245408426498) (0.6180904522613065,8.372461636351428) (0.628140703517588,8.383968702497198) (0.6381909547738693,8.402579774608933) (0.6482412060301508,8.420275130690271) (0.6582914572864322,8.4315411397966) (0.6683417085427136,8.449123773212737) (0.678391959798995,8.461322313929648) (0.6884422110552765,8.467828410701019) (0.6984924623115577,8.489470830804) (0.7085427135678392,8.522404270989673) (0.7185929648241206,8.532499411329155) (0.728643216080402,8.545748969799309) (0.7386934673366834,8.570088562766074) (0.7487437185929648,8.572078255787018) (0.7587939698492461,8.595777206110114) (0.7688442211055276,8.620475919517471) (0.7788944723618091,8.633469886770285) (0.7889447236180905,8.642204816710127) (0.7989949748743718,8.65664126369292) (0.8090452261306533,8.672129023459489) (0.8190954773869347,8.680773717340633) (0.829145728643216,8.689293452900928) (0.8391959798994975,8.695640224160893) (0.8492462311557789,8.705211809026597) (0.8592964824120601,8.714290304720706) (0.8693467336683417,8.719608566059124) (0.8793969849246231,8.722850363353269) (0.8894472361809044,8.737731942689374) (0.8994974874371858,8.750440798975669) (0.9095477386934674,8.769900256249779) (0.9195979899497487,8.77786447188876) (0.9296482412060301,8.78232458900146) (0.9396984924623116,8.771175352669239) (0.949748743718593,8.775806372255158) (0.9597989949748743,8.780114724760443) (0.9698492462311559,8.79115347044844) (0.9798994974874372,8.80259987280355) (0.9899497487437185,8.808448870638662) (1.0,8.816430079305212) };
  \addplot[ours_gau, fill, fill opacity=0.18, draw=none, mark=none, forget plot] fill between [of=U_gau_grp and L_gau_grp];
  \addplot+[ours_gau, mark=none] coordinates {(-1.0,9.816308989567414) (-0.9899497487437188,9.753169414474113) (-0.979899497487437,9.689181754731438) (-0.9698492462311559,9.638168735943887) (-0.9597989949748745,9.580981298652107) (-0.9497487437185929,9.527020239041109) (-0.9396984924623116,9.457299492887248) (-0.9296482412060304,9.417684022644014) (-0.9195979899497487,9.354577161224023) (-0.9095477386934674,9.305980929709984) (-0.899497487437186,9.277456279955135) (-0.8894472361809044,9.212971760602565) (-0.8793969849246231,9.141270824524344) (-0.8693467336683419,9.082039396743875) (-0.8592964824120602,9.020862332668697) (-0.8492462311557789,8.960520899195606) (-0.8391959798994972,8.923912656708493) (-0.829145728643216,8.872350229022071) (-0.8190954773869348,8.827306359090354) (-0.8090452261306531,8.761376529999483) (-0.7989949748743719,8.718559746238679) (-0.7889447236180905,8.647808808378644) (-0.7788944723618089,8.6304437625494) (-0.7688442211055276,8.60590012799305) (-0.7587939698492464,8.543454989624696) (-0.7487437185929647,8.52364112738291) (-0.7386934673366834,8.488463257911178) (-0.7286432160804021,8.461392312808256) (-0.7185929648241205,8.429929632048117) (-0.7085427135678393,8.389744864212485) (-0.6984924623115579,8.365275776631044) (-0.6884422110552763,8.337043348555417) (-0.678391959798995,8.301245661733018) (-0.6683417085427137,8.277754821978071) (-0.6582914572864321,8.250343893140295) (-0.6482412060301508,8.216749325928085) (-0.6381909547738696,8.158715771272236) (-0.6281407035175879,8.120712887285247) (-0.6180904522613067,8.105276382631999) (-0.6080402010050253,8.07534280514724) (-0.5979899497487436,8.056090817254942) (-0.5879396984924624,8.051752570965037) (-0.5778894472361811,8.025877913071975) (-0.5678391959798994,8.00256424179017) (-0.5577889447236181,7.986038926937897) (-0.5477386934673368,7.954625845299677) (-0.5376884422110553,7.940994182089793) (-0.5276381909547739,7.909699256639789) (-0.5175879396984923,7.886815927483132) (-0.507537688442211,7.875752458568497) (-0.4974874371859298,7.85881739990225) (-0.4874371859296481,7.833814481983755) (-0.47738693467336674,7.821203043004107) (-0.4673366834170855,7.792945492535197) (-0.45728643216080384,7.777838002173048) (-0.4472361809045226,7.7596153139053525) (-0.4371859296482414,7.743248973330728) (-0.4271356783919596,7.726909654136265) (-0.4170854271356784,7.728854993186624) (-0.40703517587939714,7.7169102361581325) (-0.39698492462311546,7.703779217218948) (-0.3869346733668342,7.6955934682976626) (-0.3768844221105529,7.685581306068636) (-0.3668341708542712,7.68857661585927) (-0.35678391959798994,7.676235813250306) (-0.3467336683417087,7.669758449512829) (-0.33668341708542704,7.661446742453896) (-0.3266331658291457,7.656709059073484) (-0.3165829145728644,7.649216357088053) (-0.30653266331658285,7.647737226818535) (-0.29648241206030157,7.642222831300258) (-0.2864321608040203,7.648497717897768) (-0.2763819095477386,7.644236828027803) (-0.26633165829145733,7.639539642301596) (-0.2562814070351761,7.631363290786292) (-0.24623115577889443,7.632487224227362) (-0.23618090452261317,7.616977653801122) (-0.2261306532663315,7.620556291079598) (-0.2160804020100502,7.616625428791727) (-0.206030150753769,7.602944559158111) (-0.19597989949748731,7.610151303584833) (-0.18592964824120595,7.622628117732435) (-0.1758793969849247,7.623055190954649) (-0.1658291457286431,7.6192480792165) (-0.15577889447236173,7.612232846008533) (-0.14572864321608053,7.605097233308134) (-0.13567839195979886,7.598946644592874) (-0.12562814070351758,7.596673373780837) (-0.1155778894472363,7.588573962726409) (-0.10552763819095466,7.5928657517649665) (-0.09547738693467334,7.6014502470102014) (-0.08542713567839214,7.597307546736158) (-0.07537688442211042,7.611977096494699) (-0.06532663316582912,7.608377114496528) (-0.05527638190954792,7.605392181334657) (-0.045226130653266264,7.598545455365976) (-0.03517587939698496,7.594928578356644) (-0.025125628140703626,7.597288181173186) (-0.015075376884422013,7.586599087188211) (-0.005025125628140791,7.581106711329138) (0.005025125628140479,7.585886065842706) (0.015075376884422259,7.579702352921412) (0.025125628140703425,7.586164289848242) (0.0351758793969847,7.5831772093959575) (0.045226130653266354,7.581920629920574) (0.0552763819095477,7.579728350916003) (0.0653266331658289,7.598999447170703) (0.07537688442211066,7.591678826377608) (0.08542713567839183,7.61394843644953) (0.09547738693467356,7.619411285818601) (0.10552763819095481,7.637066494342463) (0.11557788944723611,7.643311660767515) (0.12562814070351774,7.670237701890879) (0.13567839195979903,7.6978661988050465) (0.1457286432160803,7.710555372025983) (0.15577889447236196,7.728409413478028) (0.16582914572864324,7.713576444689492) (0.17587939698492452,7.72944343138821) (0.1859296482412062,7.749206691667752) (0.19597989949748745,7.75403115844228) (0.2060301507537687,7.756585048833978) (0.2160804020100504,7.766394893370312) (0.2261306532663317,7.781215549969741) (0.2361809045226129,7.804496409642472) (0.24623115577889454,7.8265175621927545) (0.2562814070351758,7.849524035512984) (0.2663316582914571,7.845602736361478) (0.2763819095477388,7.849793135005001) (0.28643216080402006,7.87772531462458) (0.29648241206030135,7.898614446135298) (0.30653266331658297,7.9375981009788115) (0.31658291457286436,7.949982642353185) (0.3266331658291455,7.958986994173378) (0.33668341708542726,7.96911571742483) (0.3467336683417085,7.989085650970275) (0.3567839195979898,8.01686885727386) (0.36683417085427145,8.04740139277245) (0.37688442211055273,8.066246429987702) (0.38693467336683435,8.087543499823516) (0.39698492462311563,8.088118412449857) (0.40703517587939686,8.10625768461777) (0.41708542713567864,8.150957650120453) (0.42713567839195987,8.177613968956281) (0.4371859296482411,8.187677559795222) (0.4472361809045228,8.214687909373216) (0.457286432160804,8.232619895278937) (0.4673366834170853,8.244890167820882) (0.47738693467336696,8.262730453459142) (0.48743718592964824,8.28268755084932) (0.4974874371859295,8.30210574498807) (0.5075376884422113,8.336889412965405) (0.5175879396984925,8.351143484896738) (0.5276381909547737,8.36377006550185) (0.5376884422110554,8.380786393883971) (0.5477386934673366,8.401974875857379) (0.557788944723618,8.44046870157186) (0.5678391959798996,8.45649175396892) (0.5778894472361809,8.472719900945433) (0.5879396984924622,8.500613341415242) (0.5979899497487436,8.505196634434498) (0.6080402010050252,8.536711485964192) (0.6180904522613065,8.558643043511536) (0.628140703517588,8.570090814039393) (0.6381909547738693,8.590904503835946) (0.6482412060301508,8.610009174282064) (0.6582914572864322,8.624843874524453) (0.6683417085427136,8.644593269257086) (0.678391959798995,8.656385007413585) (0.6884422110552765,8.66770293117614) (0.6984924623115577,8.689617449482371) (0.7085427135678392,8.730423174468138) (0.7185929648241206,8.742697291209671) (0.728643216080402,8.755156712481199) (0.7386934673366834,8.779160127630195) (0.7487437185929648,8.782925074591821) (0.7587939698492461,8.808004919148816) (0.7688442211055276,8.837642363580159) (0.7788944723618091,8.849770216536495) (0.7889447236180905,8.860421872232784) (0.7989949748743718,8.878320204612155) (0.8090452261306533,8.895840871350138) (0.8190954773869347,8.904291306145854) (0.829145728643216,8.91161497071375) (0.8391959798994975,8.915386627846901) (0.8492462311557789,8.926964530426984) (0.8592964824120601,8.938590526604484) (0.8693467336683417,8.94449011700109) (0.8793969849246231,8.948665282155883) (0.8894472361809044,8.965287644794882) (0.8994974874371858,8.978984946181459) (0.9095477386934674,8.996783112988698) (0.9195979899497487,9.005630033341228) (0.9296482412060301,9.010483419088366) (0.9396984924623116,9.002053277272463) (0.949748743718593,9.006816496284843) (0.9597989949748743,9.011133400419732) (0.9698492462311559,9.02215515800597) (0.9798994974874372,9.030619282711966) (0.9899497487437185,9.035734520263363) (1.0,9.042746858183255) };
  \addplot+[ours_gau, mark=none, line width=0.4pt, dashed, forget plot] coordinates {(-1.0,4.847998881134512) (1.0,4.847998881134512)};
  \addplot+[ours_gau, mark=*, only marks, mark size=\PanelArgminSize, forget plot] coordinates {(0.015075376884422259,7.579702352921412)};
  \addplot[name path=U_magsac_grp, draw=none, mark=none, forget plot] coordinates {(-0.9999999999999999,10.466339556286881) (-0.9899497487437185,10.387652677720759) (-0.979899497487437,10.342549303802372) (-0.9698492462311556,10.26520638901296) (-0.9597989949748744,10.172927900943382) (-0.9497487437185927,10.103723507822657) (-0.9396984924623115,10.06335196669818) (-0.9296482412060302,10.015476694244539) (-0.9195979899497485,9.946720868566915) (-0.9095477386934673,9.908478464826636) (-0.8994974874371859,9.856105005815285) (-0.8894472361809044,9.819837260059048) (-0.8793969849246231,9.760587999068232) (-0.8693467336683417,9.697295596829363) (-0.8592964824120601,9.62685690962283) (-0.8492462311557789,9.583521469160583) (-0.8391959798994976,9.543216944204623) (-0.8291457286432159,9.496327171854539) (-0.8190954773869347,9.465185152694707) (-0.8090452261306534,9.403038051509574) (-0.7989949748743717,9.352857503391501) (-0.7889447236180904,9.285576473735395) (-0.7788944723618089,9.260474180500347) (-0.7688442211055275,9.196660553738933) (-0.7587939698492462,9.15502094427114) (-0.7487437185929646,9.118510797843493) (-0.7386934673366833,9.077486829334198) (-0.7286432160804021,9.04275566461486) (-0.7185929648241204,9.00115845696368) (-0.708542713567839,8.936680776516425) (-0.6984924623115578,8.909626024603812) (-0.6884422110552761,8.890685462113659) (-0.6783919597989949,8.870629086482698) (-0.6683417085427136,8.814805583596415) (-0.658291457286432,8.7667882717664) (-0.6482412060301507,8.730169117910437) (-0.6381909547738694,8.659397684689008) (-0.6281407035175878,8.59908363916912) (-0.6180904522613064,8.578695211868547) (-0.6080402010050252,8.551339541293848) (-0.5979899497487435,8.516997338543888) (-0.5879396984924623,8.480645745132247) (-0.577889447236181,8.470132460320446) (-0.5678391959798994,8.419012170907878) (-0.5577889447236181,8.368867975654716) (-0.5477386934673368,8.321909786771608) (-0.5376884422110552,8.27859405428858) (-0.5276381909547738,8.232745688542236) (-0.5175879396984925,8.19293752780484) (-0.5075376884422109,8.192316807410403) (-0.4974874371859296,8.175635969419947) (-0.4874371859296484,8.16519403075955) (-0.4773869346733667,8.1360394603355) (-0.46733668341708545,8.117163158117348) (-0.4572864321608041,8.0929280709455) (-0.44723618090452255,8.065048305279202) (-0.43718592964824116,8.046266521580622) (-0.42713567839196,8.03389349102057) (-0.41708542713567825,8.013781605476254) (-0.40703517587939697,7.987719286286753) (-0.39698492462311535,7.970948657494829) (-0.38693467336683407,7.960220882347204) (-0.3768844221105528,7.938758502837127) (-0.3668341708542711,7.9078535381727) (-0.35678391959798983,7.892544425431338) (-0.34673366834170855,7.89011990966368) (-0.33668341708542693,7.885380527085755) (-0.32663316582914564,7.87435164791213) (-0.3165829145728644,7.8599421072291955) (-0.30653266331658274,7.853057026471042) (-0.29648241206030146,7.854236984692652) (-0.2864321608040202,7.845973625295958) (-0.27638190954773856,7.83229574150237) (-0.2663316582914572,7.831877523966488) (-0.256281407035176,7.838695746504018) (-0.24623115577889435,7.832078717465566) (-0.236180904522613,7.825497379215081) (-0.2261306532663318,7.813472153394696) (-0.2160804020100501,7.807153032237581) (-0.20603015075376882,7.812780494834461) (-0.19597989949748762,7.813921568943804) (-0.1859296482412059,7.802292374504792) (-0.1758793969849246,7.795627799561214) (-0.1658291457286433,7.796695075707297) (-0.15577889447236168,7.794854152888525) (-0.14572864321608042,7.782770708365612) (-0.13567839195979914,7.792585410184425) (-0.12562814070351747,7.784901401528287) (-0.11557788944723618,7.7862031821348525) (-0.10552763819095491,7.787753079254339) (-0.09547738693467328,7.7805917495656765) (-0.08542713567839201,7.779305932359707) (-0.07537688442211071,7.7650551269722445) (-0.06532663316582901,7.773452194818042) (-0.055276381909547756,7.757403377034147) (-0.045226130653266534,7.750543237975252) (-0.03517587939698485,7.752441566819316) (-0.025125628140703578,7.747790746674941) (-0.015075376884422311,7.750190510451361) (-0.0050251256281406455,7.75747755493767) (0.00502512562814067,7.764149040502589) (0.015075376884422259,7.759046498606787) (0.02512562814070352,7.763158415090749) (0.03517587939698479,7.763365558487656) (0.04522613065326653,7.7622288326847215) (0.05527638190954779,7.766507612572186) (0.06532663316582907,7.758994717682568) (0.07537688442211074,7.747578949947799) (0.08542713567839198,7.74890519498158) (0.09547738693467324,7.750394891368323) (0.10552763819095473,7.762113819646958) (0.11557788944723618,7.762142280497613) (0.1256281407035177,7.747894930273722) (0.1356783919597989,7.756362461746695) (0.1457286432160804,7.767380823896236) (0.1557788944723619,7.7853584656688835) (0.1658291457286431,7.803856839170275) (0.17587939698492464,7.806872270739431) (0.18592964824120609,7.832549314376015) (0.19597989949748737,7.837722335164691) (0.20603015075376882,7.869774199228008) (0.21608040201005027,7.889472632673431) (0.22613065326633158,7.895945321734015) (0.236180904522613,7.901874107822857) (0.24623115577889443,7.9114524715360135) (0.25628140703517577,7.902887199856969) (0.2663316582914572,7.9149885470193535) (0.2763819095477387,7.927522407985711) (0.28643216080401995,7.925679892409445) (0.29648241206030146,7.944169361947853) (0.30653266331658285,7.960573719922039) (0.31658291457286436,7.988473022170995) (0.3266331658291457,8.00753620643416) (0.3366834170854271,8.019022717657501) (0.34673366834170855,8.010499747408524) (0.35678391959798983,8.036993975173518) (0.36683417085427134,8.053343855755084) (0.3768844221105528,8.07655124177976) (0.38693467336683407,8.1092302042536) (0.3969849246231155,8.11867536709593) (0.407035175879397,8.121670714042141) (0.4170854271356783,8.128945529468684) (0.42713567839195976,8.165957674651077) (0.4371859296482412,8.191370816281749) (0.4472361809045225,8.214404190355099) (0.457286432160804,8.237823299888102) (0.4673366834170854,8.248059996805088) (0.4773869346733666,8.2622050018083) (0.4874371859296482,8.301550023971586) (0.49748743718592964,8.326743154207414) (0.5075376884422111,8.350435399053273) (0.5175879396984924,8.355974088973902) (0.5276381909547738,8.383351279330943) (0.5376884422110553,8.39695781012686) (0.5477386934673365,8.433264604006938) (0.5577889447236181,8.454141672637036) (0.5678391959798995,8.470530775645805) (0.5778894472361809,8.497689757589091) (0.5879396984924623,8.52663372963364) (0.5979899497487438,8.549366084756844) (0.608040201005025,8.554095252692255) (0.6180904522613064,8.568660756212008) (0.6281407035175879,8.58025032993652) (0.6381909547738692,8.591145271151762) (0.6482412060301507,8.616382283689223) (0.6582914572864321,8.642788597820743) (0.6683417085427134,8.653657129086737) (0.6783919597989949,8.663025445500152) (0.6884422110552764,8.709063441661783) (0.6984924623115576,8.769985892396196) (0.708542713567839,8.789514975854416) (0.7185929648241206,8.82177847051535) (0.728643216080402,8.845423512818346) (0.7386934673366833,8.890765593608322) (0.7487437185929647,8.927506892161167) (0.7587939698492462,8.97068561552294) (0.7688442211055275,9.009261038457035) (0.778894472361809,9.058154346630417) (0.7889447236180904,9.108775030464582) (0.7989949748743717,9.146717324570872) (0.8090452261306531,9.200767994237598) (0.8190954773869347,9.259000431167655) (0.8291457286432159,9.312071441507188) (0.8391959798994973,9.34874184038058) (0.8492462311557788,9.381467996166004) (0.8592964824120601,9.425757002038841) (0.8693467336683416,9.448243464724325) (0.879396984924623,9.490280271270077) (0.8894472361809043,9.519987157342545) (0.8994974874371858,9.548717169977781) (0.9095477386934673,9.598635586476034) (0.9195979899497487,9.633829480497393) (0.9296482412060301,9.651278959592574) (0.9396984924623114,9.714578316389447) (0.949748743718593,9.776298748271314) (0.9597989949748743,9.844632631183291) (0.9698492462311558,9.91741127895221) (0.979899497487437,9.968240229203994) (0.9899497487437185,10.038725411832976) (0.9999999999999999,10.087504744312115) };
  \addplot[name path=L_magsac_grp, draw=none, mark=none, forget plot] coordinates {(-0.9999999999999999,9.871773663587517) (-0.9899497487437185,9.806180163758565) (-0.979899497487437,9.762762155593137) (-0.9698492462311556,9.708287531523755) (-0.9597989949748744,9.622604134706474) (-0.9497487437185927,9.54614216124848) (-0.9396984924623115,9.50633142326492) (-0.9296482412060302,9.452743615746188) (-0.9195979899497485,9.40392629937883) (-0.9095477386934673,9.3756179343087) (-0.8994974874371859,9.332333015421137) (-0.8894472361809044,9.300084544673156) (-0.8793969849246231,9.231535519828117) (-0.8693467336683417,9.166954888513622) (-0.8592964824120601,9.102887316473794) (-0.8492462311557789,9.060080018465683) (-0.8391959798994976,9.029743298434386) (-0.8291457286432159,8.979114159243998) (-0.8190954773869347,8.934715414777576) (-0.8090452261306534,8.870006674713936) (-0.7989949748743717,8.812513044267327) (-0.7889447236180904,8.760663686066904) (-0.7788944723618089,8.737286101251792) (-0.7688442211055275,8.688142986454888) (-0.7587939698492462,8.65226418284064) (-0.7487437185929646,8.616126746700962) (-0.7386934673366833,8.570693897620007) (-0.7286432160804021,8.530913840137627) (-0.7185929648241204,8.47822898016003) (-0.708542713567839,8.423943542286128) (-0.6984924623115578,8.40001026782023) (-0.6884422110552761,8.383143785370816) (-0.6783919597989949,8.363909062373281) (-0.6683417085427136,8.334010975505604) (-0.658291457286432,8.308744161168311) (-0.6482412060301507,8.287968575314972) (-0.6381909547738694,8.255138600836599) (-0.6281407035175878,8.215045681957186) (-0.6180904522613064,8.193238925169167) (-0.6080402010050252,8.158345448968488) (-0.5979899497487435,8.11067029260349) (-0.5879396984924623,8.054856318580647) (-0.577889447236181,8.040464977416956) (-0.5678391959798994,7.984002166474142) (-0.5577889447236181,7.94087488036359) (-0.5477386934673368,7.908819123647546) (-0.5376884422110552,7.880706028873566) (-0.5276381909547738,7.852818925774988) (-0.5175879396984925,7.826521255601965) (-0.5075376884422109,7.826200532608648) (-0.4974874371859296,7.80740736745352) (-0.4874371859296484,7.797230113074587) (-0.4773869346733667,7.77962587833827) (-0.46733668341708545,7.767725172180555) (-0.4572864321608041,7.742228488959667) (-0.44723618090452255,7.722434383864587) (-0.43718592964824116,7.709576196056426) (-0.42713567839196,7.702288888416748) (-0.41708542713567825,7.689179185158681) (-0.40703517587939697,7.6721464994475905) (-0.39698492462311535,7.651654091701647) (-0.38693467336683407,7.646947889635441) (-0.3768844221105528,7.6242000199520295) (-0.3668341708542711,7.594558626002458) (-0.35678391959798983,7.568483319996026) (-0.34673366834170855,7.570261804551894) (-0.33668341708542693,7.563528705549044) (-0.32663316582914564,7.550832004313767) (-0.3165829145728644,7.528789854419876) (-0.30653266331658274,7.520632389644545) (-0.29648241206030146,7.519949979069298) (-0.2864321608040202,7.517586362906133) (-0.27638190954773856,7.502100587984622) (-0.2663316582914572,7.501985833298345) (-0.256281407035176,7.509250819420463) (-0.24623115577889435,7.496581962000399) (-0.236180904522613,7.491166565759826) (-0.2261306532663318,7.478399091219367) (-0.2160804020100501,7.473775424839296) (-0.20603015075376882,7.475378800751977) (-0.19597989949748762,7.472984401739466) (-0.1859296482412059,7.461047251358253) (-0.1758793969849246,7.448004250845194) (-0.1658291457286433,7.447227588881431) (-0.15577889447236168,7.446110112444557) (-0.14572864321608042,7.434228259369385) (-0.13567839195979914,7.439635038950256) (-0.12562814070351747,7.440886966323009) (-0.11557788944723618,7.445735736194027) (-0.10552763819095491,7.441896379729419) (-0.09547738693467328,7.438409043301085) (-0.08542713567839201,7.43678374380872) (-0.07537688442211071,7.429363603413114) (-0.06532663316582901,7.4365934699388125) (-0.055276381909547756,7.421630414925487) (-0.045226130653266534,7.412584774487579) (-0.03517587939698485,7.420643752934383) (-0.025125628140703578,7.415800673793363) (-0.015075376884422311,7.418755925824479) (-0.0050251256281406455,7.430046383444361) (0.00502512562814067,7.43847985441388) (0.015075376884422259,7.430829963919339) (0.02512562814070352,7.4406983607797) (0.03517587939698479,7.438090280952295) (0.04522613065326653,7.438258583199623) (0.05527638190954779,7.442347844129335) (0.06532663316582907,7.435997401924912) (0.07537688442211074,7.416516098504054) (0.08542713567839198,7.4208989018350096) (0.09547738693467324,7.419317492421971) (0.10552763819095473,7.427820698689681) (0.11557788944723618,7.428651074729561) (0.1256281407035177,7.414486961884281) (0.1356783919597989,7.42711192381284) (0.1457286432160804,7.445181190242552) (0.1557788944723619,7.460959849802219) (0.1658291457286431,7.477553736443869) (0.17587939698492464,7.471276460035604) (0.18592964824120609,7.4869863447710046) (0.19597989949748737,7.4892659016378245) (0.20603015075376882,7.514478118533116) (0.21608040201005027,7.5309934820961635) (0.22613065326633158,7.534223931198287) (0.236180904522613,7.541192255903346) (0.24623115577889443,7.552863930269614) (0.25628140703517577,7.553563612455185) (0.2663316582914572,7.56675946407603) (0.2763819095477387,7.579016119196734) (0.28643216080401995,7.57807360166375) (0.29648241206030146,7.594745440238049) (0.30653266331658285,7.6105861133666055) (0.31658291457286436,7.644942252301448) (0.3266331658291457,7.662385317284471) (0.3366834170854271,7.671896880130627) (0.34673366834170855,7.664372062602141) (0.35678391959798983,7.69720818024366) (0.36683417085427134,7.708852741948953) (0.3768844221105528,7.734091776931448) (0.38693467336683407,7.769512203185567) (0.3969849246231155,7.782095319840764) (0.407035175879397,7.793593203592006) (0.4170854271356783,7.794677561126762) (0.42713567839195976,7.827775563563844) (0.4371859296482412,7.845088854966891) (0.4472361809045225,7.870155016457674) (0.457286432160804,7.892118354200865) (0.4673366834170854,7.901585730356258) (0.4773869346733666,7.913693512572567) (0.4874371859296482,7.945964281076695) (0.49748743718592964,7.969423782587678) (0.5075376884422111,7.991540053176562) (0.5175879396984924,7.995897886969488) (0.5276381909547738,8.011386684975227) (0.5376884422110553,8.014243852220826) (0.5477386934673365,8.05490706408114) (0.5577889447236181,8.066799867276224) (0.5678391959798995,8.088228938419842) (0.5778894472361809,8.109282827993786) (0.5879396984924623,8.149658630589192) (0.5979899497487438,8.166337145176872) (0.608040201005025,8.172035310735302) (0.6180904522613064,8.197462002840783) (0.6281407035175879,8.205359150878374) (0.6381909547738692,8.214908640523939) (0.6482412060301507,8.234124438664004) (0.6582914572864321,8.271769255291346) (0.6683417085427134,8.290320627797763) (0.6783919597989949,8.304434711070394) (0.6884422110552764,8.338441614571448) (0.6984924623115576,8.387439058918275) (0.708542713567839,8.410559307744162) (0.7185929648241206,8.440364649902348) (0.728643216080402,8.458278638503533) (0.7386934673366833,8.487569220755741) (0.7487437185929647,8.5225560143751) (0.7587939698492462,8.568922430262534) (0.7688442211055275,8.593992076626256) (0.778894472361809,8.626634830898176) (0.7889447236180904,8.664466356043413) (0.7989949748743717,8.692585638115172) (0.8090452261306531,8.738087730468656) (0.8190954773869347,8.79478913890563) (0.8291457286432159,8.851588311068431) (0.8391959798994973,8.895962736882854) (0.8492462311557788,8.930131305070235) (0.8592964824120601,8.998627885737013) (0.8693467336683416,9.02043083076186) (0.879396984924623,9.068259514739289) (0.8894472361809043,9.111554646387493) (0.8994974874371858,9.148163620967221) (0.9095477386934673,9.201295766492917) (0.9195979899497487,9.244432546136895) (0.9296482412060301,9.268229434759991) (0.9396984924623114,9.32076230085468) (0.949748743718593,9.380377375108138) (0.9597989949748743,9.441605984374224) (0.9698492462311558,9.494719559881842) (0.979899497487437,9.534989351504263) (0.9899497487437185,9.57703740662379) (0.9999999999999999,9.605523993060311) };
  \addplot[ours_magsac, fill, fill opacity=0.18, draw=none, mark=none, forget plot] fill between [of=U_magsac_grp and L_magsac_grp];
  \addplot+[ours_magsac, mark=none] coordinates {(-0.9999999999999999,10.169056609937199) (-0.9899497487437185,10.096916420739662) (-0.979899497487437,10.052655729697754) (-0.9698492462311556,9.986746960268357) (-0.9597989949748744,9.897766017824928) (-0.9497487437185927,9.824932834535568) (-0.9396984924623115,9.78484169498155) (-0.9296482412060302,9.734110154995363) (-0.9195979899497485,9.675323583972872) (-0.9095477386934673,9.642048199567668) (-0.8994974874371859,9.594219010618211) (-0.8894472361809044,9.559960902366102) (-0.8793969849246231,9.496061759448175) (-0.8693467336683417,9.432125242671493) (-0.8592964824120601,9.364872113048312) (-0.8492462311557789,9.321800743813133) (-0.8391959798994976,9.286480121319505) (-0.8291457286432159,9.237720665549269) (-0.8190954773869347,9.199950283736142) (-0.8090452261306534,9.136522363111755) (-0.7989949748743717,9.082685273829414) (-0.7889447236180904,9.02312007990115) (-0.7788944723618089,8.99888014087607) (-0.7688442211055275,8.942401770096911) (-0.7587939698492462,8.90364256355589) (-0.7487437185929646,8.867318772272228) (-0.7386934673366833,8.824090363477103) (-0.7286432160804021,8.786834752376244) (-0.7185929648241204,8.739693718561854) (-0.708542713567839,8.680312159401277) (-0.6984924623115578,8.65481814621202) (-0.6884422110552761,8.636914623742237) (-0.6783919597989949,8.61726907442799) (-0.6683417085427136,8.57440827955101) (-0.658291457286432,8.537766216467356) (-0.6482412060301507,8.509068846612704) (-0.6381909547738694,8.457268142762803) (-0.6281407035175878,8.407064660563153) (-0.6180904522613064,8.385967068518857) (-0.6080402010050252,8.354842495131168) (-0.5979899497487435,8.31383381557369) (-0.5879396984924623,8.267751031856447) (-0.577889447236181,8.255298718868701) (-0.5678391959798994,8.20150716869101) (-0.5577889447236181,8.154871428009153) (-0.5477386934673368,8.115364455209576) (-0.5376884422110552,8.079650041581074) (-0.5276381909547738,8.042782307158612) (-0.5175879396984925,8.009729391703402) (-0.5075376884422109,8.009258670009526) (-0.4974874371859296,7.991521668436733) (-0.4874371859296484,7.981212071917068) (-0.4773869346733667,7.957832669336884) (-0.46733668341708545,7.9424441651489515) (-0.4572864321608041,7.917578279952584) (-0.44723618090452255,7.893741344571895) (-0.43718592964824116,7.877921358818524) (-0.42713567839196,7.86809118971866) (-0.41708542713567825,7.851480395317467) (-0.40703517587939697,7.829932892867172) (-0.39698492462311535,7.811301374598238) (-0.38693467336683407,7.803584385991322) (-0.3768844221105528,7.781479261394578) (-0.3668341708542711,7.751206082087579) (-0.35678391959798983,7.730513872713682) (-0.34673366834170855,7.730190857107787) (-0.33668341708542693,7.7244546163173995) (-0.32663316582914564,7.7125918261129485) (-0.3165829145728644,7.694365980824536) (-0.30653266331658274,7.6868447080577935) (-0.29648241206030146,7.687093481880975) (-0.2864321608040202,7.681779994101046) (-0.27638190954773856,7.667198164743496) (-0.2663316582914572,7.666931678632417) (-0.256281407035176,7.6739732829622405) (-0.24623115577889435,7.664330339732983) (-0.236180904522613,7.658331972487454) (-0.2261306532663318,7.645935622307031) (-0.2160804020100501,7.6404642285384385) (-0.20603015075376882,7.644079647793219) (-0.19597989949748762,7.643452985341635) (-0.1859296482412059,7.631669812931523) (-0.1758793969849246,7.621816025203204) (-0.1658291457286433,7.621961332294364) (-0.15577889447236168,7.620482132666541) (-0.14572864321608042,7.608499483867498) (-0.13567839195979914,7.616110224567341) (-0.12562814070351747,7.612894183925648) (-0.11557788944723618,7.61596945916444) (-0.10552763819095491,7.614824729491879) (-0.09547738693467328,7.609500396433381) (-0.08542713567839201,7.608044838084213) (-0.07537688442211071,7.597209365192679) (-0.06532663316582901,7.6050228323784275) (-0.055276381909547756,7.589516895979817) (-0.045226130653266534,7.581564006231416) (-0.03517587939698485,7.58654265987685) (-0.025125628140703578,7.581795710234152) (-0.015075376884422311,7.58447321813792) (-0.0050251256281406455,7.593761969191015) (0.00502512562814067,7.601314447458234) (0.015075376884422259,7.594938231263063) (0.02512562814070352,7.601928387935224) (0.03517587939698479,7.600727919719976) (0.04522613065326653,7.600243707942172) (0.05527638190954779,7.60442772835076) (0.06532663316582907,7.59749605980374) (0.07537688442211074,7.582047524225927) (0.08542713567839198,7.584902048408295) (0.09547738693467324,7.584856191895147) (0.10552763819095473,7.594967259168319) (0.11557788944723618,7.595396677613587) (0.1256281407035177,7.581190946079001) (0.1356783919597989,7.591737192779767) (0.1457286432160804,7.606281007069394) (0.1557788944723619,7.623159157735551) (0.1658291457286431,7.640705287807072) (0.17587939698492464,7.6390743653875175) (0.18592964824120609,7.65976782957351) (0.19597989949748737,7.663494118401258) (0.20603015075376882,7.692126158880562) (0.21608040201005027,7.710233057384797) (0.22613065326633158,7.715084626466151) (0.236180904522613,7.7215331818631014) (0.24623115577889443,7.732158200902814) (0.25628140703517577,7.728225406156077) (0.2663316582914572,7.740874005547692) (0.2763819095477387,7.753269263591222) (0.28643216080401995,7.751876747036597) (0.29648241206030146,7.769457401092951) (0.30653266331658285,7.7855799166443225) (0.31658291457286436,7.816707637236221) (0.3266331658291457,7.834960761859316) (0.3366834170854271,7.845459798894065) (0.34673366834170855,7.837435905005332) (0.35678391959798983,7.867101077708589) (0.36683417085427134,7.881098298852019) (0.3768844221105528,7.905321509355604) (0.38693467336683407,7.939371203719584) (0.3969849246231155,7.950385343468347) (0.407035175879397,7.957631958817073) (0.4170854271356783,7.961811545297723) (0.42713567839195976,7.99686661910746) (0.4371859296482412,8.01822983562432) (0.4472361809045225,8.042279603406387) (0.457286432160804,8.064970827044483) (0.4673366834170854,8.074822863580673) (0.4773869346733666,8.087949257190434) (0.4874371859296482,8.12375715252414) (0.49748743718592964,8.148083468397546) (0.5075376884422111,8.170987726114918) (0.5175879396984924,8.175935987971695) (0.5276381909547738,8.197368982153085) (0.5376884422110553,8.205600831173843) (0.5477386934673365,8.24408583404404) (0.5577889447236181,8.26047076995663) (0.5678391959798995,8.279379857032824) (0.5778894472361809,8.303486292791439) (0.5879396984924623,8.338146180111416) (0.5979899497487438,8.357851614966858) (0.608040201005025,8.363065281713778) (0.6180904522613064,8.383061379526396) (0.6281407035175879,8.392804740407447) (0.6381909547738692,8.40302695583785) (0.6482412060301507,8.425253361176614) (0.6582914572864321,8.457278926556045) (0.6683417085427134,8.47198887844225) (0.6783919597989949,8.483730078285273) (0.6884422110552764,8.523752528116615) (0.6984924623115576,8.578712475657236) (0.708542713567839,8.60003714179929) (0.7185929648241206,8.631071560208849) (0.728643216080402,8.65185107566094) (0.7386934673366833,8.689167407182032) (0.7487437185929647,8.725031453268134) (0.7587939698492462,8.769804022892737) (0.7688442211055275,8.801626557541645) (0.778894472361809,8.842394588764297) (0.7889447236180904,8.886620693253997) (0.7989949748743717,8.919651481343022) (0.8090452261306531,8.969427862353127) (0.8190954773869347,9.026894785036642) (0.8291457286432159,9.08182987628781) (0.8391959798994973,9.122352288631717) (0.8492462311557788,9.15579965061812) (0.8592964824120601,9.212192443887927) (0.8693467336683416,9.234337147743092) (0.879396984924623,9.279269893004683) (0.8894472361809043,9.315770901865019) (0.8994974874371858,9.348440395472501) (0.9095477386934673,9.399965676484475) (0.9195979899497487,9.439131013317144) (0.9296482412060301,9.459754197176283) (0.9396984924623114,9.517670308622064) (0.949748743718593,9.578338061689726) (0.9597989949748743,9.643119307778758) (0.9698492462311558,9.706065419417026) (0.979899497487437,9.751614790354129) (0.9899497487437185,9.807881409228383) (0.9999999999999999,9.846514368686213) };
  \addplot+[ours_magsac, mark=none, line width=0.4pt, dashed, forget plot] coordinates {(-1.0,4.702301336741773) (1.0,4.702301336741773)};
  \addplot+[ours_magsac, mark=*, only marks, mark size=\PanelArgminSize, forget plot] coordinates {(0.1256281407035177,7.581190946079001)};
  \addplot[name path=U_jeffreys_grp, draw=none, mark=none, forget plot] coordinates {(-1.0,10.24306170376416) (-0.9899497487437185,10.188858960140953) (-0.9798994974874373,10.161016058530276) (-0.9698492462311559,10.114852868140384) (-0.9597989949748744,10.066153952357523) (-0.9497487437185931,10.016401918329835) (-0.9396984924623116,9.980998475676884) (-0.9296482412060302,9.952292085544133) (-0.9195979899497487,9.927275418205532) (-0.9095477386934674,9.874553261291638) (-0.8994974874371859,9.831578948793094) (-0.8894472361809045,9.76493616279572) (-0.8793969849246233,9.719858505157127) (-0.8693467336683417,9.685969729876886) (-0.8592964824120604,9.63807345495411) (-0.849246231155779,9.57756266451604) (-0.8391959798994976,9.537830538328745) (-0.829145728643216,9.486880240596927) (-0.8190954773869348,9.450008325414245) (-0.8090452261306533,9.408136212917775) (-0.7989949748743719,9.384256523768062) (-0.7889447236180905,9.350681798179004) (-0.7788944723618091,9.310477846413624) (-0.7688442211055276,9.298522977987705) (-0.7587939698492462,9.284488724506447) (-0.748743718592965,9.250346409568907) (-0.7386934673366835,9.257897566382335) (-0.728643216080402,9.196445827249491) (-0.7185929648241207,9.171827227268269) (-0.7085427135678393,9.131498658538106) (-0.6984924623115577,9.10364129814162) (-0.6884422110552765,9.08432351976159) (-0.678391959798995,9.064954864884017) (-0.6683417085427136,9.056384809825722) (-0.6582914572864323,9.031170442565736) (-0.6482412060301508,9.035642685089076) (-0.6381909547738693,9.016656374031333) (-0.6281407035175881,8.987076433579164) (-0.6180904522613067,8.964573520648285) (-0.6080402010050251,8.930435495173107) (-0.5979899497487436,8.914663986999377) (-0.5879396984924624,8.889353404786487) (-0.5778894472361809,8.85243515131157) (-0.5678391959798995,8.818072787154733) (-0.5577889447236182,8.8040863499936) (-0.5477386934673368,8.778276480337336) (-0.5376884422110553,8.735538546058795) (-0.527638190954774,8.72319459695099) (-0.5175879396984925,8.690320184143404) (-0.507537688442211,8.672521813506682) (-0.4974874371859298,8.647668901679522) (-0.4874371859296483,8.619047658507322) (-0.4773869346733669,8.588133003709848) (-0.4673366834170855,8.56696143110824) (-0.4572864321608041,8.55393681276765) (-0.4472361809045226,8.544283180965218) (-0.43718592964824116,8.531622260546822) (-0.42713567839195987,8.517630801191606) (-0.4170854271356785,8.507088725809345) (-0.40703517587939697,8.486124958697726) (-0.39698492462311574,8.469728467008851) (-0.38693467336683424,8.4534000776378) (-0.37688442211055273,8.440817601599996) (-0.36683417085427145,8.43726082377696) (-0.35678391959799,8.41945157904425) (-0.3467336683417085,8.405593652616467) (-0.33668341708542726,8.370991554757166) (-0.3266331658291458,8.34420413895295) (-0.31658291457286436,8.327382267085605) (-0.30653266331658285,8.296562851740921) (-0.29648241206030157,8.279852999973812) (-0.2864321608040201,8.253427219616574) (-0.27638190954773884,8.243049967323133) (-0.26633165829145733,8.235411670209018) (-0.25628140703517593,8.215386220098713) (-0.2462311557788945,8.20983074475275) (-0.2361809045226131,8.18951871867787) (-0.22613065326633172,8.192012830379241) (-0.21608040201005035,8.17510093187668) (-0.20603015075376882,8.172926562892759) (-0.19597989949748754,8.159953150613584) (-0.18592964824120603,8.147471908371907) (-0.1758793969849247,8.134163013671767) (-0.1658291457286433,8.128351858245868) (-0.15577889447236182,8.127427038275698) (-0.14572864321608048,8.124311178517647) (-0.13567839195979906,8.126870888176049) (-0.12562814070351766,8.118306080126121) (-0.11557788944723625,8.103367958738405) (-0.10552763819095479,8.108473099960792) (-0.0954773869346734,8.107006958165787) (-0.08542713567839201,8.090578623605488) (-0.07537688442211053,8.07959143278635) (-0.06532663316582923,8.069161043466263) (-0.05527638190954781,8.082518759960703) (-0.04522613065326632,8.083572701195546) (-0.03517587939698496,8.090088346082661) (-0.025125628140703526,8.099212962945742) (-0.015075376884422112,8.105918808560418) (-0.005025125628140743,8.108675781776878) (0.00502512562814067,8.104488373775078) (0.015075376884422073,8.107676138829675) (0.025125628140703425,8.113728126350837) (0.035175879396984966,8.107823694893812) (0.045226130653266264,8.108094797762641) (0.0552763819095477,8.10770967697658) (0.06532663316582915,8.12244708481436) (0.0753768844221105,8.144272216370467) (0.08542713567839191,8.161044138728517) (0.09547738693467332,8.161878757712373) (0.10552763819095466,8.179524797925282) (0.11557788944723618,8.193754068546044) (0.12562814070351747,8.201330094623769) (0.13567839195979894,8.218587282441105) (0.1457286432160804,8.213925712481146) (0.15577889447236176,8.218925890701698) (0.16582914572864316,8.22821686937558) (0.17587939698492458,8.258576238658705) (0.185929648241206,8.273744093491889) (0.19597989949748737,8.29397249601712) (0.20603015075376876,8.326290675707241) (0.21608040201005016,8.349597356756908) (0.22613065326633164,8.346328299188515) (0.236180904522613,8.360725312412162) (0.24623115577889443,8.403089635440976) (0.25628140703517577,8.425449371120822) (0.2663316582914573,8.442171789685375) (0.2763819095477386,8.47197008108488) (0.28643216080402,8.497064116759661) (0.29648241206030146,8.517399438054127) (0.30653266331658285,8.535071140018337) (0.31658291457286425,8.54902463504879) (0.3266331658291457,8.574615690648974) (0.3366834170854271,8.59933392246451) (0.3467336683417085,8.656068683349599) (0.35678391959798994,8.704282841700078) (0.3668341708542713,8.740522282576853) (0.37688442211055273,8.785694569754774) (0.3869346733668341,8.837514919310383) (0.3969849246231155,8.889200911378454) (0.40703517587939697,8.959780639148908) (0.4170854271356784,9.054949209357567) (0.42713567839195976,9.106422180340516) (0.43718592964824116,9.209534200340444) (0.44723618090452255,9.30033417127166) (0.457286432160804,9.357318615164127) (0.4673366834170854,9.446270422878232) (0.47738693467336685,9.53953975829843) (0.48743718592964824,9.657286778019984) (0.4974874371859296,9.755007677186107) (0.507537688442211,9.887646478761043) (0.5175879396984924,9.994632679832721) (0.5276381909547738,10.15869141097277) (0.5376884422110553,10.346229153523534) (0.5477386934673366,10.522014771091605) (0.557788944723618,10.722286869253047) (0.5678391959798994,10.876287183031357) (0.5778894472361809,11.08886213356022) (0.5879396984924622,11.28932712552214) (0.5979899497487436,11.527481739471984) (0.6080402010050251,11.769129275802593) (0.6180904522613064,12.046229734093666) (0.6281407035175879,12.272384621231804) (0.6381909547738693,12.539250454057974) (0.6482412060301507,12.854427959521836) (0.6582914572864321,13.115602048632901) (0.6683417085427136,13.41722805084328) (0.6783919597989949,13.686887357240536) (0.6884422110552764,14.071411222333147) (0.6984924623115578,14.422838362969777) (0.7085427135678392,14.730565938473093) (0.7185929648241206,15.177174253111819) (0.728643216080402,15.648355088612094) (0.7386934673366834,16.039251583737272) (0.7487437185929647,16.39934291851234) (0.7587939698492462,16.78801296817899) (0.7688442211055275,17.144866275676073) (0.778894472361809,17.487587186243637) (0.7889447236180904,17.87468907140792) (0.7989949748743718,18.21974638945397) (0.8090452261306533,18.57291319123124) (0.8190954773869346,18.974447713061377) (0.829145728643216,19.426836218709578) (0.8391959798994975,19.73162805461645) (0.8492462311557789,20.014788390279005) (0.8592964824120602,20.25015419829674) (0.8693467336683416,20.424939071334283) (0.8793969849246231,20.835769917457114) (0.8894472361809045,21.207803615578033) (0.8994974874371859,21.424021257932516) (0.9095477386934673,21.55084993357098) (0.9195979899497487,21.866258220682898) (0.9296482412060301,22.29929456166736) (0.9396984924623115,22.447501482712486) (0.949748743718593,22.751442086338834) (0.9597989949748743,23.03494882469603) (0.9698492462311558,23.42976925552145) (0.9798994974874371,23.706284407170624) (0.9899497487437185,23.891461087808704) (0.9999999999999999,24.043915104365635) };
  \addplot[name path=L_jeffreys_grp, draw=none, mark=none, forget plot] coordinates {(-1.0,9.620412333715992) (-0.9899497487437185,9.558861218266179) (-0.9798994974874373,9.529215397679515) (-0.9698492462311559,9.484132635977998) (-0.9597989949748744,9.44028167947902) (-0.9497487437185931,9.398339541033065) (-0.9396984924623116,9.358311092818429) (-0.9296482412060302,9.319635198640924) (-0.9195979899497487,9.291899465932929) (-0.9095477386934674,9.232962149509788) (-0.8994974874371859,9.187425864976195) (-0.8894472361809045,9.133476891243795) (-0.8793969849246233,9.093188239120288) (-0.8693467336683417,9.0745278417043) (-0.8592964824120604,9.02964304209949) (-0.849246231155779,8.980326573896754) (-0.8391959798994976,8.948302764581172) (-0.829145728643216,8.89794388443285) (-0.8190954773869348,8.860104895528028) (-0.8090452261306533,8.813036473535316) (-0.7989949748743719,8.800510754192377) (-0.7889447236180905,8.776451218423876) (-0.7788944723618091,8.744487325995616) (-0.7688442211055276,8.725930391837906) (-0.7587939698492462,8.71810449250153) (-0.748743718592965,8.68754009970627) (-0.7386934673366835,8.682663636060619) (-0.728643216080402,8.623160376115752) (-0.7185929648241207,8.5990565849791) (-0.7085427135678393,8.561991035868047) (-0.6984924623115577,8.534943325813307) (-0.6884422110552765,8.518039697929751) (-0.678391959798995,8.513266398334588) (-0.6683417085427136,8.503071220685694) (-0.6582914572864323,8.478206006040905) (-0.6482412060301508,8.478512752458766) (-0.6381909547738693,8.46321063012409) (-0.6281407035175881,8.450262589466929) (-0.6180904522613067,8.43606957249258) (-0.6080402010050251,8.416323391913796) (-0.5979899497487436,8.399107431790801) (-0.5879396984924624,8.375691067963324) (-0.5778894472361809,8.351227184271705) (-0.5678391959798995,8.325095860153496) (-0.5577889447236182,8.317198409137706) (-0.5477386934673368,8.29926865476678) (-0.5376884422110553,8.269858927465311) (-0.527638190954774,8.257506316108913) (-0.5175879396984925,8.234638208623648) (-0.507537688442211,8.223470146994439) (-0.4974874371859298,8.206999999622552) (-0.4874371859296483,8.18404206211399) (-0.4773869346733669,8.163651370904477) (-0.4673366834170855,8.148575054512143) (-0.4572864321608041,8.142656995642964) (-0.4472361809045226,8.136570283080546) (-0.43718592964824116,8.124679190897254) (-0.42713567839195987,8.118844491642724) (-0.4170854271356785,8.100720797712247) (-0.40703517587939697,8.076333688351127) (-0.39698492462311574,8.056554904036568) (-0.38693467336683424,8.031475220146001) (-0.37688442211055273,8.017782725377817) (-0.36683417085427145,8.01989639987774) (-0.35678391959799,8.005889393541572) (-0.3467336683417085,7.9899528658312695) (-0.33668341708542726,7.954906002206141) (-0.3266331658291458,7.92688301079961) (-0.31658291457286436,7.912817636585446) (-0.30653266331658285,7.888355329142909) (-0.29648241206030157,7.876834106981185) (-0.2864321608040201,7.861072446130434) (-0.27638190954773884,7.863021475987671) (-0.26633165829145733,7.852238598965683) (-0.25628140703517593,7.831202869154241) (-0.2462311557788945,7.826770312531242) (-0.2361809045226131,7.814433823746533) (-0.22613065326633172,7.8124022366899295) (-0.21608040201005035,7.797343054086729) (-0.20603015075376882,7.79622182929237) (-0.19597989949748754,7.783344605168018) (-0.18592964824120603,7.770440144938668) (-0.1758793969849247,7.760635400294896) (-0.1658291457286433,7.7588104670657785) (-0.15577889447236182,7.757335585563104) (-0.14572864321608048,7.749458138226164) (-0.13567839195979906,7.754353816919998) (-0.12562814070351766,7.745586438824313) (-0.11557788944723625,7.7373811817829505) (-0.10552763819095479,7.743164322547868) (-0.0954773869346734,7.744098101125698) (-0.08542713567839201,7.730624118222039) (-0.07537688442211053,7.717758345741104) (-0.06532663316582923,7.706167290855624) (-0.05527638190954781,7.717141228371924) (-0.04522613065326632,7.711621562169738) (-0.03517587939698496,7.716206117728142) (-0.025125628140703526,7.723036214955806) (-0.015075376884422112,7.72970737927661) (-0.005025125628140743,7.739935162016844) (0.00502512562814067,7.739373181199498) (0.015075376884422073,7.741600819131607) (0.025125628140703425,7.744935242779098) (0.035175879396984966,7.7448953746805955) (0.045226130653266264,7.745767744597352) (0.0552763819095477,7.744946724063006) (0.06532663316582915,7.756259350875885) (0.0753768844221105,7.782927437991758) (0.08542713567839191,7.796470793040619) (0.09547738693467332,7.800746818222524) (0.10552763819095466,7.81121954003093) (0.11557788944723618,7.823654838946942) (0.12562814070351747,7.832175718984469) (0.13567839195979894,7.845814616882729) (0.1457286432160804,7.841870143831412) (0.15577889447236176,7.846653453965491) (0.16582914572864316,7.854867763545425) (0.17587939698492458,7.881328311665576) (0.185929648241206,7.892835480706765) (0.19597989949748737,7.9131879409909995) (0.20603015075376876,7.9361656828966565) (0.21608040201005016,7.949790714607891) (0.22613065326633164,7.951757501747395) (0.236180904522613,7.96454748224996) (0.24623115577889443,8.002181503667495) (0.25628140703517577,8.020596092130583) (0.2663316582914573,8.039784011193415) (0.2763819095477386,8.072184746685124) (0.28643216080402,8.10208723638402) (0.29648241206030146,8.129494555485469) (0.30653266331658285,8.146748290123293) (0.31658291457286425,8.167348091897997) (0.3266331658291457,8.198626317817043) (0.3366834170854271,8.222518981667665) (0.3467336683417085,8.264301355776286) (0.35678391959798994,8.299260568130212) (0.3668341708542713,8.330966194064315) (0.37688442211055273,8.361613137352748) (0.3869346733668341,8.399319169220096) (0.3969849246231155,8.43624850533479) (0.40703517587939697,8.492610239318088) (0.4170854271356784,8.581218996508127) (0.42713567839195976,8.636258793043442) (0.43718592964824116,8.731739873878501) (0.44723618090452255,8.835044076436985) (0.457286432160804,8.904235750830155) (0.4673366834170854,8.980792730757372) (0.47738693467336685,9.047743168712058) (0.48743718592964824,9.155805843164515) (0.4974874371859296,9.247845945464277) (0.507537688442211,9.346406031106183) (0.5175879396984924,9.433803067054336) (0.5276381909547738,9.561263095596948) (0.5376884422110553,9.720163262215308) (0.5477386934673366,9.887198452229246) (0.557788944723618,10.132435901697278) (0.5678391959798994,10.287548847864036) (0.5778894472361809,10.482080326911642) (0.5879396984924622,10.631627038143618) (0.5979899497487436,10.818724483641713) (0.6080402010050251,11.039443945276853) (0.6180904522613064,11.317802205428572) (0.6281407035175879,11.564652484373187) (0.6381909547738693,11.773712679779152) (0.6482412060301507,12.056752796829048) (0.6582914572864321,12.320400563682348) (0.6683417085427136,12.589616937089248) (0.6783919597989949,12.77594626387609) (0.6884422110552764,13.074962845077433) (0.6984924623115578,13.396417169178218) (0.7085427135678392,13.7223208901784) (0.7185929648241206,14.073718070060597) (0.728643216080402,14.435135783942082) (0.7386934673366834,14.872583731351336) (0.7487437185929647,15.248413283991598) (0.7587939698492462,15.52120768643316) (0.7688442211055275,15.866295549834293) (0.778894472361809,16.21735088474457) (0.7889447236180904,16.569404086603996) (0.7989949748743718,16.89667780156515) (0.8090452261306533,17.265616774720495) (0.8190954773869346,17.629733539581) (0.829145728643216,18.00425243274557) (0.8391959798994975,18.30771969685226) (0.8492462311557789,18.573593669931025) (0.8592964824120602,18.830130268031823) (0.8693467336683416,19.05709451801828) (0.8793969849246231,19.444526635902243) (0.8894472361809045,19.771181423807388) (0.8994974874371859,20.028055786351857) (0.9095477386934673,20.159448943464035) (0.9195979899497487,20.486961264835887) (0.9296482412060301,20.852897051877193) (0.9396984924623115,21.007622288947726) (0.949748743718593,21.335916263746462) (0.9597989949748743,21.63653874257679) (0.9698492462311558,21.940231958051555) (0.9798994974874371,22.19523809702393) (0.9899497487437185,22.353945576557987) (0.9999999999999999,22.529260790213808) };
  \addplot[ours_jeffreys, fill, fill opacity=0.18, draw=none, mark=none, forget plot] fill between [of=U_jeffreys_grp and L_jeffreys_grp];
  \addplot+[ours_jeffreys, mark=none] coordinates {(-1.0,9.931737018740076) (-0.9899497487437185,9.873860089203566) (-0.9798994974874373,9.845115728104895) (-0.9698492462311559,9.79949275205919) (-0.9597989949748744,9.753217815918271) (-0.9497487437185931,9.70737072968145) (-0.9396984924623116,9.669654784247657) (-0.9296482412060302,9.635963642092529) (-0.9195979899497487,9.60958744206923) (-0.9095477386934674,9.553757705400713) (-0.8994974874371859,9.509502406884645) (-0.8894472361809045,9.449206527019758) (-0.8793969849246233,9.406523372138707) (-0.8693467336683417,9.380248785790593) (-0.8592964824120604,9.3338582485268) (-0.849246231155779,9.278944619206397) (-0.8391959798994976,9.243066651454958) (-0.829145728643216,9.192412062514888) (-0.8190954773869348,9.155056610471137) (-0.8090452261306533,9.110586343226545) (-0.7989949748743719,9.09238363898022) (-0.7889447236180905,9.06356650830144) (-0.7788944723618091,9.02748258620462) (-0.7688442211055276,9.012226684912806) (-0.7587939698492462,9.001296608503988) (-0.748743718592965,8.968943254637589) (-0.7386934673366835,8.970280601221477) (-0.728643216080402,8.909803101682622) (-0.7185929648241207,8.885441906123685) (-0.7085427135678393,8.846744847203077) (-0.6984924623115577,8.819292311977463) (-0.6884422110552765,8.80118160884567) (-0.678391959798995,8.789110631609303) (-0.6683417085427136,8.779728015255708) (-0.6582914572864323,8.75468822430332) (-0.6482412060301508,8.757077718773921) (-0.6381909547738693,8.739933502077712) (-0.6281407035175881,8.718669511523046) (-0.6180904522613067,8.700321546570432) (-0.6080402010050251,8.673379443543451) (-0.5979899497487436,8.65688570939509) (-0.5879396984924624,8.632522236374905) (-0.5778894472361809,8.601831167791637) (-0.5678391959798995,8.571584323654115) (-0.5577889447236182,8.560642379565653) (-0.5477386934673368,8.538772567552058) (-0.5376884422110553,8.502698736762053) (-0.527638190954774,8.490350456529951) (-0.5175879396984925,8.462479196383526) (-0.507537688442211,8.44799598025056) (-0.4974874371859298,8.427334450651037) (-0.4874371859296483,8.401544860310656) (-0.4773869346733669,8.375892187307162) (-0.4673366834170855,8.357768242810192) (-0.4572864321608041,8.348296904205307) (-0.4472361809045226,8.340426732022882) (-0.43718592964824116,8.328150725722038) (-0.42713567839195987,8.318237646417165) (-0.4170854271356785,8.303904761760796) (-0.40703517587939697,8.281229323524427) (-0.39698492462311574,8.26314168552271) (-0.38693467336683424,8.242437648891901) (-0.37688442211055273,8.229300163488906) (-0.36683417085427145,8.22857861182735) (-0.35678391959799,8.212670486292911) (-0.3467336683417085,8.197773259223869) (-0.33668341708542726,8.162948778481654) (-0.3266331658291458,8.13554357487628) (-0.31658291457286436,8.120099951835526) (-0.30653266331658285,8.092459090441915) (-0.29648241206030157,8.078343553477499) (-0.2864321608040201,8.057249832873504) (-0.27638190954773884,8.053035721655402) (-0.26633165829145733,8.043825134587351) (-0.25628140703517593,8.023294544626477) (-0.2462311557788945,8.018300528641996) (-0.2361809045226131,8.001976271212202) (-0.22613065326633172,8.002207533534586) (-0.21608040201005035,7.986221992981704) (-0.20603015075376882,7.984574196092564) (-0.19597989949748754,7.971648877890801) (-0.18592964824120603,7.9589560266552875) (-0.1758793969849247,7.947399206983331) (-0.1658291457286433,7.943581162655823) (-0.15577889447236182,7.9423813119194016) (-0.14572864321608048,7.9368846583719055) (-0.13567839195979906,7.9406123525480234) (-0.12562814070351766,7.931946259475217) (-0.11557788944723625,7.920374570260678) (-0.10552763819095479,7.92581871125433) (-0.0954773869346734,7.925552529645743) (-0.08542713567839201,7.9106013709137635) (-0.07537688442211053,7.898674889263726) (-0.06532663316582923,7.887664167160943) (-0.05527638190954781,7.899829994166313) (-0.04522613065326632,7.897597131682643) (-0.03517587939698496,7.903147231905401) (-0.025125628140703526,7.911124588950774) (-0.015075376884422112,7.917813093918514) (-0.005025125628140743,7.924305471896861) (0.00502512562814067,7.921930777487288) (0.015075376884422073,7.924638478980641) (0.025125628140703425,7.929331684564968) (0.035175879396984966,7.926359534787204) (0.045226130653266264,7.926931271179996) (0.0552763819095477,7.926328200519793) (0.06532663316582915,7.939353217845122) (0.0753768844221105,7.963599827181113) (0.08542713567839191,7.978757465884568) (0.09547738693467332,7.981312787967448) (0.10552763819095466,7.995372168978106) (0.11557788944723618,8.008704453746493) (0.12562814070351747,8.016752906804118) (0.13567839195979894,8.032200949661917) (0.1457286432160804,8.027897928156278) (0.15577889447236176,8.032789672333594) (0.16582914572864316,8.041542316460502) (0.17587939698492458,8.06995227516214) (0.185929648241206,8.083289787099327) (0.19597989949748737,8.10358021850406) (0.20603015075376876,8.13122817930195) (0.21608040201005016,8.1496940356824) (0.22613065326633164,8.149042900467956) (0.236180904522613,8.162636397331061) (0.24623115577889443,8.202635569554236) (0.25628140703517577,8.223022731625703) (0.2663316582914573,8.240977900439395) (0.2763819095477386,8.272077413885002) (0.28643216080402,8.299575676571841) (0.29648241206030146,8.323446996769798) (0.30653266331658285,8.340909715070815) (0.31658291457286425,8.358186363473393) (0.3266331658291457,8.386621004233008) (0.3366834170854271,8.410926452066088) (0.3467336683417085,8.460185019562942) (0.35678391959798994,8.501771704915145) (0.3668341708542713,8.535744238320584) (0.37688442211055273,8.573653853553761) (0.3869346733668341,8.61841704426524) (0.3969849246231155,8.662724708356622) (0.40703517587939697,8.726195439233498) (0.4170854271356784,8.818084102932847) (0.42713567839195976,8.87134048669198) (0.43718592964824116,8.970637037109473) (0.44723618090452255,9.067689123854322) (0.457286432160804,9.130777182997141) (0.4673366834170854,9.213531576817802) (0.47738693467336685,9.293641463505244) (0.48743718592964824,9.40654631059225) (0.4974874371859296,9.501426811325192) (0.507537688442211,9.617026254933613) (0.5175879396984924,9.714217873443529) (0.5276381909547738,9.859977253284859) (0.5376884422110553,10.033196207869421) (0.5477386934673366,10.204606611660425) (0.557788944723618,10.427361385475162) (0.5678391959798994,10.581918015447696) (0.5778894472361809,10.785471230235931) (0.5879396984924622,10.96047708183288) (0.5979899497487436,11.173103111556848) (0.6080402010050251,11.404286610539723) (0.6180904522613064,11.68201596976112) (0.6281407035175879,11.918518552802496) (0.6381909547738693,12.156481566918563) (0.6482412060301507,12.455590378175442) (0.6582914572864321,12.718001306157625) (0.6683417085427136,13.003422493966264) (0.6783919597989949,13.231416810558313) (0.6884422110552764,13.57318703370529) (0.6984924623115578,13.909627766073998) (0.7085427135678392,14.226443414325747) (0.7185929648241206,14.625446161586208) (0.728643216080402,15.041745436277088) (0.7386934673366834,15.455917657544303) (0.7487437185929647,15.82387810125197) (0.7587939698492462,16.154610327306074) (0.7688442211055275,16.505580912755182) (0.778894472361809,16.852469035494103) (0.7889447236180904,17.22204657900596) (0.7989949748743718,17.55821209550956) (0.8090452261306533,17.919264982975868) (0.8190954773869346,18.30209062632119) (0.829145728643216,18.715544325727574) (0.8391959798994975,19.019673875734355) (0.8492462311557789,19.294191030105015) (0.8592964824120602,19.540142233164282) (0.8693467336683416,19.74101679467628) (0.8793969849246231,20.140148276679678) (0.8894472361809045,20.48949251969271) (0.8994974874371859,20.726038522142186) (0.9095477386934673,20.855149438517508) (0.9195979899497487,21.176609742759393) (0.9296482412060301,21.576095806772276) (0.9396984924623115,21.727561885830106) (0.949748743718593,22.043679175042648) (0.9597989949748743,22.33574378363641) (0.9698492462311558,22.685000606786502) (0.9798994974874371,22.950761252097276) (0.9899497487437185,23.122703332183345) (0.9999999999999999,23.28658794728972) };
  \addplot+[ours_jeffreys, mark=none, line width=0.4pt, dashed, forget plot] coordinates {(-1.0,5.642667516314929) (1.0,5.642667516314929)};
  \addplot+[ours_jeffreys, mark=*, only marks, mark size=\PanelArgminSize, forget plot] coordinates {(-0.06532663316582923,7.887664167160943)};
  \addplot[name path=U_ig_robust_grp, draw=none, mark=none, forget plot] coordinates {(-0.9999999999999999,8.154366255420385) (-0.9899497487437185,8.145607490211907) (-0.979899497487437,8.126919919780415) (-0.9698492462311558,8.11466408146943) (-0.9597989949748744,8.109323149763974) (-0.9497487437185927,8.097777780049254) (-0.9396984924623114,8.087370005803276) (-0.9296482412060303,8.071370087206336) (-0.9195979899497485,8.065017195951599) (-0.9095477386934673,8.050841294598996) (-0.8994974874371858,8.037075179298338) (-0.8894472361809044,8.017982307453904) (-0.8793969849246231,8.004276978826459) (-0.8693467336683416,7.993803207803356) (-0.85929648241206,7.996460641970502) (-0.8492462311557789,7.985936906598144) (-0.8391959798994973,7.977546373966519) (-0.8291457286432159,7.956186007458553) (-0.8190954773869346,7.937947808500501) (-0.8090452261306531,7.932771112040567) (-0.7989949748743718,7.925786524397144) (-0.7889447236180904,7.923523534664316) (-0.778894472361809,7.902654290452255) (-0.7688442211055275,7.898869487536266) (-0.7587939698492462,7.898939889461272) (-0.7487437185929647,7.88171884754566) (-0.7386934673366833,7.8749977357862) (-0.728643216080402,7.877802830789378) (-0.7185929648241205,7.876970815718678) (-0.708542713567839,7.861334285335007) (-0.6984924623115576,7.853387538329716) (-0.6884422110552764,7.858279184857819) (-0.6783919597989949,7.853969388286104) (-0.6683417085427135,7.856001559649608) (-0.6582914572864321,7.841990476750018) (-0.6482412060301507,7.845929450002102) (-0.6381909547738692,7.853897606413636) (-0.628140703517588,7.856668691070131) (-0.6180904522613064,7.85197942542954) (-0.608040201005025,7.856527058919981) (-0.5979899497487438,7.851164496204253) (-0.5879396984924622,7.837953336945279) (-0.5778894472361807,7.828444412677489) (-0.5678391959798995,7.825500234206178) (-0.5577889447236181,7.813121703534036) (-0.5477386934673365,7.813716940026592) (-0.5376884422110553,7.806916519796045) (-0.5276381909547738,7.7951401742282584) (-0.5175879396984924,7.78617372125922) (-0.507537688442211,7.7904058997655055) (-0.49748743718592964,7.782931463077429) (-0.4874371859296482,7.78578649379977) (-0.4773869346733667,7.788997011228959) (-0.46733668341708545,7.79650378902005) (-0.45728643216080395,7.797244166863677) (-0.4472361809045225,7.794165833657951) (-0.4371859296482412,7.785059261158039) (-0.42713567839195976,7.772599180585425) (-0.41708542713567825,7.777632058653657) (-0.40703517587939697,7.779228030669467) (-0.3969849246231155,7.784031983740899) (-0.386934673366834,7.775429443811911) (-0.3768844221105528,7.775177219566015) (-0.3668341708542713,7.781387070336417) (-0.3567839195979898,7.778145919754432) (-0.34673366834170855,7.783051285401272) (-0.3366834170854271,7.7792462256705965) (-0.32663316582914564,7.768003795837341) (-0.3165829145728644,7.759978677659861) (-0.30653266331658297,7.754128175304489) (-0.29648241206030146,7.74275045062943) (-0.28643216080402,7.731414146917063) (-0.2763819095477386,7.729004510308711) (-0.2663316582914572,7.723198753087499) (-0.2562814070351757,7.716151441680364) (-0.2462311557788945,7.716458401147718) (-0.236180904522613,7.713870103844293) (-0.22613065326633158,7.7172779771243185) (-0.21608040201005027,7.712542857606149) (-0.20603015075376882,7.703746011514389) (-0.19597989949748737,7.704478350514993) (-0.1859296482412061,7.7025183081953825) (-0.1758793969849246,7.710524159909053) (-0.16582914572864316,7.6980309263650115) (-0.15577889447236187,7.6891126149634434) (-0.14572864321608042,7.702604602585496) (-0.13567839195979894,7.712246210723784) (-0.12562814070351766,7.707270581126195) (-0.11557788944723618,7.704502525670842) (-0.10552763819095472,7.693771622949887) (-0.09547738693467328,7.696000071666709) (-0.08542713567839195,7.695537833820037) (-0.07537688442211053,7.689858359746457) (-0.06532663316582901,7.690793718520027) (-0.055276381909547756,7.680762035428519) (-0.045226130653266264,7.682731984863168) (-0.03517587939698496,7.6805199170203) (-0.02512562814070342,7.685298384781568) (-0.015075376884422061,7.6814238940393995) (-0.005025125628140743,7.695513957914678) (0.005025125628140765,7.689102802734334) (0.015075376884422073,7.689161020026033) (0.02512562814070352,7.695770876950185) (0.035175879396984966,7.705443004311386) (0.045226130653266354,7.710169622695541) (0.05527638190954779,7.7137535718803765) (0.06532663316582915,7.706563023623336) (0.0753768844221105,7.698297636192328) (0.08542713567839198,7.69547746592197) (0.09547738693467339,7.691949565774008) (0.10552763819095473,7.707079846477183) (0.11557788944723618,7.715955768280213) (0.1256281407035176,7.728489788029016) (0.1356783919597989,7.730640280491853) (0.14572864321608045,7.734625691245661) (0.15577889447236176,7.738707394220106) (0.16582914572864324,7.742526064707306) (0.17587939698492464,7.748007893587985) (0.185929648241206,7.743828331931569) (0.1959798994974875,7.7411759326077485) (0.20603015075376882,7.743038809513251) (0.21608040201005016,7.740508345770449) (0.2261306532663317,7.736053276784342) (0.23618090452261306,7.738807220528718) (0.24623115577889437,7.752713758984691) (0.2562814070351759,7.770649655910447) (0.2663316582914572,7.788858419200145) (0.2763819095477387,7.817780367041054) (0.2864321608040201,7.829674581084943) (0.2964824120603015,7.850442959856029) (0.30653266331658297,7.852062333197119) (0.31658291457286425,7.854794808826342) (0.3266331658291457,7.8706376974511345) (0.33668341708542715,7.905238832905274) (0.34673366834170855,7.932436134152281) (0.35678391959799,7.9206606274299425) (0.36683417085427134,7.933300215589411) (0.37688442211055273,7.950050652179425) (0.3869346733668341,7.958117279720479) (0.3969849246231155,7.994181686963847) (0.40703517587939697,8.006992884872579) (0.4170854271356783,8.025880338060562) (0.42713567839195976,8.046936466317065) (0.43718592964824116,8.068398339533955) (0.4472361809045226,8.083169061897777) (0.457286432160804,8.120073675040612) (0.46733668341708534,8.151348374516765) (0.47738693467336685,8.181960936385503) (0.48743718592964824,8.19687976129951) (0.4974874371859296,8.208675079791842) (0.507537688442211,8.22866029793242) (0.5175879396984924,8.281831661626388) (0.5276381909547738,8.293088956070168) (0.5376884422110552,8.348688281364666) (0.5477386934673366,8.38253056629595) (0.5577889447236181,8.427617803520256) (0.5678391959798994,8.447031742278845) (0.5778894472361809,8.47333086746987) (0.5879396984924623,8.487322834477814) (0.5979899497487436,8.5154633626828) (0.6080402010050251,8.566070356993963) (0.6180904522613064,8.629499079726047) (0.6281407035175878,8.669464732406787) (0.6381909547738693,8.718327447812944) (0.6482412060301507,8.770169051712546) (0.6582914572864321,8.836812084666791) (0.6683417085427136,8.921449259626849) (0.6783919597989949,8.98631624518629) (0.6884422110552764,9.064896441115964) (0.6984924623115577,9.142891281307561) (0.7085427135678392,9.22476969967396) (0.7185929648241206,9.315404516509313) (0.728643216080402,9.399708557526592) (0.7386934673366834,9.506060665475644) (0.7487437185929647,9.624582219836013) (0.7587939698492462,9.736052243529704) (0.7688442211055275,9.876115388095913) (0.778894472361809,10.00416513323652) (0.7889447236180903,10.172743832207324) (0.7989949748743718,10.37835150382964) (0.8090452261306532,10.566331527276127) (0.8190954773869346,10.72004631006038) (0.829145728643216,10.909559904143045) (0.8391959798994973,11.098852012762917) (0.8492462311557788,11.310666227798219) (0.8592964824120602,11.530551527539018) (0.8693467336683416,11.811945190721808) (0.879396984924623,12.089944122019899) (0.8894472361809044,12.32929080372012) (0.8994974874371858,12.553791557380665) (0.9095477386934673,12.903459104434347) (0.9195979899497487,13.129369084202427) (0.9296482412060301,13.418254728958804) (0.9396984924623115,13.7949759503415) (0.9497487437185929,14.130749409603917) (0.9597989949748743,14.490792162515701) (0.9698492462311558,14.775025151095322) (0.9798994974874371,15.278705654318246) (0.9899497487437185,15.701302346949126) (0.9999999999999999,16.087936395500634) };
  \addplot[name path=L_ig_robust_grp, draw=none, mark=none, forget plot] coordinates {(-0.9999999999999999,7.7791987282469295) (-0.9899497487437185,7.772178093635042) (-0.979899497487437,7.758084288366982) (-0.9698492462311558,7.749181686368904) (-0.9597989949748744,7.743977526997073) (-0.9497487437185927,7.733520761301134) (-0.9396984924623114,7.7157466254794045) (-0.9296482412060303,7.701305902295037) (-0.9195979899497485,7.6919049586182435) (-0.9095477386934673,7.677818679353327) (-0.8994974874371858,7.66263761103958) (-0.8894472361809044,7.645274612191432) (-0.8793969849246231,7.63742748432607) (-0.8693467336683416,7.62357634837259) (-0.85929648241206,7.622306765460717) (-0.8492462311557789,7.610374334393389) (-0.8391959798994973,7.601998117030564) (-0.8291457286432159,7.576853360314749) (-0.8190954773869346,7.5580629594356425) (-0.8090452261306531,7.557516731138199) (-0.7989949748743718,7.552783129926472) (-0.7889447236180904,7.552972787851536) (-0.778894472361809,7.540576764335488) (-0.7688442211055275,7.5364278194799335) (-0.7587939698492462,7.538087664534099) (-0.7487437185929647,7.52242122390075) (-0.7386934673366833,7.517598234395763) (-0.728643216080402,7.522612306249366) (-0.7185929648241205,7.52011908124001) (-0.708542713567839,7.508464144617421) (-0.6984924623115576,7.504296042758464) (-0.6884422110552764,7.500677953130049) (-0.6783919597989949,7.496397466870241) (-0.6683417085427135,7.499509792130329) (-0.6582914572864321,7.483375350991853) (-0.6482412060301507,7.489391780542198) (-0.6381909547738692,7.499876493316414) (-0.628140703517588,7.499315669796405) (-0.6180904522613064,7.495942516936423) (-0.608040201005025,7.4981852755657235) (-0.5979899497487438,7.494684394039717) (-0.5879396984924622,7.481884425239734) (-0.5778894472361807,7.473510451667022) (-0.5678391959798995,7.46894185507905) (-0.5577889447236181,7.455901163446636) (-0.5477386934673365,7.45764770317293) (-0.5376884422110553,7.4508652798496415) (-0.5276381909547738,7.43614032431033) (-0.5175879396984924,7.428469038708619) (-0.507537688442211,7.4352773831793035) (-0.49748743718592964,7.43192743347197) (-0.4874371859296482,7.432236187754904) (-0.4773869346733667,7.437215445377742) (-0.46733668341708545,7.441044456381168) (-0.45728643216080395,7.438105782266862) (-0.4472361809045225,7.431621292349291) (-0.4371859296482412,7.421677680921024) (-0.42713567839195976,7.409921153467545) (-0.41708542713567825,7.416687361044822) (-0.40703517587939697,7.423096903712466) (-0.3969849246231155,7.424440453023397) (-0.386934673366834,7.415901507481665) (-0.3768844221105528,7.418975640202514) (-0.3668341708542713,7.423687048910639) (-0.3567839195979898,7.426825762262429) (-0.34673366834170855,7.430408481786746) (-0.3366834170854271,7.426750955825814) (-0.32663316582914564,7.419635118361814) (-0.3165829145728644,7.414710450347661) (-0.30653266331658297,7.40232733604231) (-0.29648241206030146,7.397204298841157) (-0.28643216080402,7.384585696500879) (-0.2763819095477386,7.38108816610284) (-0.2663316582914572,7.37986471428924) (-0.2562814070351757,7.370932530101476) (-0.2462311557788945,7.373190759638246) (-0.236180904522613,7.369506237386139) (-0.22613065326633158,7.372674103336511) (-0.21608040201005027,7.365808481974151) (-0.20603015075376882,7.353078197758646) (-0.19597989949748737,7.353154379325961) (-0.1859296482412061,7.351226632686502) (-0.1758793969849246,7.363147117087877) (-0.16582914572864316,7.3554173364471795) (-0.15577889447236187,7.35467186258334) (-0.14572864321608042,7.3676022608359295) (-0.13567839195979894,7.373270610311259) (-0.12562814070351766,7.374903493682644) (-0.11557788944723618,7.37631520582382) (-0.10552763819095472,7.366858607486376) (-0.09547738693467328,7.366335396683659) (-0.08542713567839195,7.369266304529109) (-0.07537688442211053,7.363917297915818) (-0.06532663316582901,7.36587010579372) (-0.055276381909547756,7.352013485237955) (-0.045226130653266264,7.354141865432214) (-0.03517587939698496,7.3512396022941076) (-0.02512562814070342,7.357827481298653) (-0.015075376884422061,7.353418931967489) (-0.005025125628140743,7.367747037497108) (0.005025125628140765,7.360947031614337) (0.015075376884422073,7.367095863302277) (0.02512562814070352,7.373552096413864) (0.035175879396984966,7.390939188816471) (0.045226130653266354,7.395108339264714) (0.05527638190954779,7.401054113393219) (0.06532663316582915,7.389049950207865) (0.0753768844221105,7.375388341920982) (0.08542713567839198,7.371433522845274) (0.09547738693467339,7.366027169089545) (0.10552763819095473,7.3800971275997815) (0.11557788944723618,7.387256019647183) (0.1256281407035176,7.3991357979743615) (0.1356783919597989,7.402101024315524) (0.14572864321608045,7.411325743229838) (0.15577889447236176,7.414613920551399) (0.16582914572864324,7.417344671202512) (0.17587939698492464,7.418508405039473) (0.185929648241206,7.414092342965661) (0.1959798994974875,7.414563989240528) (0.20603015075376882,7.417202698250062) (0.21608040201005016,7.417116855848862) (0.2261306532663317,7.407184187401726) (0.23618090452261306,7.412181255827857) (0.24623115577889437,7.426141645758421) (0.2562814070351759,7.4498826251821795) (0.2663316582914572,7.460313604004944) (0.2763819095477387,7.479121980951678) (0.2864321608040201,7.484081906013282) (0.2964824120603015,7.494681227515951) (0.30653266331658297,7.492489125963185) (0.31658291457286425,7.495036890561695) (0.3266331658291457,7.506316744703137) (0.33668341708542715,7.534212731766823) (0.34673366834170855,7.55704911436629) (0.35678391959799,7.54750978197022) (0.36683417085427134,7.559995997554721) (0.37688442211055273,7.580201367723416) (0.3869346733668341,7.5822880875287435) (0.3969849246231155,7.612628493631869) (0.40703517587939697,7.619044079936436) (0.4170854271356783,7.635520268208657) (0.42713567839195976,7.660534742318356) (0.43718592964824116,7.682002542000837) (0.4472361809045226,7.696932940566522) (0.457286432160804,7.736666758627066) (0.46733668341708534,7.771419768981986) (0.47738693467336685,7.80595508959769) (0.48743718592964824,7.822511094524245) (0.4974874371859296,7.832638286658011) (0.507537688442211,7.849015688666151) (0.5175879396984924,7.888035043801557) (0.5276381909547738,7.898113013318246) (0.5376884422110552,7.932210702562599) (0.5477386934673366,7.962724410962312) (0.5577889447236181,8.00802114081053) (0.5678391959798994,8.030364530684924) (0.5778894472361809,8.064409513181092) (0.5879396984924623,8.088125076774665) (0.5979899497487436,8.115418381601712) (0.6080402010050251,8.167773043338476) (0.6180904522613064,8.222803982049884) (0.6281407035175878,8.262189065459527) (0.6381909547738693,8.30061360860435) (0.6482412060301507,8.332689347973368) (0.6582914572864321,8.375067465014228) (0.6683417085427136,8.438638133270237) (0.6783919597989949,8.505490613683447) (0.6884422110552764,8.581529771834282) (0.6984924623115577,8.661888478666155) (0.7085427135678392,8.74947766344756) (0.7185929648241206,8.84221668005816) (0.728643216080402,8.935333695397832) (0.7386934673366834,9.016371862973447) (0.7487437185929647,9.114387939351987) (0.7587939698492462,9.204942318556204) (0.7688442211055275,9.315563255746627) (0.778894472361809,9.411547692295585) (0.7889447236180903,9.552741418739615) (0.7989949748743718,9.738831428223877) (0.8090452261306532,9.91619841845084) (0.8190954773869346,10.109955787198079) (0.829145728643216,10.309996348215428) (0.8391959798994973,10.478166923202584) (0.8492462311557788,10.643079122965865) (0.8592964824120602,10.837006394622403) (0.8693467336683416,11.065984554662244) (0.879396984924623,11.37320349348382) (0.8894472361809044,11.595980650323138) (0.8994974874371858,11.7709764748427) (0.9095477386934673,12.081190939200948) (0.9195979899497487,12.326497055920257) (0.9296482412060301,12.612890085689376) (0.9396984924623115,12.875545149833973) (0.9497487437185929,13.133260990514954) (0.9597989949748743,13.458368429849719) (0.9698492462311558,13.772509300028931) (0.9798994974874371,14.16436747015546) (0.9899497487437185,14.49036896543287) (0.9999999999999999,14.94263010190224) };
  \addplot[ours_ig_robust, fill, fill opacity=0.18, draw=none, mark=none, forget plot] fill between [of=U_ig_robust_grp and L_ig_robust_grp];
  \addplot+[ours_ig_robust, mark=none] coordinates {(-0.9999999999999999,7.966782491833657) (-0.9899497487437185,7.958892791923475) (-0.979899497487437,7.942502104073698) (-0.9698492462311558,7.9319228839191664) (-0.9597989949748744,7.926650338380523) (-0.9497487437185927,7.915649270675194) (-0.9396984924623114,7.90155831564134) (-0.9296482412060303,7.886337994750687) (-0.9195979899497485,7.878461077284921) (-0.9095477386934673,7.864329986976162) (-0.8994974874371858,7.849856395168959) (-0.8894472361809044,7.831628459822668) (-0.8793969849246231,7.820852231576264) (-0.8693467336683416,7.808689778087973) (-0.85929648241206,7.80938370371561) (-0.8492462311557789,7.798155620495766) (-0.8391959798994973,7.789772245498542) (-0.8291457286432159,7.766519683886651) (-0.8190954773869346,7.748005383968072) (-0.8090452261306531,7.745143921589383) (-0.7989949748743718,7.739284827161808) (-0.7889447236180904,7.738248161257926) (-0.778894472361809,7.721615527393872) (-0.7688442211055275,7.7176486535081) (-0.7587939698492462,7.7185137769976855) (-0.7487437185929647,7.702070035723205) (-0.7386934673366833,7.696297985090982) (-0.728643216080402,7.700207568519372) (-0.7185929648241205,7.698544948479344) (-0.708542713567839,7.684899214976214) (-0.6984924623115576,7.67884179054409) (-0.6884422110552764,7.679478568993934) (-0.6783919597989949,7.6751834275781725) (-0.6683417085427135,7.677755675889968) (-0.6582914572864321,7.662682913870936) (-0.6482412060301507,7.66766061527215) (-0.6381909547738692,7.676887049865025) (-0.628140703517588,7.677992180433268) (-0.6180904522613064,7.673960971182981) (-0.608040201005025,7.677356167242852) (-0.5979899497487438,7.672924445121985) (-0.5879396984924622,7.659918881092507) (-0.5778894472361807,7.6509774321722555) (-0.5678391959798995,7.647221044642614) (-0.5577889447236181,7.634511433490336) (-0.5477386934673365,7.635682321599761) (-0.5376884422110553,7.628890899822843) (-0.5276381909547738,7.615640249269294) (-0.5175879396984924,7.60732137998392) (-0.507537688442211,7.6128416414724045) (-0.49748743718592964,7.6074294482747) (-0.4874371859296482,7.609011340777337) (-0.4773869346733667,7.613106228303351) (-0.46733668341708545,7.618774122700609) (-0.45728643216080395,7.617674974565269) (-0.4472361809045225,7.612893563003621) (-0.4371859296482412,7.6033684710395315) (-0.42713567839195976,7.591260167026485) (-0.41708542713567825,7.597159709849239) (-0.40703517587939697,7.601162467190966) (-0.3969849246231155,7.604236218382148) (-0.386934673366834,7.595665475646788) (-0.3768844221105528,7.597076429884265) (-0.3668341708542713,7.602537059623528) (-0.3567839195979898,7.602485841008431) (-0.34673366834170855,7.606729883594009) (-0.3366834170854271,7.602998590748205) (-0.32663316582914564,7.593819457099578) (-0.3165829145728644,7.587344564003761) (-0.30653266331658297,7.5782277556733995) (-0.29648241206030146,7.569977374735293) (-0.28643216080402,7.557999921708971) (-0.2763819095477386,7.555046338205775) (-0.2663316582914572,7.551531733688369) (-0.2562814070351757,7.54354198589092) (-0.2462311557788945,7.544824580392982) (-0.236180904522613,7.541688170615216) (-0.22613065326633158,7.544976040230415) (-0.21608040201005027,7.53917566979015) (-0.20603015075376882,7.5284121046365176) (-0.19597989949748737,7.528816364920477) (-0.1859296482412061,7.526872470440942) (-0.1758793969849246,7.536835638498465) (-0.16582914572864316,7.5267241314060955) (-0.15577889447236187,7.5218922387733915) (-0.14572864321608042,7.535103431710713) (-0.13567839195979894,7.542758410517521) (-0.12562814070351766,7.541087037404419) (-0.11557788944723618,7.540408865747331) (-0.10552763819095472,7.530315115218132) (-0.09547738693467328,7.531167734175184) (-0.08542713567839195,7.532402069174573) (-0.07537688442211053,7.526887828831137) (-0.06532663316582901,7.528331912156873) (-0.055276381909547756,7.516387760333237) (-0.045226130653266264,7.518436925147691) (-0.03517587939698496,7.515879759657204) (-0.02512562814070342,7.52156293304011) (-0.015075376884422061,7.5174214130034445) (-0.005025125628140743,7.531630497705893) (0.005025125628140765,7.525024917174336) (0.015075376884422073,7.528128441664155) (0.02512562814070352,7.534661486682024) (0.035175879396984966,7.5481910965639285) (0.045226130653266354,7.552638980980127) (0.05527638190954779,7.5574038426367975) (0.06532663316582915,7.547806486915601) (0.0753768844221105,7.536842989056655) (0.08542713567839198,7.533455494383622) (0.09547738693467339,7.528988367431777) (0.10552763819095473,7.543588487038482) (0.11557788944723618,7.551605893963698) (0.1256281407035176,7.563812793001689) (0.1356783919597989,7.5663706524036884) (0.14572864321608045,7.57297571723775) (0.15577889447236176,7.5766606573857525) (0.16582914572864324,7.579935367954909) (0.17587939698492464,7.583258149313729) (0.185929648241206,7.578960337448615) (0.1959798994974875,7.577869960924138) (0.20603015075376882,7.580120753881657) (0.21608040201005016,7.5788126008096555) (0.2261306532663317,7.571618732093034) (0.23618090452261306,7.575494238178288) (0.24623115577889437,7.589427702371556) (0.2562814070351759,7.6102661405463135) (0.2663316582914572,7.624586011602545) (0.2763819095477387,7.648451173996366) (0.2864321608040201,7.656878243549112) (0.2964824120603015,7.67256209368599) (0.30653266331658297,7.672275729580152) (0.31658291457286425,7.674915849694019) (0.3266331658291457,7.688477221077136) (0.33668341708542715,7.719725782336049) (0.34673366834170855,7.744742624259286) (0.35678391959799,7.734085204700081) (0.36683417085427134,7.746648106572066) (0.37688442211055273,7.7651260099514205) (0.3869346733668341,7.770202683624611) (0.3969849246231155,7.803405090297858) (0.40703517587939697,7.813018482404507) (0.4170854271356783,7.830700303134609) (0.42713567839195976,7.853735604317711) (0.43718592964824116,7.875200440767396) (0.4472361809045226,7.890051001232149) (0.457286432160804,7.928370216833839) (0.46733668341708534,7.961384071749375) (0.47738693467336685,7.993958012991596) (0.48743718592964824,8.009695427911877) (0.4974874371859296,8.020656683224926) (0.507537688442211,8.038837993299286) (0.5175879396984924,8.084933352713973) (0.5276381909547738,8.095600984694206) (0.5376884422110552,8.140449491963633) (0.5477386934673366,8.17262748862913) (0.5577889447236181,8.217819472165393) (0.5678391959798994,8.238698136481885) (0.5778894472361809,8.268870190325481) (0.5879396984924623,8.28772395562624) (0.5979899497487436,8.315440872142256) (0.6080402010050251,8.36692170016622) (0.6180904522613064,8.426151530887966) (0.6281407035175878,8.465826898933157) (0.6381909547738693,8.509470528208647) (0.6482412060301507,8.551429199842957) (0.6582914572864321,8.60593977484051) (0.6683417085427136,8.680043696448543) (0.6783919597989949,8.745903429434868) (0.6884422110552764,8.823213106475123) (0.6984924623115577,8.902389879986858) (0.7085427135678392,8.98712368156076) (0.7185929648241206,9.078810598283736) (0.728643216080402,9.167521126462212) (0.7386934673366834,9.261216264224545) (0.7487437185929647,9.369485079594) (0.7587939698492462,9.470497281042954) (0.7688442211055275,9.59583932192127) (0.778894472361809,9.707856412766052) (0.7889447236180903,9.86274262547347) (0.7989949748743718,10.058591466026758) (0.8090452261306532,10.241264972863483) (0.8190954773869346,10.415001048629229) (0.829145728643216,10.609778126179236) (0.8391959798994973,10.78850946798275) (0.8492462311557788,10.976872675382042) (0.8592964824120602,11.18377896108071) (0.8693467336683416,11.438964872692026) (0.879396984924623,11.73157380775186) (0.8894472361809044,11.96263572702163) (0.8994974874371858,12.162384016111682) (0.9095477386934673,12.492325021817647) (0.9195979899497487,12.727933070061342) (0.9296482412060301,13.01557240732409) (0.9396984924623115,13.335260550087737) (0.9497487437185929,13.632005200059435) (0.9597989949748743,13.97458029618271) (0.9698492462311558,14.273767225562127) (0.9798994974874371,14.721536562236853) (0.9899497487437185,15.095835656190998) (0.9999999999999999,15.515283248701436) };
  \addplot+[ours_ig_robust, mark=none, line width=0.4pt, dashed, forget plot] coordinates {(-1.0,5.353128536252401) (1.0,5.353128536252401)};
  \addplot+[ours_ig_robust, mark=*, only marks, mark size=\PanelArgminSize, forget plot] coordinates {(-0.03517587939698496,7.515879759657204)};
  \addplot[name path=U_ig_mle_grp, draw=none, mark=none, forget plot] coordinates {(-1.0,10.293926670229563) (-0.9899497487437185,10.233484115646123) (-0.9798994974874373,10.18861319082079) (-0.9698492462311559,10.166152559897615) (-0.9597989949748744,10.093604171326689) (-0.9497487437185931,10.05019973224735) (-0.9396984924623116,10.009754154864309) (-0.9296482412060302,9.9816965065634) (-0.9195979899497487,9.958118517786911) (-0.9095477386934674,9.906073315997746) (-0.8994974874371859,9.880042155863963) (-0.8894472361809045,9.815739837814947) (-0.8793969849246233,9.756501168960133) (-0.8693467336683417,9.71561299571838) (-0.8592964824120604,9.671373335096462) (-0.849246231155779,9.616944979285451) (-0.8391959798994976,9.56830955211563) (-0.829145728643216,9.526308613834154) (-0.8190954773869348,9.479636838175884) (-0.8090452261306533,9.46193307672219) (-0.7989949748743719,9.41761970040807) (-0.7889447236180905,9.374269897020197) (-0.7788944723618091,9.342001387972388) (-0.7688442211055276,9.333829774464146) (-0.7587939698492462,9.298049394254052) (-0.748743718592965,9.260989812288349) (-0.7386934673366835,9.278449376048936) (-0.728643216080402,9.249146334202537) (-0.7185929648241207,9.185039067557138) (-0.7085427135678393,9.152734900282697) (-0.6984924623115577,9.122122565682355) (-0.6884422110552765,9.102679578246532) (-0.678391959798995,9.083652021957835) (-0.6683417085427136,9.081618521923149) (-0.6582914572864323,9.0465298544718) (-0.6482412060301508,9.055318735604706) (-0.6381909547738693,9.025997004102276) (-0.6281407035175881,8.999112568953304) (-0.6180904522613067,8.991553069995728) (-0.6080402010050251,8.965653216810814) (-0.5979899497487436,8.918041949199253) (-0.5879396984924624,8.919542645384231) (-0.5778894472361809,8.8757901188807) (-0.5678391959798995,8.848537027264575) (-0.5577889447236182,8.83397908961332) (-0.5477386934673368,8.803343388907008) (-0.5376884422110553,8.784318304742158) (-0.527638190954774,8.750028678537557) (-0.5175879396984925,8.718590258781731) (-0.507537688442211,8.697637216384372) (-0.4974874371859298,8.672304043441757) (-0.4874371859296483,8.63559859706834) (-0.4773869346733669,8.60887289722435) (-0.4673366834170855,8.591328435797235) (-0.4572864321608041,8.576755205560701) (-0.4472361809045226,8.561538755116006) (-0.43718592964824116,8.554993856903371) (-0.42713567839195987,8.539020521658314) (-0.4170854271356785,8.52967714003905) (-0.40703517587939697,8.50369951092573) (-0.39698492462311574,8.487343096843428) (-0.38693467336683424,8.472262259740235) (-0.37688442211055273,8.461484425162695) (-0.36683417085427145,8.45646842724298) (-0.35678391959799,8.435233062026226) (-0.3467336683417085,8.42095624151185) (-0.33668341708542726,8.390746041068729) (-0.3266331658291458,8.36388109328411) (-0.31658291457286436,8.346272756880625) (-0.30653266331658285,8.31215727140024) (-0.29648241206030157,8.298475629505777) (-0.2864321608040201,8.275025083344772) (-0.27638190954773884,8.260856194051154) (-0.26633165829145733,8.243805518602617) (-0.25628140703517593,8.227828012417882) (-0.2462311557788945,8.22287983755811) (-0.2361809045226131,8.210379794911773) (-0.22613065326633172,8.204909209867308) (-0.21608040201005035,8.18541727588506) (-0.20603015075376882,8.177886091394333) (-0.19597989949748754,8.171260745830576) (-0.18592964824120603,8.159618881757613) (-0.1758793969849247,8.1422993013953) (-0.1658291457286433,8.1347955884619) (-0.15577889447236182,8.132472368854826) (-0.14572864321608048,8.130306667684014) (-0.13567839195979906,8.130455237307862) (-0.12562814070351766,8.121974389476147) (-0.11557788944723625,8.11706919630585) (-0.10552763819095479,8.114824033712775) (-0.0954773869346734,8.109270582715832) (-0.08542713567839201,8.098232312230833) (-0.07537688442211053,8.090607887847824) (-0.06532663316582923,8.070587040769555) (-0.05527638190954781,8.081498510307725) (-0.04522613065326632,8.091172337206363) (-0.03517587939698496,8.084341702235784) (-0.025125628140703526,8.0995179841046) (-0.015075376884422112,8.10521443794581) (-0.005025125628140743,8.108148462865506) (0.00502512562814067,8.110016436358688) (0.015075376884422073,8.110879274465596) (0.025125628140703425,8.117835783000107) (0.035175879396984966,8.111452896212844) (0.045226130653266264,8.104108197631662) (0.0552763819095477,8.114401481105379) (0.06532663316582915,8.122308830628274) (0.0753768844221105,8.14068821883516) (0.08542713567839191,8.156471466821145) (0.09547738693467332,8.155166119351511) (0.10552763819095466,8.175928882224847) (0.11557788944723618,8.186244453132645) (0.12562814070351747,8.202090131960425) (0.13567839195979894,8.22269278402255) (0.1457286432160804,8.214685889681972) (0.15577889447236176,8.224530380090583) (0.16582914572864316,8.229896408594755) (0.17587939698492458,8.257283914983834) (0.185929648241206,8.28259017400937) (0.19597989949748737,8.301495949447878) (0.20603015075376876,8.334184122525066) (0.21608040201005016,8.35835131076137) (0.22613065326633164,8.35206572724464) (0.236180904522613,8.372125969283184) (0.24623115577889443,8.407931381467517) (0.25628140703517577,8.4237522931109) (0.2663316582914573,8.445144266187175) (0.2763819095477386,8.474732942703271) (0.28643216080402,8.502969584967978) (0.29648241206030146,8.523899794490259) (0.30653266331658285,8.537883067482305) (0.31658291457286425,8.558837405804445) (0.3266331658291457,8.576985531185393) (0.3366834170854271,8.606988232751915) (0.3467336683417085,8.66354661679018) (0.35678391959798994,8.71821175052416) (0.3668341708542713,8.75203217057162) (0.37688442211055273,8.79284404157655) (0.3869346733668341,8.845044526382186) (0.3969849246231155,8.898688181096329) (0.40703517587939697,8.960022313754902) (0.4170854271356784,9.055090652787435) (0.42713567839195976,9.10470003149057) (0.43718592964824116,9.20601230683615) (0.44723618090452255,9.2978724915666) (0.457286432160804,9.353706296408834) (0.4673366834170854,9.441940999217374) (0.47738693467336685,9.542761642392811) (0.48743718592964824,9.661256353207463) (0.4974874371859296,9.757110063742372) (0.507537688442211,9.892838652153191) (0.5175879396984924,9.997053288930188) (0.5276381909547738,10.162229167673567) (0.5376884422110553,10.34769907391653) (0.5477386934673366,10.523184708312575) (0.557788944723618,10.719118387386061) (0.5678391959798994,10.877340523386605) (0.5778894472361809,11.080953708782777) (0.5879396984924622,11.295102266377747) (0.5979899497487436,11.524287292648935) (0.6080402010050251,11.778478669110378) (0.6180904522613064,12.049333948277006) (0.6281407035175879,12.28417109292928) (0.6381909547738693,12.554889917133043) (0.6482412060301507,12.859898528931913) (0.6582914572864321,13.123722436364048) (0.6683417085427136,13.411751006452588) (0.6783919597989949,13.688456422321742) (0.6884422110552764,14.058947699141148) (0.6984924623115578,14.413948793414969) (0.7085427135678392,14.725747491501512) (0.7185929648241206,15.17129973517253) (0.728643216080402,15.644647015567458) (0.7386934673366834,16.0416285894858) (0.7487437185929647,16.411416589996083) (0.7587939698492462,16.796570858631465) (0.7688442211055275,17.13650353379442) (0.778894472361809,17.473338540352746) (0.7889447236180904,17.860363670359153) (0.7989949748743718,18.217163488187932) (0.8090452261306533,18.56491614486774) (0.8190954773869346,18.978993424448866) (0.829145728643216,19.42035152948556) (0.8391959798994975,19.752261565999454) (0.8492462311557789,20.036400147828363) (0.8592964824120602,20.264815147731785) (0.8693467336683416,20.447097632886603) (0.8793969849246231,20.87253652633315) (0.8894472361809045,21.22170726862629) (0.8994974874371859,21.43321558312304) (0.9095477386934673,21.569159280068924) (0.9195979899497487,21.893507661394743) (0.9296482412060301,22.322874950561406) (0.9396984924623115,22.474464143039025) (0.949748743718593,22.7559800693313) (0.9597989949748743,23.067204401885483) (0.9698492462311558,23.430919651915307) (0.9798994974874371,23.730923059951127) (0.9899497487437185,23.90945193457296) (0.9999999999999999,24.059158245119747) };
  \addplot[name path=L_ig_mle_grp, draw=none, mark=none, forget plot] coordinates {(-1.0,9.66206862903912) (-0.9899497487437185,9.60419605440226) (-0.9798994974874373,9.551844314218846) (-0.9698492462311559,9.531122536524524) (-0.9597989949748744,9.473283948278592) (-0.9497487437185931,9.438155629442885) (-0.9396984924623116,9.38295147983253) (-0.9296482412060302,9.348914363648094) (-0.9195979899497487,9.327731405510978) (-0.9095477386934674,9.268704565421322) (-0.8994974874371859,9.239002493772949) (-0.8894472361809045,9.186968090617674) (-0.8793969849246233,9.128042332357554) (-0.8693467336683417,9.102512225623542) (-0.8592964824120604,9.061069936649773) (-0.849246231155779,9.012919148013436) (-0.8391959798994976,8.9787018592872) (-0.829145728643216,8.93892315438099) (-0.8190954773869348,8.895502504276406) (-0.8090452261306533,8.871915048694348) (-0.7989949748743719,8.830765738177933) (-0.7889447236180905,8.804591035300588) (-0.7788944723618091,8.775134683756267) (-0.7688442211055276,8.770940239750239) (-0.7587939698492462,8.738989887851423) (-0.748743718592965,8.71328574480521) (-0.7386934673366835,8.718665107188617) (-0.728643216080402,8.684458107270602) (-0.7185929648241207,8.624553187222734) (-0.7085427135678393,8.59908014778082) (-0.6984924623115577,8.568206000202343) (-0.6884422110552765,8.548200053105125) (-0.678391959798995,8.534903788832043) (-0.6683417085427136,8.531017144061442) (-0.6582914572864323,8.505525716499925) (-0.6482412060301508,8.5091759127238) (-0.6381909547738693,8.489660475706575) (-0.6281407035175881,8.477678440177318) (-0.6180904522613067,8.469811891331979) (-0.6080402010050251,8.450949079465495) (-0.5979899497487436,8.41309400295225) (-0.5879396984924624,8.407297208210915) (-0.5778894472361809,8.370118089767969) (-0.5678391959798995,8.351007402343134) (-0.5577889447236182,8.339975185433316) (-0.5477386934673368,8.321347813145165) (-0.5376884422110553,8.307976265659725) (-0.527638190954774,8.279840218753144) (-0.5175879396984925,8.257572368558094) (-0.507537688442211,8.237869649917776) (-0.4974874371859298,8.226214795063305) (-0.4874371859296483,8.202527488475868) (-0.4773869346733669,8.182100563244417) (-0.4673366834170855,8.170672850228621) (-0.4572864321608041,8.165613041893543) (-0.4472361809045226,8.155392856937494) (-0.43718592964824116,8.14987915684314) (-0.42713567839195987,8.13974191243849) (-0.4170854271356785,8.123487263933063) (-0.40703517587939697,8.098345317067489) (-0.39698492462311574,8.081827975774331) (-0.38693467336683424,8.05999553658319) (-0.37688442211055273,8.045755828827854) (-0.36683417085427145,8.042557810981274) (-0.35678391959799,8.028467758171525) (-0.3467336683417085,8.015340172603997) (-0.33668341708542726,7.990235322376802) (-0.3266331658291458,7.958123372687518) (-0.31658291457286436,7.94243323559854) (-0.30653266331658285,7.911956918244224) (-0.29648241206030157,7.899815928529798) (-0.2864321608040201,7.882994435936734) (-0.27638190954773884,7.878453015180757) (-0.26633165829145733,7.8680578233326655) (-0.25628140703517593,7.8502363783703055) (-0.2462311557788945,7.844981890297524) (-0.2361809045226131,7.840044285772898) (-0.22613065326633172,7.834029956411816) (-0.21608040201005035,7.816983735922582) (-0.20603015075376882,7.812160496652212) (-0.19597989949748754,7.80308479914071) (-0.18592964824120603,7.786715354448266) (-0.1758793969849247,7.778877208307413) (-0.1658291457286433,7.773473864879759) (-0.15577889447236182,7.771991139478406) (-0.14572864321608048,7.76647743919422) (-0.13567839195979906,7.769050149186791) (-0.12562814070351766,7.7601926545639195) (-0.11557788944723625,7.758641697869106) (-0.10552763819095479,7.763784075380675) (-0.0954773869346734,7.759091285349576) (-0.08542713567839201,7.7493733729326015) (-0.07537688442211053,7.739372319324665) (-0.06532663316582923,7.718610180686331) (-0.05527638190954781,7.729610959763211) (-0.04522613065326632,7.731153125829171) (-0.03517587939698496,7.727040044279483) (-0.025125628140703526,7.7395124639347195) (-0.015075376884422112,7.743172929089331) (-0.005025125628140743,7.755420956882267) (0.00502512562814067,7.7590338176030516) (0.015075376884422073,7.757508436597243) (0.025125628140703425,7.7657160603902895) (0.035175879396984966,7.765024635725415) (0.045226130653266264,7.758418724025288) (0.0552763819095477,7.766557464764539) (0.06532663316582915,7.772594886972669) (0.0753768844221105,7.799086630592844) (0.08542713567839191,7.813134049014839) (0.09547738693467332,7.816916650275757) (0.10552763819095466,7.828289154836413) (0.11557788944723618,7.833911476200444) (0.12562814070351747,7.847674318396488) (0.13567839195979894,7.859398983371967) (0.1457286432160804,7.851848806604241) (0.15577889447236176,7.860226973263548) (0.16582914572864316,7.861669541453058) (0.17587939698492458,7.8842863119871005) (0.185929648241206,7.902174628855538) (0.19597989949748737,7.920832017917534) (0.20603015075376876,7.944151545089749) (0.21608040201005016,7.959053567093603) (0.22613065326633164,7.958423751221002) (0.236180904522613,7.974149197419703) (0.24623115577889443,8.01030345172017) (0.25628140703517577,8.024844947474117) (0.2663316582914573,8.051126107652635) (0.2763819095477386,8.084278696597925) (0.28643216080402,8.116366353327901) (0.29648241206030146,8.140495751871363) (0.30653266331658285,8.155845176028073) (0.31658291457286425,8.176355034239952) (0.3266331658291457,8.200999752478591) (0.3366834170854271,8.22963042075279) (0.3467336683417085,8.268581503041661) (0.35678391959798994,8.307268713158178) (0.3668341708542713,8.337777215166998) (0.37688442211055273,8.368932040453675) (0.3869346733668341,8.406360847026296) (0.3969849246231155,8.446074336268701) (0.40703517587939697,8.50118519058152) (0.4170854271356784,8.593124192605) (0.42713567839195976,8.647969318114374) (0.43718592964824116,8.743258738406212) (0.44723618090452255,8.846091682201685) (0.457286432160804,8.917199487709098) (0.4673366834170854,8.988697425362771) (0.47738693467336685,9.056643006630324) (0.48743718592964824,9.169058838467297) (0.4974874371859296,9.254592779115002) (0.507537688442211,9.36194370516493) (0.5175879396984924,9.448599993150319) (0.5276381909547738,9.574791541457861) (0.5376884422110553,9.738643048736261) (0.5477386934673366,9.906881492858771) (0.557788944723618,10.158000797481456) (0.5678391959798994,10.31873224548366) (0.5778894472361809,10.503234645769057) (0.5879396984924622,10.659947698170853) (0.5979899497487436,10.83989992754151) (0.6080402010050251,11.064952400871896) (0.6180904522613064,11.333895435340224) (0.6281407035175879,11.584699545479149) (0.6381909547738693,11.792834647005886) (0.6482412060301507,12.072812466608417) (0.6582914572864321,12.33769532134076) (0.6683417085427136,12.605035176732821) (0.6783919597989949,12.802316082099491) (0.6884422110552764,13.100836226360661) (0.6984924623115578,13.423887503903392) (0.7085427135678392,13.768346600713869) (0.7185929648241206,14.115289647728643) (0.728643216080402,14.471506721896862) (0.7386934673366834,14.915489639214911) (0.7487437185929647,15.278510674888068) (0.7587939698492462,15.585880894555219) (0.7688442211055275,15.926329882775423) (0.778894472361809,16.294423129686518) (0.7889447236180904,16.63315635730154) (0.7989949748743718,16.95928675614157) (0.8090452261306533,17.321058456547757) (0.8190954773869346,17.677809529880655) (0.829145728643216,18.050593443888406) (0.8391959798994975,18.359876896753782) (0.8492462311557789,18.629090795174616) (0.8592964824120602,18.88026601073525) (0.8693467336683416,19.103166745756695) (0.8793969849246231,19.461206937578375) (0.8894472361809045,19.795790540751675) (0.8994974874371859,20.055743337730394) (0.9095477386934673,20.201884562744738) (0.9195979899497487,20.514048143252907) (0.9296482412060301,20.873486934444145) (0.9396984924623115,21.029616932062407) (0.949748743718593,21.35676631746155) (0.9597989949748743,21.660189726082777) (0.9698492462311558,21.970256041967264) (0.9798994974874371,22.229553258876027) (0.9899497487437185,22.391322072943698) (0.9999999999999999,22.564179419286795) };
  \addplot[ours_ig_mle, fill, fill opacity=0.18, draw=none, mark=none, forget plot] fill between [of=U_ig_mle_grp and L_ig_mle_grp];
  \addplot+[ours_ig_mle, mark=none] coordinates {(-1.0,9.977997649634341) (-0.9899497487437185,9.918840085024192) (-0.9798994974874373,9.870228752519818) (-0.9698492462311559,9.84863754821107) (-0.9597989949748744,9.78344405980264) (-0.9497487437185931,9.744177680845118) (-0.9396984924623116,9.69635281734842) (-0.9296482412060302,9.665305435105747) (-0.9195979899497487,9.642924961648944) (-0.9095477386934674,9.587388940709534) (-0.8994974874371859,9.559522324818456) (-0.8894472361809045,9.50135396421631) (-0.8793969849246233,9.442271750658843) (-0.8693467336683417,9.409062610670961) (-0.8592964824120604,9.366221635873117) (-0.849246231155779,9.314932063649444) (-0.8391959798994976,9.273505705701416) (-0.829145728643216,9.232615884107572) (-0.8190954773869348,9.187569671226145) (-0.8090452261306533,9.166924062708269) (-0.7989949748743719,9.124192719293001) (-0.7889447236180905,9.089430466160392) (-0.7788944723618091,9.058568035864328) (-0.7688442211055276,9.052385007107192) (-0.7587939698492462,9.018519641052738) (-0.748743718592965,8.98713777854678) (-0.7386934673366835,8.998557241618776) (-0.728643216080402,8.96680222073657) (-0.7185929648241207,8.904796127389936) (-0.7085427135678393,8.875907524031758) (-0.6984924623115577,8.84516428294235) (-0.6884422110552765,8.825439815675828) (-0.678391959798995,8.80927790539494) (-0.6683417085427136,8.806317832992296) (-0.6582914572864323,8.776027785485862) (-0.6482412060301508,8.782247324164253) (-0.6381909547738693,8.757828739904426) (-0.6281407035175881,8.738395504565311) (-0.6180904522613067,8.730682480663853) (-0.6080402010050251,8.708301148138155) (-0.5979899497487436,8.665567976075751) (-0.5879396984924624,8.663419926797573) (-0.5778894472361809,8.622954104324334) (-0.5678391959798995,8.599772214803854) (-0.5577889447236182,8.586977137523318) (-0.5477386934673368,8.562345601026086) (-0.5376884422110553,8.546147285200941) (-0.527638190954774,8.51493444864535) (-0.5175879396984925,8.488081313669912) (-0.507537688442211,8.467753433151074) (-0.4974874371859298,8.449259419252531) (-0.4874371859296483,8.419063042772104) (-0.4773869346733669,8.395486730234383) (-0.4673366834170855,8.381000643012928) (-0.4572864321608041,8.371184123727122) (-0.4472361809045226,8.35846580602675) (-0.43718592964824116,8.352436506873255) (-0.42713567839195987,8.339381217048402) (-0.4170854271356785,8.326582201986056) (-0.40703517587939697,8.30102241399661) (-0.39698492462311574,8.28458553630888) (-0.38693467336683424,8.266128898161712) (-0.37688442211055273,8.253620126995274) (-0.36683417085427145,8.249513119112127) (-0.35678391959799,8.231850410098875) (-0.3467336683417085,8.218148207057924) (-0.33668341708542726,8.190490681722766) (-0.3266331658291458,8.161002232985814) (-0.31658291457286436,8.144352996239583) (-0.30653266331658285,8.112057094822232) (-0.29648241206030157,8.099145779017787) (-0.2864321608040201,8.079009759640753) (-0.27638190954773884,8.069654604615955) (-0.26633165829145733,8.055931670967642) (-0.25628140703517593,8.039032195394094) (-0.2462311557788945,8.033930863927818) (-0.2361809045226131,8.025212040342335) (-0.22613065326633172,8.019469583139562) (-0.21608040201005035,8.001200505903821) (-0.20603015075376882,7.995023294023273) (-0.19597989949748754,7.987172772485643) (-0.18592964824120603,7.97316711810294) (-0.1758793969849247,7.960588254851356) (-0.1658291457286433,7.95413472667083) (-0.15577889447236182,7.952231754166616) (-0.14572864321608048,7.948392053439117) (-0.13567839195979906,7.949752693247326) (-0.12562814070351766,7.941083522020033) (-0.11557788944723625,7.937855447087478) (-0.10552763819095479,7.939304054546725) (-0.0954773869346734,7.934180934032704) (-0.08542713567839201,7.923802842581717) (-0.07537688442211053,7.914990103586244) (-0.06532663316582923,7.894598610727943) (-0.05527638190954781,7.905554735035468) (-0.04522613065326632,7.911162731517766) (-0.03517587939698496,7.905690873257633) (-0.025125628140703526,7.919515224019659) (-0.015075376884422112,7.92419368351757) (-0.005025125628140743,7.931784709873886) (0.00502512562814067,7.93452512698087) (0.015075376884422073,7.9341938555314195) (0.025125628140703425,7.941775921695198) (0.035175879396984966,7.938238765969129) (0.045226130653266264,7.931263460828475) (0.0552763819095477,7.940479472934959) (0.06532663316582915,7.947451858800472) (0.0753768844221105,7.969887424714002) (0.08542713567839191,7.984802757917992) (0.09547738693467332,7.9860413848136345) (0.10552763819095466,8.00210901853063) (0.11557788944723618,8.010077964666545) (0.12562814070351747,8.024882225178457) (0.13567839195979894,8.041045883697258) (0.1457286432160804,8.033267348143106) (0.15577889447236176,8.042378676677066) (0.16582914572864316,8.045782975023906) (0.17587939698492458,8.070785113485467) (0.185929648241206,8.092382401432454) (0.19597989949748737,8.111163983682706) (0.20603015075376876,8.139167833807408) (0.21608040201005016,8.158702438927486) (0.22613065326633164,8.15524473923282) (0.236180904522613,8.173137583351444) (0.24623115577889443,8.209117416593843) (0.25628140703517577,8.224298620292508) (0.2663316582914573,8.248135186919905) (0.2763819095477386,8.279505819650598) (0.28643216080402,8.30966796914794) (0.29648241206030146,8.33219777318081) (0.30653266331658285,8.346864121755189) (0.31658291457286425,8.367596220022198) (0.3266331658291457,8.388992641831992) (0.3366834170854271,8.418309326752352) (0.3467336683417085,8.46606405991592) (0.35678391959798994,8.51274023184117) (0.3668341708542713,8.54490469286931) (0.37688442211055273,8.580888041015113) (0.3869346733668341,8.62570268670424) (0.3969849246231155,8.672381258682515) (0.40703517587939697,8.73060375216821) (0.4170854271356784,8.824107422696217) (0.42713567839195976,8.876334674802472) (0.43718592964824116,8.974635522621181) (0.44723618090452255,9.071982086884143) (0.457286432160804,9.135452892058966) (0.4673366834170854,9.215319212290073) (0.47738693467336685,9.299702324511568) (0.48743718592964824,9.41515759583738) (0.4974874371859296,9.505851421428687) (0.507537688442211,9.62739117865906) (0.5175879396984924,9.722826641040253) (0.5276381909547738,9.868510354565714) (0.5376884422110553,10.043171061326396) (0.5477386934673366,10.215033100585673) (0.557788944723618,10.438559592433759) (0.5678391959798994,10.598036384435133) (0.5778894472361809,10.792094177275917) (0.5879396984924622,10.9775249822743) (0.5979899497487436,11.182093610095222) (0.6080402010050251,11.421715534991137) (0.6180904522613064,11.691614691808615) (0.6281407035175879,11.934435319204214) (0.6381909547738693,12.173862282069464) (0.6482412060301507,12.466355497770165) (0.6582914572864321,12.730708878852404) (0.6683417085427136,13.008393091592705) (0.6783919597989949,13.245386252210617) (0.6884422110552764,13.579891962750905) (0.6984924623115578,13.91891814865918) (0.7085427135678392,14.24704704610769) (0.7185929648241206,14.643294691450587) (0.728643216080402,15.05807686873216) (0.7386934673366834,15.478559114350356) (0.7487437185929647,15.844963632442076) (0.7587939698492462,16.191225876593343) (0.7688442211055275,16.531416708284922) (0.778894472361809,16.88388083501963) (0.7889447236180904,17.246760013830347) (0.7989949748743718,17.58822512216475) (0.8090452261306533,17.94298730070775) (0.8190954773869346,18.32840147716476) (0.829145728643216,18.735472486686984) (0.8391959798994975,19.05606923137662) (0.8492462311557789,19.33274547150149) (0.8592964824120602,19.572540579233518) (0.8693467336683416,19.77513218932165) (0.8793969849246231,20.166871731955762) (0.8894472361809045,20.508748904688982) (0.8994974874371859,20.744479460426717) (0.9095477386934673,20.88552192140683) (0.9195979899497487,21.203777902323825) (0.9296482412060301,21.598180942502776) (0.9396984924623115,21.752040537550716) (0.949748743718593,22.056373193396425) (0.9597989949748743,22.36369706398413) (0.9698492462311558,22.700587846941286) (0.9798994974874371,22.980238159413577) (0.9899497487437185,23.15038700375833) (0.9999999999999999,23.31166883220327) };
  \addplot+[ours_ig_mle, mark=none, line width=0.4pt, dashed, forget plot] coordinates {(-1.0,5.65183986220968) (1.0,5.65183986220968)};
  \addplot+[ours_ig_mle, mark=*, only marks, mark size=\PanelArgminSize, forget plot] coordinates {(-0.06532663316582923,7.894598610727943)};
  \addplot+[ours_oracle, mark=none, forget plot] coordinates {(-1.0,1.930526897094334) (1.0,1.930526897094334)};
%
			\let\nextgroupplot\origNGP%
		\end{groupplot}
	\end{tikzpicture}
	\par\vspace{0.4em}
	% Paper-side legend: single row.  4 threshold baselines, then an
	% "Ours:" group label before the three proposed variants, then Oracle.
	% 1.5pt swatches for methods, 1.0pt solid black for Oracle, \footnotesize.
	\centerline{%
		\begin{tikzpicture}[font=\footnotesize]
			\matrix [matrix of nodes, ampersand replacement=\&,
				column sep=0.3em, nodes={anchor=west, inner sep=1pt}] {
				\tikz[baseline=-0.6ex]{\draw[color=ColorRANSAC, line width=1.5pt] (0,0) -- (0.9em,0);} \& RANSAC \&
				\tikz[baseline=-0.6ex]{\draw[color=ColorMSAC, line width=1.5pt] (0,0) -- (0.9em,0);} \& MSAC \&
				\tikz[baseline=-0.6ex]{\draw[color=ColorGaU, line width=1.5pt] (0,0) -- (0.9em,0);} \& GaU \&
				\tikz[baseline=-0.6ex]{\draw[color=ColorMAGSAC, line width=1.5pt] (0,0) -- (0.9em,0);} \& MAGSAC++ \&
				\tikz[baseline=-0.6ex]{\draw[color=black, line width=1.0pt] (0,0) -- (0.9em,0);} \& Oracle \&
				\& \emph{Ours:} \&
				\tikz[baseline=-0.6ex]{\draw[color=ColorJeffreys, line width=1.5pt] (0,0) -- (0.9em,0);} \& Jeffreys \&
				\tikz[baseline=-0.6ex]{\draw[color=ColorIGRobust, line width=1.5pt] (0,0) -- (0.9em,0);} \& IG-Robust \&
				\tikz[baseline=-0.6ex]{\draw[color=ColorIGMLE, line width=1.5pt] (0,0) -- (0.9em,0);} \& IG-MLE \\
			};
		\end{tikzpicture}%
	}
	\caption{\textbf{Profile} across three datasets, one per
		geometric problem: PT~/~$\bE$ (left), ETH3D~/~$\bF$
		(middle), HEB~/~$\bH$ (right).  Each panel plots, as
		solid lines, the median rotation error across all pairs
		in the dataset versus each method's hyperparameter on a
		log axis, swept exhaustively over a two-decade interval.
		Color-matched dashed horizontal lines mark each method's
		achievable performance when its hyperparameter is free
		to be tuned per pair to select the model closest to the
		ground truth.  Shaded bands are 95\% bootstrap confidence
		intervals on the median.
		\textbf{In the data-rich regime (PT, ETH3D) the proposed
			methods are remarkably robust to hyperparameter
			miscalibration}, with curves nearly flat across the full
		two-decade window where threshold-based baselines degrade
		markedly.  On HEB---the data-scarce
		regime---\textbf{the IG-Robust prior regularizes the
			score and outperforms the baselines over most of the
			hyperparameter interval}
		(cf.~\Cref{par:pseudo-obs}).}
	\label{fig:profile-grid}
\end{figure}

Each method is swept on its own native hyperparameter over a
two-decade search interval centered on its anchor~$\kappa^\star$
(\Cref{sec:eval-protocol}); the panels of
\Cref{fig:profile-grid} plot median error against the ratio
$\kappa/\kappa^\star$, so the same horizontal decade
represents comparable parameter neighbourhoods across
methods.  Shaded bands are bootstrapped 95\% confidence
intervals on the median, computed from $1\,000$
pair-resamples (with replacement) of the dataset; because
validation and test share the same pool, these bands measure
the precision of the fixed-pool median, not a generalization
gap (the latter is covered by the Sensitivity panels,
\Cref{sec:sensitivity-expt}).  The solid black line is the
GT-Oracle---the lowest pose error in each pair's candidate
model pool if error against the ground-truth model is used
for scoring, which represents the lower error bound for a
scoring function on the dataset.  Each
color-matched dashed horizontal line is the corresponding
method's per-pair-best envelope: the median error if its
hyperparameter were tuned per pair to pick the model closest
to ground truth, i.e., the method's achievable performance
under oracle per-pair hyperparameter.

\subsubsection{Results}
In the data-rich regime (PhotoTourism and ETH3D), the proposed
methods are remarkably robust to miscalibration of their
hyperparameter---the outlier half-width $a$: the median-error
curves are nearly flat across the full two-decade window, in
contrast to the threshold-based baselines, whose error grows
visibly as the hyperparameter is perturbed away from its
optimum.  The three IG variants (Jeffreys, IG-MLE, IG-Robust)
overlap throughout the window---varying the prior does not change
performance.  Here the conjugate update~\eqref{eq:ig-posterior} reads the
informative prior as $2\alpha_0$ pseudo-residuals carrying
$2\beta_0$ of pseudo-RSS.  Recall from \Cref{par:pseudo-obs} that
$\alpha_0$ is the prior's strength of belief about the noise
scale (in pseudo-observations): it controls how much the
prior pulls, while $\beta_0/\alpha_0$ controls toward what
scale.
IG-Robust fixes $\alpha_0 = 1$ (2 pseudo-residuals); IG-MLE
fits $\alpha_0$ from the validation $\hat\sigma^2_p$ set and
on these datasets it likewise lands at an $\alpha_0$ small
relative to $n_p$.  PT and ETH3D per-pair inlier
counts are sufficiently larger than $2\alpha_0$ that the data
term in the posterior outweighs the prior contribution, and
the informative variants approach the non-informative
Jeffreys limit.  Note that a large per-pair inlier count $n_p$ does
not imply a large inlier \emph{ratio}: PT in particular has
$n_p \approx 140$ but a low ratio.  The ratio is a separate
axis of difficulty---it controls RANSAC's expected iteration
count, which grows as $\epsilon^{-k}$ in the inlier ratio
$\epsilon$ and minimal-sample size $k$.

\paragraph{Convergence at the optimum.}
The proposed methods and the threshold-based baselines reach
the \emph{same} ceiling on every dataset, despite using
different inlier distributions: the marginalized model has a
Student-$t$ inlier (from the Gaussian + IG hierarchy after
integrating $\sigma^2$ analytically), while the
threshold-based baselines have a Gaussian inlier at a
globally pinned $\sigma$.  Two effects
collapse the gap as $n_p$ grows.  First, the per-partition
posterior on $\sigma^2$,
$\mathrm{IG}\!\bigl(\alpha_0 + \tfrac{\nu_p}{2},\,
	\beta_0 + \tfrac{\mathrm{RSS}_p}{2}\bigr)$, has shape
$\alpha_0 + \nu_p/2 \to \infty$, so the Student-$t$ inlier has
$2\alpha_0 + \nu_p$ degrees of freedom and converges to a
Gaussian centered at the per-partition scale estimate
$\hat s_p$.  Second, $\hat s_p$ has $\sqrt{2/\nu_p}$ relative
standard error across pairs, so it concentrates around a
population-typical value.  An optimally-tuned $\kappa$ pins
$\sigma$ globally to that same population-typical value, and
the model selected per pair coincides with what the marginal
score selects.  What the threshold-based baselines
\emph{cannot} reproduce, however, is the proposed methods'
robustness to miscalibration \emph{in the data-rich regime}:
on PhotoTourism and ETH3D the proposed curves stay essentially
flat across the full two-decade hyperparameter window---varying
by ${<}0.02^\circ$ on PhotoTourism---while the threshold-based
baselines degrade under the same perturbation, most sharply for
vanilla RANSAC and on PhotoTourism, where the baselines span
${\sim}0.8^\circ$--$2.9^\circ$ over the window---a property no
choice of fixed $\sigma$ can recover.  In the low-cardinality
HEB regime the prior no longer recedes and this flatness does
not hold; that regime is governed instead by the
adaptive-regularization behavior analyzed next.

\paragraph{Low-cardinality regime.}
HEB stresses a second axis of difficulty: in addition to a
low inlier ratio (which controls RANSAC's iteration count),
it has a low cardinality of inlying matches per pair, which
controls the statistical reliability of any per-partition
scale estimate.  Mechanically, its annotations are per-plane
rather than per-pair, so a typical pair carries hundreds of
geometrically valid correspondences but only those on the
annotated dominant plane (median $n_p \approx 10$, a
${\sim}4\%$ ratio) are GT inliers.  The remaining ${\sim}96\%$ are \emph{not}
mismatches; rather they are correct matches from competing
surfaces, but from the perspective of a particular dominant
plane they appear as coherently structured outliers, which
violate the uniform-outlier assumption shared by the
marginal-scoring and threshold-based families.  Imperfect assignment of spatially-verified point matches to
dominant planes injects further label noise, raising the
floor of the per-pair noise-variance estimate
$\hat\sigma^2_p$ (the variance of inlier residuals under the
oracle pose, used to anchor the IG prior).  In this context, IG-Robust pulls away from the baselines and
the other proposed methods over most of the two-decade
search region.  With $\alpha_0 = 1$ contributing two
pseudo-residuals and a median-anchored $\beta_0$, the prior
regularizes per-partition scale estimates when $n_p$ is
comparable to the model's degrees of freedom.  The median
anchor is itself insensitive to the inflated $\hat\sigma^2_p$
tails produced by structured-outlier contamination and by
annotation errors in the dominant-plane labels.  MLE regression of the IG parameters is not robust to these
corruptions, which negates IG-MLE's regularization advantage.
Jeffreys gives up the regularization entirely, and with only
$n_p \approx 10$ inlying measurements per pair its
per-partition scale estimate is too noisy to discriminate
models reliably.  Threshold-based methods cannot adapt at
all, since a single global $\sigma$ has no per-partition
correction.

\paragraph{Oracle-based scoring.}
The per-pair-best envelope (dashed line) reports each method's
achievable error if its hyperparameter were retuned on every
pair to select the model closest to the ground truth.  These
envelopes invert the ordering of the solid curves' minima:
RANSAC gives the best lower bound, threshold-based baselines
next, the proposed methods highest.
This is not a measure of attainable performance but an
artifact of how each family's hyperparameter interacts with
the per-pair argmax.  For threshold-based methods the
hyperparameter is a \emph{model-selection} knob: different
thresholds induce different inlier sets and select different
models, so a per-pair sweep explores a wide range of
selections.  For our IG marginal, $a$ enters only the
outlier-penalty term $-(N - n_p)\,d_g\log(2a)$; in the data-rich
regime the sort-and-sweep picks an inlier partition at the
natural residual gap, largely insensitive to $a$, so sweeping
$a$ shifts the score's value but not which model is selected,
and the IG per-pair-best collapses to a single fixed model.
This decoupling is the desired property: $a$ calibrates
score values while the selected model is governed by the
data, not the hyperparameter.  The threshold-based baselines'
lower dashed envelope is therefore not a deployable
benefit---it is the visible footprint of the same dual-role
mechanism that makes their solid curves so sensitive to
hyperparameter perturbation.  In deployment one
hyperparameter is used per method, so the relevant
comparison is between the solid-curve minima; the dashed
envelope is a per-pair oracle no real scoring function has
access to.

% ---------------------------------------------------------------------
\subsection{Sensitivity}
\label{sec:sensitivity-expt}

A RANSAC user typically has only a few ground-truth examples
available for tuning and selecting their method's
hyperparameter.  Small validation sets do not fully capture
the test distribution, so one scoring method may appear
better than another by chance, depending on which validation
pairs were drawn~\cite{Shekhovtsov2025}.  We measure
\emph{expected} test performance by randomly drawing small
validation sets from the data.

We compute this expected test performance via bootstrap.  For
each image pair we precompute the rotation error of every
method at every grid point in the two-decade hyperparameter
window from \Cref{sec:eval-protocol}: $1\,000$ minimal-sample
models are drawn per pair, scored at every hyperparameter
$\kappa$, the best-scoring model is selected at each $\kappa$,
and its error against the GT pose is recorded.  With these errors
cached, the bootstrap loop is fast: draw a validation pool
$\mathcal{V}$ of cardinality $|\mathcal{V}|$ (on HEB, sampled
from a random subset of scenes held out for validation, with the
disjoint remaining scenes forming the test pool, so trial
variance reflects cross-scene generalization; split at the pair
level on PhotoTourism and ETH3D) and pick the
validation-optimal hyperparameter
$\kappa^\star(\mathcal{V}) = \arg\min_\kappa
	\mathrm{median}_{p \in \mathcal{V}}\, e_p(\kappa)$, and
evaluate $e^{\mathrm{test}}(\kappa^\star(\mathcal{V})) =
	\mathrm{median}_{p \in \mathcal{T}}\, e_p(\kappa^\star(\mathcal{V}))$
on the held-out test pool $\mathcal{T}$.  Repeating $1\,000$
times gives the expected median test error
$\mathbb{E}[e_{\mathbf{R}}]$ at each validation budget
$|\mathcal{V}| \in \{2, 4, 8, \ldots, 1024\}$ (extended to
$2048$ on HEB).

\begin{figure}[t]
	\centering
	% paper_styles_shared.tex — colour palette + Sensitivity per-method
% named styles.  Identical across cells/experiments; \input once
% into your paper preamble (or copy verbatim).  Override per-method
% via `\pgfplotsset{ours_<m>/.style={...}}` in your preamble.

% --- Colour palette (\providecolor: paper-side override wins) ---
\providecolor{ColorGaU}{HTML}{42D4F4}
\providecolor{ColorIGMLE}{HTML}{4363D8}
\providecolor{ColorIGRobust}{HTML}{3CB44B}
\providecolor{ColorJeffreys}{HTML}{E6194B}
\providecolor{ColorMAGSAC}{HTML}{F58231}
\providecolor{ColorMSAC}{HTML}{FFE119}
\providecolor{ColorOracle}{HTML}{808080}
\providecolor{ColorRANSAC}{HTML}{FFB6C1}
\providecolor{ColorgauPerPairBest}{HTML}{42D4F4}
\providecolor{ColormagsacPerPairBest}{HTML}{F58231}
\providecolor{ColormsacPerPairBest}{HTML}{FFE119}
\providecolor{ColorransacPerPairBest}{HTML}{FFB6C1}

% --- Per-method addplot styles ---
\pgfplotsset{ours_ransac/.style={color=ColorRANSAC, line width=0.8pt}}
\pgfplotsset{ours_ransac_marker/.style={ours_ransac, mark=x, mark size=2pt, mark options={solid}}}
\pgfplotsset{ours_msac/.style={color=ColorMSAC, line width=0.8pt}}
\pgfplotsset{ours_msac_marker/.style={ours_msac, mark=diamond*, mark size=2pt, mark options={solid}}}
\pgfplotsset{ours_gau/.style={color=ColorGaU, line width=0.8pt}}
\pgfplotsset{ours_gau_marker/.style={ours_gau, mark=triangle*, mark size=2pt, mark options={solid}}}
\pgfplotsset{ours_magsac/.style={color=ColorMAGSAC, line width=0.8pt}}
\pgfplotsset{ours_magsac_marker/.style={ours_magsac, mark=star, mark size=2pt, mark options={solid}}}
\pgfplotsset{ours_jeffreys/.style={color=ColorJeffreys, line width=0.8pt}}
\pgfplotsset{ours_jeffreys_marker/.style={ours_jeffreys, mark=pentagon, mark size=2pt, mark options={solid}}}
\pgfplotsset{ours_ig_robust/.style={color=ColorIGRobust, line width=0.8pt}}
\pgfplotsset{ours_ig_robust_marker/.style={ours_ig_robust, mark=o, mark size=2pt, mark options={solid}}}
\pgfplotsset{ours_ig_mle/.style={color=ColorIGMLE, line width=0.8pt}}
\pgfplotsset{ours_ig_mle_marker/.style={ours_ig_mle, mark=square, mark size=2pt, mark options={solid}}}
%
	% paper_styles.tex — cell-specific defaults: \Panel* macros + axis cosmetics.
% Standalone preview: \input{paper_styles_shared}\input{paper_styles}
% BEFORE the panel \input.  In main.tex, \input{paper_styles_shared}
% globally, omit this file, and \renewcommand the \Panel* macros
% per panel-input as needed (use `{}` to suppress a label).

\providecommand{\PanelXLabel}{Validation set size}
\providecommand{\PanelYLabeleRmean}{$\mathbb{E}[e_{\mathbf{R}}]$ (deg)}
\providecommand{\PanelYLabeleRstd}{$\mathrm{std}[e_{\mathbf{R}}]$ (deg)}
\providecommand{\PanelYLabelmAAmean}{$\mathbb{E}[\mathrm{mAA}]$}
\providecommand{\PanelYLabelmAAstd}{$\mathrm{std}[\mathrm{mAA}]$}
\providecommand{\PanelXTicks}{2,4,8,16,32,64,128,256,512,1024,2048}

% Named axis style — panels opt-in via `panel_axis` in their
% axis options.  Legends/insets/color-bars stay clean.
% NOTE: visual properties only — NO width/height.  Layout (size,
% subfigure dimensions, groupplot geometry) is the paper's call;
% baking width/height here would silently override paper-side
% sizing via `every axis/.append style{}` last-wins.  Standalone
% preview supplies the default size in `preview.tex`.
\pgfplotsset{panel_axis/.style={
    grid style={gray!20},
    label style={font=\footnotesize},
    tick label style={font=\scriptsize},
}}
%
	\pgfplotsset{panel_axis/.append style={
				label style={font=\scriptsize},
				title style={font=\scriptsize},
				tick label style={font=\tiny}}}%
	% Strip the gen's `fill=Color<m>` from mark options so the "Ours"
	% markers stay open (regardless of which open shape the gen chooses).
	\pgfplotsset{ours_jeffreys_marker/.append style={mark options={solid}}}%
	\pgfplotsset{ours_ig_robust_marker/.append style={mark options={solid}}}%
	\pgfplotsset{ours_ig_mle_marker/.append style={mark options={solid}}}%
	\begin{tikzpicture}
		\begin{groupplot}[panel_axis,
				width=0.40\linewidth, height=1.9cm, scale only axis,
				group style={group size=2 by 3, horizontal sep=1.4cm,
						vertical sep=1.2cm, xlabels at=edge bottom},
				xmode=log, log basis x=2,
				xmajorgrids,
				xtick={4,16,64,256,1024},
				xlabel={Validation set size},
				scaled y ticks=false,
				yticklabel style={/pgf/number format/.cd, fixed, precision=1},
				every axis y label/.style={
						at={(axis description cs:-0.13,0.5)},
						anchor=center, rotate=90, font=\scriptsize}]
			\let\origNGP\nextgroupplot%
			% Row 1 (PT)
			\renewcommand\nextgroupplot[1][]{\origNGP[#1, ylabel={$\mathbb{E}[e_{\mathbf{R}}]$ (deg)}, ymax=1, ytick={0.8,0.9,1}]}%
			% panel ..._group.tex — \nextgroupplot fragment.
% Host preamble MUST load: \usepgfplotslibrary{groupplots},
%                          \usepgfplotslibrary{fillbetween}.
% Wrap N such fragments inside one \begin{groupplot}[...] that
% sets `panel_axis, xmode=log, log basis x=2, xmajorgrids,
% ymajorgrids, xlabel={…}, ylabel={…}, xtick={…},
% ytick distance=…, yminorgrids, scaled y ticks=false,
% yticklabel style={…}, group style={group size=N by M, …}].

\nextgroupplot[title={PhotoTourism / RootSIFT / \(\bE\)}]
    \addplot[name path={UeR_mean_median_1}, draw={none}, mark={none}, forget plot]
        coordinates {(2.0,1.1070490281341718) (4.0,1.0333906284995615) (8.0,0.9838754576249147) (16.0,0.9363953377466808) (32.0,0.913857932959856) (64.0,0.8965922488910179) (128.0,0.8850571147015021) (256.0,0.8774553786645836) (512.0,0.8711349610293699) (1024.0,0.8676867720880008)}
        ;
    \addplot[name path={LeR_mean_median_1}, draw={none}, mark={none}, forget plot]
        coordinates {(2.0,1.0600828982564128) (4.0,0.9999801998696647) (8.0,0.9605340938024355) (16.0,0.9221798160256827) (32.0,0.9039839942706194) (64.0,0.8901178890939874) (128.0,0.8803213286431087) (256.0,0.8744119448209509) (512.0,0.8691302755348331) (1024.0,0.8662064216343628)}
        ;
    \addplot[fill opacity={0.18}, draw={none}, mark={none}, forget plot, ours_ransac, fill]
        fill between [of=UeR_mean_median_1 and LeR_mean_median_1]
        ;
    \addplot[ours_ransac_marker]
        coordinates {(2.0,1.0813993292025077) (4.0,1.0150192212234272) (8.0,0.9711206266464197) (16.0,0.9285177552362548) (32.0,0.9086291658762065) (64.0,0.8932175617101635) (128.0,0.8825310171292678) (256.0,0.8758437993384915) (512.0,0.8700824721053934) (1024.0,0.8669128632910591)}
        ;
    \addplot[name path={UeR_mean_median_2}, draw={none}, mark={none}, forget plot]
        coordinates {(2.0,0.9083698301727812) (4.0,0.8861190831685596) (8.0,0.8618979231267565) (16.0,0.8382041370170974) (32.0,0.8163843044249659) (64.0,0.8027115118009892) (128.0,0.789100921889138) (256.0,0.7810002864410525) (512.0,0.7782305176662303) (1024.0,0.7751101342903364)}
        ;
    \addplot[name path={LeR_mean_median_2}, draw={none}, mark={none}, forget plot]
        coordinates {(2.0,0.889515532430158) (4.0,0.8689783473430979) (8.0,0.8476818271929922) (16.0,0.8280771935514312) (32.0,0.809668140061611) (64.0,0.7977588681541807) (128.0,0.7860314431140772) (256.0,0.7793211755252165) (512.0,0.7770226541543002) (1024.0,0.7744549639669179)}
        ;
    \addplot[fill opacity={0.18}, draw={none}, mark={none}, forget plot, ours_msac, fill]
        fill between [of=UeR_mean_median_2 and LeR_mean_median_2]
        ;
    \addplot[ours_msac_marker]
        coordinates {(2.0,0.8982403891421761) (4.0,0.8768294455999059) (8.0,0.8544007567852029) (16.0,0.8326459593456753) (32.0,0.8129497814548909) (64.0,0.8001176659366652) (128.0,0.7875122413328902) (256.0,0.7801196955972585) (512.0,0.7775859776426637) (1024.0,0.7747661875137731)}
        ;
    \addplot[name path={UeR_mean_median_3}, draw={none}, mark={none}, forget plot]
        coordinates {(2.0,0.9101229188927784) (4.0,0.8861551460271839) (8.0,0.8687148043057135) (16.0,0.8380145266577256) (32.0,0.8171195544782448) (64.0,0.7979023572226711) (128.0,0.7851163806889946) (256.0,0.7781867799680118) (512.0,0.7743624108993951) (1024.0,0.7718601344172893)}
        ;
    \addplot[name path={LeR_mean_median_3}, draw={none}, mark={none}, forget plot]
        coordinates {(2.0,0.8931502794573792) (4.0,0.8708455715905774) (8.0,0.8543386510263542) (16.0,0.8274202676780508) (32.0,0.8091197315730131) (64.0,0.792895601838139) (128.0,0.7821536184608238) (256.0,0.7764691337874172) (512.0,0.7732320196548996) (1024.0,0.7711366457085841)}
        ;
    \addplot[fill opacity={0.18}, draw={none}, mark={none}, forget plot, ours_gau, fill]
        fill between [of=UeR_mean_median_3 and LeR_mean_median_3]
        ;
    \addplot[ours_gau_marker]
        coordinates {(2.0,0.9016604238077025) (4.0,0.8782418832836449) (8.0,0.8613619400376817) (16.0,0.8324390268538071) (32.0,0.8129132341631401) (64.0,0.7952023668878584) (128.0,0.7835519408928203) (256.0,0.7772712930428742) (512.0,0.7737551235787647) (1024.0,0.77147676420646)}
        ;
    \addplot[name path={UeR_mean_median_4}, draw={none}, mark={none}, forget plot]
        coordinates {(2.0,0.91506237189572) (4.0,0.8828742037188014) (8.0,0.8603615970927316) (16.0,0.834340654690511) (32.0,0.8125728497651983) (64.0,0.7969299441504197) (128.0,0.7836019440664547) (256.0,0.7764824966756939) (512.0,0.7732382352347673) (1024.0,0.7705784277546504)}
        ;
    \addplot[name path={LeR_mean_median_4}, draw={none}, mark={none}, forget plot]
        coordinates {(2.0,0.8944844861261633) (4.0,0.8665387148784836) (8.0,0.8462291439557804) (16.0,0.8242408999735746) (32.0,0.8056013810433793) (64.0,0.7920279490292966) (128.0,0.7806928475648756) (256.0,0.7749524488903998) (512.0,0.772131648281894) (1024.0,0.7699064338523274)}
        ;
    \addplot[fill opacity={0.18}, draw={none}, mark={none}, forget plot, ours_magsac, fill]
        fill between [of=UeR_mean_median_4 and LeR_mean_median_4]
        ;
    \addplot[ours_magsac_marker]
        coordinates {(2.0,0.9041416810741899) (4.0,0.8743673468412441) (8.0,0.8527583038993093) (16.0,0.828904461176687) (32.0,0.8089285571893952) (64.0,0.7942646202425553) (128.0,0.7820313480847291) (256.0,0.7756588251648093) (512.0,0.7726541141699813) (1024.0,0.7702268461161167)}
        ;
    \addplot[name path={UeR_mean_median_5}, draw={none}, mark={none}, forget plot]
        coordinates {(2.0,0.7755789457487948) (4.0,0.7753362668976919) (8.0,0.7751698351556234) (16.0,0.774919813836707) (32.0,0.7745045425367579) (64.0,0.774162432758541) (128.0,0.7737983192412816) (256.0,0.7734898540416159) (512.0,0.7731804606585923) (1024.0,0.7728240926674161)}
        ;
    \addplot[name path={LeR_mean_median_5}, draw={none}, mark={none}, forget plot]
        coordinates {(2.0,0.7753525166744001) (4.0,0.7750949619422975) (8.0,0.7749228190705739) (16.0,0.7746674024965099) (32.0,0.7742520372160219) (64.0,0.7739172888609865) (128.0,0.7735616515650692) (256.0,0.7732790093142596) (512.0,0.7729772384236333) (1024.0,0.7726433369665925)}
        ;
    \addplot[fill opacity={0.18}, draw={none}, mark={none}, forget plot, ours_jeffreys, fill]
        fill between [of=UeR_mean_median_5 and LeR_mean_median_5]
        ;
    \addplot[ours_jeffreys_marker]
        coordinates {(2.0,0.7754717946789311) (4.0,0.775216278682111) (8.0,0.7750478296939765) (16.0,0.7747978107647243) (32.0,0.7743817788064894) (64.0,0.7740373258156114) (128.0,0.7736753700144668) (256.0,0.7733827830206872) (512.0,0.7730780297983018) (1024.0,0.7727298998007993)}
        ;
    \addplot[name path={UeR_mean_median_6}, draw={none}, mark={none}, forget plot]
        coordinates {(2.0,0.7792310998508729) (4.0,0.7787802791145997) (8.0,0.778388151093989) (16.0,0.7779107979086894) (32.0,0.7769927795702545) (64.0,0.7761239012347011) (128.0,0.7754734240505493) (256.0,0.7747516773164581) (512.0,0.7743511950458859) (1024.0,0.7736006559072777)}
        ;
    \addplot[name path={LeR_mean_median_6}, draw={none}, mark={none}, forget plot]
        coordinates {(2.0,0.7787242236468057) (4.0,0.7782429570070118) (8.0,0.7778341555653917) (16.0,0.7773684748353454) (32.0,0.7764512079864243) (64.0,0.7756082188972739) (128.0,0.7749825406555043) (256.0,0.7743256152898296) (512.0,0.7739534100531675) (1024.0,0.7732801222853881)}
        ;
    \addplot[fill opacity={0.18}, draw={none}, mark={none}, forget plot, ours_ig_robust, fill]
        fill between [of=UeR_mean_median_6 and LeR_mean_median_6]
        ;
    \addplot[ours_ig_robust_marker]
        coordinates {(2.0,0.7789875395786501) (4.0,0.778514200166671) (8.0,0.7781112000894216) (16.0,0.77764172975574) (32.0,0.7767164335239425) (64.0,0.7758626011154242) (128.0,0.7752176462147151) (256.0,0.7745314704558963) (512.0,0.7741439081790176) (1024.0,0.7734308485735185)}
        ;
    \addplot[name path={UeR_mean_median_7}, draw={none}, mark={none}, forget plot]
        coordinates {(2.0,0.7757184796098235) (4.0,0.7766318989648024) (8.0,0.7769029416446185) (16.0,0.7766708750678075) (32.0,0.7760532186615717) (64.0,0.7753074631601454) (128.0,0.7749603308297818) (256.0,0.7744434300426967) (512.0,0.7739888437804419) (1024.0,0.7735524642752899)}
        ;
    \addplot[name path={LeR_mean_median_7}, draw={none}, mark={none}, forget plot]
        coordinates {(2.0,0.7751188149356902) (4.0,0.776223577489732) (8.0,0.7765011047241251) (16.0,0.7762538059635933) (32.0,0.7756435609541468) (64.0,0.7749223513871717) (128.0,0.7745834635305422) (256.0,0.7740916878700669) (512.0,0.7736754789370132) (1024.0,0.7732514061460328)}
        ;
    \addplot[fill opacity={0.18}, draw={none}, mark={none}, forget plot, ours_ig_mle, fill]
        fill between [of=UeR_mean_median_7 and LeR_mean_median_7]
        ;
    \addplot[ours_ig_mle_marker]
        coordinates {(2.0,0.7753794972074334) (4.0,0.7764288803309831) (8.0,0.7767023821479637) (16.0,0.7764649032508066) (32.0,0.7758491995002309) (64.0,0.7751134832853456) (128.0,0.7747639002455805) (256.0,0.7742667441358791) (512.0,0.7738275566254672) (1024.0,0.7733953051515597)}
        ;
%
			\renewcommand\nextgroupplot[1][]{\origNGP[#1, ylabel={$\mathrm{std}[e_{\mathbf{R}}]$ (deg)}, ymax=0.2, ytick={0,0.1,0.2}]}%
			% panel ..._group.tex — \nextgroupplot fragment.
% Host preamble MUST load: \usepgfplotslibrary{groupplots},
%                          \usepgfplotslibrary{fillbetween}.
% Wrap N such fragments inside one \begin{groupplot}[...] that
% sets `panel_axis, xmode=log, log basis x=2, xmajorgrids,
% ymajorgrids, xlabel={…}, ylabel={…}, xtick={…},
% ytick distance=…, yminorgrids, scaled y ticks=false,
% yticklabel style={…}, group style={group size=N by M, …}].

\nextgroupplot[title={PhotoTourism / RootSIFT / \(\bE\)}]
    \addplot[name path={UeR_std_median_1}, draw={none}, mark={none}, forget plot]
        coordinates {(2.0,0.4250638374411128) (4.0,0.31978745476549675) (8.0,0.23461355695999972) (16.0,0.13823164727981968) (32.0,0.09115484626368632) (64.0,0.058702206698013265) (128.0,0.04402394315364475) (256.0,0.027603033473517995) (512.0,0.0182809312828274) (1024.0,0.012897029106491199)}
        ;
    \addplot[name path={LeR_std_median_1}, draw={none}, mark={none}, forget plot]
        coordinates {(2.0,0.33256255322756145) (4.0,0.23014283876415764) (8.0,0.16188659165580027) (16.0,0.09924567056629226) (32.0,0.07063467622586025) (64.0,0.04685315515548139) (128.0,0.0340921228673042) (256.0,0.022980198573522296) (512.0,0.01524675852451468) (1024.0,0.01093006295169404)}
        ;
    \addplot[fill opacity={0.18}, draw={none}, mark={none}, forget plot, ours_ransac, fill]
        fill between [of=UeR_std_median_1 and LeR_std_median_1]
        ;
    \addplot[ours_ransac_marker]
        coordinates {(2.0,0.3761738336923212) (4.0,0.26808070872535605) (8.0,0.18677963787009103) (16.0,0.11315850758241369) (32.0,0.07880719056376961) (64.0,0.05211966677428214) (128.0,0.03836411500248642) (256.0,0.025007890795739885) (512.0,0.016493643414667237) (1024.0,0.01179157899240496)}
        ;
    \addplot[name path={UeR_std_median_2}, draw={none}, mark={none}, forget plot]
        coordinates {(2.0,0.17240119890180652) (4.0,0.16033155956529657) (8.0,0.13566417044797002) (16.0,0.09666751085245631) (32.0,0.05996911356207924) (64.0,0.044167028087133464) (128.0,0.02737971623917924) (256.0,0.015037034370578722) (512.0,0.01109380229809698) (1024.0,0.005886929199922818)}
        ;
    \addplot[name path={LeR_std_median_2}, draw={none}, mark={none}, forget plot]
        coordinates {(2.0,0.137548418301029) (4.0,0.11993786163317845) (8.0,0.10078815207234454) (16.0,0.07215880513260238) (32.0,0.05069633996039619) (64.0,0.03662527689305623) (128.0,0.022159406403962495) (256.0,0.01206884732657094) (512.0,0.00863635300270119) (1024.0,0.004886869648929523)}
        ;
    \addplot[fill opacity={0.18}, draw={none}, mark={none}, forget plot, ours_msac, fill]
        fill between [of=UeR_std_median_2 and LeR_std_median_2]
        ;
    \addplot[ours_msac_marker]
        coordinates {(2.0,0.15233230286164506) (4.0,0.13591659454612012) (8.0,0.11434723856139889) (16.0,0.08069366938487828) (32.0,0.054870035212452144) (64.0,0.03993649393554722) (128.0,0.024427164047111552) (256.0,0.013357223944634854) (512.0,0.009603234324099229) (1024.0,0.005321085555902054)}
        ;
    \addplot[name path={UeR_std_median_3}, draw={none}, mark={none}, forget plot]
        coordinates {(2.0,0.14408352048014322) (4.0,0.12957307425209638) (8.0,0.12296159243948176) (16.0,0.09210780903781093) (32.0,0.07172227424058307) (64.0,0.04603029838877345) (128.0,0.0268034241036161) (256.0,0.015231306703243252) (512.0,0.010183407458789416) (1024.0,0.006414304649186255)}
        ;
    \addplot[name path={LeR_std_median_3}, draw={none}, mark={none}, forget plot]
        coordinates {(2.0,0.13201631909656392) (4.0,0.11723315540550948) (8.0,0.1090226977349953) (16.0,0.07833128730920833) (32.0,0.05880018530624422) (64.0,0.03601256806654567) (128.0,0.021888128049026418) (256.0,0.012557653397568582) (512.0,0.008339390796973278) (1024.0,0.0052264764259547)}
        ;
    \addplot[fill opacity={0.18}, draw={none}, mark={none}, forget plot, ours_gau, fill]
        fill between [of=UeR_std_median_3 and LeR_std_median_3]
        ;
    \addplot[ours_gau_marker]
        coordinates {(2.0,0.1377034858771786) (4.0,0.1230696075766373) (8.0,0.1157188951078065) (16.0,0.08468075459213217) (32.0,0.06476034086017944) (64.0,0.04017337237281868) (128.0,0.024026548472893042) (256.0,0.013704664941209881) (512.0,0.009156012155508141) (1024.0,0.0057530957212257455)}
        ;
    \addplot[name path={UeR_std_median_4}, draw={none}, mark={none}, forget plot]
        coordinates {(2.0,0.18577095427111806) (4.0,0.1495272398989278) (8.0,0.12788564015816337) (16.0,0.09249615819193688) (32.0,0.06177413393373871) (64.0,0.044467483327758475) (128.0,0.0263220843863512) (256.0,0.014032987420886087) (512.0,0.009901560118246427) (1024.0,0.0059467479306292235)}
        ;
    \addplot[name path={LeR_std_median_4}, draw={none}, mark={none}, forget plot]
        coordinates {(2.0,0.15268181093144437) (4.0,0.11856270239691666) (8.0,0.10213142763659229) (16.0,0.07441337661326247) (32.0,0.05188900226596692) (64.0,0.03641836817845375) (128.0,0.021191426058024314) (256.0,0.011331767186801005) (512.0,0.008162891758652568) (1024.0,0.00491618931579838)}
        ;
    \addplot[fill opacity={0.18}, draw={none}, mark={none}, forget plot, ours_magsac, fill]
        fill between [of=UeR_std_median_4 and LeR_std_median_4]
        ;
    \addplot[ours_magsac_marker]
        coordinates {(2.0,0.16738033171379058) (4.0,0.13124297796574397) (8.0,0.11275847350286441) (16.0,0.08164501099469855) (32.0,0.056508503087928645) (64.0,0.039964152646690315) (128.0,0.02329073458363271) (256.0,0.012451869374108245) (512.0,0.008965197258415927) (1024.0,0.005345945912510181)}
        ;
    \addplot[name path={UeR_std_median_5}, draw={none}, mark={none}, forget plot]
        coordinates {(2.0,0.001944634814909155) (4.0,0.0020160554927484) (8.0,0.002035278322914633) (16.0,0.002093866415698901) (32.0,0.0020716055825515754) (64.0,0.002008403309756039) (128.0,0.0019389924652442163) (256.0,0.0017703140190946569) (512.0,0.001704103804931452) (1024.0,0.0015430419341425005)}
        ;
    \addplot[name path={LeR_std_median_5}, draw={none}, mark={none}, forget plot]
        coordinates {(2.0,0.0017907306201545054) (4.0,0.0018879957766723336) (8.0,0.0019212696587295366) (16.0,0.001993520341455437) (32.0,0.0019838902859917362) (64.0,0.0019137966402536233) (128.0,0.0018349807833766153) (256.0,0.001657885903680609) (512.0,0.0015795448397457345) (1024.0,0.0014010154635903868)}
        ;
    \addplot[fill opacity={0.18}, draw={none}, mark={none}, forget plot, ours_jeffreys, fill]
        fill between [of=UeR_std_median_5 and LeR_std_median_5]
        ;
    \addplot[ours_jeffreys_marker]
        coordinates {(2.0,0.0018681302856835506) (4.0,0.0019534525488411367) (8.0,0.0019785659756991892) (16.0,0.0020428105253539617) (32.0,0.0020279598531314333) (64.0,0.0019598991553775253) (128.0,0.0018867001927865127) (256.0,0.001713356678665478) (512.0,0.0016411800355272922) (1024.0,0.0014715378407413866)}
        ;
    \addplot[name path={UeR_std_median_6}, draw={none}, mark={none}, forget plot]
        coordinates {(2.0,0.004223271479999092) (4.0,0.004455044952897498) (8.0,0.004512790707705413) (16.0,0.004513450485681955) (32.0,0.004480916030310053) (64.0,0.004226442800329668) (128.0,0.004029486779513063) (256.0,0.0036084319758416398) (512.0,0.003350453260287942) (1024.0,0.0027354280205810037)}
        ;
    \addplot[name path={LeR_std_median_6}, draw={none}, mark={none}, forget plot]
        coordinates {(2.0,0.003995785741167552) (4.0,0.004264635960715984) (8.0,0.004356154429384332) (16.0,0.004375056595591633) (32.0,0.004323990933021929) (64.0,0.003999217211861323) (128.0,0.003762836686756834) (256.0,0.0032949309293558732) (512.0,0.0030135870267630777) (1024.0,0.0023630825739699743)}
        ;
    \addplot[fill opacity={0.18}, draw={none}, mark={none}, forget plot, ours_ig_robust, fill]
        fill between [of=UeR_std_median_6 and LeR_std_median_6]
        ;
    \addplot[ours_ig_robust_marker]
        coordinates {(2.0,0.004111559888377618) (4.0,0.004362474297606713) (8.0,0.004437207700889582) (16.0,0.004443426469168238) (32.0,0.004401591628106889) (64.0,0.004117143116799544) (128.0,0.003897640861358305) (256.0,0.003453228416396415) (512.0,0.003184366650049167) (1024.0,0.0025479059231767087)}
        ;
    \addplot[name path={UeR_std_median_7}, draw={none}, mark={none}, forget plot]
        coordinates {(2.0,0.006849730886987917) (4.0,0.003369032758860056) (8.0,0.0033139203603372237) (16.0,0.0034255217404411) (32.0,0.0033530487352610547) (64.0,0.0032007545702274235) (128.0,0.0031059622642684434) (256.0,0.002987351292108891) (512.0,0.002662149999886915) (1024.0,0.002587980132288756)}
        ;
    \addplot[name path={LeR_std_median_7}, draw={none}, mark={none}, forget plot]
        coordinates {(2.0,0.003980690170200885) (4.0,0.003193469020230689) (8.0,0.0031563450743166167) (16.0,0.003284058977257632) (32.0,0.0032152042687680093) (64.0,0.0030278759018286198) (128.0,0.002914671901620378) (256.0,0.002764702381265516) (512.0,0.0024157576740768388) (1024.0,0.002315421287838653)}
        ;
    \addplot[fill opacity={0.18}, draw={none}, mark={none}, forget plot, ours_ig_mle, fill]
        fill between [of=UeR_std_median_7 and LeR_std_median_7]
        ;
    \addplot[ours_ig_mle_marker]
        coordinates {(2.0,0.0047408977868862116) (4.0,0.0032775676437723734) (8.0,0.0032365188837514425) (16.0,0.0033538239993247823) (32.0,0.0032826096514307025) (64.0,0.0031155024661815873) (128.0,0.0030119061722728804) (256.0,0.0028781093023903755) (512.0,0.002539344948447794) (1024.0,0.0024506516513952783)}
        ;
%
			% Row 2 (ETH3D)
			\renewcommand\nextgroupplot[1][]{\origNGP[#1, ylabel={$\mathbb{E}[e_{\mathbf{R}}]$ (deg)}, ymax=0.6, ytick={0.2,0.4,0.5,0.6}]}%
			% panel ..._group.tex — \nextgroupplot fragment.
% Host preamble MUST load: \usepgfplotslibrary{groupplots},
%                          \usepgfplotslibrary{fillbetween}.
% Wrap N such fragments inside one \begin{groupplot}[...] that
% sets `panel_axis, xmode=log, log basis x=2, xmajorgrids,
% ymajorgrids, xlabel={…}, ylabel={…}, xtick={…},
% ytick distance=…, yminorgrids, scaled y ticks=false,
% yticklabel style={…}, group style={group size=N by M, …}].

\nextgroupplot[title={ETH3D / SuperPoint+LightGlue / \(\bF\)}]
    \addplot[name path={UeR_mean_median_1}, draw={none}, mark={none}, forget plot]
        coordinates {(2.0,0.7130894555740613) (4.0,0.6067577526119927) (8.0,0.5648510683138851) (16.0,0.5248907317308567) (32.0,0.5087611792876578) (64.0,0.4995474663701258) (128.0,0.49138953536495616) (256.0,0.48646742842203183) (512.0,0.4852318592312562) (1024.0,0.4846162011092281)}
        ;
    \addplot[name path={LeR_mean_median_1}, draw={none}, mark={none}, forget plot]
        coordinates {(2.0,0.669617173195409) (4.0,0.5817485862315438) (8.0,0.5475381800133896) (16.0,0.5178538046898489) (32.0,0.5041994952675836) (64.0,0.49638044653930047) (128.0,0.48961491672484114) (256.0,0.48544343009367724) (512.0,0.48440731264084547) (1024.0,0.4839314011652888)}
        ;
    \addplot[fill opacity={0.18}, draw={none}, mark={none}, forget plot, ours_ransac, fill]
        fill between [of=UeR_mean_median_1 and LeR_mean_median_1]
        ;
    \addplot[ours_ransac_marker]
        coordinates {(2.0,0.6900661338260969) (4.0,0.5931380906246605) (8.0,0.555065692166264) (16.0,0.5210515964529595) (32.0,0.5063105027605873) (64.0,0.4978512030615848) (128.0,0.49043960947147486) (256.0,0.48594748790837344) (512.0,0.4848257246910174) (1024.0,0.4842831758318618)}
        ;
    \addplot[name path={UeR_mean_median_2}, draw={none}, mark={none}, forget plot]
        coordinates {(2.0,0.528619668172554) (4.0,0.5192975021072611) (8.0,0.5112130209573171) (16.0,0.499413863448037) (32.0,0.4911845989097201) (64.0,0.4850735987078901) (128.0,0.4805068475205802) (256.0,0.4775822071664563) (512.0,0.4750694639429207) (1024.0,0.4712293855976524)}
        ;
    \addplot[name path={LeR_mean_median_2}, draw={none}, mark={none}, forget plot]
        coordinates {(2.0,0.5221715956159088) (4.0,0.5135311762405543) (8.0,0.5058675416889196) (16.0,0.4952974208443974) (32.0,0.48810842627335904) (64.0,0.48325415601664456) (128.0,0.47940464381404485) (256.0,0.4767904766055918) (512.0,0.4743443929509067) (1024.0,0.47068298827520616)}
        ;
    \addplot[fill opacity={0.18}, draw={none}, mark={none}, forget plot, ours_msac, fill]
        fill between [of=UeR_mean_median_2 and LeR_mean_median_2]
        ;
    \addplot[ours_msac_marker]
        coordinates {(2.0,0.5254778110598177) (4.0,0.5163769375198668) (8.0,0.5084273699228893) (16.0,0.4972411766720832) (32.0,0.48955037352051795) (64.0,0.4841047509867784) (128.0,0.479943040086518) (256.0,0.477164933981659) (512.0,0.4747061995768128) (1024.0,0.47096311114554684)}
        ;
    \addplot[name path={UeR_mean_median_3}, draw={none}, mark={none}, forget plot]
        coordinates {(2.0,0.5207676211957863) (4.0,0.5137949083992311) (8.0,0.5089535272151258) (16.0,0.49915073131075394) (32.0,0.49275388848896123) (64.0,0.4882325349345764) (128.0,0.48288163479640506) (256.0,0.48042868122274207) (512.0,0.47892317043626864) (1024.0,0.4779215025032322)}
        ;
    \addplot[name path={LeR_mean_median_3}, draw={none}, mark={none}, forget plot]
        coordinates {(2.0,0.5158368817684594) (4.0,0.5092332543540438) (8.0,0.5047070834642982) (16.0,0.49572919195625437) (32.0,0.4901787845904408) (64.0,0.4864069554379722) (128.0,0.48189718693373174) (256.0,0.47983063765910794) (512.0,0.4785657786583545) (1024.0,0.4777330679005084)}
        ;
    \addplot[fill opacity={0.18}, draw={none}, mark={none}, forget plot, ours_gau, fill]
        fill between [of=UeR_mean_median_3 and LeR_mean_median_3]
        ;
    \addplot[ours_gau_marker]
        coordinates {(2.0,0.5182777321119925) (4.0,0.511489468080361) (8.0,0.5068120345934449) (16.0,0.4973834772617429) (32.0,0.4914147591056085) (64.0,0.4873001372927456) (128.0,0.4823580193412153) (256.0,0.48010369270254116) (512.0,0.4787348076122779) (1024.0,0.4778234288131388)}
        ;
    \addplot[name path={UeR_mean_median_4}, draw={none}, mark={none}, forget plot]
        coordinates {(2.0,0.5290907858684297) (4.0,0.5199348143411169) (8.0,0.5131353018179798) (16.0,0.5026426710331433) (32.0,0.49421553886122166) (64.0,0.48826790552126387) (128.0,0.48309068664360766) (256.0,0.4808303195501493) (512.0,0.47956325067987704) (1024.0,0.47894297260788604)}
        ;
    \addplot[name path={LeR_mean_median_4}, draw={none}, mark={none}, forget plot]
        coordinates {(2.0,0.5231983597126088) (4.0,0.5145479398876395) (8.0,0.5080049177946917) (16.0,0.49836763636915504) (32.0,0.49114506627687315) (64.0,0.4862906083518573) (128.0,0.4820479053384042) (256.0,0.4802107966445609) (512.0,0.4791923418744472) (1024.0,0.4788128533386959)}
        ;
    \addplot[fill opacity={0.18}, draw={none}, mark={none}, forget plot, ours_magsac, fill]
        fill between [of=UeR_mean_median_4 and LeR_mean_median_4]
        ;
    \addplot[ours_magsac_marker]
        coordinates {(2.0,0.5261543766103297) (4.0,0.5172657969084432) (8.0,0.5104978382312595) (16.0,0.5004057163808373) (32.0,0.4925586463605091) (64.0,0.48720853651769547) (128.0,0.48252790541024115) (256.0,0.4804912285223536) (512.0,0.4793658469268019) (1024.0,0.478877660125436)}
        ;
    \addplot[name path={UeR_mean_median_5}, draw={none}, mark={none}, forget plot]
        coordinates {(2.0,0.4978093145421719) (4.0,0.49584903606338293) (8.0,0.49392582289187226) (16.0,0.48904200442337437) (32.0,0.4860830570816155) (64.0,0.4826266997187777) (128.0,0.47976875048342094) (256.0,0.4773637417215331) (512.0,0.4758013785151908) (1024.0,0.4750444437096246)}
        ;
    \addplot[name path={LeR_mean_median_5}, draw={none}, mark={none}, forget plot]
        coordinates {(2.0,0.49549715877312456) (4.0,0.4935769353680035) (8.0,0.49177309348327936) (16.0,0.4871012636706825) (32.0,0.48438586316170523) (64.0,0.4812952604820183) (128.0,0.47879062484692075) (256.0,0.47673590305879915) (512.0,0.47539287318057505) (1024.0,0.4747109954016536)}
        ;
    \addplot[fill opacity={0.18}, draw={none}, mark={none}, forget plot, ours_jeffreys, fill]
        fill between [of=UeR_mean_median_5 and LeR_mean_median_5]
        ;
    \addplot[ours_jeffreys_marker]
        coordinates {(2.0,0.49667773540539795) (4.0,0.49472784927315566) (8.0,0.49283588907653725) (16.0,0.4880455181004976) (32.0,0.4851924914017542) (64.0,0.4819299302566645) (128.0,0.4792479651162912) (256.0,0.47703852295259885) (512.0,0.47559095701670956) (1024.0,0.474872269997515)}
        ;
    \addplot[name path={UeR_mean_median_6}, draw={none}, mark={none}, forget plot]
        coordinates {(2.0,0.49695563713761576) (4.0,0.4948476395503463) (8.0,0.4922631042041021) (16.0,0.4879827198200142) (32.0,0.48535136731653966) (64.0,0.48289935534419415) (128.0,0.47984972083829447) (256.0,0.4774939437881324) (512.0,0.47604907985421835) (1024.0,0.4748685458255201)}
        ;
    \addplot[name path={LeR_mean_median_6}, draw={none}, mark={none}, forget plot]
        coordinates {(2.0,0.4942008831602435) (4.0,0.4925904865738521) (8.0,0.4903428566812287) (16.0,0.4863054898382214) (32.0,0.4838571849631848) (64.0,0.48167502401324713) (128.0,0.4789873243584597) (256.0,0.47691582550718803) (512.0,0.4756544052004063) (1024.0,0.47461798762586416)}
        ;
    \addplot[fill opacity={0.18}, draw={none}, mark={none}, forget plot, ours_ig_robust, fill]
        fill between [of=UeR_mean_median_6 and LeR_mean_median_6]
        ;
    \addplot[ours_ig_robust_marker]
        coordinates {(2.0,0.49556062414734203) (4.0,0.4937133329757681) (8.0,0.49128141299565603) (16.0,0.48712155894724185) (32.0,0.4845703658898578) (64.0,0.4822575626649993) (128.0,0.47939525814921885) (256.0,0.47719086969112273) (512.0,0.47584662241169523) (1024.0,0.4747366225700393)}
        ;
    \addplot[name path={UeR_mean_median_7}, draw={none}, mark={none}, forget plot]
        coordinates {(2.0,0.49696314483357407) (4.0,0.4946342949255977) (8.0,0.4913273277694342) (16.0,0.48752451254804274) (32.0,0.4854252763443529) (64.0,0.48274777393591783) (128.0,0.4798683774809821) (256.0,0.47756657675810876) (512.0,0.4761219522709746) (1024.0,0.47500477886742165)}
        ;
    \addplot[name path={LeR_mean_median_7}, draw={none}, mark={none}, forget plot]
        coordinates {(2.0,0.4917295919948339) (4.0,0.49128474391047783) (8.0,0.4892757020994776) (16.0,0.4858831758706056) (32.0,0.483920710592003) (64.0,0.48152875945158) (128.0,0.47900557589026677) (256.0,0.4769568117485662) (512.0,0.47572716288843353) (1024.0,0.47475248670085496)}
        ;
    \addplot[fill opacity={0.18}, draw={none}, mark={none}, forget plot, ours_ig_mle, fill]
        fill between [of=UeR_mean_median_7 and LeR_mean_median_7]
        ;
    \addplot[ours_ig_mle_marker]
        coordinates {(2.0,0.49440757957808634) (4.0,0.4929555029871284) (8.0,0.4903134777829979) (16.0,0.4866677833955305) (32.0,0.4846450920373491) (64.0,0.48211432571686264) (128.0,0.4794150951490465) (256.0,0.4772423310394906) (512.0,0.4759179194917093) (1024.0,0.47487406692219986)}
        ;
%
			\renewcommand\nextgroupplot[1][]{\origNGP[#1, ylabel={$\mathrm{std}[e_{\mathbf{R}}]$ (deg)}, ymax=0.2, ytick={0,0.1,0.2}]}%
			% panel ..._group.tex — \nextgroupplot fragment.
% Host preamble MUST load: \usepgfplotslibrary{groupplots},
%                          \usepgfplotslibrary{fillbetween}.
% Wrap N such fragments inside one \begin{groupplot}[...] that
% sets `panel_axis, xmode=log, log basis x=2, xmajorgrids,
% ymajorgrids, xlabel={…}, ylabel={…}, xtick={…},
% ytick distance=…, yminorgrids, scaled y ticks=false,
% yticklabel style={…}, group style={group size=N by M, …}].

\nextgroupplot[title={ETH3D / SuperPoint+LightGlue / \(\bF\)}]
    \addplot[name path={UeR_std_median_1}, draw={none}, mark={none}, forget plot]
        coordinates {(2.0,0.3803221144234912) (4.0,0.23411401237169796) (8.0,0.17561388716043558) (16.0,0.06646171181913499) (32.0,0.04246445834616308) (64.0,0.028028154394594294) (128.0,0.015949367864838124) (256.0,0.00909957792512571) (512.0,0.007294423602649588) (1024.0,0.005781524129761898)}
        ;
    \addplot[name path={LeR_std_median_1}, draw={none}, mark={none}, forget plot]
        coordinates {(2.0,0.31690870477591) (4.0,0.17337645933701276) (8.0,0.11610521443240886) (16.0,0.050904160845066035) (32.0,0.03352021035412845) (64.0,0.022985304487840686) (128.0,0.012920130471155294) (256.0,0.007613166819940715) (512.0,0.006223740352889002) (1024.0,0.005039225069477116)}
        ;
    \addplot[fill opacity={0.18}, draw={none}, mark={none}, forget plot, ours_ransac, fill]
        fill between [of=UeR_std_median_1 and LeR_std_median_1]
        ;
    \addplot[ours_ransac_marker]
        coordinates {(2.0,0.3476995667023428) (4.0,0.2009901356635738) (8.0,0.13980851790580665) (16.0,0.05686834505499979) (32.0,0.03726180336260101) (64.0,0.025249594812930126) (128.0,0.014304342464866365) (256.0,0.00823913926645836) (512.0,0.006662203028145353) (1024.0,0.005387752647008922)}
        ;
    \addplot[name path={UeR_std_median_2}, draw={none}, mark={none}, forget plot]
        coordinates {(2.0,0.05273819657760633) (4.0,0.04838699124531745) (8.0,0.04534000159583285) (16.0,0.035685284696200976) (32.0,0.02798194115270608) (64.0,0.016628544943267546) (128.0,0.009802977750063214) (256.0,0.00667454890291027) (512.0,0.006101442858357225) (1024.0,0.004611298562349495)}
        ;
    \addplot[name path={LeR_std_median_2}, draw={none}, mark={none}, forget plot]
        coordinates {(2.0,0.04993328548008215) (4.0,0.04498851604465807) (8.0,0.04120774385929711) (16.0,0.031097415105962707) (32.0,0.022704028288343798) (64.0,0.013265613516590651) (128.0,0.008345651837593226) (256.0,0.0059970236042311675) (512.0,0.005446925140104416) (1024.0,0.004276980865520519)}
        ;
    \addplot[fill opacity={0.18}, draw={none}, mark={none}, forget plot, ours_msac, fill]
        fill between [of=UeR_std_median_2 and LeR_std_median_2]
        ;
    \addplot[ours_msac_marker]
        coordinates {(2.0,0.051387454436126426) (4.0,0.04665939793395458) (8.0,0.04323043930275437) (16.0,0.03331990689885394) (32.0,0.025166262934194404) (64.0,0.014680562977545355) (128.0,0.009000703974206305) (256.0,0.0063237881761861224) (512.0,0.005741766180284516) (1024.0,0.004406204698603088)}
        ;
    \addplot[name path={UeR_std_median_3}, draw={none}, mark={none}, forget plot]
        coordinates {(2.0,0.04074687961247517) (4.0,0.03804184540130856) (8.0,0.035408464672630376) (16.0,0.029213715319384895) (32.0,0.022882582811412748) (64.0,0.016192444201530123) (128.0,0.008669839488638614) (256.0,0.005756711358167751) (512.0,0.003263230091027392) (1024.0,0.0018702598559254614)}
        ;
    \addplot[name path={LeR_std_median_3}, draw={none}, mark={none}, forget plot]
        coordinates {(2.0,0.03879454901324211) (4.0,0.03556080905444895) (8.0,0.03259478850436757) (16.0,0.025837461250931913) (32.0,0.019298601963872758) (64.0,0.013731437523663443) (128.0,0.007183447504994401) (256.0,0.0042519892067657715) (512.0,0.002545290444356502) (1024.0,0.0014381106721864942)}
        ;
    \addplot[fill opacity={0.18}, draw={none}, mark={none}, forget plot, ours_gau, fill]
        fill between [of=UeR_std_median_3 and LeR_std_median_3]
        ;
    \addplot[ours_gau_marker]
        coordinates {(2.0,0.03980435855253438) (4.0,0.036814492379379206) (8.0,0.034009004429112526) (16.0,0.027521761576151105) (32.0,0.02099262427591586) (64.0,0.01484497820980218) (128.0,0.00787212619151829) (256.0,0.004801447568929025) (512.0,0.0028690498900697473) (1024.0,0.0015438848441240004)}
        ;
    \addplot[name path={UeR_std_median_4}, draw={none}, mark={none}, forget plot]
        coordinates {(2.0,0.04883263084156099) (4.0,0.04569351521390991) (8.0,0.042928230142475265) (16.0,0.03636800347484935) (32.0,0.027154585076054463) (64.0,0.01837093043776064) (128.0,0.010174045102330652) (256.0,0.005940522576352733) (512.0,0.0033363826236082767) (1024.0,0.001275700684613992)}
        ;
    \addplot[name path={LeR_std_median_4}, draw={none}, mark={none}, forget plot]
        coordinates {(2.0,0.04642450906932718) (4.0,0.04261895919757354) (8.0,0.03923520251791654) (16.0,0.03196140989022652) (32.0,0.022346744243717007) (64.0,0.014324950263972977) (128.0,0.007471817130378956) (256.0,0.004389218275156436) (512.0,0.002638743378069194) (1024.0,0.0009445304339881829)}
        ;
    \addplot[fill opacity={0.18}, draw={none}, mark={none}, forget plot, ours_magsac, fill]
        fill between [of=UeR_std_median_4 and LeR_std_median_4]
        ;
    \addplot[ours_magsac_marker]
        coordinates {(2.0,0.04765303114588453) (4.0,0.04419270285838805) (8.0,0.041057956613911796) (16.0,0.03416013431964441) (32.0,0.024661003644776597) (64.0,0.016035355790404795) (128.0,0.008337899411994752) (256.0,0.0049293279053722475) (512.0,0.002952622279597421) (1024.0,0.001058085607398551)}
        ;
    \addplot[name path={UeR_std_median_5}, draw={none}, mark={none}, forget plot]
        coordinates {(2.0,0.018999394674633373) (4.0,0.01850477536427486) (8.0,0.017983077873652654) (16.0,0.016152597078670245) (32.0,0.014413586284557463) (64.0,0.011537107951010836) (128.0,0.008637259464443792) (256.0,0.005635954017834496) (512.0,0.0035810822181179813) (1024.0,0.002802763878237355)}
        ;
    \addplot[name path={LeR_std_median_5}, draw={none}, mark={none}, forget plot]
        coordinates {(2.0,0.01846997341999419) (4.0,0.017842717431031427) (8.0,0.017201381568561193) (16.0,0.014996525686470738) (32.0,0.013097047365656325) (64.0,0.010072385404641417) (128.0,0.007153142485849912) (256.0,0.004679267568938223) (512.0,0.0031344189166839604) (1024.0,0.00258922480173439)}
        ;
    \addplot[fill opacity={0.18}, draw={none}, mark={none}, forget plot, ours_jeffreys, fill]
        fill between [of=UeR_std_median_5 and LeR_std_median_5]
        ;
    \addplot[ours_jeffreys_marker]
        coordinates {(2.0,0.01873024025718189) (4.0,0.018178790941980275) (8.0,0.01759892417338628) (16.0,0.015574630454634798) (32.0,0.013759897789652973) (64.0,0.010773786712462175) (128.0,0.007861241729340216) (256.0,0.00509249902692951) (512.0,0.003292305687247468) (1024.0,0.0026952267441987007)}
        ;
    \addplot[name path={UeR_std_median_6}, draw={none}, mark={none}, forget plot]
        coordinates {(2.0,0.016283796302454554) (4.0,0.015733840430485008) (8.0,0.015415121148795778) (16.0,0.013857017368734022) (32.0,0.012561244470173163) (64.0,0.010529596562422562) (128.0,0.007634586552274358) (256.0,0.005214870808545569) (512.0,0.00351758774146975) (1024.0,0.002166431344208142)}
        ;
    \addplot[name path={LeR_std_median_6}, draw={none}, mark={none}, forget plot]
        coordinates {(2.0,0.015576144722678036) (4.0,0.015105391125568798) (8.0,0.014739920997540982) (16.0,0.01291617684958342) (32.0,0.011433388139730733) (64.0,0.009277672286316381) (128.0,0.006391919348793561) (256.0,0.0043203785401002065) (512.0,0.0030068713106925563) (1024.0,0.0019075313305087897)}
        ;
    \addplot[fill opacity={0.18}, draw={none}, mark={none}, forget plot, ours_ig_robust, fill]
        fill between [of=UeR_std_median_6 and LeR_std_median_6]
        ;
    \addplot[ours_ig_robust_marker]
        coordinates {(2.0,0.015924886219231486) (4.0,0.015412347929180748) (8.0,0.01507736599977425) (16.0,0.01341072000013056) (32.0,0.012012904553737134) (64.0,0.00988601745393794) (128.0,0.0069664299150951695) (256.0,0.004727690328943184) (512.0,0.003192259365389359) (1024.0,0.002034224114157228)}
        ;
    \addplot[name path={UeR_std_median_7}, draw={none}, mark={none}, forget plot]
        coordinates {(2.0,0.014614083716587584) (4.0,0.015424877457645522) (8.0,0.014388300497577729) (16.0,0.013505145631918276) (32.0,0.012565801729754834) (64.0,0.010478564837891462) (128.0,0.007714596713877692) (256.0,0.005482255825958038) (512.0,0.0034833616952680254) (1024.0,0.002219932962201251)}
        ;
    \addplot[name path={LeR_std_median_7}, draw={none}, mark={none}, forget plot]
        coordinates {(2.0,0.012575264360672689) (4.0,0.014321865870200335) (8.0,0.013615436092937526) (16.0,0.01253771456327273) (32.0,0.011465879422815552) (64.0,0.00923605447895927) (128.0,0.006438180148744924) (256.0,0.004491418414320243) (512.0,0.0030076089127906614) (1024.0,0.001954179838750523)}
        ;
    \addplot[fill opacity={0.18}, draw={none}, mark={none}, forget plot, ours_ig_mle, fill]
        fill between [of=UeR_std_median_7 and LeR_std_median_7]
        ;
    \addplot[ours_ig_mle_marker]
        coordinates {(2.0,0.013563036169231652) (4.0,0.014856997609416013) (8.0,0.013996444162087496) (16.0,0.013021490073960584) (32.0,0.012030969698538714) (64.0,0.00983244135914551) (128.0,0.007030959163658759) (256.0,0.004933345081509247) (512.0,0.0031911802670837913) (1024.0,0.0020843028318508683)}
        ;
%
			% Row 3 (HEB)
			\renewcommand\nextgroupplot[1][]{\origNGP[#1, ylabel={$\mathbb{E}[e_{\mathbf{R}}]$ (deg)}, ymax=10, ytick={7,8,9,10}]}%
			% panel ..._group.tex — \nextgroupplot fragment.
% Host preamble MUST load: \usepgfplotslibrary{groupplots},
%                          \usepgfplotslibrary{fillbetween}.
% Wrap N such fragments inside one \begin{groupplot}[...] that
% sets `panel_axis, xmode=log, log basis x=2, xmajorgrids,
% ymajorgrids, xlabel={…}, ylabel={…}, xtick={…},
% ytick distance=…, yminorgrids, scaled y ticks=false,
% yticklabel style={…}, group style={group size=N by M, …}].

\nextgroupplot[title={HEB / SIFT / \(\bH\)}]
    \addplot[name path={UeR_mean_median_1}, draw={none}, mark={none}, forget plot]
        coordinates {(2.0,9.375609366129964) (4.0,9.251873700824559) (8.0,9.122050451828091) (16.0,9.009047390425003) (32.0,8.882401708956877) (64.0,8.703620622729108) (128.0,8.596673158065306) (256.0,8.604287253813148) (512.0,8.511660728955363) (1024.0,8.468435806517673) (2048.0,8.465954667733332)}
        ;
    \addplot[name path={LeR_mean_median_1}, draw={none}, mark={none}, forget plot]
        coordinates {(2.0,9.212488582518944) (4.0,9.099566890888907) (8.0,8.967978482443163) (16.0,8.863806951344252) (32.0,8.740269891530755) (64.0,8.576293060911507) (128.0,8.481050157677855) (256.0,8.481488866106565) (512.0,8.39517331585784) (1024.0,8.358327148554663) (2048.0,8.34722955668986)}
        ;
    \addplot[fill opacity={0.18}, draw={none}, mark={none}, forget plot, ours_ransac, fill]
        fill between [of=UeR_mean_median_1 and LeR_mean_median_1]
        ;
    \addplot[ours_ransac_marker]
        coordinates {(2.0,9.29107615619851) (4.0,9.173358404633476) (8.0,9.043909476060618) (16.0,8.93306615977832) (32.0,8.807021434335885) (64.0,8.63844796842829) (128.0,8.53650120456056) (256.0,8.539474749480402) (512.0,8.450779442039467) (1024.0,8.410484843266882) (2048.0,8.401605910392343)}
        ;
    \addplot[name path={UeR_mean_median_2}, draw={none}, mark={none}, forget plot]
        coordinates {(2.0,8.708577509006398) (4.0,8.621570440866192) (8.0,8.510005389077548) (16.0,8.469120784924296) (32.0,8.321436089193002) (64.0,8.210771808171039) (128.0,8.157544478679771) (256.0,8.114361568289837) (512.0,8.033088034581555) (1024.0,8.043579570044145) (2048.0,8.05037953857694)}
        ;
    \addplot[name path={LeR_mean_median_2}, draw={none}, mark={none}, forget plot]
        coordinates {(2.0,8.563125233638967) (4.0,8.493815279706501) (8.0,8.376312960287176) (16.0,8.341738941333636) (32.0,8.203979498202422) (64.0,8.101588827647618) (128.0,8.048655972261495) (256.0,8.008771084652725) (512.0,7.935407969773153) (1024.0,7.944671932615853) (2048.0,7.9474817601451635)}
        ;
    \addplot[fill opacity={0.18}, draw={none}, mark={none}, forget plot, ours_msac, fill]
        fill between [of=UeR_mean_median_2 and LeR_mean_median_2]
        ;
    \addplot[ours_msac_marker]
        coordinates {(2.0,8.634436906483272) (4.0,8.557961896926473) (8.0,8.441176507252438) (16.0,8.403316871071802) (32.0,8.260979690719628) (64.0,8.153040583159187) (128.0,8.100935043226817) (256.0,8.06016421723944) (512.0,7.982969940637057) (1024.0,7.9925895303218475) (2048.0,7.996816209699357)}
        ;
    \addplot[name path={UeR_mean_median_3}, draw={none}, mark={none}, forget plot]
        coordinates {(2.0,8.37910689564648) (4.0,8.29815363736837) (8.0,8.26111938139011) (16.0,8.233280399280988) (32.0,8.099500027406956) (64.0,8.06274205801016) (128.0,7.990923357032992) (256.0,7.964645106292856) (512.0,7.914009875937439) (1024.0,7.913575871942516) (2048.0,7.889337838348618)}
        ;
    \addplot[name path={LeR_mean_median_3}, draw={none}, mark={none}, forget plot]
        coordinates {(2.0,8.25584793810004) (4.0,8.189520074716775) (8.0,8.145010238167693) (16.0,8.126138921817141) (32.0,7.998662062342371) (64.0,7.965749043627387) (128.0,7.897480495690275) (256.0,7.8744834272318265) (512.0,7.827156861294008) (1024.0,7.826389969371114) (2048.0,7.802889787110061)}
        ;
    \addplot[fill opacity={0.18}, draw={none}, mark={none}, forget plot, ours_gau, fill]
        fill between [of=UeR_mean_median_3 and LeR_mean_median_3]
        ;
    \addplot[ours_gau_marker]
        coordinates {(2.0,8.315545265775961) (4.0,8.242583211592347) (8.0,8.200931610071686) (16.0,8.178407809455654) (32.0,8.048517737964765) (64.0,8.011271934619925) (128.0,7.941474463094286) (256.0,7.918891438551958) (512.0,7.8695659985530755) (1024.0,7.868821664453724) (2048.0,7.845207025127676)}
        ;
    \addplot[name path={UeR_mean_median_4}, draw={none}, mark={none}, forget plot]
        coordinates {(2.0,8.513960753638594) (4.0,8.456539961645035) (8.0,8.367642624369761) (16.0,8.33607496084311) (32.0,8.165864466549449) (64.0,8.110462899162911) (128.0,8.06285339848042) (256.0,8.003068483878884) (512.0,7.932544337196706) (1024.0,7.956488714972748) (2048.0,7.949590199261708)}
        ;
    \addplot[name path={LeR_mean_median_4}, draw={none}, mark={none}, forget plot]
        coordinates {(2.0,8.376031732490665) (4.0,8.332766483339496) (8.0,8.244851861167168) (16.0,8.21687971996129) (32.0,8.058019560649251) (64.0,8.003450601585564) (128.0,7.958447346691978) (256.0,7.9070464688358335) (512.0,7.842685803288674) (1024.0,7.861615322279638) (2048.0,7.853822910415687)}
        ;
    \addplot[fill opacity={0.18}, draw={none}, mark={none}, forget plot, ours_magsac, fill]
        fill between [of=UeR_mean_median_4 and LeR_mean_median_4]
        ;
    \addplot[ours_magsac_marker]
        coordinates {(2.0,8.443698851842498) (4.0,8.393111824888136) (8.0,8.303554187473699) (16.0,8.275322034597439) (32.0,8.11153561573624) (64.0,8.054299378125453) (128.0,8.008622953452276) (256.0,7.954170246975897) (512.0,7.887186395922164) (1024.0,7.907783771576466) (2048.0,7.89901457660238)}
        ;
    \addplot[name path={UeR_mean_median_5}, draw={none}, mark={none}, forget plot]
        coordinates {(2.0,9.770832860671701) (4.0,9.472142644305483) (8.0,9.115118102550454) (16.0,8.876974777078281) (32.0,8.607713604379787) (64.0,8.536750114860462) (128.0,8.434856389581924) (256.0,8.37287763763695) (512.0,8.281558995615338) (1024.0,8.247085124142014) (2048.0,8.20807337164064)}
        ;
    \addplot[name path={LeR_mean_median_5}, draw={none}, mark={none}, forget plot]
        coordinates {(2.0,9.417676754958034) (4.0,9.180612139698185) (8.0,8.882523075981313) (16.0,8.693820382991293) (32.0,8.481390876339045) (64.0,8.415595774102716) (128.0,8.319937520537406) (256.0,8.262749800146826) (512.0,8.176951376723753) (1024.0,8.146204831175801) (2048.0,8.106280175955654)}
        ;
    \addplot[fill opacity={0.18}, draw={none}, mark={none}, forget plot, ours_jeffreys, fill]
        fill between [of=UeR_mean_median_5 and LeR_mean_median_5]
        ;
    \addplot[ours_jeffreys_marker]
        coordinates {(2.0,9.5825503881191) (4.0,9.313736049650792) (8.0,8.985794779400418) (16.0,8.776781629026175) (32.0,8.54211229579993) (64.0,8.474652161520433) (128.0,8.375342195020943) (256.0,8.31720349273107) (512.0,8.228916176636035) (1024.0,8.195741657134343) (2048.0,8.155081508533032)}
        ;
    \addplot[name path={UeR_mean_median_6}, draw={none}, mark={none}, forget plot]
        coordinates {(2.0,8.276752704540591) (4.0,8.164278283140387) (8.0,8.06503110462473) (16.0,7.9860954052170925) (32.0,7.857438340814213) (64.0,7.813218876499476) (128.0,7.77382145449076) (256.0,7.750113865104355) (512.0,7.723290124580357) (1024.0,7.712338188209683) (2048.0,7.689485164311206)}
        ;
    \addplot[name path={LeR_mean_median_6}, draw={none}, mark={none}, forget plot]
        coordinates {(2.0,8.094550502787348) (4.0,8.021407929100995) (8.0,7.9337769057527945) (16.0,7.877784973062472) (32.0,7.769672213240448) (64.0,7.733497950971524) (128.0,7.698943702428771) (256.0,7.677987543529758) (512.0,7.651088341921271) (1024.0,7.642373914563045) (2048.0,7.616985487092856)}
        ;
    \addplot[fill opacity={0.18}, draw={none}, mark={none}, forget plot, ours_ig_robust, fill]
        fill between [of=UeR_mean_median_6 and LeR_mean_median_6]
        ;
    \addplot[ours_ig_robust_marker]
        coordinates {(2.0,8.180198176633311) (4.0,8.086819127561752) (8.0,7.993327062731261) (16.0,7.9292357797492645) (32.0,7.811357905805368) (64.0,7.771838187888898) (128.0,7.735187210772751) (256.0,7.713615309230363) (512.0,7.687943598498725) (1024.0,7.67698898842929) (2048.0,7.652080033166571)}
        ;
    \addplot[name path={UeR_mean_median_7}, draw={none}, mark={none}, forget plot]
        coordinates {(2.0,8.771100398374388) (4.0,8.55600887166098) (8.0,8.247004897219195) (16.0,8.170189965515178) (32.0,8.02248702153555) (64.0,8.118466949890257) (128.0,8.185957879216367) (256.0,8.258174191700457) (512.0,8.26883887416377) (1024.0,8.250173966553339) (2048.0,8.211408033333967)}
        ;
    \addplot[name path={LeR_mean_median_7}, draw={none}, mark={none}, forget plot]
        coordinates {(2.0,8.364300542774046) (4.0,8.277116673340725) (8.0,8.046464500366845) (16.0,8.001247203789237) (32.0,7.917116550208506) (64.0,8.008175072266136) (128.0,8.078885037312569) (256.0,8.153864365828438) (512.0,8.163918950364652) (1024.0,8.148644050582824) (2048.0,8.109528705801969)}
        ;
    \addplot[fill opacity={0.18}, draw={none}, mark={none}, forget plot, ours_ig_mle, fill]
        fill between [of=UeR_mean_median_7 and LeR_mean_median_7]
        ;
    \addplot[ours_ig_mle_marker]
        coordinates {(2.0,8.537319092066571) (4.0,8.398839022850785) (8.0,8.134233878199973) (16.0,8.073875722605578) (32.0,7.966718811965107) (64.0,8.059680609306213) (128.0,8.13078856766258) (256.0,8.20527966439354) (512.0,8.21609279741437) (1024.0,8.198369751080367) (2048.0,8.158742709047687)}
        ;
%
			\renewcommand\nextgroupplot[1][]{\origNGP[#1, ylabel={$\mathrm{std}[e_{\mathbf{R}}]$ (deg)}, ytick={1,2,3}]}%
			% panel ..._group.tex — \nextgroupplot fragment.
% Host preamble MUST load: \usepgfplotslibrary{groupplots},
%                          \usepgfplotslibrary{fillbetween}.
% Wrap N such fragments inside one \begin{groupplot}[...] that
% sets `panel_axis, xmode=log, log basis x=2, xmajorgrids,
% ymajorgrids, xlabel={…}, ylabel={…}, xtick={…},
% ytick distance=…, yminorgrids, scaled y ticks=false,
% yticklabel style={…}, group style={group size=N by M, …}].

\nextgroupplot[title={HEB / SIFT / \(\bH\)}]
    \addplot[name path={UeR_std_median_1}, draw={none}, mark={none}, forget plot]
        coordinates {(2.0,1.3707409949945495) (4.0,1.2897031271645931) (8.0,1.3096495476074654) (16.0,1.227554870811919) (32.0,1.218683040897872) (64.0,1.0965442350619454) (128.0,1.0047293349903628) (256.0,1.0554793564923637) (512.0,1.0429592503480847) (1024.0,0.966744410975953) (2048.0,1.0524774861121182)}
        ;
    \addplot[name path={LeR_std_median_1}, draw={none}, mark={none}, forget plot]
        coordinates {(2.0,1.2565700253088892) (4.0,1.1792086357356988) (8.0,1.1867727785008344) (16.0,1.1080696380283954) (32.0,1.0849891307427484) (64.0,0.9465788651936473) (128.0,0.8815654412947551) (256.0,0.9151154102568446) (512.0,0.8794454572132442) (1024.0,0.8218226411442249) (2048.0,0.8835627141267864)}
        ;
    \addplot[fill opacity={0.18}, draw={none}, mark={none}, forget plot, ours_ransac, fill]
        fill between [of=UeR_std_median_1 and LeR_std_median_1]
        ;
    \addplot[ours_ransac_marker]
        coordinates {(2.0,1.3109464474007995) (4.0,1.2305754178999537) (8.0,1.2422357801428987) (16.0,1.1657912437425466) (32.0,1.148715337692474) (64.0,1.0148177285568167) (128.0,0.9378209035818704) (256.0,0.981016813011893) (512.0,0.9514805132932037) (1024.0,0.8890419713880137) (2048.0,0.9616173165007792)}
        ;
    \addplot[name path={UeR_std_median_2}, draw={none}, mark={none}, forget plot]
        coordinates {(2.0,1.2381713493644197) (4.0,1.095367930169106) (8.0,1.1302992863602344) (16.0,1.0833064844269207) (32.0,0.9943674348561737) (64.0,0.9481281577051925) (128.0,0.9313207343710974) (256.0,0.8965463415613151) (512.0,0.8363009230960217) (1024.0,0.8583200839897793) (2048.0,0.8695032201800078)}
        ;
    \addplot[name path={LeR_std_median_2}, draw={none}, mark={none}, forget plot]
        coordinates {(2.0,1.126485286183687) (4.0,1.0000367966881285) (8.0,1.0144747481278915) (16.0,0.972720632172859) (32.0,0.895856049307688) (64.0,0.8337768353823451) (128.0,0.8336580087020528) (256.0,0.7994342705402531) (512.0,0.7487593732427997) (1024.0,0.7670106199106379) (2048.0,0.7896281207468889)}
        ;
    \addplot[fill opacity={0.18}, draw={none}, mark={none}, forget plot, ours_msac, fill]
        fill between [of=UeR_std_median_2 and LeR_std_median_2]
        ;
    \addplot[ours_msac_marker]
        coordinates {(2.0,1.1795984749980257) (4.0,1.0451210241501558) (8.0,1.0676553992943247) (16.0,1.0238092232604314) (32.0,0.9413368841732798) (64.0,0.885278350667172) (128.0,0.8793134189512422) (256.0,0.8442540226074862) (512.0,0.7895162319306748) (1024.0,0.8088869698971284) (2048.0,0.8268096032970168)}
        ;
    \addplot[name path={UeR_std_median_3}, draw={none}, mark={none}, forget plot]
        coordinates {(2.0,1.0644103089512469) (4.0,0.9217398636892598) (8.0,0.9803570392471989) (16.0,0.9105168087181759) (32.0,0.8436354875035988) (64.0,0.8253423912497178) (128.0,0.8026778931570164) (256.0,0.7717274798990492) (512.0,0.7385665364881484) (1024.0,0.7446670478068286) (2048.0,0.7358568151496672)}
        ;
    \addplot[name path={LeR_std_median_3}, draw={none}, mark={none}, forget plot]
        coordinates {(2.0,0.9614945037423815) (4.0,0.8397989201768753) (8.0,0.8738418690748126) (16.0,0.8241091994891152) (32.0,0.7620223139419254) (64.0,0.7394800893007483) (128.0,0.7156776268932953) (256.0,0.6969348088131889) (512.0,0.6726339401865025) (1024.0,0.6785020831406003) (2048.0,0.6746007705057535)}
        ;
    \addplot[fill opacity={0.18}, draw={none}, mark={none}, forget plot, ours_gau, fill]
        fill between [of=UeR_std_median_3 and LeR_std_median_3]
        ;
    \addplot[ours_gau_marker]
        coordinates {(2.0,1.0102660561827619) (4.0,0.8775084741006248) (8.0,0.9238790206733046) (16.0,0.8645440828794524) (32.0,0.79872968984324) (64.0,0.7773414538379531) (128.0,0.7527725836591146) (256.0,0.7312765547265019) (512.0,0.7037386555246156) (1024.0,0.709944192838224) (2048.0,0.7033063379888816)}
        ;
    \addplot[name path={UeR_std_median_4}, draw={none}, mark={none}, forget plot]
        coordinates {(2.0,1.1724913747733612) (4.0,1.0544744553264844) (8.0,1.0487653891541053) (16.0,1.001937658059871) (32.0,0.9357200727670264) (64.0,0.9197458796724041) (128.0,0.8793602668402898) (256.0,0.8254594406682104) (512.0,0.7699121005724114) (1024.0,0.8032068709036222) (2048.0,0.8113714972348659)}
        ;
    \addplot[name path={LeR_std_median_4}, draw={none}, mark={none}, forget plot]
        coordinates {(2.0,1.0594507838739156) (4.0,0.9586533427988448) (8.0,0.9415239612380067) (16.0,0.901769744112455) (32.0,0.8379315710042918) (64.0,0.8155117632373005) (128.0,0.7940096938958067) (256.0,0.7395424581922478) (512.0,0.6969656346100497) (1024.0,0.7307751493505419) (2048.0,0.7396724487066353)}
        ;
    \addplot[fill opacity={0.18}, draw={none}, mark={none}, forget plot, ours_magsac, fill]
        fill between [of=UeR_std_median_4 and LeR_std_median_4]
        ;
    \addplot[ours_magsac_marker]
        coordinates {(2.0,1.113475012779567) (4.0,1.0033498012212814) (8.0,0.9908641493520502) (16.0,0.9478389695088729) (32.0,0.8822283998032329) (64.0,0.8625825073576555) (128.0,0.8328376901629249) (256.0,0.7788499921201657) (512.0,0.7312188657568987) (1024.0,0.764222144079287) (2048.0,0.7717958258127753)}
        ;
    \addplot[name path={UeR_std_median_5}, draw={none}, mark={none}, forget plot]
        coordinates {(2.0,3.169501372473362) (4.0,2.667722727047213) (8.0,2.2315291750283452) (16.0,1.827140312905471) (32.0,1.114148656379042) (64.0,1.0378931178886326) (128.0,0.9774711129666143) (256.0,0.9537417158676438) (512.0,0.8909335806440324) (1024.0,0.8656806984500836) (2048.0,0.8584662603998336)}
        ;
    \addplot[name path={LeR_std_median_5}, draw={none}, mark={none}, forget plot]
        coordinates {(2.0,2.5466036503892955) (4.0,2.0472548225992875) (8.0,1.5999497428272969) (16.0,1.2591524338850204) (32.0,0.9800432452948643) (64.0,0.9351748425675679) (128.0,0.8815012815409833) (256.0,0.854386418206039) (512.0,0.8068189297867412) (1024.0,0.7839185167161915) (2048.0,0.7806866590555446)}
        ;
    \addplot[fill opacity={0.18}, draw={none}, mark={none}, forget plot, ours_jeffreys, fill]
        fill between [of=UeR_std_median_5 and LeR_std_median_5]
        ;
    \addplot[ours_jeffreys_marker]
        coordinates {(2.0,2.839477989810774) (4.0,2.3195466571962426) (8.0,1.8574309543191263) (16.0,1.4562280433962609) (32.0,1.038671271407863) (64.0,0.9837844513522667) (128.0,0.9268215379098875) (256.0,0.9000212813991475) (512.0,0.8473548377521976) (1024.0,0.8219780624127189) (2048.0,0.8188694436530203)}
        ;
    \addplot[name path={UeR_std_median_6}, draw={none}, mark={none}, forget plot]
        coordinates {(2.0,1.5770670521458066) (4.0,1.2932507487278861) (8.0,1.1938101377166523) (16.0,0.9980879685684954) (32.0,0.8297280488817029) (64.0,0.6858130266780297) (128.0,0.6416155874414512) (256.0,0.6065073260885331) (512.0,0.6061468313779546) (1024.0,0.5949588597484011) (2048.0,0.6026922109392029)}
        ;
    \addplot[name path={LeR_std_median_6}, draw={none}, mark={none}, forget plot]
        coordinates {(2.0,1.243937286840556) (4.0,1.0001608405735176) (8.0,0.9285120417109055) (16.0,0.7853625032623441) (32.0,0.6514651532678811) (64.0,0.6094683093971764) (128.0,0.5834769121439809) (256.0,0.5588064803472487) (512.0,0.5596046263064386) (1024.0,0.5482973098875898) (2048.0,0.5578393113639692)}
        ;
    \addplot[fill opacity={0.18}, draw={none}, mark={none}, forget plot, ours_ig_robust, fill]
        fill between [of=UeR_std_median_6 and LeR_std_median_6]
        ;
    \addplot[ours_ig_robust_marker]
        coordinates {(2.0,1.397520848158763) (4.0,1.122392419665474) (8.0,1.0462684754072633) (16.0,0.8696601692290526) (32.0,0.7088832794963871) (64.0,0.6417243232328209) (128.0,0.6095230785601408) (256.0,0.5814203498419312) (512.0,0.5821827977513748) (1024.0,0.5703670221788019) (2048.0,0.5796789166414801)}
        ;
    \addplot[name path={UeR_std_median_7}, draw={none}, mark={none}, forget plot]
        coordinates {(2.0,2.8366343928573294) (4.0,2.4839932013519634) (8.0,1.9951750358702531) (16.0,1.7742165062981519) (32.0,0.9506975357229608) (64.0,0.9487289601025888) (128.0,0.9105659176895525) (256.0,0.9030084299039529) (512.0,0.8867088636406684) (1024.0,0.8672447607020352) (2048.0,0.8620532183374328)}
        ;
    \addplot[name path={LeR_std_median_7}, draw={none}, mark={none}, forget plot]
        coordinates {(2.0,1.8275316397153032) (4.0,1.7980907498171843) (8.0,1.3265633412466495) (16.0,1.0766982092345576) (32.0,0.7875175539362146) (64.0,0.8374317487525185) (128.0,0.8152868091773241) (256.0,0.8022668683733394) (512.0,0.8040079023110085) (1024.0,0.7844201463256321) (2048.0,0.7840099160910314)}
        ;
    \addplot[fill opacity={0.18}, draw={none}, mark={none}, forget plot, ours_ig_mle, fill]
        fill between [of=UeR_std_median_7 and LeR_std_median_7]
        ;
    \addplot[ours_ig_mle_marker]
        coordinates {(2.0,2.2553758210976254) (4.0,2.1026591966055492) (8.0,1.593966339435274) (16.0,1.3419196580427566) (32.0,0.8482858903939612) (64.0,0.8874505638446181) (128.0,0.8597434923507102) (256.0,0.8475497434122939) (512.0,0.8436491484955049) (1024.0,0.822274077170287) (2048.0,0.8217707205147038)}
        ;
%
			\let\nextgroupplot\origNGP%
		\end{groupplot}
	\end{tikzpicture}
	\par\vspace{0.4em}
	\newcommand{\SensSwatch}[2]{%
		\tikz[baseline=-0.6ex]{%
			\draw[color=#1, line width=1.5pt] (0,0) -- (0.9em,0);%
			\draw[color=#1, line width=1.0pt] plot[mark=#2, mark size=3pt, mark options={solid}] coordinates {(0.45em,0)};%
		}%
	}%
	\noindent\hbox to \linewidth{\footnotesize
		\SensSwatch{ColorRANSAC}{x}\,RANSAC\hfill
		\SensSwatch{ColorMSAC}{diamond*}\,MSAC\hfill
		\SensSwatch{ColorGaU}{triangle*}\,GaU\hfill
		\SensSwatch{ColorMAGSAC}{star}\,MAGSAC++\hfill
		\emph{Ours:}\hfill
		\SensSwatch{ColorJeffreys}{pentagon}\,Jeffreys\hfill
		\SensSwatch{ColorIGRobust}{o}\,IG-Robust\hfill
		\SensSwatch{ColorIGMLE}{square}\,IG-MLE%
	}
	\caption{\textbf{Sensitivity}: expected test rotation error
	(left) and its standard deviation (right) as a function
	of validation-set size, for three cells spanning the
	geometric-problem axis---PT~/~$\bE$ (top), ETH3D~/~$\bF$
	(middle), HEB~/~$\bH$ (bottom).  Shaded bands are 95\%
	BCa bootstrap confidence intervals over the
	cross-validation trials.
	\textbf{Our methods are nearly insensitive to validation-set
		size}: even with as few as 2--4 validation pairs the IG prior
	regularizes the per-partition scale enough to recover
	near-optimal hyperparameters, while threshold-based baselines
	need $\ConvRatioPT{\times}$ more validation data to converge on PhotoTourism
	($\ConvRatioETH{\times}$ on ETH3D at a 5\% tolerance).
	On HEB---where per-pair inlier counts are too sparse to
	support a stable joint-MLE prior fit---\textbf{IG-Robust's
		median-only anchor} is the appropriate variant.}
	\label{fig:sensitivity-grid}
\end{figure}

\subsubsection{Results}
The proposed methods reach near-optimal expected test error
at validation budgets as small as $|\mathcal{V}| = 2$, while
the threshold-based baselines need $\ConvRatioPT\times$ more
validation data to converge on PhotoTourism ($\ConvRatioETH\times$ on ETH3D
at a 5\% tolerance).  On PhotoTourism the three IG
variants (Jeffreys, IG-MLE, IG-Robust) sit within
$0.006^\circ$ of their large-$|\mathcal{V}|$ asymptote of
$0.77^\circ$ already at $|\mathcal{V}| = 2$, whereas
MSAC, GaU, and MAGSAC++ start at ${\sim}0.90^\circ$
and need $|\mathcal{V}| \approx 1024$ to reach the same
floor; vanilla RANSAC starts higher ($1.08^\circ$) and never
reaches it, plateauing near $0.87^\circ$.  ETH3D shows the same convergence pattern.  The proposed methods are also more reproducible
at small $|\mathcal{V}|$: their cross-trial standard deviation
$\mathrm{std}[e_{\mathbf{R}}]$ sits at
$0.002^\circ$--$0.019^\circ$ at $|\mathcal{V}| = 2$, while
threshold-based baselines disperse over
$\BaselinesStdVtwoLo^\circ$--$\BaselinesStdVtwoHi^\circ$ at the same budget---every baseline at least $\StdGapVtwoLo\times$ (ETH3D) to
$\StdGapVtwoHi\times$ (PhotoTourism) more than any proposed variant.

This is the deployment counterpart of the Profile robustness
result.  In \Cref{fig:profile-grid} the proposed methods'
curves are nearly flat across the two-decade
$\kappa/\kappa^\star$ window---a single hyperparameter is
nearly as good as any other.  Under sampling noise that
property translates directly into validation-budget
insensitivity: when many $\kappa$ values are nearly equally
good, the particular $\kappa^\star(\mathcal{V})$ chosen on a
tiny validation pool lands inside the ``good $\kappa$''
region with high probability.  Threshold-based methods, by
contrast, have steeper curves and therefore narrower good
regions, so a small $\mathcal{V}$ is unlikely to put
$\kappa^\star(\mathcal{V})$ inside one.  For the RANSAC
practitioner this means the proposed methods can be deployed
with a handful of ground-truth pairs---or, in the data-rich
regime, with the prior alone---while threshold-based scoring
demands a substantial labelled validation set to be tuned
reliably.

\paragraph{HEB inverts the prior-irrelevance story.}
On PT and ETH3D the three IG variants are
indistinguishable---Profile's data-rich regime, where the
prior recedes.  HEB undoes this.  Jeffreys is now the
\emph{worst} of the three at small $|\mathcal{V}|$
($9.58^\circ$ at $|\mathcal{V}| = 2$, vs.\ $8.18^\circ$ for
IG-Robust), and its trial-to-trial $\mathrm{std}$ is the
largest of any method ($2.84^\circ$).  This is the same
low-$\nu_p$ regime we identified in Profile: with only
$n_p \approx 10$ inlying matches per pair, the per-partition
scale estimate is too noisy to discriminate models without
prior regularization, and Jeffreys gives that regularization
up.  IG-Robust's median-anchored prior supplies the missing
regularization without being dragged by outlier
$\hat\sigma^2_p$ tails, and dominates HEB at every
$|\mathcal{V}|$ in both mean ($7.65^\circ$ asymptote vs.\
$7.85^\circ$ for the best threshold-based baseline) and
std (lowest among the proposed variants at every budget, and
lowest of any method for $|\mathcal{V}| \ge 32$).  IG-MLE's joint
fit is corrupted by the same heavy-tailed $\hat\sigma^2_p$
produced by HEB's structured-outlier contamination, costing
it the regularization advantage---exactly the Profile
diagnosis.

% =====================================================================
\section{Conclusion}
\label{sec:conclusion}
% =====================================================================

We marginalize the noise scale $\sigma^2$ analytically out of
the Gaussian--Uniform inlier model under a conjugate
Inverse-Gamma prior, yielding a closed-form, scale-free RANSAC
score that spans the non-informative Jeffreys limit and informative
empirical-Bayes priors within a single expression, computable
in $O(N \log N)$ via sort-and-sweep.
Across three datasets covering three difficult two-view geometry
problems, with model proposals sampled from both engineered
(SIFT/RootSIFT) and learned (SuperPoint+LightGlue) feature
pipelines, in the data-rich regime the proposed methods are
nearly insensitive to hyperparameter miscalibration where
threshold-based scoring degrades markedly
(\Cref{fig:profile-grid}).  They reach
near-optimal expected test error with as few as two
validation pairs, while threshold-based baselines need
$\ConvRatioPT\times$ more on PhotoTourism (\Cref{fig:sensitivity-grid}).  And in
the data-scarce regime they supply principled regularization
that is automatically downweighted as data grows.  For the
practitioner, validation-size insensitivity means a handful
of ground-truth pairs (or none, in the data-rich regime
where Jeffreys suffices) is enough to deploy near-optimally,
which distinguishes the proposed family from the
threshold-based state of the art.  To our knowledge, this is
the first RANSAC score that genuinely marginalizes the noise
scale $\sigma$ in closed form.

% ---- Bibliography ----
\bibliographystyle{plainnat}
\bibliography{main}

% =====================================================================
%  SUPPLEMENTARY MATERIAL
% =====================================================================
\newpage
\appendix
\setcounter{section}{0}
\renewcommand{\thesection}{\Alph{section}}

\section*{Supplementary Material}

% =====================================================================
\section{Model Dimensions: Parameters, Codimension, and Minimal Sample Size}
\label{sec:supp-dims}
% ---------------------------------------------------------------------

\Cref{tab:model-dims} summarizes the key dimensional quantities for
common geometric models.  The model parameter
count~$n_\theta = \dim(\theta)$ is the number of free parameters after
fixing gauge freedoms.  The codimension~$d_g$ is the number of
independent algebraic constraints each observation imposes on the
model.  The minimal sample size $m = \lceil n_\theta / d_g \rceil$ is the
fewest observations that determine~$\theta$; equivalently, the
constraint-space design matrix
$\pdv{\bg_\cI}!{\theta} \in \Reals^{k\,d_g \times n_\theta}$ has rank~$n_\theta$ only when
$k \geq m$.  The observation dimension~$d$ is the ambient
dimension of each measurement~$\bx_i$.

\begin{table}[ht]
	\centering
	\caption{Dimensional quantities for standard two-view models.
		$n_\theta$: model parameters (free DOF);
		$d_g$: codimension (constraints per point);
		$m = \lceil n_\theta/d_g\rceil$: minimal sample size;
		$d$: observation dimension.}
	\label{tab:model-dims}
	\smallskip
	\begin{tabular}{lcccc}
		\toprule
		\textbf{Model}                    & $n_\theta$ & $d_g$ & $m$ & $d$ \\
		\midrule
		Fundamental matrix ($\mathbf{F}$) & 7          & 1     & 7   & 4   \\
		Homography ($\mathbf{H}$)         & 8          & 2     & 4   & 4   \\
		Essential matrix ($\mathbf{E}$)   & 5          & 1     & 5   & 4   \\
		\bottomrule
	\end{tabular}
\end{table}

\noindent The distinction between~$d_g$ and~$d$ is critical.
A na\"{\i}ve Mahalanobis approach that uses~$d$ in the inlier
likelihood over-counts degrees of freedom (this is one of the
errors identified in MAGSAC++~\cite{Shekhovtsov2025}).
Our formulation uses the
codimension~$d_g$---the number of independent constraints per observation---which
is the correct number of degrees of freedom per observation in
the projected likelihood.
Note that $m\,d_g = n_\theta$ for all models in the table; in
general the effective degrees of freedom for $n_I$~inliers is
$n_I\,d_g - n_\theta$.

% =====================================================================
\section{Shekhovtsov's Unified Threshold Framework}
\label{sec:supp-unified}
% =====================================================================

We review the unified comparison framework
of~\cite{Shekhovtsov2025}, which parameterizes the most widely
used RANSAC scoring functions by a single decision
threshold~$\tau$ (in pixels).  This framework is the basis of the experimental
comparisons in Section~\ref{sec:experiments}.

\subsection{From Mixture Model to Decision Threshold}
\label{sec:supp-mixture}

For the scalar case $d_g = 1$, the residual~$r_i$
is a scalar drawn from a two-component mixture: an inlier density
$\pin(r;\sigma) =
	\tfrac{1}{\sqrt{2\pi}\,\sigma}\exp(-r^2/2\sigma^2)$
and a uniform outlier density
$\pout = \frac{1}{2a}$ over $[-a, a]$, where $a$ is
the outlier domain half-width~\eqref{eq:component-densities}
(same units as~$\sigma$).
This 1D mixture is the starting point for
Shekhovtsov's~\cite{Shekhovtsov2025} $\tau$-parameterized framework.
With prior inlier probability
$\gamma$, the MAP classification rule labels point~$i$ as an inlier
when
\begin{equation}
	\gamma\, \pin(r_i;\sigma)
	\;>\; (1{-}\gamma)\, \pout\,.
	\label{eq:supp-map}
\end{equation}
Geometrically, the decision threshold~$\tau$ is the
\emph{intersection} of the prior-weighted inlier and outlier
densities:
\begin{equation}
	\gamma\, \pin(\tau;\,\sigma)
	\;=\; (1{-}\gamma)\, \pout\,,
	\label{eq:supp-tau-def}
\end{equation}
i.e., the residual value at which the two components of the mixture
are equally likely.  Since the Gaussian density
$\pin(r;\sigma)$ is monotone decreasing for $r \geq 0$
while the uniform density is constant, the intersection is unique,
and the MAP rule reduces to the simple threshold test
$|r_i| < \tau$.  Every method that uses this mixture
model---regardless of how it parameterizes itself---implicitly
defines such a crossing point~$\tau$.

\iffalse  % unreferenced; commented out for space
	\subsection{Scoring Functions Parameterized by $\tau$}
	\label{sec:supp-scoring}

	Shekhovtsov~\cite{Shekhovtsov2025} shows that the profile and
	marginal likelihood scoring functions can both be written in terms
	of~$\tau$ through the reference level
	$\log \pin(\tau;\sigma)$
	(we suppress the~$\sigma$ argument where it is absorbed
	into~$\tau$):
	\begin{align}
		\rho^{\mathrm{prof}}(r)
		 & = \max\!\bigl(\log \pin(r;\sigma)
		- \log \pin(\tau;\sigma),\; 0\bigr),
		\label{eq:supp-rho-prof}              \\[4pt]
		\rho^{\mathrm{MARG}}(r)
		 & = \smax\!\bigl(\log \pin(r;\sigma)
		- \log \pin(\tau;\sigma),\; 0\bigr),
		\label{eq:supp-rho-marg}
	\end{align}
	where
	$\smax(x,y) \coloneqq \log(e^x + e^y)$ is the
	smooth (soft) maximum.  In both cases the total score is
	$Q(\theta) = \sum_i \rho(r_i)$, an additive function of the
	residuals.

	\subsection{Gaussian Inliers: Normalization and $\sigma$-Cancellation}
	\label{sec:supp-gauss}

	For Gaussian inlier density
	$\pin(r;\sigma) =
		\tfrac{1}{\sqrt{2\pi}\,\sigma}\exp(-r^2/2\sigma^2)$,
	the scoring functions become:
	\begin{align}
		\rho^{\mathrm{prof}}(r)
		 & = \max\!\biggl(\frac{\tau^2 - r^2}{2\sigma^2},\; 0\biggr),
		\label{eq:supp-msac-unnorm}                                   \\[4pt]
		\rho^{\mathrm{MARG}}(r)
		 & = \smax\!\biggl(
		\frac{\tau^2 - r^2}{2\sigma^2},\; 0\biggr).
		\label{eq:supp-gau-unnorm}
	\end{align}
	\emph{Normalizing} the profile score so that
	$\tilde\rho(0) = 1$ and $\tilde\rho(\tau) = 0$:
	\begin{equation}
		\tilde\rho^{\,\mathrm{MSAC}}(r)
		= \max\!\bigl(1 - r^2/\tau^2,\; 0\bigr).
		\label{eq:supp-msac}
	\end{equation}
	The scale~$\sigma$ cancels completely---the normalized MSAC
	scoring function depends \emph{only} on~$\tau$.  Two different
	$(\sigma, \gamma)$ pairs that produce the same~$\tau$
	via~\eqref{eq:supp-tau-def} yield the same scoring function
	and the same inlier set.

	For the marginal (GaU) score, normalization gives
	\begin{equation}
		\tilde\rho^{\,\mathrm{GaU}}(r)
		= \frac{\smax\!\bigl(
		(\tau^2 {-} r^2)/(2\sigma^2),\; 0\bigr)}
		{\smax\!\bigl(
		\tau^2/(2\sigma^2),\; 0\bigr)}\,.
		\label{eq:supp-gau}
	\end{equation}
	Here $\sigma$ does not fully cancel: GaU depends on both $\tau$
	and~$\sigma$.  However, Shekhovtsov shows empirically that for
	fitted GaU models $\sigma \approx \tau$, so even GaU
	effectively has a single degree of freedom.

	\subsection{Mapping Each Method's Native Parameter to $\tau$}
	\label{sec:supp-mapping}

	Different methods parameterize themselves differently, but each
	implicitly defines a decision threshold~$\tau$:

	\begin{center}
		\begin{tabular}{@{}lll@{}}
			\textbf{Method} & \textbf{Native parameter}
			                & \textbf{Mapping to $\tau$}                          \\
			\midrule
			MSAC~\cite{Torr1998}
			                & $\tau$ directly
			                & identity                                            \\[2pt]
			MLESAC / GaU~\cite{Torr2000}
			                & $\sigma$
			                & root of $\pin(\tau;\sigma) = \mu$                   \\[2pt]
			MAGSAC++~\cite{Barath2020}
			                & $\sigma_{\max}$
			                & $\tau \approx \kappa\,\sigma_{\max}$
			\;\text{($\kappa$: 99th \%ile of $\chi_d$, $d$: obs.\ dim.)}          \\[2pt]
			Ours\textsuperscript{$\dagger$}
			                & $a$
			                & $\tau^* = \hat s\sqrt{2\log(a\sqrt{2/\pi}/\hat s)}$
			\;\text{(Eq.~\ref{eq:tau-star})}
		\end{tabular}
	\end{center}
	\noindent\textsuperscript{$\dagger$}The mapping
	$\tau^* = \hat s\sqrt{2\log(a\sqrt{2/\pi}/\hat s)}$ is a per-model
	asymptotic approximation: $\hat s$ depends on the candidate model, so
	$\tau^*$ varies across models and is not a global parameter.  This
	mapping is used only for exposition; the evaluation protocol
	(Section~\ref{sec:supp-protocol}) does not rely on it.

	\noindent For MAGSAC++, Shekhovtsov demonstrates that the
	scale-marginalized $\chi_d$ density closely approximates the GaU
	posterior, giving $\sigma \approx 0.96\,\tau$ when $d = 4$
	(the observation dimension for two-view correspondences).  The
	parameter $\sigma_{\max}$ therefore functions as an alternative
	parameterization of $\tau$, and MAGSAC++ reduces to a member of the
	GaU family~\cite{Shekhovtsov2025}.

	\subsection{Experimental Protocol}
	\label{sec:supp-protocol}

	With all methods expressed in terms of~$\tau$, Shekhovtsov's
	evaluation proceeds as follows:
	\begin{enumerate}
		\item Discretize $\tau$ over a grid of 200 values (e.g.,
		      $[0.1, 30]$ pixels).
		\item For each $\tau$ and each method, use
		      $\cI(\tau) = \{i : |r_i| < \tau\}$ as the inlier set and
		      evaluate the method's scoring function.  Additive scores
		      $Q = \sum_i \rho(r_i;\tau)$ can be computed efficiently via
		      a residual histogram shared across all threshold
		      values~\cite{Shekhovtsov2025}.
		\item Select the model~$\theta$ with the highest score, measure
		      pose error, and record the result.
		\item Plot pose error \emph{vs.}\ $\tau$ on a common axis.
	\end{enumerate}
	This protocol ensures a fair comparison: at each operating
	point the same inlier set is used, and the only difference
	between methods is the weighting of residuals within that set.
	The key conclusion of~\cite{Shekhovtsov2025} is that when $\tau$
	is properly tuned, all scoring functions---MSAC, GaU, and
	MAGSAC++---perform identically for model selection; the
	bottleneck is not the score but the choice of~$\tau$.

	\paragraph{Generalized protocol for non-additive methods.}
	The protocol above assumes that every scoring function is additive,
	i.e., $Q = \sum_i \rho(r_i;\tau)$, so that step~2 can be carried out
	via the shared residual histogram.
	Our scale-free score is non-additive
	(Section~\ref{sec:experiments}), so step~2 must be replaced by
	method-specific scoring: the method receives the candidate models and
	correspondences and returns the best model under its own rule.
	In practice the protocol becomes: (1)~generate a shared model pool as
	before; (2)~each method scores the pool using its own function and
	hyperparameter ($\tau$ for the additive baselines, $a$ for ours);
	(3)~measure pose error of the selected model; (4)~plot pose error
	vs.\ each method's native hyperparameter.
	The $x$-axis is therefore method-specific rather than a shared~$\tau$,
	but the comparison remains fair because every method scores the same
	models and is evaluated by the same metric, with its hyperparameter
	chosen on a held-out validation set
	(see Section~\ref{sec:eval-protocol}).
\fi  % end commented-out (sec:supp-scoring, sec:supp-gauss, sec:supp-mapping, sec:supp-protocol)

% =====================================================================
\section{Derivation of the Scale-Free Marginal Score}
\label{sec:supp-score-deriv}
% =====================================================================

For a fixed partition~$\cI$, applying the Inverse-Gamma prior
$\pi(\sigma^2) \propto (\sigma^2)^{-(\alpha_0+1)}\exp(-\beta_0/\sigma^2)$
to the partition likelihood~\eqref{eq:joint-partition} multiplies
the inlier Gaussian factor with a $\sigma^{-2(\alpha_0+1)}$ tail and an
exponential dampener at scale~$\beta_0$:
\begin{equation}
  p(\br \mid \theta, \cI)
  \;\propto\;
  \frac{1}{(2a)^{n_O d_g}\,(2\pi)^{n_I d_g/2}}
  \int_0^\infty
  \frac{
    \exp\!\bigl({-}\tfrac{\RSS_\cI/2 + \beta_0}{\sigma^2}\bigr)
  }{
    (\sigma^2)^{n_I d_g/2 + \alpha_0 + 1}
  }\,
  d\sigma^2,
  \label{eq:marg-sigma}
\end{equation}
where $\RSS_\cI \coloneqq \sum_{i \in \cI}\|\br_i\|^2$
is the residual sum of squares over inliers and the
proportionality absorbs the prior normalizer
$\beta_0^{\alpha_0}/\Gamma(\alpha_0)$, which is constant in
$(\theta, \cI)$ and drops out of the partition optimization.
The integral evaluates to the Inverse-Gamma normalizing
constant
\begin{equation}
  \int_0^\infty
  (\sigma^2)^{-(n_I d_g/2 + \alpha_0 + 1)}
  \exp\!\Bigl({-}\frac{\RSS_\cI/2 + \beta_0}{\sigma^2}\Bigr)\,
  d\sigma^2
  = \Gamma\!\Bigl(\frac{n_I\,d_g}{2} + \alpha_0\Bigr)\,
  \Bigl(\frac{\RSS_\cI}{2} + \beta_0\Bigr)^{\!-(n_I d_g/2 + \alpha_0)}.
  \label{eq:invgamma-integral}
\end{equation}
At $(\alpha_0,\beta_0) = (0,0)$ this collapses to
$\Gamma(n_I d_g/2)\,(\RSS_\cI/2)^{-n_I d_g/2}$, the standard
Jeffreys-marginalized form.

Because the model~$\theta$ is estimated from an
$m$-point minimal sample~$\cS$, the $n_\theta$~parameters
consumed by the fit reduce the effective degrees of freedom
from~$n_I\,d_g$ to~$n_I\,d_g - n_\theta$
in~\eqref{eq:marg-sigma}--\eqref{eq:invgamma-integral}.
When $m\,d_g = n_\theta$ (as for all models in
\Cref{tab:model-dims}), the minimal
sample constraints are satisfied exactly
($\bg_i = \bzero$ for $i \in \cS$) and
$n_I\,d_g - n_\theta = (n_I{-}m)\,d_g$; otherwise the
minimal sample is overdetermined and its residuals contribute
to~$\RSS_\cI$.  Applying this correction to the
marginal~\eqref{eq:marg-sigma} (replacing $n_I\,d_g$ with~$\nu$)
and taking logarithms gives the boxed scale-free marginal
score~\eqref{eq:score}.

\iffalse  % unreferenced; commented out for space
	% ---------------------------------------------------------------------
	\section{Specialization to Scalar Constraints ($d_g = 1$)}
	\label{sec:supp-scalar}
	% ---------------------------------------------------------------------

	For a single scalar constraint
	$g_i \coloneqq g(\bx_i,\theta)$,
	the projected covariance shape reduces to the scalar
	\begin{equation}
		c_i \;\coloneqq\;
		\pdv{g_i}!{\bx}\, \btSigma_{\bx_i}\,
		(\pdv{g_i}!{\bx})^\top,
		\label{eq:proj-var}
	\end{equation}
	and the weighted squared residual
	$r_i^2 = g_i^2/c_i \sim \sigma^2\,\chi^2_1$
	reduces the problem to the isotropic setting of
	Section~\ref{sec:score}.
	With $d_g = 1$, the score~\eqref{eq:score} becomes
	\begin{equation}
		S(\theta, \cI)
		= \log\Gamma\!\Bigl(\frac{n_I{-}n_\theta}{2}\Bigr)
		- \frac{n_I{-}n_\theta}{2}\,\log(\pi\,\RSS_\cI)
		- n_O\, \log(2a),
		\label{eq:weighted-score}
	\end{equation}
	where $\RSS_\cI = \sum_{i\in\cI} r_i^2 = \sum_{i\in\cI} g_i^2 / c_i$.
	When $\btSigma_{\bx_i} = \bI$ for all~$i$,
	$c_i = \|\pdv{g_i}!{\bx}\|^2$ and $r_i^2$ is the squared
	Sampson residual.

	\paragraph{Isotropic special case.}
	For isotropic observation noise ($\btSigma_{\bx_i} = \bI$) and a
	scalar constraint ($d_g = 1$),
	$\btSigma_{\bg_i}$ reduces to the scalar
	$c_i = \|\pdv{g_i}!{\bx}\|^2$, the residual
	$g_i^2/c_i$ is the squared Sampson
	residual~$r_i^2$~\cite{Hartley2003,Chum2005}, and
	$\RSS_\cI = \sum_{i\in\cI} r_i^2$;
	the score~\eqref{eq:weighted-score} is unchanged.

	% ---------------------------------------------------------------------
	\section{Notation Reference}
	\label{sec:supp-notation}
	% ---------------------------------------------------------------------

	\Cref{tab:notation} collects the principal symbols used
	throughout the paper.

	\begin{table}[ht]
		\centering
		\caption{Notation reference.}
		\label{tab:notation}
		\smallskip
		\begin{tabular}{@{}cl@{}}
			\toprule
			\textbf{Symbol}     & \textbf{Description}                                                                                            \\
			\midrule
			\multicolumn{2}{@{}l}{\emph{Data and model}}                                                                                          \\[2pt]
			$N$                 & number of observations (correspondences)                                                                        \\
			$\bx_i$             & $i$-th observation, $\bx_i \in \Reals^d$                                                                        \\
			$\bar\bx_i$         & true (noise-free) observation                                                                                   \\
			$d$                 & observation dimension                                                                                           \\
			$\theta$            & geometric model parameters, $\theta \in \Reals^{n_\theta}$                                                      \\
			$n_\theta$          & number of free model parameters                                                                                 \\
			$d_g$               & codimension of $\cM_\theta$ ($= d - \dim\cM_\theta$; constraints per point)                                     \\
			$m$                 & minimal sample size, $m = \lceil n_\theta / d_g \rceil$                                                         \\
			\midrule
			\multicolumn{2}{@{}l}{\emph{Constraints, residuals, and partitions}}                                                                  \\[2pt]
			$\bg_i$             & constraint vector at point $i$,
			$\bg_i \in \Reals^{d_g}$                                                                                                              \\
			$\bL_i$             & Cholesky factor of $\btSigma_{\bg_i}$:
			$\bL_i\bL_i^\top = \btSigma_{\bg_i}$                                                                                                  \\
			$\br_i$             & shape-whitened residual,
			$\br_i = \bL_i^{-1}\bg_i$;
			$\|\br_i\|^2 = \bg_i^\top\btSigma_{\bg_i}^{-1}\bg_i$                                                                                  \\
			$r_i$               & signed scalar residual ($d_g = 1$):
			$r_i = g_i/\sqrt{c_i}$;
			Sampson residual when $\btSigma_{\bx_i} = \bI$                                                                                        \\
			$z_i$               & latent inlier label, $z_i \in \{0,1\}$                                                                          \\
			$\bz$               & label vector, $\bz = (z_1,\ldots,z_N)$                                                                          \\
			$\cI,\;\cO$         & inlier / outlier sets;
			$\cO = \{1,\ldots,N\}\setminus\cI$                                                                                                    \\
			$n_I,\;n_O$         & $|\cI|$ and $|\cO|$                                                                                             \\
			$k$                 & running inlier count in the sweep ($k^* = n_I$ at the optimum)                                                  \\
			$q_i$               & per-point Mahalanobis residual,
			$q_i = \|\br_i\|^2
			= \bg_i^\top\btSigma_{\bg_i}^{-1}\bg_i$                                                                                               \\
			$\RSS_\cI$          & $\sum_{i\in\cI} q_i$;
			weighted residual sum of squares                                                                                                      \\
			$\hat s$            & inlier RMS, $\hat s^2 = \RSS_\cI / (n_I\,d_g - n_\theta)$;
			$n_I\,d_g - n_\theta$ effective DOF                                                                                                   \\
			\midrule
			\multicolumn{2}{@{}l}{\emph{Model parameters and priors}}                                                                             \\[2pt]
			$\sigma$            & inlier noise scale                                                                                              \\
			$a$                 & outlier domain half-width in shape-whitened constraint space~\eqref{eq:component-densities}; same units as~$\sigma$ \\
			$\gamma$            & prior inlier probability; set to $\tfrac{1}{2}$                                                                 \\
			$\tau$              & decision threshold (GaU / MSAC); implicit in our score                                                          \\
			\midrule
			\multicolumn{2}{@{}l}{\emph{Scores}}                                                                                                  \\[2pt]
			$S(\theta,\cI)$     & scale-free marginal score~\eqref{eq:score}                                                                      \\
			$S^*(\theta)$       & $\max_\cI S(\theta,\cI)$                                                                                        \\
			\midrule
			\multicolumn{2}{@{}l}{\emph{Jacobians}}                                                                                               \\[2pt]
			$\pdv{\bg_i}!{\bx}$ & constraint Jacobian w.r.t.\ observation,
			$\in \Reals^{d_g \times d}$                                                                                                           \\
			\midrule
			\multicolumn{2}{@{}l}{\emph{Covariances}}                                                                                             \\[2pt]
			$\bvarepsilon_i$    & observation-space error vector,
			$\bvarepsilon_i \in \Reals^d$                                                                                                         \\
			$\bSigma_{\bx_i}$   & observation covariance,
			$\bSigma_{\bx_i} = \sigma^2\btSigma_{\bx_i}$                                                                                          \\
			$\btSigma_{\bx_i}$  & observation covariance shape (known)                                                                            \\
			$\bSigma_{\bg_i}$   & constraint covariance,
			$= \sigma^2\btSigma_{\bg_i}$                                                                                                          \\
			$\btSigma_{\bg_i}$  & constraint covariance shape from
			measurement propagation                                                                                                               \\
			$c_i$               & projected scalar variance~\eqref{eq:proj-var}
			($d_g = 1$)                                                                                                                           \\
			\bottomrule
		\end{tabular}
	\end{table}
\fi  % end commented-out (sec:supp-scalar, sec:supp-notation)

\iffalse  % unreferenced; commented out for space
	\section{Derivation of the Constraint Covariance}
	\label{sec:supp-proj-cov}
	% ---------------------------------------------------------------------

	Treating $\theta$ as fixed, a first-order Taylor expansion of
	$\bg_i = \bg(\bx_i,\theta)$ about the noise-free
	observation~$\bar\bx_i$ gives
	\begin{equation}
		\bg_i
		\;\approx\;
		\underbrace{\bg(\bar\bx_i,\theta)}_{=\,\bzero}
		\;+\;
		\pdv{\bg_i}!{\bx}\,
		(\bx_i - \bar\bx_i)
		\;=\;
		J_{\bx}\,\bvarepsilon_i,
		\label{eq:taylor-g}
	\end{equation}
	where $\bvarepsilon_i \coloneqq \bx_i - \bar\bx_i$
	is the observation noise.  The covariance of $\bg_i$ is then
	$\bSigma_{\bg_i} = \Cov(\bg_i)
		= J_{\bx}\,\bSigma_{\bx_i}\,J_{\bx}^\top
		= \sigma^2\, J_{\bx}\,\btSigma_{\bx_i}\,J_{\bx}^\top$,
	so $\btSigma_{\bg_i} = J_{\bx}\,\btSigma_{\bx_i}\,J_{\bx}^\top$.
\fi  % end commented-out (sec:supp-proj-cov)



\iffalse  % unreferenced; commented out for space
	% =====================================================================
	\section{Fit Quality versus Consensus Size}
	\label{sec:supp-quality-vs-size}
	% =====================================================================

	We show that the marginal score can prefer a model with fewer
	inliers over one with more, provided the former has a tighter fit.
	For the isotropic $d_g = 1$ case, collecting the constant
	$-N\log(2a)$ (constant across
	partitions at fixed~$\theta$) and applying the Stirling
	approximation
	$\log\Gamma(n/2) \approx \tfrac{n}{2}\log\tfrac{n}{2} -
		\tfrac{n}{2}$ to the $n_I{-}n_\theta$ effective degrees of freedom
	(since $m = n_\theta$ for $d_g = 1$)
	gives
	\begin{equation}
		S + N\log(2a)
		\;=\; n_I\,\log(2a)
		- \tfrac{n_I{-}n_\theta}{2}\,\log(2\pi e\,\hat s^2)\,.
		\label{eq:per-inlier}
	\end{equation}
	For two models with the same~$n_\theta$, the
	$n_\theta\,\log(2a)$ contribution cancels and the comparison reduces to
	the effective inlier counts $n_I' = n_I - n_\theta$.
	Consider two models~$A$ and~$B$
	with $n_A' < n_B'$ effective inliers but $s_A < s_B$.
	Model~$A$ outscores model~$B$ whenever
	\begin{equation}
		\frac{n_A'}{n_B'}
		\;>\;
		\frac{\log(2a) - \tfrac{1}{2}\log(2\pi e\,s_B^2)}
		{\log(2a) - \tfrac{1}{2}\log(2\pi e\,s_A^2)}\,.
		\label{eq:quality-vs-size}
	\end{equation}
	Since $s_A < s_B$, the denominator exceeds the numerator, so the
	right-hand side is strictly less than~1.  Hence a sufficiently
	tight fit can compensate for a smaller consensus set.  For example,
	with $a = 50$\,px, $s_A = 1$\,px ($n_A' = 193$), and
	$s_B = 3$\,px ($n_B' = 243$): the right-hand side
	is ${\approx}\,0.66$, while $n_A'/n_B' = 0.79$, so model~$A$ wins
	despite having fewer inliers.

	% =====================================================================
	\section{Hyperparameter Grid: Coarse Probe and Fine Sweep}
	\label{sec:supp-a-grid}
	% =====================================================================

	The outlier half-width $a$ in~\eqref{eq:score} is the only
	hyperparameter of our score and is selected by held-out
	cross-validation, exactly as the $\tau$-baselines select their
	threshold $\tau$.  We apply a uniform two-stage tuning
	procedure to every method (ours and baselines), so that no
	method is advantaged by a wider search range, a more refined
	grid, or a better-calibrated centre point.

	\paragraph{Stage 1 -- coarse range probe (one-time).}
	For each (dataset, problem, protocol) cell and each of the six
	methods $\{$\textsc{IG-median}, \textsc{IG-mAA},
	\textsc{ransac}, \textsc{msac}, \textsc{gau},
	\textsc{magsac}$\}$, we sweep a 21-point log-spaced grid over
	\begin{equation*}
		x \in [\,\sigma_{\mathrm{med}}\!\cdot\!10^{-2.5},\;
		\sigma_{\mathrm{med}}\!\cdot\!10^{+2.5}\,],
		\quad \text{$0.25$-decade spacing,}
	\end{equation*}
	where $x = a$ for the IG family and $x = \tau$ for the
	$\tau$-baselines, and
	$\sigma_{\mathrm{med}} = \sqrt{\mathrm{median}_p
			\hat\sigma^2_p}$ is the cell's empirical noise scale measured
	from the GT inlier sets
	(\Cref{sec:supp-per-pair-sigma}).  The probe minimises the
	same per-protocol scalar objective that the Sensitivity
	panels later report---per-scene reduction then
	mean across scenes for hierarchical, flatten then reduce for
	no-split, flatten then 50/50 val split for flat---under either
	median rotation error or $-\mathrm{mAA}@10^\circ$, so the
	window is correctly centred for whichever criterion the cell
	ultimately tunes against.  The probe is GPU-accelerated for
	the IG family via a backend-portable
	\textsc{KernelAbstractions} kernel that is the \emph{same}
	kernel the per-trial Sensitivity tuning uses, so probe
	$x^\star$ and per-trial argmin agree to floating-point
	rounding by construction.

	\paragraph{Boundary extension.}
	If the coarse argmin lands within one grid step of either
	edge, we extend the grid outward by $4$~additional
	$0.25$-decade points on the violating side, then re-evaluate
	the argmin.  Up to two such extensions are allowed per side;
	beyond that the run is flagged ($\textsf{boundary\_hit{=}true}$)
	for review rather than silently centring on a possibly biased
	endpoint.  In production this safeguard fires only on
	methods/protocols where the optimal hyperparameter falls
	outside the data-driven $\sigma_{\mathrm{med}}$ envelope by
	more than two decades, which has not occurred on any
	production cell to date.

	\paragraph{Stage 2 -- fine evaluation grid (per trial).}
	Stage~1's per-(method, protocol) optimum~$x^\star$ defines a
	$2$-decade log-spaced fine grid
	\begin{equation*}
		x \in [\,x^\star/10,\; x^\star\cdot 10\,],
		\quad \text{$200$~log-spaced points,}
	\end{equation*}
	identical in width and resolution across all six methods.
	Stage~2 is the actual hyperparameter selection: every
	bootstrap trial samples $n$~validation pairs, evaluates the
	chosen objective at all $200$~grid points (one batched GPU
	sweep for the IG kernel; a histogram-trick CPU pass for the
	$\tau$-baselines), and reports the argmin.

	\paragraph{Why log-spaced fine grids.}
	Earlier drafts of this protocol used a linear-spaced fine
	grid (mirroring Shekhovtsov's $\tau$-grid
	convention~\cite{Shekhovtsov2025}); we switched to log-spaced
	because the per-method coarse probe already locates~$x^\star$
	to within a decade, and a linear grid over a $2$-decade
	log-window devotes most of its resolution to a single decade
	near the upper edge.  Log spacing within the same window
	gives uniform coverage in the units that actually matter for
	the score---scaling the grid by a factor of~$2$ moves equally
	many evaluation points across all decades---without changing
	the total per-trial tuning budget.  The fairness invariant is
	preserved: every method has the same $200$ trial-time
	evaluations across the same $2$-decade window, just with
	points placed where they distinguish the score's behaviour
	most.

	\paragraph{Per-cell reproducibility.}
	Stage~1's output is cached per
	$(\text{cell}, \text{method}, \text{protocol})$ and content-
	addressed by the SHA of the relevant bench source files plus
	$(\sigma^2_{\mathrm{med}}, \text{grid bounds},
		\text{seed})$, so re-running the pipeline reuses prior probes
	unchanged unless the scoring code or the data has changed.
	Sidecar files
	\texttt{<cell>/coarse\_ranges\_<protocol>.jld2} record the
	per-method curve $\{(x_i, y_i)\}_{i=1}^{21+\text{ext}}$,
	$x^\star$, the window endpoints, $\textsf{boundary\_hit}$, and
	$\sigma^2_{\mathrm{med}}$ so that downstream analyses can
	audit the centring decision without re-probing.
\fi  % end commented-out (sec:supp-quality-vs-size, sec:supp-a-grid)

% =====================================================================
\section{Posterior and Pseudo-Observation Interpretation}
\label{sec:supp-pseudo-obs}
% =====================================================================

The pseudo-observation reading of \Cref{par:pseudo-obs}---the
prior as $2\alpha_0$ pseudo-residuals carrying $2\beta_0$ of
pseudo-RSS via the posterior~\eqref{eq:ig-posterior}, combined
with the data as if the two samples were concatenated---yields
three regimes in the score:
\begin{itemize}
	\item \emph{Data-dominated} ($\nu \gg 2\alpha_0$): the prior is
	      negligible and the score~\eqref{eq:score} reduces to the
	      Jeffreys form~\eqref{eq:score-jeffreys} up to additive constants.
	\item \emph{Prior-dominated} ($\nu \ll 2\alpha_0$): the
	      partition's RSS contributes little and the score behaves as
	      if $\sigma^2$ were locked at $\beta_0/\alpha_0$.
	\item \emph{Balanced} ($\nu \sim 2\alpha_0$): the prior smooths
	      the per-partition scale estimate by the factor
	      $(\nu+2\alpha_0)/\nu$ toward $\beta_0/\alpha_0$, equivalent to
	      a weighted combination of the data RSS and the pseudo-RSS.
\end{itemize}
This makes $\alpha_0$ interpretable as the strength of the
regularization---$\alpha_0 = 1$ contributes two pseudo-residuals
(weak regularization, dominated by data once $\nu \gg 2$);
$\alpha_0 = 10$ contributes twenty (strong, comparable to a
typical inlier count and therefore more committed). The shape
$\alpha_0$ controls \emph{how much} the prior pulls; the ratio
$\beta_0/\alpha_0$ controls \emph{toward what scale}.

% =====================================================================
\section{The Score's Implicit Threshold}
\label{sec:implicit-threshold}
% =====================================================================

Although our score~\eqref{eq:score} eliminates~$\sigma$ as an
	explicit parameter, Algorithm~\ref{alg:score} selects an optimal
	inlier prefix $\cI^* = \{(1),\ldots,(k^*)\}$ whose boundary
	implicitly defines a threshold.  We derive the relationship
	between this implicit threshold, the outlier half-width~$a$, and
	the estimated scale~$\hat s$---first exactly, then via a
	large-sample approximation that reveals the connection to the GaU
	framework.  The analysis below uses the scalar case
	$d_g = 1$ (score~\eqref{eq:score}),
	where $\|\br_i\|^2$ reduces to the squared
	residual~$r_i^2$ and the boundary residual is
	$\tau^* = |r_{(k^*)}|$.
	For $d_g = 1$ the residual~$r_i$ is a scalar with
	$r_i \sim \NN(0,\,\sigma^2)$ under~$H_0$, and the
	outlier density reduces to the 1D uniform
	$\pout = 1/(2a)$ on~$[-a,a]$.  This
	scalar mixture model is the setting of
	Shekhovtsov's~\cite{Shekhovtsov2025} $\tau$-parameterized
	framework (Appendix~\ref{sec:supp-mixture}); the connection
	below relies on this $d_g{=}1$ reduction.

	\paragraph{Exact discrete condition.}
	Algorithm~\ref{alg:score} sweeps over sorted residuals and selects
	the prefix length~$k^*$ that maximizes~\eqref{eq:score}.  The
	optimal~$k^*$ is the largest~$k$ for which adding the $(k{+}1)$-th
	point still improves the score, i.e.,
	$\Delta S_k \coloneqq S(\theta,\cI_{k+1}) - S(\theta,\cI_k) \geq 0$:
	\begin{equation}
		\Delta S_k
		= \log\frac{\Gamma\!\bigl(\tfrac{k{-}n_\theta{+}1}{2}\bigr)}
		{\Gamma\!\bigl(\tfrac{k{-}n_\theta}{2}\bigr)}
		- \tfrac{1}{2}\log\pi
		+ \frac{k{-}n_\theta}{2}\,\log\RSS_k
		- \frac{k{-}n_\theta{+}1}{2}\,\log\!\bigl(\RSS_k + r_{(k+1)}^2\bigr)
		+ \log(2a)\,,
		\label{eq:delta-S}
	\end{equation}
	where $\RSS_k = \sum_{j=1}^{k} r_{(j)}^2$.
	The first term is
	the inlier reward from the $\log\Gamma$ factor; the middle terms
	capture the degradation in fit quality from including a larger
	residual; the last term is the outlier penalty
	saved.
	Equation~\eqref{eq:delta-S} is exact and valid for any inlier
	count, including the small sets that arise from minimal samples.

	\paragraph{Relation to MAP classification.}
	Given the predictive Student-$t$--uniform
	mixture, the MAP rule classifies
	point~$i$ as an inlier whenever
	$\pin^*(r_i) > \pout$.  Because the sweep
	operates on a discrete, sorted sequence, the optimal prefix
	boundary~$k^*$ does not in general land exactly on this MAP
	threshold: the true MAP boundary---where the Student-$t$ inlier
	density equals the uniform outlier density---lies between
	$|r_{(k^*)}|$ and $|r_{(k^*+1)}|$.  The sweep therefore
	\emph{bounds} the MAP threshold rather than computing it
	exactly.  Moreover, the MAP threshold itself is not a fixed
	quantity: it depends on the sufficient statistics
	$(\nu,\,\RSS_\cI)$, which change with the partition, creating
	a fixed-point problem.  The sweep resolves this circularity in a
	single pass by jointly optimizing the partition and the implicit
	scale.  In the large-sample limit (below), the discrete
	condition~\eqref{eq:delta-S} converges to the continuous MAP
	boundary, and the bound becomes tight.

	\paragraph{Large-sample approximation.}
	Using Stirling's approximation
	$\log\Gamma(n/2) \approx \tfrac{n}{2}\log\tfrac{n}{2} -
		\tfrac{n}{2}$ and treating the threshold
	continuously (\Cref{sec:supp-stirling}),
	the condition $\Delta S_k = 0$ reduces to
	\begin{equation}
		\boxed{
			\frac{1}{\sqrt{2\pi}\,\hat s}\,
			\exp\!\Bigl({-}\frac{{\tau^*}^2}{2\hat s^2}\Bigr)
			\;=\; \frac{1}{2a}\,,
		}
		\label{eq:implicit-tau}
	\end{equation}
	where $\hat s = \sqrt{\RSS_{k^*}/(k^*{-}n_\theta)}$ is the inlier
	RMS over the $k^*{-}n_\theta$ effective degrees of freedom
	(since $m = n_\theta$ for $d_g = 1$).  The
	left-hand side is the Gaussian inlier density
	$\pin(\tau^*;\hat s)$---the large-$\nu$ limit of the
	Student-$t$ predictive density
	$\pin^*(\tau^*)$---and the right-hand side is the
	uniform outlier density~$\pout$.
	This is precisely the MAP condition
	$\pin^*(\tau^*) = \pout$, confirming that
	the sweep's discrete bound converges to the MAP boundary in the
	large-sample limit.

	\paragraph{Connection to GaU scoring.}
	Recall that Shekhovtsov's GaU model classifies point~$i$ as an
	inlier when
	$\pin(r_i;\sigma) > \mu$, with
	$\mu = \tfrac{1-\gamma}{\gamma}\,\tfrac{1}{2a}$~\cite{Shekhovtsov2025}.
	The threshold $\tau_{\mathrm{GaU}}$ is the root of
	$\pin(\tau;\sigma) = \mu$.
	Comparing with~\eqref{eq:implicit-tau}, our score's implicit
	threshold satisfies the \emph{same} equation with
	\begin{equation}
		\frac{1{-}\gamma}{\gamma}
		\;\longrightarrow\;
		1\,,
		\qquad\text{i.e.,}\quad
		\gamma_{\mathrm{eff}}
		= \tfrac{1}{2}\,.
		\label{eq:gamma-eff}
	\end{equation}
	The effective prior is $\gamma_{\mathrm{eff}} = 1/2$,
	corresponding to maximum ignorance about inlier/outlier status.
	The non-informative choice $\gamma = \tfrac{1}{2}$ adopted in
	Section~\ref{sec:reverse-order} is therefore self-consistent:
	the score's implicit threshold reproduces the GaU threshold at
	exactly $\gamma = \tfrac{1}{2}$.  Crucially, GaU
	requires~$\sigma$ (and hence~$\tau$) to be supplied externally,
	whereas our method jointly estimates~$\hat s$ and selects the
	optimal partition from the data
	(Section~\ref{sec:score}).

	Solving~\eqref{eq:implicit-tau} for~$\tau^*$ in closed form:
	\begin{equation}
		\tau^*
		= \hat s\,\sqrt{2\,\log\!\bigl(a\sqrt{2/\pi}\,/\,\hat s\bigr)}\,.
		\label{eq:tau-star}
	\end{equation}
	For typical values (e.g., $a = 10\hat s$), this gives
	$\tau^* \approx 2.0\,\hat s$; for $a = 100\hat s$, we obtain
	$\tau^* \approx 3.0\,\hat s$---squarely in the 2--3$\sigma$ range
	commonly used in practice.

	\paragraph{Per-model adaptation.}
	Equation~\eqref{eq:tau-star} makes explicit a fundamental
	difference from the baseline framework: the implicit
	threshold~$\tau^*$ is \emph{not} a global constant.  Because
	$\hat s = \sqrt{\RSS_{k^*}/(k^*{-}n_\theta)}$ depends on the residuals of each
	candidate model~$\theta$, the threshold adapts automatically to
	the noise level of each hypothesis.  A well-fitting model with
	small inlier residuals produces a tight~$\hat s$ and hence a
	tight~$\tau^*$; a poor model with large residuals produces a
	loose one.  By contrast, the baselines apply a single, externally
	supplied~$\tau$ uniformly to all candidate models.  This
	per-model adaptation is a direct consequence of jointly estimating
	the scale and partition from the data, and it cannot be replicated
	by any global threshold.

	\paragraph{Asymptotic structural superiority over threshold-based scoring.}
	The per-model adaptation above translates into a strict
	ranking-quality advantage in the high-$\nu$ limit.  As $\nu \to
		\infty$, $\hat s \to \sigma_{\mathrm{true}}$ and the IG-marginal
	score (with any prior, including Jeffreys) converges to the GaU
	log-likelihood evaluated at the per-partition MLE
	$\sigma = \hat s$.  By construction, the per-partition MLE
	maximises the GaU likelihood for that partition.  Any threshold-based
	method evaluating GaU at a globally tuned $\tau$ (and hence at a
	fixed $\sigma$ that does not depend on the partition) achieves
	$\leq$ that value, with strict inequality whenever the data-implied
	$\hat s$ for the optimal partition differs from the global $\tau$.
	Equivalently, in the high-$\nu$ limit the IG-marginal score is the
	Bayes-optimal model-ranking objective under a Gaussian-Uniform
	mixture, and threshold-based scoring is its constrained
	specialisation that fixes $\sigma$ globally.

	This is the predicted structural superiority that the
	high-$\nu$ datasets (\Cref{fig:profile-grid}) verify
	empirically.
	The framework wins because it is the unconstrained version of the
	objective the threshold methods are constrained to.  In the
	low-$\nu$ regime the asymptotic argument breaks down because
	$\hat s$ is no longer a low-variance estimate of
	$\sigma_{\mathrm{true}}$; this is the regime where the informative
	prior re-establishes ranking quality by replacing the noisy
	data-only $\hat s$ with the posterior mean
	$(\RSS + 2\beta_0)/(\nu + 2\alpha_0)$.

% =====================================================================
\section{Derivation of the Large-Sample Optimality Condition}
	\label{sec:supp-stirling}
	% ---------------------------------------------------------------------

	We derive~\eqref{eq:implicit-tau} from the exact discrete
	condition~\eqref{eq:delta-S}.  Write $k' \coloneqq k - n_\theta$ for
	the effective degrees of freedom (recalling $m = n_\theta$ for
	$d_g = 1$).  For large~$k'$, the log-Gamma ratio
	satisfies
	\begin{equation}
		\log\frac{\Gamma\!\bigl(\tfrac{k'{+}1}{2}\bigr)}
		{\Gamma\!\bigl(\tfrac{k'}{2}\bigr)}
		\;\approx\; \tfrac{1}{2}\log(k'/2)\,,
		\label{eq:supp-loggamma}
	\end{equation}
	by Stirling's formula.  For the RSS terms, write
	$\RSS_{k+1} = \RSS_k + \tau^2$ and expand:
	\begin{align}
		 & \frac{k'}{2}\,\log\RSS_k
		- \frac{k'{+}1}{2}\,\log(\RSS_k + \tau^2) \notag                     \\
		 & \quad= -\frac{k'}{2}\,\log\!\Bigl(1 + \frac{\tau^2}{\RSS_k}\Bigr)
		- \frac{1}{2}\,\log(\RSS_k + \tau^2) \notag                          \\
		 & \quad\approx -\frac{\tau^2}{2\hat s^2}
		- \frac{1}{2}\,\log(k'\,\hat s^2)\,,
		\label{eq:supp-rss-expand}
	\end{align}
	where $\hat s^2 = \RSS_k/k'$ and we used
	$\log(1 + \tau^2/\RSS_k) \approx \tau^2/\RSS_k = \tau^2/(k'\hat s^2)$
	for $\tau^2 \ll \RSS_k$ (the boundary residual is small compared to
	the total), and $\log(\RSS_k + \tau^2) \approx \log(k'\hat s^2)$.
	Substituting~\eqref{eq:supp-loggamma}
	and~\eqref{eq:supp-rss-expand} into $\Delta S = 0$
	(including the $-\tfrac{1}{2}\log\pi$ term
	from~\eqref{eq:delta-S}):
	\begin{equation}
		\tfrac{1}{2}\log(k'/2)
		- \tfrac{1}{2}\log\pi
		- \frac{\tau^2}{2\hat s^2}
		- \tfrac{1}{2}\log(k'\,\hat s^2)
		+ \log(2a) = 0\,.
		\label{eq:supp-combined}
	\end{equation}
	Simplifying $\tfrac{1}{2}\log(k'/2) - \tfrac{1}{2}\log\pi
		- \tfrac{1}{2}\log(k'\,\hat s^2)
		= -\tfrac{1}{2}\log(2\pi \hat s^2)$:
	\begin{equation}
		-\tfrac{1}{2}\log(2\pi \hat s^2)
		- \frac{\tau^2}{2\hat s^2} + \log(2a) = 0\,.
		\label{eq:supp-simplified}
	\end{equation}
	Recognizing that $-\tfrac{1}{2}\log(2\pi \hat s^2) - \tau^2/(2\hat s^2)
		= \log \pin(\tau; \hat s)$ for a Gaussian density, we
	rewrite~\eqref{eq:supp-simplified} as
	\begin{equation}
		\log \pin(\tau; \hat s)
		= -\log(2a)\,,
	\end{equation}
	which is~\eqref{eq:implicit-tau}.

% =====================================================================
\section{Estimating the Prior Parameters $(\alpha_0, \beta_0)$}
\label{sec:supp-prior-estimation}
% =====================================================================

We consider two procedures for setting the prior parameters
from a validation set of per-pair scale estimates
$\{\hat\sigma^2_p\}_p$.  Per-pair $\hat\sigma^2_p$ may be
obtained either from an unconstrained model fit to the
ground-truth inlier set (whose residuals reflect within-plane
texture and the model class's flexibility) or from a fit
constrained by the oracle pose (whose residuals reflect any
deviation from a single rigid plane under that pose).  The
two estimators differ in which sources of variability they
absorb into the model versus into the residuals; we use the
unconstrained form for problems whose ground-truth inlier sets
are not strictly co-planar (e.g., facade-relief in
homography estimation) and the pose-constrained form when the
oracle pose tightly determines the model
(see Appendix~\ref{sec:supp-per-pair-sigma}).

\paragraph{Type-II maximum likelihood.}  Using the identity
$\sigma^2 \!\sim\! \IG(\alpha,\beta) \iff
	1/\sigma^2 \!\sim\! \mathrm{Gamma}(\alpha,\,1/\beta)$,
we fit a Gamma distribution to $\{1/\hat\sigma^2_p\}_p$ and
read off $(\alpha_0, \beta_0)$ from its shape and inverse
scale.  Both parameters are jointly identified, so the prior
captures both the location and the dispersion of the
validation $\hat\sigma^2_p$.  This estimator is preferred when
the per-pair $\hat\sigma^2_p$ are well concentrated, but is
sensitive to outliers: pairs with few inliers produce
$\hat\sigma^2_p$ estimates with $\sqrt{2/\nu_p}$ relative
standard error, and a small number of extreme values can drag
the joint MLE far from the bulk of the distribution.

\paragraph{Robust median anchor.}  We fix $\alpha_0$ at a weak
default (we use $\alpha_0 = 1$, i.e.\ two pseudo-observations)
and set $\beta_0 = \alpha_0 \cdot \mathrm{median}(\hat\sigma^2_p)$
so that the prior pseudo-scale $\beta_0/\alpha_0$ matches the
empirical median of $\{\hat\sigma^2_p\}$.  This sacrifices the
dispersion information captured by joint MLE in exchange for
robustness to heavy tails in $\{\hat\sigma^2_p\}$, which arise
on datasets with widely-varying per-pair inlier counts.

Both estimators are recomputed once per cross-validation trial
and held fixed for both the $a$-selection sweep and the test
evaluation.

% =====================================================================
\section{Per-Pair Scale Estimation for the Empirical-Bayes Prior}
\label{sec:supp-per-pair-sigma}
% =====================================================================

The IG prior parameters $(\alpha_0, \beta_0)$ are fit (via
either type-II MLE or the median anchor; see
\Cref{sec:supp-prior-estimation}) from a sample of per-pair scale
estimates $\{\hat\sigma^2_p\}_p$ obtained from the
ground-truth inlier set of each validation pair.  How
$\hat\sigma^2_p$ is computed is itself a methodological
choice: different ways of bringing the ground-truth pose into
the residual definition trade off bias from non-planarity
against variance from data scarcity.

\paragraph{Two natural estimators.}
Let $\cI_p^{\mathrm{gt}}$ be the labelled inlier set of pair
$p$ and let $H$ denote a homography (the analogous derivation
for relative pose replaces $H$ with $E$ throughout).  The two
natural choices are

\emph{Unconstrained model fit.}
Solve for $H$ by an unconstrained model fit on
$\cI_p^{\mathrm{gt}}$ (e.g., DLT for $H$, with $n_\theta = 8$
free parameters), giving $\hat H_{\mathrm{DLT}}$.  The per-pair
estimate is then
$\hat\sigma^2_p = \mathrm{RSS}_p / (d_g\,n_p^{\mathrm{in}} - n_\theta)$,
where $\mathrm{RSS}_p$ is the Sampson sum-of-squares of the
inliers under $\hat H_{\mathrm{DLT}}$.  All $n_\theta$ DOF of
the model are absorbed into the fit, so the residuals reflect
only what the unconstrained $H$ class cannot explain---measurement
noise plus any \emph{within-class} structure such as facade
relief that is not parameterizable by a single homography.

\emph{Pose-constrained fit.}
With the oracle relative pose $(R_p, t_p)$ known, the homography
of any plane in the scene takes the form
$H = R_p + t_p\,\tilde n^\top$ for some $\tilde n$
(the plane normal scaled by inverse plane depth).  Fitting only
the 3-DOF $\tilde n$ on $\cI_p^{\mathrm{gt}}$ gives a residual
DOF of $d_g\,n_p^{\mathrm{in}} - 3$ (more degrees of freedom for
the same inlier count) but a residual that includes any
deviation from the assumption that $\cI_p^{\mathrm{gt}}$ lies on
\emph{exactly one} rigid plane under $(R_p, t_p)$.

\paragraph{Trade-off.}
Each estimator absorbs different sources of variability into
the model versus the residual:
\begin{itemize}
	\item Unconstrained fit: absorbs within-plane geometric
	      variability into $\hat H$.  Residuals $\to$ measurement
	      noise.  Robust to imperfect coplanarity of the labelled
	      inliers.  Bias: tends to underestimate the noise scale,
	      because some of what we'd want to call ``inlier noise'' is
	      absorbed by the extra model DOF.
	\item Pose-constrained fit: locks the pose to the oracle.
	      Residuals $\to$ measurement noise \emph{plus} any
	      inconsistency with the single-plane assumption.  More DOF
	      per pair (so lower variance for fixed
	      $n_p^{\mathrm{in}}$).  Bias: overestimates the noise scale
	      in the presence of facade relief or labelling noise.
\end{itemize}
On the homography benchmarks we use, ground-truth inlier sets
are not strictly co-planar---facades exhibit several pixels of
relief and a small fraction of the labelled inliers are
geometrically inconsistent with a single plane at the oracle
pose---so the pose-constrained $\hat\sigma^2_p$ inflates by one
to several orders of magnitude on a non-trivial fraction of
pairs.  The resulting heavy tails in $\{\hat\sigma^2_p\}_p$
destabilize both estimators of $(\alpha_0, \beta_0)$.  We
therefore use the unconstrained DLT estimator for homography
estimation.  For relative-pose problems where the inlier set is
governed by exactly the model class being fitted (e.g., 5-pt
essential), the pose-constrained form coincides with the
unconstrained model fit and the distinction is moot.

\paragraph{Insufficient inliers.}
Either estimator requires $d_g\,n_p^{\mathrm{in}} > n_\theta'$
where $n_\theta'$ is the fitted-model DOF count (8 for
unconstrained $H$; 3 for pose-constrained).  Pairs failing
this floor produce no $\hat\sigma^2_p$ and are excluded from
the prior aggregation.  Their inclusion as test pairs is
unaffected: the per-partition score~\eqref{eq:score} requires
only that the partition itself satisfy $\nu = (k - m)\,d_g > 0$,
not that the GT inlier count be large enough to estimate
$\sigma^2$ robustly.



\iffalse  % unreferenced; commented out for space
	% =====================================================================
	\section{Calibrating the Ground-Truth Inlier Tolerance}
	\label{sec:supp-gt-calibration}
	% =====================================================================

	The protocol of Section~\ref{sec:eval-protocol} requires
	ground-truth inlier labels $\cI_p^{\mathrm{gt}}$ for the per-pair
	scale estimate $\hat\sigma^2_p$ that feeds the empirical-Bayes
	prior (Appendix~\ref{sec:supp-per-pair-sigma}) and, more
	broadly, for every per-cspond accuracy diagnostic on
	real-data benchmarks.  Producing these labels from a noisy
	ground-truth pose plus per-view depth maps requires a tolerance
	threshold whose value is dataset-dependent: laser-derived
	ETH3D depth has sub-millimetre noise, BundleFusion-tracked
	ScanNet1500 has centimetre-class noise, and so on.  Choosing
	this threshold by hand---or worse, sharing one across
	datasets---introduces a hidden lever that can silently bias all
	downstream results.  This appendix describes an empirical,
	matcher- and dataset-agnostic procedure for setting the threshold
	from the data, in physically meaningful units (metres at the
	surface).

	\paragraph{Why metric, not pixel, units.}
	Earlier iterations of the labelling pipeline used a pixel-residual
	threshold $\tau_{\mathrm{px}}$ on the forward depth-warp error.
	This conflates resolution and focal length with structural
	noise: a 3-pixel disagreement at $f{=}3408$, $z{=}10\text{m}$
	corresponds to $\sim$$9\text{mm}$ at the surface, but the same
			$3\text{px}$ at $f{=}585$, $z{=}1\text{m}$ corresponds to
			$\sim$$5\text{mm}$.  Cross-dataset cspond labels constructed
			this way are not apples-to-apples.  Instead we work directly with
			the 3D disagreement vector at the surface,
			\begin{equation}
				\label{eq:gt-residual}
				r^{A\to B}_p
				= \mathrm{unproj}\bigl(K_{\mathrm{tgt}}, \hat p_B, \hat d_B\bigr)
				- \mathrm{unproj}\bigl(K_{\mathrm{tgt}}, p_B,      d^{\mathrm{obs}}_B\bigr),
			\end{equation}
			where $\hat p_B$ and $\hat d_B$ are the forward-warp prediction
			of the cspond's source endpoint $(p_A, d^{\mathrm{obs}}_A)$ into
			the target view, $\mathrm{unproj}(K, u, d) = K^{-1}\,[u; 1]\,d$, and
		$d^{\mathrm{obs}}_A,\,d^{\mathrm{obs}}_B$ are the per-view
			depth-map values at the cspond endpoints.  Both terms live in the
			target camera frame in metres; their difference is the
			\emph{metric} disagreement at the surface, independent of
		$K$ and image resolution.  The intrinsic $K$ is in pixel units
			but the depth map is in metres, so the depth map provides the
			physical scale that $K$ alone lacks.

			\paragraph{Bidirectional residual.}
			Because the depth-warp prediction can be biased asymmetrically
			across views (one depth map may have a sharper edge than the
			other), we compute the metric residual in both directions and
			take the worse:
			\begin{equation}
				\|r_p\|
				= \max\bigl(\|r_p^{A\to B}\|,\;
				\|r_p^{B\to A}\|\bigr).
			\end{equation}
			A cspond is accepted only if both directions agree at the
			chosen tolerance, which guards against single-view occlusion
			edges that leak past the forward check alone.

			\paragraph{Calibration algorithm.}
			The tolerance is set empirically from the data via the procedure
			\textsc{calibrate-metric-tol}:

			\begin{enumerate}
				\item Sample $N_{\mathrm{sample}}$ pairs (default 20--30) evenly
				      from the dataset.
				\item For each sampled pair, run the unthresholded version
				      \textsc{gt-validate-metric}: the forward and backward
				      depth-consistency checks at relative depth tolerance
				      $\eta = 0.05$ (no metric threshold), keeping only csponds whose
				      predicted depth agrees with the observed depth on both views.
				      This is the structurally independent inlier filter, removing
				      occlusion edges, NaN-depth pixels, and gross outliers without
				      itself using any $r_p$ tolerance---so it does not bias the
				      noise distribution we are trying to measure.
				\item For each surviving cspond, compute~$\|r_p\|$ in metres.
				\item Estimate the noise scale robustly by the median absolute
				      deviation:
				      \begin{equation}
					      \sigma_{\mathrm{match}}
					      = 1.4826 \cdot
					      \mathrm{median}\!\left(
					      \bigl|\;\|r_p\| - \mathrm{median}_q(\|r_q\|)\;\bigr|
					      \right).
				      \end{equation}
				      The factor $1.4826 = 1/\Phi^{-1}(3/4)$ rescales the MAD to the
				      standard deviation of a Gaussian.  Unlike $\widehat{\mathrm{std}}$,
				      this estimator has a 50\% breakdown point: any small fraction of
				      residual outliers that slipped past the depth-consistency check
				      cannot inflate~$\sigma_{\mathrm{match}}$.
				\item Set the tolerance to a multiple of the robust scale:
				      \begin{equation}
					      \tau_{\mathrm{metric}} = K_{\sigma}\,\sigma_{\mathrm{match}},
					      \qquad K_{\sigma} = 3
				      \end{equation}
				      capturing $\sim$99.7\% of the true-inlier mass under a Gaussian
				      approximation.  We use a Unitful representation
				      ($\mathrm{Quantity}\langle\mathrm{Float64}, \mathrm{Length}\rangle$;
				      e.g.\ $14.1\text{mm}$) so the threshold is portable across
				      consumers without unit confusion.
			\end{enumerate}

			\paragraph{Why MAD, not std or a percentile.}
			Standard deviation is biased high by any residual outliers
			\textsc{gt-validate-metric} fails to filter (especially on
			occlusion edges, where the depth flips smoothly across two
			surfaces).  The 95th-percentile of $\|r_p\|$ is a valid
			alternative but is harder to interpret as a noise floor: it
			answers ``where does 95\% of the residual mass lie'' rather
			than ``what is~$\sigma$.''  The MAD-based estimator is the
			matched convention to the threshold our scoring loop uses
			(also $K_\sigma\cdot\hat\sigma$ with $K_\sigma{=}3$), so the
			tolerance is in the same scale family as everything else.

			\paragraph{Empirical values.}
			\Cref{tab:gt-tolerance-empirical} reports the calibrated
		$\sigma_{\mathrm{match}}$ and $\tau_{\mathrm{metric}}$ for the two
			real-data benchmarks used in this paper, run on a 20--30 pair
			sample with our default RoMa-V2 csponds.  ETH3D's noise
			floor is sub-cm (laser depth + COLMAP-BA poses); ScanNet1500
			sits roughly $4\!\times$ higher (Kinect depth + BundleFusion
			poses), exactly the ratio expected from sensor specifications.
			The validation rate---fraction of csponds that pass the
			unthresholded forward$+$backward depth-consistency check---is
			much lower for ETH3D ($4.2\%$) than ScanNet ($66.7\%$) because
			laser-scanned depth is sparse where Kinect depth is dense; both
			remain large enough samples for stable scale estimation
			($\geq 4\,000$ csponds).
			%
			\begin{table}[t]
				\centering\small
				\begin{tabular}{lcccc}
					\toprule
					Dataset     & $\sigma_{\mathrm{match}}$ & $\tau_{\mathrm{metric}}$
					            & $n_{\mathrm{validated}}$  & validation rate          \\
					\midrule
					ETH3D       & $4.7\,$mm                 & $14.1\,$mm
					            & $4{,}245$                 & $4.2\%$                  \\
					ScanNet1500 & $20.9\,$mm                & $62.6\,$mm
					            & $100{,}033$               & $66.7\%$                 \\
					\bottomrule
				\end{tabular}
				\caption{Calibrated GT inlier tolerances per dataset.  Both
					values produced by \textsc{calibrate-metric-tol} on a small
					dataset-stratified sample with $K_\sigma{=}3$ on the worst-direction
					3D residual~\eqref{eq:gt-residual}.}
				\label{tab:gt-tolerance-empirical}
			\end{table}

			\paragraph{Matcher dependence.}
			The structural component of $\sigma_{\mathrm{match}}$ (sensor
			depth noise plus pose drift) is independent of which matcher
			produced the csponds; only the cspond-localization noise depends
			on the matcher.  At typical depths and focal lengths considered
			here, the matcher contribution is sub-millimetre while the
			structural contribution is tens of millimetres
			(\Cref{tab:gt-tolerance-empirical}), so reusing
		$\tau_{\mathrm{metric}}$ across matchers
			(RoMa-V2 versus SuperPoint+LightGlue) introduces at most a
			few-percent error.  We nonetheless re-run
			\textsc{calibrate-metric-tol} per matcher for rigor, and report
			matcher-specific values when the recalibration moves the
			tolerance by more than~$\sim$10\%.

			\paragraph{Refinements.}
			A few refinements not used in the headline numbers but worth
			noting:
			\emph{(i) Stratified sampling.}  Even-by-index sampling
			matches ScanNet1500's roughly uniform pair density but
			over-weights covisibility-rich scenes in ETH3D; a
			scene-stratified sampler would tighten $\sigma_{\mathrm{match}}$
			by a few percent on the latter.
			\emph{(ii) Anisotropic noise.}  Decomposing $r_p$ into
			components parallel and perpendicular to the line of sight
			($\sigma_\parallel,\sigma_\perp$) and thresholding per-axis
			yields a tighter inlier set for sensors whose depth and
			lateral noise differ in scale (true for both Kinect and laser),
			at the cost of two thresholds rather than one.
			\emph{(iii) Depth-relative tolerance.}  When residuals scale
			with~$z$ (e.g.\ KITTI's lidar at long range), a relative
			tolerance $\tau_{\mathrm{rel}} = \kappa\,z$ removes the systematic
		$z$-dependence at the cost of an extra hyperparameter; we have
	not needed this on ETH3D or ScanNet1500.
\fi  % end commented-out (sec:supp-gt-calibration)

% =====================================================================
\section{Dataset Details}
\label{sec:supp-datasets}
% =====================================================================

\Cref{tab:datasets} summarizes the three datasets, which span
three axes of variation---dataset, geometric problem, and
feature pipeline.

\begin{table}[h]
	\centering
	\small
	\caption{Datasets used in our evaluation: geometric problem,
		feature pipeline, bootstrap protocol, number of scenes, and
		total image-pair count.}
	\label{tab:datasets}
	\begin{tabular}{l@{~}r c l@{~}r l c c}
		\toprule
		\multicolumn{2}{l}{Dataset} & Problem & \multicolumn{2}{l}{Features} & Protocol & Scenes & Pairs \\
		\midrule
		PhotoTourism & \cite{Jin2021}    & $\bE$ & RootSIFT & \cite{Arandjelovic2012}              & Flat (resample) & 21 & 44\,211 \\
		ETH3D        & \cite{Schoeps2017} & $\bF$ & SP+LG    & \cite{DeTone2018,Lindenberger2023}   & Flat (resample) & 13 &  3\,039 \\
		HEB          & \cite{Barath2023}  & $\bH$ & SIFT     & \cite{Lowe2004}                      & Hierarchical    & 10 & 21\,999 \\
		\bottomrule
	\end{tabular}
\end{table}

The \emph{Homography Estimation
	Benchmark} (HEB)~\cite{Barath2023} provides 10 architectural
sites of building facades (21\,999 pairs) for
homography~($\bH$) estimation.  \emph{PhotoTourism}
(PT)~\cite{Jin2021,RansacTutorial2020} provides 21 landmark
scenes from the Image Matching
Challenge\footnote{No canonical PT split exists; we use the
	publicly-available scenes with GT poses and SfM matches.}
(44\,211 pairs) captured under varying illumination and
viewpoint, for essential-matrix ($\bE$) estimation.
\emph{ETH3D}~\cite{Schoeps2017} provides 13 indoor and outdoor
sub-scenes (3\,039 pairs) with laser-scanned per-view depth
for fundamental-matrix ($\bF$) estimation.

\paragraph{Annotations.}
Each dataset provides stereo pairs with ground-truth absolute
pose and camera calibration, from which we construct the
ground-truth relative pose.  HEB additionally annotates the
homographies induced by the dominant planes in each scene.

Each scoring method ranks a set of \emph{tentative
	correspondences} per pair; IG-MLE and IG-Robust additionally
consume \emph{ground-truth point matches} during validation
(\Cref{sec:eval-protocol}).  HEB and PhotoTourism ship both a
tentative correspondence set---SIFT~\cite{Lowe2004} on HEB,
RootSIFT~\cite{Arandjelovic2012} on PhotoTourism---and
SfM bundle-adjusted GT matches (for HEB, restricted to the
annotated plane).  These GT matches are
pseudo-ground-truth---a multi-view SfM reference treated as
gold-standard for the two-view sub-problem---and their
reprojection residuals also seed our per-pair variance
estimates.  ETH3D ships only per-view laser depth (training
subset only; the test set is unreleased); we run
SuperPoint~\cite{DeTone2018}\,+\,LightGlue~\cite{Lindenberger2023}
(SP+LG), a learned detect-and-match pipeline, as the
tentative set, and verify each candidate by 3D consistency:
pixels are unprojected via the depth, brought into a common
frame by the GT poses, and accepted when the resulting 3D
points lie within a dataset-specific threshold (reflecting
feature, depth-sensor, and absolute-pose accuracy).  Pairing
SP+LG with SIFT/RootSIFT on the other two datasets exercises
our methods against feature noise from both hand-designed and
learned pipelines.  One caveat for HEB: the mapping from
SfM-derived ground-truth correspondences to the annotated
dominant planes is not always clean---a pair's matches can
span multiple planes, so a scoring method may select a model
that fits one plane while the GT inliers we score against
belong to another.  This shows up downstream as anomalous
residual structure when a selected model is evaluated against
the GT inlier set, and surfaces in some HEB panels of the
experiments.

% =====================================================================
%  NeurIPS Paper Checklist (omitted from the arXiv version)
% =====================================================================
% \newpage
% \input{checklist.tex}

\end{document}

